\chapt{अयोध्याकाण्डः}

\sect{प्रथमः सर्गः}

\twolineshloka
{कस्यचित्त्वथ कालस्य राजा दशरथः सुतम्}
{भरतं केकयीपुत्रमब्रवीद्रघुनन्दनः} %||2-1-1||

\twolineshloka
{अयं केकयराजस्य पुत्रो वसति पुत्रक}
{त्वां नेतुमागतो वीर युधाजिन्मातुलस्तव} %||2-1-2||

\twolineshloka
{श्रुत्वा दशरथस्यैतद्भरतः केकयीसुतः}
{गमनायाभिचक्राम शत्रुघ्नसहितस्तदा} %||2-1-3||

\twolineshloka
{आपृच्छ्य पितरं शूरो रामं चाक्लिष्टकारिणम्}
{मातॄंश्चापि नरश्रेष्ठः शत्रुघ्नसहितो ययौ} %||2-1-4||

\twolineshloka
{युधाजित्प्राप्य भरतं सशत्रुघ्नं प्रहर्षितः}
{स्वपुरं प्राविशद्वीरः पिता तस्य तुतोष ह} %||2-1-5||

\twolineshloka
{स तत्र न्यवसद्भ्रात्रा सह सत्कारसत्कृतः}
{मातुलेनाश्वपतिना पुत्रस्नेहेन लालितः} %||2-1-6||

\twolineshloka
{तत्रापि निवसन्तौ तौ तर्प्यमाणौ च कामतः}
{भ्रातरौ स्मरतां वीरौ वृद्धं दशरथं नृपम्} %||2-1-7||

\twolineshloka
{राजापि तौ महातेजाः सस्मार प्रोषितौ सुतौ}
{उभौ भरतशत्रुघ्नौ महेन्द्रवरुणोपमौ} %||2-1-8||

\twolineshloka
{सर्व एव तु तस्येष्टाश्चत्वारः पुरुषर्षभाः}
{स्वशरीराद्विनिर्वृत्ताश्चत्वार इव बाहवः} %||2-1-9||

\twolineshloka
{तेषामपि महातेजा रामो रतिकरः पितुः}
{स्वयम्भूरिव भूतानां बभूव गुणवत्तरः} %||2-1-10||

\twolineshloka
{गते च भरते रामो लक्ष्मणश्च महाबलः}
{पितरं देवसङ्काशं पूजयामासतुस्तदा} %||2-1-11||

\twolineshloka
{पितुराज्ञां पुरस्कृत्य पौरकार्याणि सर्वशः}
{चकार रामो धर्मात्मा प्रियाणि च हितानि च} %||2-1-12||

\twolineshloka
{मातृभ्यो मातृकार्याणि कृत्वा परमयन्त्रितः}
{गुरूणां गुरुकार्याणि काले कालेऽन्ववैक्षत} %||2-1-13||

\twolineshloka
{एवं दशरथः प्रीतो ब्राह्मणा नैगमास्तथा}
{रामस्य शीलवृत्तेन सर्वे विषयवासिनः} %||2-1-14||

\twolineshloka
{स हि नित्यं प्रशान्तात्मा मृदुपूर्वं च भाषते}
{उच्यमानोऽपि परुषं नोत्तरं प्रतिपद्यते} %||2-1-15||

\twolineshloka
{कथञ्चिदुपकारेण कृतेनैकेन तुष्यति}
{न स्मरत्यपकाराणां शतमप्यात्मवत्तया} %||2-1-16||

\twolineshloka
{शीलवृद्धैर्ज्ञानवृद्धैर्वयोवृद्धैश्च सज्जनैः}
{कथयन्नास्त वै नित्यमस्त्रयोग्यान्तरेष्वपि} %||2-1-17||

\twolineshloka
{कल्याणाभिजनः साधुरदीनः सत्यवागृजुः}
{वृद्धैरभिविनीतश्च द्विजैर्धर्मार्थदर्शिभिः} %||2-1-18||

\twolineshloka
{धर्मार्थकामतत्त्वज्ञः स्मृतिमान्प्रतिभावनान्}
{लौकिके समयाचरे कृतकल्पो विशारदः} %||2-1-19||

\twolineshloka
{शास्त्रज्ञश्च कृतज्ञश्च पुरुषान्तरकोविदः}
{यः प्रग्रहानुग्रहयोर्यथान्यायं विचक्षणः} %||2-1-20||

\twolineshloka
{आयकर्मण्युपायज्ञः सन्दृष्टव्ययकर्मवित्}
{श्रैष्ठ्यं शास्त्रसमूहेषु प्राप्तो व्यामिश्रकेष्वपि} %||2-1-21||

\twolineshloka
{अर्थधर्मौ च सङ्गृह्य सुखतन्त्रो न चालसः}
{वैहारिकाणां शिल्पानां विज्ञातार्थविभागवित्} %||2-1-22||

\twolineshloka
{आरोहे विनये चैव युक्तो वारणवाजिनाम्}
{धनुर्वेदविदां श्रेष्ठो लोकेऽतिरथसम्मतः} %||2-1-23||

\twolineshloka
{अभियाता प्रहर्ता च सेनानयविशारदः}
{अप्रधृष्यश्च सङ्ग्रामे क्रुद्धैरपि सुरासुरैः} %||2-1-24||

\twolineshloka
{अनसूयो जितक्रोधो न दृप्तो न च मत्सरी}
{न चावमन्ता भूतानां न च कालवशानुगः} %||2-1-25||

\threelineshloka
{एवं श्रैष्ठैर्गुणैर्युक्तः प्रजानां पार्थिवात्मजः}
{सम्मतस्त्रिषु लोकेषु वसुधायाः क्षमागुणैः}
{बुद्ध्या बृहस्पतेस्तुल्यो वीर्येणापि शचीपतेः} %||2-1-26||

\twolineshloka
{तथा सर्वप्रजाकान्तैः प्रीतिसंजननैः पितुः}
{गुणैर्विरुरुचे रामो दीप्तः सूर्य इवांशुभिः} %||2-1-27||

\twolineshloka
{तमेवंवृत्तसम्पन्नमप्रधृष्य पराक्रमम्}
{लोकपालोपमं नाथमकामयत मेदिनी} %||2-1-28||

\twolineshloka
{एतैस्तु बहुभिर्युक्तं गुणैरनुपमैः सुतम्}
{दृष्ट्वा दशरथो राजा चक्रे चिन्तां परन्तपः} %||2-1-29||

\twolineshloka
{एषा ह्यस्य परा प्रीतिर्हृदि सम्परिवर्तते}
{कदा नाम सुतं द्रक्ष्याम्यभिषिक्तमहं प्रियम्} %||2-1-30||

\twolineshloka
{वृद्धिकामो हि लोकस्य सर्वभूतानुकम्पनः}
{मत्तः प्रियतरो लोके पर्जन्य इव वृष्टिमान्} %||2-1-31||

\twolineshloka
{यमशक्रसमो वीर्ये बृहस्पतिसमो मतौ}
{महीधरसमो धृत्यां मत्तश्च गुणवत्तरः} %||2-1-32||

\twolineshloka
{महीमहमिमां कृत्स्नामधितिष्ठन्तमात्मजम्}
{अनेन वयसा दृष्ट्वा यथा स्वर्गमवाप्नुयाम्} %||2-1-33||

\twolineshloka
{तं समीक्ष्य महाराजो युक्तं समुदितैर्गुणैः}
{निश्चित्य सचिवैः सार्धं युवराजममन्यत} %||2-1-34||

\twolineshloka
{नानानगरवास्तव्यान्पृथग्जानपदानपि}
{समानिनाय मेदिन्याः प्रधानान्पृथिवीपतिः} %||2-1-35||

\twolineshloka
{अथ राजवितीर्णेषु विविधेष्वासनेषु च}
{राजानमेवाभिमुखा निषेदुर्नियता नृपाः} %||2-1-36||

\fourlineindentedshloka
{स लब्धमानैर्विनयान्वितैर्नृपैः}
{ पुरालयैर्जानपदैश्च मानवैः}
{उपोपविष्टैर्नृपतिर्वृतो बभौ}
{ सहस्रचक्षुर्भगवानिवामरैः} %||2-2-37||


{॥इत्यार्षे श्रीमद्रामायणे वाल्मीकीये आदिकाव्ये अयोध्याकाण्डे प्रथमः सर्गः॥}

\sect{द्वितीयः सर्गः}

\twolineshloka
{ततः परिषदं सर्वामामन्त्र्य वसुधाधिपः}
{हितमुद्धर्षणं चेदमुवाचाप्रतिमं वचः} %||2-2-1||

\twolineshloka
{दुन्दुभिस्वनकल्पेन गम्भीरेणानुनादिना}
{स्वरेण महता राजा जीग्मूत इव नादयन्} %||2-2-2||

\twolineshloka
{सोऽहमिक्ष्वाकुभिः पूर्वैर्नरेन्द्रैः परिपालितम्}
{श्रेयसा योक्तुकामोऽस्मि सुखार्हमखिलं जगत्} %||2-2-3||

\twolineshloka
{मयाप्याचरितं पूर्वैः पन्थानमनुगच्छता}
{प्रजा नित्यमतन्द्रेण यथाशक्त्यभिरक्षता} %||2-2-4||

\twolineshloka
{इदं शरीरं कृत्स्नस्य लोकस्य चरता हितम्}
{पाण्डुरस्यातपत्रस्यच्छायायां जरितं मया} %||2-2-5||

\twolineshloka
{प्राप्य वर्षसहस्राणि बहून्यायूंषि जीवितः}
{जीर्णस्यास्य शरीरस्य विश्रान्तिमभिरोचये} %||2-2-6||

\twolineshloka
{राजप्रभावजुष्टां हि दुर्वहामजितेन्द्रियैः}
{परिश्रान्तोऽस्मि लोकस्य गुर्वीं धर्मधुरं वहन्} %||2-2-7||

\twolineshloka
{सोऽहं विश्रममिच्छामि पुत्रं कृत्वा प्रजाहिते}
{संनिकृष्टानिमान्सर्वाननुमान्य द्विजर्षभान्} %||2-2-8||

\twolineshloka
{अनुजातो हि मे सर्वैर्गुणैर्ज्येष्ठो ममात्मजः}
{पुरन्दरसमो वीर्ये रामः परपुरंजयः} %||2-2-9||

\twolineshloka
{तं चन्द्रमिव पुष्येण युक्तं धर्मभृतां वरम्}
{यौवराज्येन योक्तास्मि प्रीतः पुरुषपुङ्गवम्} %||2-2-10||

\twolineshloka
{अनुरूपः स वो नाथो लक्ष्मीवाँल्लक्ष्मणाग्रजः}
{त्रैलोक्यमपि नाथेन येन स्यान्नाथवत्तरम्} %||2-2-11||

\twolineshloka
{अनेन श्रेयसा सद्यः संयोज्याहमिमां महीम्}
{गतक्लेशो भविष्यामि सुते तस्मिन्निवेश्य वै} %||2-2-12||

\twolineshloka
{इति ब्रुवन्तं मुदिताः प्रत्यनन्दन्नृपा नृपम्}
{वृष्टिमन्तं महामेघं नर्दन्तमिव बर्हिणः} %||2-2-13||

\twolineshloka
{तस्य धर्मार्थविदुषो भावमाज्ञाय सर्वशः}
{ऊचुश्च मनसा ज्ञात्वा वृद्धं दशरथं नृपम्} %||2-2-14||

\twolineshloka
{अनेकवर्षसाहस्रो वृद्धस्त्वमसि पार्थिव}
{स रामं युवराजानमभिषिञ्चस्व पार्थिवम्} %||2-2-15||

\twolineshloka
{इति तद्वचनं श्रुत्वा राजा तेषां मनःप्रियम्}
{अजानन्निव जिज्ञासुरिदं वचनमब्रवीत्} %||2-2-16||

\twolineshloka
{कथं नु मयि धर्मेण पृथिवीमनुशासति}
{भवन्तो द्रष्टुमिच्छन्ति युवराजं ममात्मजम्} %||2-2-17||

\twolineshloka
{ते तमूचुर्महात्मानं पौरजानपदैः सह}
{बहवो नृप कल्याणा गुणाः पुत्रस्य सन्ति ते} %||2-2-18||

\twolineshloka
{दिव्यैर्गुणैः शक्रसमो रामः सत्यपराक्रमः}
{इक्ष्वाकुभ्यो हि सर्वेभ्योऽप्यतिरक्तो विशाम्पते} %||2-2-19||

\twolineshloka
{रामः सत्पुरुषो लोके सत्यधर्मपरायणः}
{धर्मज्ञः सत्यसन्धश्च शीलवाननसूयकः} %||2-2-20||

\twolineshloka
{क्षान्तः सान्त्वयिता श्लक्ष्णः कृतज्ञो विजितेन्द्रियः}
{मृदुश्च स्थिरचित्तश्च सदा भव्योऽनसूयकः} %||2-2-21||

\twolineshloka
{प्रियवादी च भूतानां सत्यवादी च राघवः}
{बहुश्रुतानां वृद्धानां ब्राह्मणानामुपासिता} %||2-2-22||

\twolineshloka
{तेनास्येहातुला कीर्तिर्यशस्तेजश्च वर्धते}
{देवासुरमनुष्याणां सर्वास्त्रेषु विशारदः} %||2-2-23||

\twolineshloka
{यदा व्रजति सङ्ग्रामं ग्रामार्थे नगरस्य वा}
{गत्वा सौमित्रिसहितो नाविजित्य निवर्तते} %||2-2-24||

\twolineshloka
{सङ्ग्रामात्पुनरागम्य कुञ्जरेण रथेन वा}
{पौरान्स्वजनवन्नित्यं कुशलं परिपृच्छति} %||2-2-25||

\twolineshloka
{पुत्रेष्वग्निषु दारेषु प्रेष्यशिष्यगणेषु च}
{निखिलेनानुपूर्व्या च पिता पुत्रानिवौरसान्} %||2-2-26||

\twolineshloka
{शुश्रूषन्ते च वः शिष्याः कच्चित्कर्मसु दंशिताः}
{इति नः पुरुषव्याघ्रः सदा रामोऽभिभाषते} %||2-2-27||

\twolineshloka
{व्यसनेषु मनुष्याणां भृशं भवति दुःखितः}
{उत्सवेषु च सर्वेषु पितेव परितुष्यति} %||2-2-28||

\threelineshloka
{सत्यवादी महेष्वासो वृद्धसेवी जितेन्द्रियः}
{वत्सः श्रेयसि जातस्ते दिष्ट्यासौ तव राघवः}
{दिष्ट्या पुत्रगुणैर्युक्तो मारीच इव कश्यपः} %||2-2-29||

\twolineshloka
{बलमारोग्यमायुश्च रामस्य विदितात्मनः}
{आशंसते जनः सर्वो राष्ट्रे पुरवरे तथा} %||2-2-30||

\twolineshloka
{अभ्यन्तरश्च बाह्यश्च पौरजानपदो जनः}
{स्त्रियो वृद्धास्तरुण्यश्च सायम्प्रातः समाहिताः} %||2-2-31||

\twolineshloka
{सर्वान्देवान्नमस्यन्ति रामस्यार्थे यशस्विनः}
{तेषामायाचितं देव त्वत्प्रसादात्समृध्यताम्} %||2-2-32||

\twolineshloka
{राममिन्दीवरश्यामं सर्वशत्रुनिबर्हणम्}
{पश्यामो यौवराज्यस्थं तव राजोत्तमात्मजम्} %||2-2-33||

\fourlineindentedshloka
{तं देवदेवोपममात्मजं ते}
{ सर्वस्य लोकस्य हिते निविष्टम्}
{हिताय नः क्षिप्रमुदारजुष्टम्}
{ मुदाभिषेक्तुं वरद त्वमर्हसि} %||2-3-34||


{॥इत्यार्षे श्रीमद्रामायणे वाल्मीकीये आदिकाव्ये अयोध्याकाण्डे द्वितीयः सर्गः॥}

\sect{तृतीयः सर्गः}

\twolineshloka
{तेषामज्ञलिपद्मानि प्रगृहीतानि सर्वशः}
{प्रतिगृह्याब्रवीद्राजा तेभ्यः प्रियहितं वचः} %||2-3-1||

\twolineshloka
{अहोऽस्मि परमप्रीतः प्रभावश्चातुलो मम}
{यन्मे ज्येष्ठं प्रियं पुत्रं यौवराज्यस्थमिच्छथ} %||2-3-2||

\twolineshloka
{इति प्रत्यर्च्य तान्राजा ब्राह्मणानिदमब्रवीत्}
{वसिष्ठं वामदेवं च तेषामेवोपशृण्वताम्} %||2-3-3||

\twolineshloka
{चैत्रः श्रीमानयं मासः पुण्यः पुष्पितकाननः}
{यौवराज्याय रामस्य सर्वमेवोपकल्प्यताम्} %||2-3-4||

\twolineshloka
{कृतमित्येव चाब्रूतामभिगम्य जगत्पतिम्}
{यथोक्तवचनं प्रीतौ हर्षयुक्तौ द्विजर्षभौ} %||2-3-5||

\twolineshloka
{ततः सुमन्त्रं द्युतिमान्राजा वचनमब्रवीत्}
{रामः कृतात्मा भवता शीघ्रमानीयतामिति} %||2-3-6||

\twolineshloka
{स तथेति प्रतिज्ञाय सुमन्त्रो राजशासनात्}
{रामं तत्रानयां चक्रे रथेन रथिनां वरम्} %||2-3-7||

\twolineshloka
{अथ तत्र समासीनास्तदा दशरथं नृपम्}
{प्राच्योदीच्याः प्रतीच्याश्च दाक्षिणात्याश्च भूमिपाः} %||2-3-8||

\twolineshloka
{म्लेच्छाश्चार्याश्च ये चान्ये वनशैलान्तवासिनः}
{उपासां चक्रिरे सर्वे तं देवा इव वासवम्} %||2-3-9||

\twolineshloka
{तेषां मध्ये स राजर्षिर्मरुतामिव वासवः}
{प्रासादस्थो रथगतं ददर्शायान्तमात्मजम्} %||2-3-10||

\twolineshloka
{गन्धर्वराजप्रतिमं लोके विख्यातपौरुषम्}
{दीर्घबाहुं महासत्त्वं मत्तमातङ्गगामिनम्} %||2-3-11||

\twolineshloka
{चन्द्रकान्ताननं राममतीव प्रियदर्शनम्}
{रूपौदार्यगुणैः पुंसां दृष्टिचित्तापहारिणम्} %||2-3-12||

\twolineshloka
{घर्माभितप्ताः पर्जन्यं ह्लादयन्तमिव प्रजाः}
{न ततर्प समायान्तं पश्यमानो नराधिपः} %||2-3-13||

\twolineshloka
{अवतार्य सुमन्त्रस्तं राघवं स्यन्दनोत्तमात्}
{पितुः समीपं गच्छन्तं प्राञ्जलिः पृष्ठतोऽन्वगात्} %||2-3-14||

\twolineshloka
{स तं कैलासशृङ्गाभं प्रासादं नरपुङ्गवः}
{आरुरोह नृपं द्रष्टुं सह सूतेन राघवः} %||2-3-15||

\twolineshloka
{स प्राञ्जलिरभिप्रेत्य प्रणतः पितुरन्तिके}
{नाम स्वं श्रावयन्रामो ववन्दे चरणौ पितुः} %||2-3-16||

\twolineshloka
{तं दृष्ट्वा प्रणतं पार्श्वे कृताञ्जलिपुटं नृपः}
{गृह्याञ्जलौ समाकृष्य सस्वजे प्रियमात्मजम्} %||2-3-17||

\twolineshloka
{तस्मै चाभ्युद्यतं श्रीमान्मणिकाञ्चनभूषितम्}
{दिदेश राजा रुचिरं रामाय परमासनम्} %||2-3-18||

\twolineshloka
{तदासनवरं प्राप्य व्यदीपयत राघवः}
{स्वयेव प्रभया मेरुमुदये विमलो रविः} %||2-3-19||

\twolineshloka
{तेन विभ्राजिता तत्र सा सभाभिव्यरोचत}
{विमलग्रहनक्षत्रा शारदी द्यौरिवेन्दुना} %||2-3-20||

\twolineshloka
{तं पश्यमानो नृपतिस्तुतोष प्रियमात्मजम्}
{अलङ्कृतमिवात्मानमादर्शतलसंस्थितम्} %||2-3-21||

\twolineshloka
{स तं सस्मितमाभाष्य पुत्रं पुत्रवतां वरः}
{उवाचेदं वचो राजा देवेन्द्रमिव कश्यपः} %||2-3-22||

\twolineshloka
{ज्येष्ठायामसि मे पत्न्यां सदृश्यां सदृशः सुतः}
{उत्पन्नस्त्वं गुणश्रेष्ठो मम रामात्मजः प्रियः} %||2-3-23||

\twolineshloka
{त्वया यतः प्रजाश्चेमाः स्वगुणैरनुरञ्जिताः}
{तस्मात्त्वं पुष्ययोगेन यौवराज्यमवाप्नुहि} %||2-3-24||

\twolineshloka
{कामतस्त्वं प्रकृत्यैव विनीतो गुणवानसि}
{गुणवत्यपि तु स्नेहात्पुत्र वक्ष्यामि ते हितम्} %||2-3-25||

\twolineshloka
{भूयो विनयमास्थाय भव नित्यं जितेन्द्रियः}
{कामक्रोधसमुत्थानि त्यजेथा व्यसनानि च} %||2-3-26||

\twolineshloka
{परोक्षया वर्तमानो वृत्त्या प्रत्यक्षया तथा}
{अमात्यप्रभृतीः सर्वाः प्रकृतीश्चानुरञ्जय} %||2-3-27||

\threelineshloka
{तुष्टानुरक्तप्रकृतिर्यः पालयति मेदिनीम्}
{तस्य नन्दन्ति मित्राणि लब्ध्वामृतमिवामराः}
{तस्मात्पुत्र त्वमात्मानं नियम्यैव समाचर} %||2-3-28||

\twolineshloka
{तच्छ्रुत्वा सुहृदस्तस्य रामस्य प्रियकारिणः}
{त्वरिताः शीघ्रमभ्येत्य कौसल्यायै न्यवेदयन्} %||2-3-29||

\twolineshloka
{सा हिरण्यं च गाश्चैव रत्नानि विविधानि च}
{व्यादिदेश प्रियाख्येभ्यः कौसल्या प्रमदोत्तमा} %||2-3-30||

\twolineshloka
{अथाभिवाद्य राजानं रथमारुह्य राघवः}
{ययौ स्वं द्युतिमद्वेश्म जनौघैः प्रतिपूजितः} %||2-3-31||

\fourlineindentedshloka
{ते चापि पौरा नृपतेर्वचस्त}
{च्छ्रुत्वा तदा लाभमिवेष्टमाप्य}
{नरेन्द्रमामन्त्य गृहाणि गत्वा}
{ देवान्समानर्चुरतीव हृष्टाः} %||2-4-32||


{॥इत्यार्षे श्रीमद्रामायणे वाल्मीकीये आदिकाव्ये अयोध्याकाण्डे तृतीयः सर्गः॥}

\sect{चतुर्थः सर्गः}

\twolineshloka
{गतेष्वथ नृपो भूयः पौरेषु सह मन्त्रिभिः}
{मन्त्रयित्वा ततश्चक्रे निश्चयज्ञः स निश्चयम्} %||2-4-1||

\twolineshloka
{श्व एव पुष्यो भविता श्वोऽभिषेच्येत मे सुतः}
{रामो राजीवताम्राक्षो यौवराज्य इति प्रभुः} %||2-4-2||

\twolineshloka
{अथान्तर्गृहमाविश्य राजा दशरथस्तदा}
{सूतमाज्ञापयामास रामं पुनरिहानय} %||2-4-3||

\twolineshloka
{प्रतिगृह्य स तद्वाक्यं सूतः पुनरुपाययौ}
{रामस्य भवनं शीघ्रं राममानयितुं पुनः} %||2-4-4||

\twolineshloka
{द्वाःस्थैरावेदितं तस्य रामायागमनं पुनः}
{श्रुत्वैव चापि रामस्तं प्राप्तं शङ्कान्वितोऽभवत्} %||2-4-5||

\twolineshloka
{प्रवेश्य चैनं त्वरितं रामो वचनमब्रवीत्}
{यदागमनकृत्यं ते भूयस्तद्ब्रूह्यशेषतः} %||2-4-6||

\twolineshloka
{तमुवाच ततः सूतो राजा त्वां द्रष्टुमिच्छति}
{श्रुत्वा प्रमाणमत्र त्वं गमनायेतराय वा} %||2-4-7||

\twolineshloka
{इति सूतवचः श्रुत्वा रामोऽथ त्वरयान्वितः}
{प्रययौ राजभवनं पुनर्द्रष्टुं नरेश्वरम्} %||2-4-8||

\twolineshloka
{तं श्रुत्वा समनुप्राप्तं रामं दशरथो नृपः}
{प्रवेशयामास गृहं विविक्षुः प्रियमुत्तमम्} %||2-4-9||

\twolineshloka
{प्रविशन्नेव च श्रीमान्राघवो भवनं पितुः}
{ददर्श पितरं दूरात्प्रणिपत्य कृताञ्जलिः} %||2-4-10||

\twolineshloka
{प्रणमन्तं समुत्थाप्य तं परिष्वज्य भूमिपः}
{प्रदिश्य चास्मै रुचिरमासनं पुनरब्रवीत्} %||2-4-11||

\twolineshloka
{राम वृद्धोऽस्मि दीर्घायुर्भुक्ता भोगा मयेप्सिताः}
{अन्नवद्भिः क्रतुशतैस्तथेष्टं भूरिदक्षिणैः} %||2-4-12||

\twolineshloka
{जातमिष्टमपत्यं मे त्वमद्यानुपमं भुवि}
{दत्तमिष्टमधीतं च मया पुरुषसत्तम} %||2-4-13||

\twolineshloka
{अनुभूतानि चेष्टानि मया वीर सुखानि च}
{देवर्षि पितृविप्राणामनृणोऽस्मि तथात्मनः} %||2-4-14||

\twolineshloka
{न किञ्चिन्मम कर्तव्यं तवान्यत्राभिषेचनात्}
{अतो यत्त्वामहं ब्रूयां तन्मे त्वं कर्तुमर्हसि} %||2-4-15||

\twolineshloka
{अद्य प्रकृतयः सर्वास्त्वामिच्छन्ति नराधिपम्}
{अतस्त्वां युवराजानमभिषेक्ष्यामि पुत्रक} %||2-4-16||

\twolineshloka
{अपि चाद्याशुभान्राम स्वप्नान्पश्यामि दारुणान्}
{सनिर्घाता महोल्काश्च पतन्तीह महास्वनाः} %||2-4-17||

\twolineshloka
{अवष्टब्धं च मे राम नक्षत्रं दारुणैर्ग्रहैः}
{आवेदयन्ति दैवज्ञाः सूर्याङ्गारकराहुभिः} %||2-4-18||

\twolineshloka
{प्रायेण हि निमित्तानामीदृशानां समुद्भवे}
{राजा वा मृत्युमाप्नोति घोरां वापदमृच्छति} %||2-4-19||

\twolineshloka
{तद्यावदेव मे चेतो न विमुह्यति राघव}
{तावदेवाभिषिञ्चस्व चला हि प्राणिनां मतिः} %||2-4-20||

\twolineshloka
{अद्य चन्द्रोऽभ्युपगतः पुष्यात्पूर्वं पुनर्वसुम्}
{श्वः पुष्य योगं नियतं वक्ष्यन्ते दैवचिन्तकाः} %||2-4-21||

\twolineshloka
{तत्र पुष्येऽभिषिञ्चस्व मनस्त्वरयतीव माम्}
{श्वस्त्वाहमभिषेक्ष्यामि यौवराज्ये परन्तप} %||2-4-22||

\twolineshloka
{तस्मात्त्वयाद्य व्रतिना निशेयं नियतात्मना}
{सह वध्वोपवस्तव्या दर्भप्रस्तरशायिना} %||2-4-23||

\twolineshloka
{सुहृदश्चाप्रमत्तास्त्वां रक्षन्त्वद्य समन्ततः}
{भवन्ति बहुविघ्नानि कार्याण्येवंविधानि हि} %||2-4-24||

\twolineshloka
{विप्रोषितश्च भरतो यावदेव पुरादितः}
{तावदेवाभिषेकस्ते प्राप्तकालो मतो मम} %||2-4-25||

\twolineshloka
{कामं खलु सतां वृत्ते भ्राता ते भरतः स्थितः}
{ज्येष्ठानुवर्ती धर्मात्मा सानुक्रोशो जितेन्द्रियः} %||2-4-26||

\twolineshloka
{किं तु चित्तं मनुष्याणामनित्यमिति मे मतिः}
{सतां च धर्मनित्यानां कृतशोभि च राघव} %||2-4-27||

\twolineshloka
{इत्युक्तः सोऽभ्यनुज्ञातः श्वोभाविन्यभिषेचने}
{व्रजेति रामः पितरमभिवाद्याभ्ययाद्गृहम्} %||2-4-28||

\twolineshloka
{प्रविश्य चात्मनो वेश्म राज्ञोद्दिष्टेऽभिषेचने}
{तस्मिन्क्षणे विनिर्गत्य मातुरन्तःपुरं ययौ} %||2-4-29||

\twolineshloka
{तत्र तां प्रवणामेव मातरं क्षौमवासिनीम्}
{वाग्यतां देवतागारे ददर्श याचतीं श्रियम्} %||2-4-30||

\twolineshloka
{प्रागेव चागता तत्र सुमित्रा लक्ष्मणस्तथा}
{सीता चानायिता श्रुत्वा प्रियं रामाभिषेचनम्} %||2-4-31||

\twolineshloka
{तस्मिन्काले हि कौसल्या तस्थावामीलितेक्षणा}
{सुमित्रयान्वास्यमाना सीतया लक्ष्मणेन च} %||2-4-32||

\twolineshloka
{श्रुत्वा पुष्येण पुत्रस्य यौवराज्याभिषेचनम्}
{प्राणायामेन पुरुषं ध्यायमाना जनार्दनम्} %||2-4-33||

\twolineshloka
{तथा सनियमामेव सोऽभिगम्याभिवाद्य च}
{उवाच वचनं रामो हर्षयंस्तामिदं तदा} %||2-4-34||

\twolineshloka
{अम्ब पित्रा नियुक्तोऽस्मि प्रजापालनकर्मणि}
{भविता श्वोऽभिषेको मे यथा मे शासनं पितुः} %||2-4-35||

\twolineshloka
{सीतयाप्युपवस्तव्या रजनीयं मया सह}
{एवमृत्विगुपाध्यायैः सह मामुक्तवान्पिता} %||2-4-36||

\twolineshloka
{यानि यान्यत्र योग्यानि श्वोभाविन्यभिषेचने}
{तानि मे मङ्गलान्यद्य वैदेह्याश्चैव कारय} %||2-4-37||

\twolineshloka
{एतच्छ्रुत्वा तु कौसल्या चिरकालाभिकाङ्क्षितम्}
{हर्षबाष्पकलं वाक्यमिदं राममभाषत} %||2-4-38||

\twolineshloka
{वत्स राम चिरं जीव हतास्ते परिपन्थिनः}
{ज्ञातीन्मे त्वं श्रिया युक्तः सुमित्रायाश्च नन्दय} %||2-4-39||

\twolineshloka
{कल्याणे बत नक्षत्रे मयि जातोऽसि पुत्रक}
{येन त्वया दशरथो गुणैराराधितः पिता} %||2-4-40||

\twolineshloka
{अमोघं बत मे क्षान्तं पुरुषे पुष्करेक्षणे}
{येयमिक्ष्वाकुराज्यश्रीः पुत्र त्वां संश्रयिष्यति} %||2-4-41||

\twolineshloka
{इत्येवमुक्तो मात्रेदं रामो भारतमब्रवीत्}
{प्राञ्जलिं प्रह्वमासीनमभिवीक्ष्य स्मयन्निव} %||2-4-42||

\twolineshloka
{लक्ष्मणेमां मया सार्धं प्रशाधि त्वं वसुन्धराम्}
{द्वितीयं मेऽन्तरात्मानं त्वामियं श्रीरुपस्थिता} %||2-4-43||

\twolineshloka
{सौमित्रे भुङ्क्ष्व भोगांस्त्वमिष्टान्राज्यफलानि च}
{जीवितं च हि राज्यं च त्वदर्थमभिकामये} %||2-4-44||

\twolineshloka
{इत्युक्त्वा लक्ष्मणं रामो मातरावभिवाद्य च}
{अभ्यनुज्ञाप्य सीतां च जगाम स्वं निवेशनम्} %||2-5-45||


{॥इत्यार्षे श्रीमद्रामायणे वाल्मीकीये आदिकाव्ये अयोध्याकाण्डे चतुर्थः सर्गः॥}

\sect{पञ्चमः सर्गः}

\twolineshloka
{सन्दिश्य रामं नृपतिः श्वोभाविन्यभिषेचने}
{पुरोहितं समाहूय वसिष्ठमिदमब्रवीत्} %||2-5-1||

\twolineshloka
{गच्छोपवासं काकुत्स्थं कारयाद्य तपोधन}
{श्रीयशोराज्यलाभाय वध्वा सह यतव्रतम्} %||2-5-2||

\twolineshloka
{तथेति च स राजानमुक्त्वा वेदविदां वरः}
{स्वयं वसिष्ठो भगवान्ययौ रामनिवेशनम्} %||2-5-3||

\twolineshloka
{स रामभवनं प्राप्य पाण्डुराभ्रघनप्रभम्}
{तिस्रः कक्ष्या रथेनैव विवेश मुनिसत्तमः} %||2-5-4||

\twolineshloka
{तमागतमृषिं रामस्त्वरन्निव ससम्भ्रमः}
{मानयिष्यन्स मानार्हं निश्चक्राम निवेशनात्} %||2-5-5||

\twolineshloka
{अभ्येत्य त्वरमाणश्च रथाभ्याशं मनीषिणः}
{ततोऽवतारयामास परिगृह्य रथात्स्वयम्} %||2-5-6||

\twolineshloka
{स चैनं प्रश्रितं दृष्ट्वा सम्भाष्याभिप्रसाद्य च}
{प्रियार्हं हर्षयन्राममित्युवाच पुरोहितः} %||2-5-7||

\twolineshloka
{प्रसन्नस्ते पिता राम यौवराज्यमवाप्स्यसि}
{उपवासं भवानद्य करोतु सह सीतया} %||2-5-8||

\twolineshloka
{प्रातस्त्वामभिषेक्ता हि यौवराज्ये नराधिपः}
{पिता दशरथः प्रीत्या ययातिं नहुषो यथा} %||2-5-9||

\twolineshloka
{इत्युक्त्वा स तदा राममुपवासं यतव्रतम्}
{मन्त्रवत्कारयामास वैदेह्या सहितं मुनिः} %||2-5-10||

\twolineshloka
{ततो यथावद्रामेण स राज्ञो गुरुरर्चितः}
{अभ्यनुज्ञाप्य काकुत्स्थं ययौ रामनिवेशनात्} %||2-5-11||

\twolineshloka
{सुहृद्भिस्तत्र रामोऽपि ताननुज्ञाप्य सर्वशः}
{सभाजितो विवेशाथ ताननुज्ञाप्य सर्वशः} %||2-5-12||

\twolineshloka
{हृष्टनारी नरयुतं रामवेश्म तदा बभौ}
{यथा मत्तद्विजगणं प्रफुल्लनलिनं सरः} %||2-5-13||

\twolineshloka
{स राजभवनप्रख्यात्तस्माद्रामनिवेशनात्}
{निर्गत्य ददृशे मार्गं वसिष्ठो जनसंवृतम्} %||2-5-14||

\twolineshloka
{वृन्दवृन्दैरयोध्यायां राजमार्गाः समन्ततः}
{बभूवुरभिसम्बाधाः कुतूहलजनैर्वृताः} %||2-5-15||

\twolineshloka
{जनवृन्दोर्मिसङ्घर्षहर्षस्वनवतस्तदा}
{बभूव राजमार्गस्य सागरस्येव निस्वनः} %||2-5-16||

\twolineshloka
{सिक्तसम्मृष्टरथ्या हि तदहर्वनमालिनी}
{आसीदयोध्या नगरी समुच्छ्रितगृहध्वजा} %||2-5-17||

\twolineshloka
{तदा ह्ययोध्या निलयः सस्त्रीबालाबलो जनः}
{रामाभिषेकमाकाङ्क्षन्नाकाङ्क्षन्नुदयं रवेः} %||2-5-18||

\twolineshloka
{प्रजालङ्कारभूतं च जनस्यानन्दवर्धनम्}
{उत्सुकोऽभूज्जनो द्रष्टुं तमयोध्या महोत्सवम्} %||2-5-19||

\twolineshloka
{एवं तं जनसम्बाधं राजमार्गं पुरोहितः}
{व्यूहन्निव जनौघं तं शनै राज कुलं ययौ} %||2-5-20||

\twolineshloka
{सिताभ्रशिखरप्रख्यं प्रासदमधिरुह्य सः}
{समियाय नरेन्द्रेण शक्रेणेव बृहस्पतिः} %||2-5-21||

\twolineshloka
{तमागतमभिप्रेक्ष्य हित्वा राजासनं नृपः}
{पप्रच्छ स च तस्मै तत्कृतमित्यभ्यवेदयत्} %||2-5-22||

\twolineshloka
{गुरुणा त्वभ्यनुज्ञातो मनुजौघं विसृज्य तम्}
{विवेशान्तःपुरं राजा सिंहो गिरिगुहामिव} %||2-5-23||

\fourlineindentedshloka
{तदग्र्यवेषप्रमदाजनाकुलम्}
{ महेन्द्रवेश्मप्रतिमं निवेशनम्}
{व्यदीपयंश्चारु विवेश पार्थिवः}
{ शशीव तारागणसङ्कुलं नभः} %||2-6-24||


{॥इत्यार्षे श्रीमद्रामायणे वाल्मीकीये आदिकाव्ये अयोध्याकाण्डे पञ्चमः सर्गः॥}

\sect{षष्ठः सर्गः}

\twolineshloka
{गते पुरोहिते रामः स्नातो नियतमानसः}
{सह पत्न्या विशालाक्ष्या नारायणमुपागमत्} %||2-6-1||

\twolineshloka
{प्रगृह्य शिरसा पात्रीं हविषो विधिवत्तदा}
{महते दैवतायाज्यं जुहाव ज्वलितेऽनले} %||2-6-2||

\twolineshloka
{शेषं च हविषस्तस्य प्राश्याशास्यात्मनः प्रियम्}
{ध्यायन्नारायणं देवं स्वास्तीर्णे कुशसंस्तरे} %||2-6-3||

\twolineshloka
{वाग्यतः सह वैदेह्या भूत्वा नियतमानसः}
{श्रीमत्यायतने विष्णोः शिश्ये नरवरात्मजः} %||2-6-4||

\twolineshloka
{एकयामावशिष्टायां रात्र्यां प्रतिविबुध्य सः}
{अलङ्कारविधिं कृत्स्नं कारयामास वेश्मनः} %||2-6-5||

\twolineshloka
{तत्र शृण्वन्सुखा वाचः सूतमागधबन्दिनाम्}
{पूर्वां सन्ध्यामुपासीनो जजाप यतमानसः} %||2-6-6||

\twolineshloka
{तुष्टाव प्रणतश्चैव शिरसा मधुसूदनम्}
{विमलक्षौमसंवीतो वाचयामास च द्विजान्} %||2-6-7||

\twolineshloka
{तेषां पुण्याहघोषोऽथ गम्भीरमधुरस्तदा}
{अयोध्यां पूरयामास तूर्यघोषानुनादितः} %||2-6-8||

\twolineshloka
{कृतोपवासं तु तदा वैदेह्या सह राघवम्}
{अयोध्या निलयः श्रुत्वा सर्वः प्रमुदितो जनः} %||2-6-9||

\twolineshloka
{ततः पौरजनः सर्वः श्रुत्वा रामाभिषेचनम्}
{प्रभातां रजनीं दृष्ट्वा चक्रे शोभां परां पुनः} %||2-6-10||

\twolineshloka
{सिताभ्रशिखराभेषु देवतायतनेषु च}
{चतुष्पथेषु रथ्यासु चैत्येष्वट्टालकेषु च} %||2-6-11||

\twolineshloka
{नानापण्यसमृद्धेषु वणिजामापणेषु च}
{कुटुम्बिनां समृद्धेषु श्रीमत्सु भवनेषु च} %||2-6-12||

\twolineshloka
{सभासु चैव सर्वासु वृक्षेष्वालक्षितेषु च}
{ध्वजाः समुच्छ्रिताश्चित्राः पताकाश्चाभवंस्तदा} %||2-6-13||

\twolineshloka
{नटनर्तकसङ्घानां गायकानां च गायताम्}
{मनःकर्णसुखा वाचः शुश्रुवुश्च ततस्ततः} %||2-6-14||

\twolineshloka
{रामाभिषेकयुक्ताश्च कथाश्चक्रुर्मिथो जनाः}
{रामाभिषेके सम्प्राप्ते चत्वरेषु गृहेषु च} %||2-6-15||

\twolineshloka
{बाला अपि क्रीडमाना गृहद्वारेषु सङ्घशः}
{रामाभिषेकसंयुक्ताश्चक्रुरेव मिथः कथाः} %||2-6-16||

\twolineshloka
{कृतपुष्पोपहारश्च धूपगन्धाधिवासितः}
{राजमार्गः कृतः श्रीमान्पौरै रामाभिषेचने} %||2-6-17||

\twolineshloka
{प्रकाशीकरणार्थं च निशागमनशङ्कया}
{दीपवृक्षांस्तथा चक्रुरनु रथ्यासु सर्वशः} %||2-6-18||

\twolineshloka
{अलङ्कारं पुरस्यैवं कृत्वा तत्पुरवासिनः}
{आकाङ्क्षमाणा रामस्य यौवराज्याभिषेचनम्} %||2-6-19||

\twolineshloka
{समेत्य सङ्घशः सर्वे चत्वरेषु सभासु च}
{कथयन्तो मिथस्तत्र प्रशशंसुर्जनाधिपम्} %||2-6-20||

\twolineshloka
{अहो महात्मा राजायमिक्ष्वाकुकुलनन्दनः}
{ज्ञात्वा यो वृद्धमात्मानं रामं राज्येऽह्बिषेक्ष्यति} %||2-6-21||

\twolineshloka
{सर्वे ह्यनुगृहीताः स्म यन्नो रामो महीपतिः}
{चिराय भविता गोप्ता दृष्टलोकपरावरः} %||2-6-22||

\twolineshloka
{अनुद्धतमना विद्वान्धर्मात्मा भ्रातृवत्सलः}
{यथा च भ्रातृषु स्निग्धस्तथास्मास्वपि राघवः} %||2-6-23||

\twolineshloka
{चिरं जीवतु धर्मात्मा राजा दशरथोऽनघः}
{यत्प्रसादेनाभिषिक्तं रामं द्रक्ष्यामहे वयम्} %||2-6-24||

\twolineshloka
{एवंविधं कथयतां पौराणां शुश्रुवुस्तदा}
{दिग्भ्योऽपि श्रुतवृत्तान्ताः प्राप्ता जानपदा जनाः} %||2-6-25||

\twolineshloka
{ते तु दिग्भ्यः पुरीं प्राप्ता द्रष्टुं रामाभिषेचनम्}
{रामस्य पूरयामासुः पुरीं जानपदा जनाः} %||2-6-26||

\twolineshloka
{जनौघैस्तैर्विसर्पद्भिः शुश्रुवे तत्र निस्वनः}
{पर्वसूदीर्णवेगस्य सागरस्येव निस्वनः} %||2-6-27||

\fourlineindentedshloka
{ततस्तदिन्द्रक्षयसंनिभं पुरम्}
{ दिदृक्षुभिर्जानपदैरुपागतैः}
{समन्ततः सस्वनमाकुलं बभौ}
{ समुद्रयादोभिरिवार्णवोदकम्} %||2-7-28||


{॥इत्यार्षे श्रीमद्रामायणे वाल्मीकीये आदिकाव्ये अयोध्याकाण्डे षष्ठः सर्गः॥}

\sect{सप्तमः सर्गः}

\twolineshloka
{ज्ञातिदासी यतो जाता कैकेय्यास्तु सहोषिता}
{प्रासादं चन्द्रसङ्काशमारुरोह यदृच्छया} %||2-7-1||

\twolineshloka
{सिक्तराजपथां कृत्स्नां प्रकीर्णकमलोत्पलाम्}
{अयोध्यां मन्थरा तस्मात्प्रासादादन्ववैक्षत} %||2-7-2||

\twolineshloka
{पताकाभिर्वरार्हाभिर्ध्वजैश्च समलङ्कृताम्}
{सिक्तां चन्दनतोयैश्च शिरःस्नातजनैर्वृताम्} %||2-7-3||

\twolineshloka
{अविदूरे स्थितां दृष्ट्वा धात्रीं पप्रच्छ मन्थरा}
{उत्तमेनाभिसंयुक्ता हर्षेणार्थपरा सती} %||2-7-4||

\threelineshloka
{राममाता धनं किं नु जनेभ्यः सम्प्रयच्छति}
{अतिमात्रं प्रहर्षोऽयं किं जनस्य च शंस मे}
{कारयिष्यति किं वापि सम्प्रहृष्टो महीपतिः} %||2-7-5||

\twolineshloka
{विदीर्यमाणा हर्षेण धात्री परमया मुदा}
{आचचक्षेऽथ कुब्जायै भूयसीं राघवे श्रियम्} %||2-7-6||

\twolineshloka
{श्वः पुष्येण जितक्रोधं यौवराज्येन राघवम्}
{राजा दशरथो राममभिषेचयितानघम्} %||2-7-7||

\twolineshloka
{धात्र्यास्तु वचनं श्रुत्वा कुब्जा क्षिप्रममर्षिता}
{कैलास शिखराकारात्प्रासादादवरोहत} %||2-7-8||

\twolineshloka
{सा दह्यमाना कोपेन मन्थरा पापदर्शिनी}
{शयानामेत्य कैकेयीमिदं वचनमब्रवीत्} %||2-7-9||

\twolineshloka
{उत्तिष्ठ मूढे किं शेषे भयं त्वामभिवर्तते}
{उपप्लुतमहौघेन किमात्मानं न बुध्यसे} %||2-7-10||

\twolineshloka
{अनिष्टे सुभगाकारे सौभाग्येन विकत्थसे}
{चलं हि तव सौभाग्यं नद्यः स्रोत इवोष्णगे} %||2-7-11||

\twolineshloka
{एवमुक्ता तु कैकेयी रुष्टया परुषं वचः}
{कुब्जया पापदर्शिन्या विषादमगमत्परम्} %||2-7-12||

\twolineshloka
{कैकेयी त्वब्रवीत्कुब्जां कच्चित्क्षेमं न मन्थरे}
{विषण्णवदनां हि त्वां लक्षये भृशदुःखिताम्} %||2-7-13||

\twolineshloka
{मन्थरा तु वचः श्रुत्वा कैकेय्या मधुराक्षरम्}
{उवाच क्रोधसंयुक्ता वाक्यं वाक्यविशारदा} %||2-7-14||

\twolineshloka
{सा विषण्णतरा भूत्वा कुब्जा तस्या हितैषिणी}
{विषादयन्ती प्रोवाच भेदयन्ती च राघवम्} %||2-7-15||

\twolineshloka
{अक्षेमं सुमहद्देवि प्रवृत्तं त्वद्विनाशनम्}
{रामं दशरथो राजा यौवराज्येऽभिषेक्ष्यति} %||2-7-16||

\twolineshloka
{सास्म्यगाधे भये मग्ना दुःखशोकसमन्विता}
{दह्यमानानलेनेव त्वद्धितार्थमिहागता} %||2-7-17||

\twolineshloka
{तव दुःखेन कैकेयि मम दुःखं महद्भवेत्}
{त्वद्वृद्धौ मम वृद्धिश्च भवेदत्र न संशयः} %||2-7-18||

\twolineshloka
{नराधिपकुले जाता महिषी त्वं महीपतेः}
{उग्रत्वं राजधर्माणां कथं देवि न बुध्यसे} %||2-7-19||

\twolineshloka
{धर्मवादी शठो भर्ता श्लक्ष्णवादी च दारुणः}
{शुद्धभावे न जानीषे तेनैवमतिसन्धिता} %||2-7-20||

\twolineshloka
{उपस्थितं पयुञ्जानस्त्वयि सान्त्वमनर्थकम्}
{अर्थेनैवाद्य ते भर्ता कौसल्यां योजयिष्यति} %||2-7-21||

\twolineshloka
{अपवाह्य स दुष्टात्मा भरतं तव बन्धुषु}
{काल्यं स्थापयिता रामं राज्ये निहतकण्टके} %||2-7-22||

\twolineshloka
{शत्रुः पतिप्रवादेन मात्रेव हितकाम्यया}
{आशीविष इवाङ्केन बाले परिधृतस्त्वया} %||2-7-23||

\twolineshloka
{यथा हि कुर्यात्सर्पो वा शत्रुर्वा प्रत्युपेक्षितः}
{राज्ञा दशरथेनाद्य सपुत्रा त्वं तथा कृता} %||2-7-24||

\twolineshloka
{पापेनानृतसन्त्वेन बाले नित्यं सुखोचिते}
{रामं स्थापयता राज्ये सानुबन्धा हता ह्यसि} %||2-7-25||

\twolineshloka
{सा प्राप्तकालं कैकेयि क्षिप्रं कुरु हितं तव}
{त्रायस्व पुत्रमात्मानं मां च विस्मयदर्शने} %||2-7-26||

\twolineshloka
{मन्थराया वचः श्रुत्वा शयनात्स शुभानना}
{एवमाभरणं तस्यै कुब्जायै प्रददौ शुभम्} %||2-7-27||

\twolineshloka
{दत्त्वा त्वाभरणं तस्यै कुब्जायै प्रमदोत्तमा}
{कैकेयी मन्थरां हृष्टा पुनरेवाब्रवीदिदम्} %||2-7-28||

\twolineshloka
{इदं तु मन्थरे मह्यमाख्यासि परमं प्रियम्}
{एतन्मे प्रियमाख्यातुः किं वा भूयः करोमि ते} %||2-7-29||

\twolineshloka
{रामे वा भरते वाहं विशेषं नोपलक्षये}
{तस्मात्तुष्टास्मि यद्राजा रामं राज्येऽभिषेक्ष्यति} %||2-7-30||

\fourlineindentedshloka
{न मे परं किञ्चिदितस्त्वया पुनः}
{ प्रियं प्रियार्हे सुवचं वचो वरम्}
{तथा ह्यवोचस्त्वमतः प्रियोत्तरम्}
{ वरं परं ते प्रददामि तं वृणु} %||2-8-31||


{॥इत्यार्षे श्रीमद्रामायणे वाल्मीकीये आदिकाव्ये अयोध्याकाण्डे सप्तमः सर्गः॥}

\sect{अष्टमः सर्गः}

\twolineshloka
{मन्थरा त्वभ्यसूय्यैनामुत्सृज्याभरणं च तत्}
{उवाचेदं ततो वाक्यं कोपदुःखसमन्विता} %||2-8-1||

\twolineshloka
{हर्षं किमिदमस्थाने कृतवत्यसि बालिशे}
{शोकसागरमध्यस्थमात्मानं नावबुध्यसे} %||2-8-2||

\twolineshloka
{सुभगा खलु कौसल्या यस्याः पुत्रोऽभिषेक्ष्यते}
{यौवराज्येन महता श्वः पुष्येण द्विजोत्तमैः} %||2-8-3||

\twolineshloka
{प्राप्तां सुमहतीं प्रीतिं प्रतीतां तां हतद्विषम्}
{उपस्थास्यसि कौसल्यां दासीव त्वं कृताञ्जलिः} %||2-8-4||

\twolineshloka
{हृष्टाः खलु भविष्यन्ति रामस्य परमाः स्त्रियः}
{अप्रहृष्टा भविष्यन्ति स्नुषास्ते भरतक्षये} %||2-8-5||

\twolineshloka
{तां दृष्ट्वा परमप्रीतां ब्रुवन्तीं मन्थरां ततः}
{रामस्यैव गुणान्देवी कैकेयी प्रशशंस ह} %||2-8-6||

\twolineshloka
{धर्मज्ञो गुरुभिर्दान्तः कृतज्ञः सत्यवाक्शुचिः}
{रामो राज्ञः सुतो ज्येष्ठो यौवराज्यमतोऽर्हति} %||2-8-7||

\twolineshloka
{भ्रातॄन्भृत्यांश्च दीर्घायुः पितृवत्पालयिष्यति}
{सन्तप्यसे कथं कुब्जे श्रुत्वा रामाभिषेचनम्} %||2-8-8||

\twolineshloka
{भरतश्चापि रामस्य ध्रुवं वर्षशतात्परम्}
{पितृपैतामहं राज्यमवाप्स्यति नरर्षभः} %||2-8-9||

\threelineshloka
{सा त्वमभ्युदये प्राप्ते वर्तमाने च मन्थरे}
{भविष्यति च कल्याणे किमर्थं परितप्यसे}
{कौसल्यातोऽतिरिक्तं च स तु शुश्रूषते हि माम्} %||2-8-10||

\twolineshloka
{कैकेय्या वचनं श्रुत्वा मन्थरा भृशदुःखिता}
{दीर्घमुष्णं विनिःश्वस्य कैकेयीमिदमब्रवीत्} %||2-8-11||

\twolineshloka
{अनर्थदर्शिनी मौर्ख्यान्नात्मानमवबुध्यसे}
{शोकव्यसनविस्तीर्णे मज्जन्ती दुःखसागरे} %||2-8-12||

\twolineshloka
{भविता राघवो राजा राघवस्य च यः सुतः}
{राजवंशात्तु भरतः कैकेयि परिहास्यते} %||2-8-13||

\twolineshloka
{न हि राज्ञः सुताः सर्वे राज्ये तिष्ठन्ति भामिनि}
{स्थाप्यमानेषु सर्वेषु सुमहाननयो भवेत्} %||2-8-14||

\twolineshloka
{तस्माज्ज्येष्ठे हि कैकेयि राज्यतन्त्राणि पार्थिवाः}
{स्थापयन्त्यनवद्याङ्गि गुणवत्स्वितरेष्वपि} %||2-8-15||

\twolineshloka
{असावत्यन्तनिर्भग्नस्तव पुत्रो भविष्यति}
{अनाथवत्सुखेभ्यश्च राजवंशाच्च वत्सले} %||2-8-16||

\twolineshloka
{साहं त्वदर्थे सम्प्राप्ता त्वं तु मां नावबुध्यसे}
{सपत्निवृद्धौ या मे त्वं प्रदेयं दातुमिच्छसि} %||2-8-17||

\twolineshloka
{ध्रुवं तु भरतं रामः प्राप्य राज्यमकण्टकम्}
{देशान्तरं नाययित्वा लोकान्तरमथापि वा} %||2-8-18||

\twolineshloka
{बाल एव हि मातुल्यं भरतो नायितस्त्वया}
{संनिकर्षाच्च सौहार्दं जायते स्थावरेष्वपि} %||2-8-19||

\twolineshloka
{गोप्ता हि रामं सौमित्रिर्लक्ष्मणं चापि राघवः}
{अश्विनोरिव सौभ्रात्रं तयोर्लोकेषु विश्रुतम्} %||2-8-20||

\twolineshloka
{तस्मान्न लक्ष्मणे रामः पापं किञ्चित्करिष्यति}
{रामस्तु भरते पापं कुर्यादिति न संशयः} %||2-8-21||

\twolineshloka
{तस्माद्राजगृहादेव वनं गच्छतु ते सुतः}
{एतद्धि रोचते मह्यं भृशं चापि हितं तव} %||2-8-22||

\twolineshloka
{एवं ते ज्ञातिपक्षस्य श्रेयश्चैव भविष्यति}
{यदि चेद्भरतो धर्मात्पित्र्यं राज्यमवाप्स्यति} %||2-8-23||

\twolineshloka
{स ते सुखोचितो बालो रामस्य सहजो रिपुः}
{समृधार्थस्य नष्टार्थो जीविष्यति कथं वशे} %||2-8-24||

\twolineshloka
{अभिद्रुतमिवारण्ये सिंहेन गजयूथपम्}
{प्रच्छाद्यमानं रामेण भरतं त्रातुमर्हसि} %||2-8-25||

\twolineshloka
{दर्पान्निराकृता पूर्वं त्वया सौभाग्यवत्तया}
{राममाता सपत्नी ते कथं वैरं न यातयेत्} %||2-8-26||

\fourlineindentedshloka
{यदा हि रामः पृथिवीमवाप्स्यति}
{ ध्रुवं प्रनष्टो भरतो भविष्यति}
{अतो हि सञ्चिन्तय राज्यमात्मजे}
{ परस्य चाद्यैव विवास कारणम्} %||2-9-27||


{॥इत्यार्षे श्रीमद्रामायणे वाल्मीकीये आदिकाव्ये अयोध्याकाण्डे अष्टमः सर्गः॥}

\sect{नवमः सर्गः}

\twolineshloka
{एवमुक्ता तु कैकेयी क्रोधेन ज्वलितानना}
{दीर्घमुष्णं विनिःश्वस्य मन्थरामिदमब्रवीत्} %||2-9-1||

\twolineshloka
{अद्य राममितः क्षिप्रं वनं प्रस्थापयाम्यहम्}
{यौवराज्येन भरतं क्षिप्रमेवाभिषेचये} %||2-9-2||

\twolineshloka
{इदं त्विदानीं सम्पश्य केनोपायेन मन्थरे}
{भरतः प्राप्नुयाद्राज्यं न तु रामः कथञ्चन} %||2-9-3||

\twolineshloka
{एवमुक्ता तया देव्या मन्थरा पापदर्शिनी}
{रामार्थमुपहिंसन्ती कैकेयीमिदमब्रवीत्} %||2-9-4||

\twolineshloka
{हन्तेदानीं प्रवक्ष्यामि कैकेयि श्रूयतां च मे}
{यथा ते भरतो राज्यं पुत्रः प्राप्स्यति केवलम्} %||2-9-5||

\twolineshloka
{श्रुत्वैवं वचनं तस्या मन्थरायास्तु कैकयी}
{किञ्चिदुत्थाय शयनात्स्वास्तीर्णादिदमब्रवीत्} %||2-9-6||

\twolineshloka
{कथय त्वं ममोपायं केनोपायेन मन्थरे}
{भरतः प्राप्नुयाद्राज्यं न तु रामः कथञ्चन} %||2-9-7||

\twolineshloka
{एवमुक्ता तया देव्या मन्थरा पापदर्शिनी}
{रामार्थमुपहिंसन्ती कुब्जा वचनमब्रवीत्} %||2-9-8||

\twolineshloka
{तव देवासुरे युद्धे सह राजर्षिभिः पतिः}
{अगच्छत्त्वामुपादाय देवराजस्य साह्यकृत्} %||2-9-9||

\twolineshloka
{दिशमास्थाय कैकेयि दक्षिणां दण्डकान्प्रति}
{वैजयन्तमिति ख्यातं पुरं यत्र तिमिध्वजः} %||2-9-10||

\twolineshloka
{स शम्बर इति ख्यातः शतमायो महासुरः}
{ददौ शक्रस्य सङ्ग्रामं देवसङ्घैरनिर्जितः} %||2-9-11||

\twolineshloka
{तस्मिन्महति सङ्ग्रामे राजा दशरथस्तदा}
{अपवाह्य त्वया देवि सङ्ग्रामान्नष्टचेतनः} %||2-9-12||

\twolineshloka
{तत्रापि विक्षतः शस्त्रैः पतिस्ते रक्षितस्त्वया}
{तुष्टेन तेन दत्तौ ते द्वौ वरौ शुभदर्शने} %||2-9-13||

\threelineshloka
{स त्वयोक्तः पतिर्देवि यदेच्छेयं तदा वरौ}
{गृह्णीयामिति तत्तेन तथेत्युक्तं महात्मना}
{अनभिज्ञा ह्यहं देवि त्वयैव कथितं पुरा} %||2-9-14||

\twolineshloka
{तौ वरौ याच भर्तारं भरतस्याभिषेचनम्}
{प्रव्राजनं च रामस्य त्वं वर्षाणि चतुर्दश} %||2-9-15||

\threelineshloka
{क्रोधागारं प्रविश्याद्य क्रुद्धेवाश्वपतेः सुते}
{शेष्वानन्तर्हितायां त्वं भूमौ मलिनवासिनी}
{मा स्मैनं प्रत्युदीक्षेथा मा चैनमभिभाषथाः} %||2-9-16||

\twolineshloka
{दयिता त्वं सदा भर्तुरत्र मे नास्ति संशयः}
{त्वत्कृते च महाराजो विशेदपि हुताशनम्} %||2-9-17||

\twolineshloka
{न त्वां क्रोधयितुं शक्तो न क्रुद्धां प्रत्युदीक्षितुम्}
{तव प्रियार्थं राजा हि प्राणानपि परित्यजेत्} %||2-9-18||

\twolineshloka
{न ह्यतिक्रमितुं शक्तस्तव वाक्यं महीपतिः}
{मन्दस्वभावे बुध्यस्व सौभाग्यबलमात्मनः} %||2-9-19||

\twolineshloka
{मणिमुक्तासुवर्णानि रत्नानि विविधानि च}
{दद्याद्दशरथो राजा मा स्म तेषु मनः कृथाः} %||2-9-20||

\twolineshloka
{यौ तौ देवासुरे युद्धे वरौ दशरथोऽददात्}
{तौ स्मारय महाभागे सोऽर्थो मा त्वामतिक्रमेत्} %||2-9-21||

\twolineshloka
{यदा तु ते वरं दद्यात्स्वयमुत्थाप्य राघवः}
{व्यवस्थाप्य महाराजं त्वमिमं वृणुया वरम्} %||2-9-22||

\twolineshloka
{रामं प्रव्राजयारण्ये नव वर्षाणि पञ्च च}
{भरतः क्रियतां राजा पृथिव्यां पार्थिवर्षभः} %||2-9-23||

\twolineshloka
{एवं प्रव्राजितश्चैव रामोऽरामो भविष्यति}
{भरतश्च हतामित्रस्तव राजा भविष्यति} %||2-9-24||

\threelineshloka
{येन कालेन रामश्च वनात्प्रत्यागमिष्यति}
{तेन कालेन पुत्रस्ते कृतमूलो भविष्यति}
{सङ्गृहीतमनुष्यश्च सुहृद्भिः सार्धमात्मवान्} %||2-9-25||

\twolineshloka
{प्राप्तकालं तु ते मन्ये राजानं वीतसाध्वसा}
{रामाभिषेकसङ्कल्पान्निगृह्य विनिवर्तय} %||2-9-26||

\twolineshloka
{अनर्थमर्थरूपेण ग्राहिता सा ततस्तया}
{हृष्टा प्रतीता कैकेयी मन्थरामिदमब्रवीत्} %||2-9-27||

\twolineshloka
{कुब्जे त्वां नाभिजानामि श्रेष्ठां श्रेष्ठाभिधायिनीम्}
{पृथिव्यामसि कुब्जानामुत्तमा बुद्धिनिश्चये} %||2-9-28||

\twolineshloka
{त्वमेव तु ममार्थेषु नित्ययुक्ता हितैषिणी}
{नाहं समवबुध्येयं कुब्जे राज्ञश्चिकीर्षितम्} %||2-9-29||

\twolineshloka
{सन्ति दुःसंस्थिताः कुब्जा वक्राः परमपापिकाः}
{त्वं पद्ममिव वातेन संनता प्रियदर्शना} %||2-9-30||

\twolineshloka
{उरस्तेऽभिनिविष्टं वै यावत्स्कन्धात्समुन्नतम्}
{अधस्ताच्चोदरं शान्तं सुनाभमिव लज्जितम्} %||2-9-31||

\twolineshloka
{जघनं तव निर्घुष्टं रशनादामशोभितम्}
{जङ्घे भृशमुपन्यस्ते पादौ चाप्यायतावुभौ} %||2-9-32||

\twolineshloka
{त्वमायताभ्यां सक्थिभ्यां मन्थरे क्षौमवासिनि}
{अग्रतो मम गच्छन्ती राजहंसीव राजसे} %||2-9-33||

\twolineshloka
{तवेदं स्थगु यद्दीर्घं रथघोणमिवायतम्}
{मतयः क्षत्रविद्याश्च मायाश्चात्र वसन्ति ते} %||2-9-34||

\twolineshloka
{अत्र ते प्रतिमोक्ष्यामि मालां कुब्जे हिरण्मयीम्}
{अभिषिक्ते च भरते राघवे च वनं गते} %||2-9-35||

\twolineshloka
{जात्येन च सुवर्णेन सुनिष्टप्तेन सुन्दरि}
{लब्धार्था च प्रतीता च लेपयिष्यामि ते स्थगु} %||2-9-36||

\twolineshloka
{मुखे च तिलकं चित्रं जातरूपमयं शुभम्}
{कारयिष्यामि ते कुब्जे शुभान्याभरणानि च} %||2-9-37||

\threelineshloka
{परिधाय शुभे वस्त्रे देवदेव चरिष्यसि}
{चन्द्रमाह्वयमानेन मुखेनाप्रतिमानना}
{गमिष्यसि गतिं मुख्यां गर्वयन्ती द्विषज्जनम्} %||2-9-38||

\twolineshloka
{तवापि कुब्जाः कुब्जायाः सर्वाभरणभूषिताः}
{पादौ परिचरिष्यन्ति यथैव त्वं सदा मम} %||2-9-39||

\twolineshloka
{इति प्रशस्यमाना सा कैकेयीमिदमब्रवीत्}
{शयानां शयने शुभ्रे वेद्यामग्निशिखामिव} %||2-9-40||

\twolineshloka
{गतोदके सेतुबन्धो न कल्याणि विधीयते}
{उत्तिष्ठ कुरु कल्याणं राजानमनुदर्शय} %||2-9-41||

\twolineshloka
{तथा प्रोत्साहिता देवी गत्वा मन्थरया सह}
{क्रोधागारं विशालाक्षी सौभाग्यमदगर्विता} %||2-9-42||

\twolineshloka
{अनेकशतसाहस्रं मुक्ताहारं वराङ्गना}
{अवमुच्य वरार्हाणि शुभान्याभरणानि च} %||2-9-43||

\twolineshloka
{ततो हेमोपमा तत्र कुब्जा वाक्यं वशं गता}
{संविश्य भूमौ कैकेयी मन्थरामिदमब्रवीत्} %||2-9-44||

\twolineshloka
{इह वा मां मृतां कुब्जे नृपायावेदयिष्यसि}
{वनं तु राघवे प्राप्ते भरतः प्राप्स्यति क्षितिम्} %||2-9-45||

\fourlineindentedshloka
{अथैतदुक्त्वा वचनं सुदारुणम्}
{ निधाय सर्वाभरणानि भामिनी}
{असंवृतामास्तरणेन मेदिनीम्}
{ तदाधिशिश्ये पतितेव किन्नरी} %||2-9-46||

\fourlineindentedshloka
{उदीर्णसंरम्भतमोवृतानना}
{ तथावमुक्तोत्तममाल्यभूषणा}
{नरेन्द्रपत्नी विमना बभूव सा}
{ तमोवृता द्यौरिव मग्नतारका} %||2-10-47||


{॥इत्यार्षे श्रीमद्रामायणे वाल्मीकीये आदिकाव्ये अयोध्याकाण्डे नवमः सर्गः॥}

\sect{दशमः सर्गः}

\twolineshloka
{आज्ञाप्य तु महाराजो राघवस्याभिषेचनम्}
{प्रियार्हां प्रियमाख्यातुं विवेशान्तःपुरं वशी} %||2-10-1||

\twolineshloka
{तां तत्र पतितां भूमौ शयानामतथोचिताम्}
{प्रतप्त इव दुःखेन सोऽपश्यज्जगतीपतिः} %||2-10-2||

\twolineshloka
{स वृद्धस्तरुणीं भार्यां प्राणेभ्योऽपि गरीयसीम्}
{अपापः पापसङ्कल्पां ददर्श धरणीतले} %||2-10-3||

\twolineshloka
{करेणुमिव दिग्धेन विद्धां मृगयुणा वने}
{महागज इवारण्ये स्नेहात्परिममर्श ताम्} %||2-10-4||

\twolineshloka
{परिमृश्य च पाणिभ्यामभिसन्त्रस्तचेतनः}
{कामी कमलपत्राक्षीमुवाच वनितामिदम्} %||2-10-5||

\twolineshloka
{न तेऽहमभिजानामि क्रोधमात्मनि संश्रितम्}
{देवि केनाभियुक्तासि केन वासि विमानिता} %||2-10-6||

\threelineshloka
{यदिदं मम दुःखाय शेषे कल्याणि पांसुषु}
{भूमौ शेषे किमर्थं त्वं मयि कल्याण चेतसि}
{भूतोपहतचित्तेव मम चित्तप्रमाथिनी} %||2-10-7||

\twolineshloka
{सन्ति मे कुशला वैद्या अभितुष्टाश्च सर्वशः}
{सुखितां त्वां करिष्यन्ति व्याधिमाचक्ष्व भामिनि} %||2-10-8||

\twolineshloka
{कस्य वा ते प्रियं कार्यं केन वा विप्रियं कृतम्}
{कः प्रियं लभतामद्य को वा सुमहदप्रियम्} %||2-10-9||

\twolineshloka
{अवध्यो वध्यतां को वा वध्यः को वा विमुच्यताम्}
{दरिद्रः को भवत्वाढ्यो द्रव्यवान्वाप्यकिञ्चनः} %||2-10-10||

\twolineshloka
{अहं चैव मदीयाश्च सर्वे तव वशानुगाः}
{न ते कञ्चिदभिप्रायं व्याहन्तुमहमुत्सहे} %||2-10-11||

\twolineshloka
{आत्मनो जीवितेनापि ब्रूहि यन्मनसेच्छसि}
{यावदावर्तते चक्रं तावती मे वसुन्धरा} %||2-10-12||

\twolineshloka
{तथोक्ता सा समाश्वस्ता वक्तुकामा तदप्रियम्}
{परिपीडयितुं भूयो भर्तारमुपचक्रमे} %||2-10-13||

\twolineshloka
{नास्मि विप्रकृता देव केनचिन्न विमानिता}
{अभिप्रायस्तु मे कश्चित्तमिच्छामि त्वया कृतम्} %||2-10-14||

\twolineshloka
{प्रतिज्ञां प्रतिजानीष्व यदि त्वं कर्तुमिच्छसि}
{अथ तद्व्याहरिष्यामि यदभिप्रार्थितं मया} %||2-10-15||

\twolineshloka
{एवमुक्तस्तया राजा प्रियया स्त्रीवशं गतः}
{तामुवाच महातेजाः कैकेयीमीषदुत्स्मितः} %||2-10-16||

\twolineshloka
{अवलिप्ते न जानासि त्वत्तः प्रियतरो मम}
{मनुजो मनुजव्याघ्राद्रामादन्यो न विद्यते} %||2-10-17||

\twolineshloka
{भद्रे हृदयमप्येतदनुमृश्श्योद्धरस्व मे}
{एतत्समीक्ष्य कैकेयि ब्रूहि यत्साधु मन्यसे} %||2-10-18||

\twolineshloka
{बलमात्मनि पश्यन्ती न मां शङ्कितुमर्हसि}
{करिष्यामि तव प्रीतिं सुकृतेनापि ते शपे} %||2-10-19||

\twolineshloka
{तेन वाक्येन संहृष्टा तमभिप्रायमात्मनः}
{व्याजहार महाघोरमभ्यागतमिवान्तकम्} %||2-10-20||

\twolineshloka
{यथाक्रमेण शपसि वरं मम ददासि च}
{तच्छृण्वन्तु त्रयस्त्रिंशद्देवाः सेन्द्रपुरोगमाः} %||2-10-21||

\twolineshloka
{चन्द्रादित्यौ नभश्चैव ग्रहा रात्र्यहनी दिशः}
{जगच्च पृथिवी चैव सगन्धर्वा सराक्षसा} %||2-10-22||

\twolineshloka
{निशाचराणि भूतानि गृहेषु गृहदेवताः}
{यानि चान्यानि भूतानि जानीयुर्भाषितं तव} %||2-10-23||

\twolineshloka
{सत्यसन्धो महातेजा धर्मज्ञः सुसमाहितः}
{वरं मम ददात्येष तन्मे शृण्वन्तु देवताः} %||2-10-24||

\twolineshloka
{इति देवी महेष्वासं परिगृह्याभिशस्य च}
{ततः परमुवाचेदं वरदं काममोहितम्} %||2-10-25||

\twolineshloka
{वरौ यौ मे त्वया देव तदा दत्तौ महीपते}
{तौ तावदहमद्यैव वक्ष्यामि शृणु मे वचः} %||2-10-26||

\twolineshloka
{अभिषेक समारम्भो राघवस्योपकल्पितः}
{अनेनैवाभिषेकेण भरतो मेऽभिषिच्यताम्} %||2-10-27||

\twolineshloka
{नव पञ्च च वर्षाणि दण्डकारण्यमाश्रितः}
{चीराजिनजटाधारी रामो भवतु तापसः} %||2-10-28||

\twolineshloka
{भरतो भजतामद्य यौवराज्यमकण्टकम्}
{अद्य चैव हि पश्येयं प्रयान्तं राघवं वने} %||2-10-29||

\twolineshloka
{ततः श्रुत्वा महाराज कैकेय्या दारुणं वचः}
{व्यथितो विलवश्चैव व्याघ्रीं दृष्ट्वा यथा मृगः} %||2-10-30||

\threelineshloka
{असंवृतायामासीनो जगत्यां दीर्घमुच्छ्वसन्}
{अहो धिगिति सामर्षो वाचमुक्त्वा नराधिपः}
{मोहमापेदिवान्भूयः शोकोपहतचेतनः} %||2-10-31||

\twolineshloka
{चिरेण तु नृपः संज्ञां प्रतिलभ्य सुदुःखितः}
{कैकेयीमब्रवीत्क्रुद्धः प्रदहन्निव चक्षुषा} %||2-10-32||

\twolineshloka
{नृशंसे दुष्टचारित्रे कुलस्यास्य विनाशिनि}
{किं कृतं तव रामेण पापे पापं मयापि वा} %||2-10-33||

\twolineshloka
{सदा ते जननी तुल्यां वृत्तिं वहति राघवः}
{तस्यैव त्वमनर्थाय किंनिमित्तमिहोद्यता} %||2-10-34||

\twolineshloka
{त्वं मयात्मविनाशाय भवनं स्वं प्रवेशिता}
{अविज्ञानान्नृपसुता व्याली तीक्ष्णविषा यथा} %||2-10-35||

\twolineshloka
{जीवलोको यदा सर्वो रामस्येह गुणस्तवम्}
{अपराधं कमुद्दिश्य त्यक्ष्यामीष्टमहं सुतम्} %||2-10-36||

\twolineshloka
{कौसल्यां वा सुमित्रां वा त्यजेयमपि वा श्रियम्}
{जीवितं वात्मनो रामं न त्वेव पितृवत्सलम्} %||2-10-37||

\twolineshloka
{परा भवति मे प्रीतिर्दृष्ट्वा तनयमग्रजम्}
{अपश्यतस्तु मे रामं नष्टा भवति चेतना} %||2-10-38||

\twolineshloka
{तिष्ठेल्लोको विना सूर्यं सस्यं वा सलिलं विना}
{न तु रामं विना देहे तिष्ठेत्तु मम जीवितम्} %||2-10-39||

\twolineshloka
{तदलं त्यज्यतामेष निश्चयः पापनिश्चये}
{अपि ते चरणौ मूर्ध्ना स्पृशाम्येष प्रसीद मे} %||2-10-40||

\fourlineindentedshloka
{स भूमिपालो विलपन्ननाथव}
{त्स्त्रिया गृहीतो दृहयेऽतिमात्रता}
{पपात देव्याश्चरणौ प्रसारिता}
{वुभावसंस्पृश्य यथातुरस्तथा} %||2-11-41||


{॥इत्यार्षे श्रीमद्रामायणे वाल्मीकीये आदिकाव्ये अयोध्याकाण्डे दशमः सर्गः॥}

\sect{एकादशः सर्गः}

\twolineshloka
{अतदर्हं महाराजं शयानमतथोचितम्}
{ययातिमिव पुण्यान्ते देवलोकात्परिच्युतम्} %||2-11-1||

\twolineshloka
{अनर्थरूपा सिद्धार्था अभीता भयदर्शिनी}
{पुनराकारयामास तमेव वरमङ्गना} %||2-11-2||

\twolineshloka
{त्वं कत्थसे महाराज सत्यवादी दृढव्रतः}
{मम चेमं वरं कस्माद्विधारयितुमिच्छसि} %||2-11-3||

\twolineshloka
{एवमुक्तस्तु कैकेय्या राजा दशरथस्तदा}
{प्रत्युवाच ततः क्रुद्धो मुहूर्तं विह्वलन्निव} %||2-11-4||

\twolineshloka
{मृते मयि गते रामे वनं मनुजपुङ्गवे}
{हन्तानार्ये ममामित्रे रामः प्रव्राजितो वनम्} %||2-11-5||

\twolineshloka
{यदि सत्यं ब्रवीम्येतत्तदसत्यं भविष्यति}
{अकीर्तिरतुला लोके ध्रुवं परिभवश्च मे} %||2-11-6||

\twolineshloka
{तथा विलपतस्तस्य परिभ्रमितचेतसः}
{अस्तमभ्यगमत्सूर्यो रजनी चाभ्यवर्तत} %||2-11-7||

\twolineshloka
{स त्रियामा तथार्तस्य चन्द्रमण्डलमण्डिता}
{राज्ञो विलपमानस्य न व्यभासत शर्वरी} %||2-11-8||

\twolineshloka
{तथैवोष्णं विनिःश्वस्य वृद्धो दशरथो नृपः}
{विललापार्तवद्दुःखं गगनासक्तलोचनः} %||2-11-9||

\threelineshloka
{न प्रभातं त्वयेच्छामि मयायं रचितोऽञ्जलिः}
{अथ वा गम्यतां शीघ्रं नाहमिच्छामि निर्घृणाम्}
{नृशंसां कैकेयीं द्रष्टुं यत्कृते व्यसनं महत्} %||2-11-10||

\twolineshloka
{एवमुक्त्वा ततो राजा कैकेयीं संयताञ्जलिः}
{प्रसादयामास पुनः कैकेयीं चेदमब्रवीत्} %||2-11-11||

\twolineshloka
{साधुवृत्तस्य दीनस्य त्वद्गतस्य गतायुषः}
{प्रसादः क्रियतां देवि भद्रे राज्ञो विशेषतः} %||2-11-12||

\twolineshloka
{शून्येन खलु सुश्रोणि मयेदं समुदाहृतम्}
{कुरु साधु प्रसादं मे बाले सहृदया ह्यसि} %||2-11-13||

\fourlineindentedshloka
{विशुद्धभावस्य हि दुष्टभावा}
{ ताम्रेक्षणस्याश्रुकलस्य राज्ञः}
{श्रुत्वा विचित्रं करुणं विलापम्}
{ भर्तुर्नृशंसा न चकार वाक्यम्} %||2-11-14||

\fourlineindentedshloka
{ततः स राजा पुनरेव मूर्छितः}
{ प्रियामतुष्टां प्रतिकूलभाषिणीम्}
{समीक्ष्य पुत्रस्य विवासनं प्रति}
{ क्षितौ विसंज्ञो निपपात दुःखितः} %||2-12-15||


{॥इत्यार्षे श्रीमद्रामायणे वाल्मीकीये आदिकाव्ये अयोध्याकाण्डे एकादशः सर्गः॥}

\sect{द्वादशः सर्गः}

\twolineshloka
{पुत्रशोकार्दितं पापा विसंज्ञं पतितं भुवि}
{विवेष्टमानमुदीक्ष्य सैक्ष्वाकमिदमब्रवीत्} %||2-12-1||

\twolineshloka
{पापं कृत्वेव किमिदं मम संश्रुत्य संश्रवम्}
{शेषे क्षितितले सन्नः स्थित्यां स्थातुं त्वमर्हसि} %||2-12-2||

\twolineshloka
{आहुः सत्यं हि परमं धर्मं धर्मविदो जनाः}
{सत्यमाश्रित्य हि मया त्वं च धर्मं प्रचोदितः} %||2-12-3||

\twolineshloka
{संश्रुत्य शैब्यः श्येनाय स्वां तनुं जगतीपतिः}
{प्रदाय पक्षिणो राजञ्जगाम गतिमुत्तमाम्} %||2-12-4||

\twolineshloka
{तथ ह्यलर्कस्तेजस्वी ब्राह्मणे वेदपारगे}
{याचमाने स्वके नेत्रे उद्धृत्याविमना ददौ} %||2-12-5||

\twolineshloka
{सरितां तु पतिः स्वल्पां मर्यादां सत्यमन्वितः}
{सत्यानुरोधात्समये वेलां खां नातिवर्तते} %||2-12-6||

\twolineshloka
{समयं च ममार्येमं यदि त्वं न करिष्यसि}
{अग्रतस्ते परित्यक्ता परित्यक्ष्यामि जीवितम्} %||2-12-7||

\twolineshloka
{एवं प्रचोदितो राजा कैकेय्या निर्विशङ्कया}
{नाशकत्पाशमुन्मोक्तुं बलिरिन्द्रकृतं यथा} %||2-12-8||

\twolineshloka
{उद्भ्रान्तहृदयश्चापि विवर्णवनदोऽभवत्}
{स धुर्यो वै परिस्पन्दन्युगचक्रान्तरं यथा} %||2-12-9||

\twolineshloka
{विह्वलाभ्यां च नेत्राभ्यामपश्यन्निव भूमिपः}
{कृच्छ्राद्धैर्येण संस्तभ्य कैकेयीमिदमब्रवीत्} %||2-12-10||

\twolineshloka
{यस्ते मन्त्रकृतः पाणिरग्नौ पापे मया धृतः}
{तं त्यजामि स्वजं चैव तव पुत्रं सह त्वया} %||2-12-11||

\twolineshloka
{ततः पापसमाचारा कैकेयी पार्थिवं पुनः}
{उवाच परुषं वाक्यं वाक्यज्ञा रोषमूर्छिता} %||2-12-12||

\twolineshloka
{किमिदं भाषसे राजन्वाक्यं गररुजोपमम्}
{आनाययितुमक्लिष्टं पुत्रं राममिहार्हसि} %||2-12-13||

\twolineshloka
{स्थाप्य राज्ये मम सुतं कृत्वा रामं वनेचरम्}
{निःसपत्नां च मां कृत्वा कृतकृत्यो भविष्यसि} %||2-12-14||

\twolineshloka
{स नुन्न इव तीक्षेण प्रतोदेन हयोत्तमः}
{राजा प्रदोचितोऽभीक्ष्णं कैकेयीमिदमब्रवीत्} %||2-12-15||

\twolineshloka
{धर्मबन्धेन बद्धोऽस्मि नष्टा च मम चेतना}
{ज्येष्ठं पुत्रं प्रियं रामं द्रष्टुमिच्छामि धार्मिकम्} %||2-12-16||

\twolineshloka
{इति राज्ञो वचः श्रुत्वा कैकेयी तदनन्तरम्}
{स्वयमेवाब्रवीत्सूतं गच्छ त्वं राममानय} %||2-12-17||

\twolineshloka
{ततः स राजा तं सूतं सन्नहर्षः सुतं प्रति}
{शोकारक्तेक्षणः श्रीमानुद्वीक्ष्योवाच धार्मिकः} %||2-12-18||

\twolineshloka
{सुमन्त्रः करुणं श्रुत्वा दृष्ट्वा दीनं च पार्थिवम्}
{प्रगृहीताञ्जलिः किञ्चित्तस्माद्देशादपाक्रमन्} %||2-12-19||

\twolineshloka
{यदा वक्तुं स्वयं दैन्यान्न शशाक महीपतिः}
{तदा सुमन्त्रं मन्त्रज्ञा कैकेयी प्रत्युवाच ह} %||2-12-20||

\twolineshloka
{सुमन्त्र रामं द्रक्ष्यामि शीघ्रमानय सुन्दरम्}
{स मन्यमानः कल्याणं हृदयेन ननन्द च} %||2-12-21||

\twolineshloka
{सुमन्त्रश्चिन्तयामास त्वरितं चोदितस्तया}
{व्यक्तं रामोऽभिषेकार्थमिहायास्यति धर्मवित्} %||2-12-22||

\twolineshloka
{इति सूतो मतिं कृत्वा हर्षेण महता पुनः}
{निर्जगाम महातेजा राघवस्य दिदृक्षया} %||2-12-23||

\fourlineindentedshloka
{ततः पुरस्तात्सहसा विनिर्गतो}
{ महीपतीन्द्वारगतान्विलोकयन्}
{ददर्श पौरान्विविधान्महाधना}
{नुपस्थितान्द्वारमुपेत्य विष्ठितान्} %||2-13-24||


{॥इत्यार्षे श्रीमद्रामायणे वाल्मीकीये आदिकाव्ये अयोध्याकाण्डे द्वादशः सर्गः॥}

\sect{त्रयोदशः सर्गः}

\twolineshloka
{ते तु तां रजनीमुष्य ब्राह्मणा वेदपारगाः}
{उपतस्थुरुपस्थानं सहराजपुरोहिताः} %||2-13-1||

\twolineshloka
{अमात्या बलमुख्याश्च मुख्या ये निगमस्य च}
{राघवस्याभिषेकार्थे प्रीयमाणास्तु सङ्गताः} %||2-13-2||

\twolineshloka
{उदिते विमले सूर्ये पुष्ये चाभ्यागतेऽहनि}
{अभिषेकाय रामस्य द्विजेन्द्रैरुपकल्पितम्} %||2-13-3||

\twolineshloka
{काञ्चना जलकुम्भाश्च भद्रपीठं स्वलङ्कृतम्}
{रामश्च सम्यगास्तीर्णो भास्वरा व्याघ्रचर्मणा} %||2-13-4||

\twolineshloka
{गङ्गायमुनयोः पुण्यात्सङ्गमादाहृतं जलम्}
{याश्चान्याः सरितः पुण्या ह्रदाः कूपाः सरांसि च} %||2-13-5||

\twolineshloka
{प्राग्वाहाश्चोर्ध्ववाहाश्च तिर्यग्वाहाः समाहिताः}
{ताभ्यश्चैवाहृतं तोयं समुद्रेभ्यश्च सर्वशः} %||2-13-6||

\threelineshloka
{क्षौद्रं दधिघृतं लाजा धर्भाः सुमनसः पयः}
{सलाजाः क्षीरिभिश्छन्ना घटाः काञ्चनराजताः}
{पद्मोत्पलयुता भान्ति पूर्णाः परमवारिणा} %||2-13-7||

\twolineshloka
{चन्द्रांशुविकचप्रख्यं पाण्डुरं रत्नभूषितम्}
{सज्जं तिष्ठति रामस्य वालव्यजनमुत्तमम्} %||2-13-8||

\twolineshloka
{चन्द्रमण्डलसङ्काशमातपत्रं च पाण्डुरम्}
{सज्जं द्युतिकरं श्रीमदभिषेकपुरस्कृतम्} %||2-13-9||

\twolineshloka
{पाण्डुरश्च वृषः सज्जः पाण्डुराश्वश्च सुस्थितः}
{प्रस्रुतश्च गजः श्रीमानौपवाह्यः प्रतीक्षते} %||2-13-10||

\twolineshloka
{अष्टौ कन्याश्च मङ्गल्याः सर्वाभरणभूषिताः}
{वादित्राणि च सर्वाणि बन्दिनश्च तथापरे} %||2-13-11||

\twolineshloka
{इक्ष्वाकूणां यथा राज्ये सम्भ्रियेताभिषेचनम्}
{तथा जातीयामादाय राजपुत्राभिषेचनम्} %||2-13-12||

\twolineshloka
{ते राजवचनात्तत्र समवेता महीपतिम्}
{अपश्यन्तोऽब्रुवन्को नु राज्ञो नः प्रतिवेदयेत्} %||2-13-13||

\twolineshloka
{न पश्यामश्च राजानमुदितश्च दिवाकरः}
{यौवराज्याभिषेकश्च सज्जो रामस्य धीमतः} %||2-13-14||

\twolineshloka
{इति तेषु ब्रुवाणेषु सार्वभौमान्महीपतीन्}
{अब्रवीत्तानिदं सर्वान्सुमन्त्रो राजसत्कृतः} %||2-13-15||

\twolineshloka
{अयं पृच्छामि वचनात्सुखमायुष्मतामहम्}
{राज्ञः सम्प्रतिबुद्धस्य यच्चागमनकारणम्} %||2-13-16||

\twolineshloka
{इत्युक्त्वान्तःपुरद्वारमाजगाम पुराणवित्}
{आशीर्भिर्गुणयुक्ताभिरभितुष्टाव राघवम्} %||2-13-17||

\twolineshloka
{गता भगवती रात्रिरहः शिवमुपस्थितम्}
{बुध्यस्व नृपशार्दूल कुरु कार्यमनन्तरम्} %||2-13-18||

\twolineshloka
{ब्राह्मणा बलमुख्याश्च नैगमाश्चागता नृप}
{दर्शनं प्रतिकाङ्क्षन्ते प्रतिबुध्यस्व राघव} %||2-13-19||

\twolineshloka
{स्तुवन्तं तं तदा सूतं सुमन्त्रं मन्त्रकोविदम्}
{प्रतिबुध्य ततो राजा इदं वचनमब्रवीत्} %||2-13-20||

\twolineshloka
{न चैव सम्प्रसुतोऽहमानयेदाशु राघवम्}
{इति राजा दशरथः सूतं तत्रान्वशात्पुनः} %||2-13-21||

\twolineshloka
{स राजवचनं श्रुत्वा शिरसा प्रतिपूज्य तम्}
{निर्जगाम नृपावासान्मन्यमानः प्रियं महत्} %||2-13-22||

\twolineshloka
{प्रपन्नो राजमार्गं च पताका ध्वजशोभितम्}
{स सूतस्तत्र शुश्राव रामाधिकरणाः कथाः} %||2-13-23||

\twolineshloka
{ततो ददर्श रुचिरं कैलाससदृशप्रभम्}
{रामवेश्म सुमन्त्रस्तु शक्रवेश्मसमप्रभम्} %||2-13-24||

\twolineshloka
{महाकपाटपिहितं वितर्दिशतशोभितम्}
{काञ्चनप्रतिमैकाग्रं मणिविद्रुमतोरणम्} %||2-13-25||

\twolineshloka
{शारदाभ्रघनप्रख्यं दीप्तं मेरुगुहोपमम्}
{दामभिर्वरमाल्यानां सुमहद्भिरलङ्कृतम्} %||2-13-26||

\fourlineindentedshloka
{स वाजियुक्तेन रथेन सारथि}
{र्नराकुलं राजकुलं विलोकयन्}
{ततः समासाद्य महाधनं मह}
{त्प्रहृष्टरोमा स बभूव सारथिः} %||2-13-27||

\fourlineindentedshloka
{तदद्रिकूटाचलमेघसंनिभम्}
{ महाविमानोत्तमवेश्मसङ्घवत्}
{अवार्यमाणः प्रविवेश सारथिः}
{ प्रभूतरत्नं मकरो यथार्णवम्} %||2-14-28||


{॥इत्यार्षे श्रीमद्रामायणे वाल्मीकीये आदिकाव्ये अयोध्याकाण्डे त्रयोदशः सर्गः॥}

\sect{चतुर्दशः सर्गः}

\twolineshloka
{स तदन्तःपुरद्वारं समतीत्य जनाकुलम्}
{प्रविविक्तां ततः कक्ष्यामाससाद पुराणवित्} %||2-14-1||

\twolineshloka
{प्रासकार्मुकबिभ्रद्भिर्युवभिर्मृष्टकुण्डलैः}
{अप्रमादिभिरेकाग्रैः स्वनुरक्तैरधिष्ठिताम्} %||2-14-2||

\twolineshloka
{तत्र काषायिणो वृद्धान्वेत्रपाणीन्स्वलङ्कृतान्}
{ददर्श विष्ठितान्द्वारि स्त्र्यध्यक्षान्सुसमाहितान्} %||2-14-3||

\twolineshloka
{ते समीक्ष्य समायान्तं रामप्रियचिकीर्षवः}
{सहभार्याय रामाय क्षिप्रमेवाचचक्षिरे} %||2-14-4||

\twolineshloka
{प्रतिवेदितमाज्ञाय सूतमभ्यन्तरं पितुः}
{तत्रैवानाययामास राघवः प्रियकाम्यया} %||2-14-5||

\twolineshloka
{तं वैश्रवणसङ्काशमुपविष्टं स्वलङ्कृतम्}
{दादर्श सूतः पर्यङ्के सौवणो सोत्तरच्छदे} %||2-14-6||

\twolineshloka
{वराहरुधिराभेण शुचिना च सुगन्धिना}
{अनुलिप्तं परार्ध्येन चन्दनेन परन्तपम्} %||2-14-7||

\twolineshloka
{स्थितया पार्श्वतश्चापि वालव्यजनहस्तया}
{उपेतं सीतया भूयश्चित्रया शशिनं यथा} %||2-14-8||

\twolineshloka
{तं तपन्तमिवादित्यमुपपन्नं स्वतेजसा}
{ववन्दे वरदं बन्दी नियमज्ञो विनीतवत्} %||2-14-9||

\twolineshloka
{प्राञ्जलिस्तु सुखं पृष्ट्वा विहारशयनासने}
{राजपुत्रमुवाचेदं सुमन्त्रो राजसत्कृतः} %||2-14-10||

\twolineshloka
{कौसल्या सुप्रभा देव पिता त्वं द्रष्टुमिच्छति}
{महिष्या सह कैकेय्या गम्यतां तत्र माचिरम्} %||2-14-11||

\twolineshloka
{एवमुक्तस्तु संहृष्टो नरसिंहो महाद्युतिः}
{ततः सम्मानयामास सीतामिदमुवाच ह} %||2-14-12||

\twolineshloka
{देवि देवश्च देवी च समागम्य मदन्तरे}
{मन्त्रेयेते ध्रुवं किञ्चिदभिषेचनसंहितम्} %||2-14-13||

\twolineshloka
{लक्षयित्वा ह्यभिप्रायं प्रियकामा सुदक्षिणा}
{सञ्चोदयति राजानं मदर्थं मदिरेक्षणा} %||2-14-14||

\twolineshloka
{यादृशी परिषत्तत्र तादृशो दूत आगतः}
{ध्रुवमद्यैव मां राजा यौवराज्येऽभिषेक्ष्यति} %||2-14-15||

\twolineshloka
{हन्त शीघ्रमितो गत्वा द्रक्ष्यामि च महीपतिः}
{सह त्वं परिवारेण सुखमास्स्व रमस्य च} %||2-14-16||

\twolineshloka
{पतिसम्मानिता सीता भर्तारमसितेक्षणा}
{आद्वारमनुवव्राज मङ्गलान्यभिदध्युषी} %||2-14-17||

\twolineshloka
{स सर्वानर्थिनो दृष्ट्वा समेत्य प्रतिनन्द्य च}
{ततः पावकसङ्काशमारुरोह रथोत्तमम्} %||2-14-18||

\twolineshloka
{मुष्णन्तमिव चक्षूंषि प्रभया हेमवर्चसम्}
{करेणुशिशुकल्पैश्च युक्तं परमवाजिभिः} %||2-14-19||

\twolineshloka
{हरियुक्तं सहस्राक्षो रथमिन्द्र इवाशुगम्}
{प्रययौ तूर्णमास्थाय राघवो ज्वलितः श्रिया} %||2-14-20||

\twolineshloka
{स पर्जन्य इवाकाशे स्वनवानभिनादयन्}
{निकेतान्निर्ययौ श्रीमान्महाभ्रादिव चन्द्रमाः} %||2-14-21||

\twolineshloka
{छत्रचामरपाणिस्तु लक्ष्मणो राघवानुजः}
{जुगोप भ्रातरं भ्राता रथमास्थाय पृष्ठतः} %||2-14-22||

\twolineshloka
{ततो हलहलाशब्दस्तुमुलः समजायत}
{तस्य निष्क्रममाणस्य जनौघस्य समन्ततः} %||2-14-23||

\fourlineindentedshloka
{स राघवस्तत्र कथाप्रलापम्}
{ शुश्राव लोकस्य समागतस्य}
{आत्माधिकारा विविधाश्च वाचः}
{ प्रहृष्टरूपस्य पुरे जनस्य} %||2-14-24||

\threelineshloka
{एष श्रियं गच्छति राघवोऽद्य; राजप्रसादाद्विपुलां गमिष्यन्}
{एते वयं सर्वसमृद्धकामा; येषामयं नो भविता प्रशास्ता}
{लाभो जनस्यास्य यदेष सर्वम्; प्रपत्स्यते राष्ट्रमिदं चिराय} %||2-14-25||

\fourlineindentedshloka
{स घोषवद्भिश्च हयैः सनागैः}
{ पुरःसरैः स्वस्तिकसूतमागधैः}
{महीयमानः प्रवरैश्च वादकै}
{रभिष्टुतो वैश्रवणो यथा ययौ} %||2-14-26||

\fourlineindentedshloka
{करेणुमातङ्गरथाश्वसङ्कुलम्}
{ महाजनौघैः परिपूर्णचत्वरम्}
{प्रभूतरत्नं बहुपण्यसञ्चयम्}
{ ददर्श रामो रुचिरं महापथम्} %||2-15-27||


{॥इत्यार्षे श्रीमद्रामायणे वाल्मीकीये आदिकाव्ये अयोध्याकाण्डे चतुर्दशः सर्गः॥}

\sect{पञ्चदशः सर्गः}

\twolineshloka
{स रामो रथमास्थाय सम्प्रहृष्टसुहृज्जनः}
{अपश्यन्नगरं श्रीमान्नानाजनसमाकुलम्} %||2-15-1||

\twolineshloka
{स गृहैरभ्रसङ्काशैः पाण्डुरैरुपशोभितम्}
{राजमार्गं ययौ रामो मध्येनागरुधूपितम्} %||2-15-2||

\twolineshloka
{शोभमानमसम्बाधं तं राजपथमुत्तमम्}
{संवृतं विविधैः पण्यैर्भक्ष्यैरुच्चावचैरपि} %||2-15-3||

\twolineshloka
{आशीर्वादान्बहूञ्शृण्वन्सुहृद्भिः समुदीरितान्}
{यथार्हं चापि सम्पूज्य सर्वानेव नरान्ययौ} %||2-15-4||

\twolineshloka
{पितामहैराचरितं तथैव प्रपितामहैः}
{अद्योपादाय तं मार्गमभिषिक्तोऽनुपालय} %||2-15-5||

\twolineshloka
{यथा स्म लालिताः पित्रा यथा पूर्वैः पितामहैः}
{ततः सुखतरं सर्वे रामे वत्स्याम राजनि} %||2-15-6||

\twolineshloka
{अलमद्य हि भुक्तेन परमार्थैरलं च नः}
{यथा पश्याम निर्यान्तं रामं राज्ये प्रतिष्ठितम्} %||2-15-7||

\twolineshloka
{अतो हि न प्रियतरं नान्यत्किञ्चिद्भविष्यति}
{यथाभिषेको रामस्य राज्येनामिततेजसः} %||2-15-8||

\twolineshloka
{एताश्चान्याश्च सुहृदामुदासीनः कथाः शुभाः}
{आत्मसम्पूजनीः शृण्वन्ययौ रामो महापथम्} %||2-15-9||

\twolineshloka
{न हि तस्मान्मनः कश्चिच्चक्षुषी वा नरोत्तमात्}
{नरः शक्नोत्यपाक्रष्टुमतिक्रान्तेऽपि राघवे} %||2-15-10||

\twolineshloka
{सर्वेषां स हि धर्मात्मा वर्णानां कुरुते दयाम्}
{चतुर्णां हि वयःस्थानां तेन ते तमनुव्रताः} %||2-15-11||

\twolineshloka
{स राजकुलमासाद्य महेन्द्रभवनोपमम्}
{राजपुत्रः पितुर्वेश्म प्रविवेश श्रिया ज्वलन्} %||2-15-12||

\twolineshloka
{स सर्वाः समतिक्रम्य कक्ष्या दशरथात्मजः}
{संनिवर्त्य जनं सर्वं शुद्धान्तःपुरमभ्यगात्} %||2-15-13||

\fourlineindentedshloka
{ततः प्रविष्टे पितुरन्तिकं तदा}
{ जनः स सर्वो मुदितो नृपात्मजे}
{प्रतीक्षते तस्य पुनः स्म निर्गमम्}
{ यथोदयं चन्द्रमसः सरित्पतिः} %||2-16-14||


{॥इत्यार्षे श्रीमद्रामायणे वाल्मीकीये आदिकाव्ये अयोध्याकाण्डे पञ्चदशः सर्गः॥}

\sect{षोडशः सर्गः}

\twolineshloka
{स ददर्शासने रामो निषण्णं पितरं शुभे}
{कैकेयीसहितं दीनं मुखेन परिशुष्यता} %||2-16-1||

\twolineshloka
{स पितुश्चरणौ पूर्वमभिवाद्य विनीतवत्}
{ततो ववन्दे चरणौ कैकेय्याः सुसमाहितः} %||2-16-2||

\twolineshloka
{रामेत्युक्त्वा च वचनं वाष्पपर्याकुलेक्षणः}
{शशाक नृपतिर्दीनो नेक्षितुं नाभिभाषितुम्} %||2-16-3||

\twolineshloka
{तदपूर्वं नरपतेर्दृष्ट्वा रूपं भयावहम्}
{रामोऽपि भयमापन्नः पदा स्पृष्ट्वेव पन्नगम्} %||2-16-4||

\twolineshloka
{इन्द्रियैरप्रहृष्टैस्तं शोकसन्तापकर्शितम्}
{निःश्वसन्तं महाराजं व्यथिताकुलचेतसम्} %||2-16-5||

\twolineshloka
{ऊर्मि मालिनमक्षोभ्यं क्षुभ्यन्तमिव सागरम्}
{उपप्लुतमिवादित्यमुक्तानृतमृषिं यथा} %||2-16-6||

\twolineshloka
{अचिन्त्यकल्पं हि पितुस्तं शोकमुपधारयन्}
{बभूव संरब्धतरः समुद्र इव पर्वणि} %||2-16-7||

\twolineshloka
{चिन्तयामास च तदा रामः पितृहिते रतः}
{किंस्विदद्यैव नृपतिर्न मां प्रत्यभिनन्दति} %||2-16-8||

\twolineshloka
{अन्यदा मां पिता दृष्ट्वा कुपितोऽपि प्रसीदति}
{तस्य मामद्य सम्प्रेक्ष्य किमायासः प्रवर्तते} %||2-16-9||

\twolineshloka
{स दीन इव शोकार्तो विषण्णवदनद्युतिः}
{कैकेयीमभिवाद्यैव रामो वचनमब्रवीत्} %||2-16-10||

\twolineshloka
{कच्चिन्मया नापराधमज्ञानाद्येन मे पिता}
{कुपितस्तन्ममाचक्ष्व त्वं चैवैनं प्रसादय} %||2-16-11||

\threelineshloka
{विवर्णवदनो दीनो न हि मामभिभाषते}
{शारीरो मानसो वापि कच्चिदेनं न बाधते}
{सन्तापो वाभितापो वा दुर्लभं हि सदा सुखम्} %||2-16-12||

\twolineshloka
{कच्चिन्न किञ्चिद्भरते कुमारे प्रियदर्शने}
{शत्रुघ्ने वा महासत्त्वे मातॄणां वा ममाशुभम्} %||2-16-13||

\twolineshloka
{अतोषयन्महाराजमकुर्वन्वा पितुर्वचः}
{मुहूर्तमपि नेच्छेयं जीवितुं कुपिते नृपे} %||2-16-14||

\twolineshloka
{यतोमूलं नरः पश्येत्प्रादुर्भावमिहात्मनः}
{कथं तस्मिन्न वर्तेत प्रत्यक्षे सति दैवते} %||2-16-15||

\twolineshloka
{कच्चित्ते परुषं किञ्चिदभिमानात्पिता मम}
{उक्तो भवत्या कोपेन यत्रास्य लुलितं मनः} %||2-16-16||

\twolineshloka
{एतदाचक्ष्व मे देवि तत्त्वेन परिपृच्छतः}
{किंनिमित्तमपूर्वोऽयं विकारो मनुजाधिपे} %||2-16-17||

\threelineshloka
{अहं हि वचनाद्राज्ञः पतेयमपि पावके}
{भक्षयेयं विषं तीक्ष्णं मज्जेयमपि चार्णवे}
{नियुक्तो गुरुणा पित्रा नृपेण च हितेन च} %||2-16-18||

\twolineshloka
{तद्ब्रूहि वचनं देवि राज्ञो यदभिकाङ्क्षितम्}
{करिष्ये प्रतिजाने च रामो द्विर्नाभिभाषते} %||2-16-19||

\twolineshloka
{तमार्जवसमायुक्तमनार्या सत्यवादिनम्}
{उवाच रामं कैकेयी वचनं भृशदारुणम्} %||2-16-20||

\twolineshloka
{पुरा देवासुरे युद्धे पित्रा ते मम राघव}
{रक्षितेन वरौ दत्तौ सशल्येन महारणे} %||2-16-21||

\twolineshloka
{तत्र मे याचितो राजा भरतस्याभिषेचनम्}
{गमनं दण्डकारण्ये तव चाद्यैव राघव} %||2-16-22||

\twolineshloka
{यदि सत्यप्रतिज्ञं त्वं पितरं कर्तुमिच्छसि}
{आत्मानं च नररेष्ठ मम वाक्यमिदं शृणु} %||2-16-23||

\twolineshloka
{स निदेशे पितुस्तिष्ठ यथा तेन प्रतिश्रुतम्}
{त्वयारण्यं प्रवेष्टव्यं नव वर्षाणि पञ्च च} %||2-16-24||

\twolineshloka
{सप्त सप्त च वर्षाणि दण्डकारण्यमाश्रितः}
{अभिषेकमिमं त्यक्त्वा जटाचीरधरो वस} %||2-16-25||

\twolineshloka
{भरतः कोसलपुरे प्रशास्तु वसुधामिमाम्}
{नानारत्नसमाकीर्णां सवाजिरथकुञ्जराम्} %||2-16-26||

\twolineshloka
{तदप्रियममित्रघ्नो वचनं मरणोपमम्}
{श्रुत्वा न विव्यथे रामः कैकेयीं चेदमब्रवीत्} %||2-16-27||

\twolineshloka
{एवमस्तु गमिष्यामि वनं वस्तुमहं त्वतः}
{जटाचीरधरो राज्ञः प्रतिज्ञामनुपालयन्} %||2-16-28||

\twolineshloka
{इदं तु ज्ञातुमिच्छामि किमर्थं मां महीपतिः}
{नाभिनन्दति दुर्धर्षो यथापुरमरिन्दमः} %||2-16-29||

\twolineshloka
{मन्युर्न च त्वया कार्यो देवि ब्रूहि तवाग्रतः}
{यास्यामि भव सुप्रीता वनं चीरजटाधरः} %||2-16-30||

\twolineshloka
{हितेन गुरुणा पित्रा कृतज्ञेन नृपेण च}
{नियुज्यमानो विश्रब्धं किं न कुर्यादहं प्रियम्} %||2-16-31||

\twolineshloka
{अलीकं मानसं त्वेकं हृदयं दहतीव मे}
{स्वयं यन्नाह मां राजा भरतस्याभिषेचनम्} %||2-16-32||

\twolineshloka
{अहं हि सीतां राज्यं च प्राणानिष्टान्धनानि च}
{हृष्टो भ्रात्रे स्वयं दद्यां भरतायाप्रचोदितः} %||2-16-33||

\twolineshloka
{किं पुनर्मनुजेन्द्रेण स्वयं पित्रा प्रचोदितः}
{तव च प्रियकामार्थं प्रतिज्ञामनुपालयन्} %||2-16-34||

\twolineshloka
{तदाश्वासय हीमं त्वं किं न्विदं यन्महीपतिः}
{वसुधासक्तनयनो मन्दमश्रूणि मुञ्चति} %||2-16-35||

\twolineshloka
{गच्छन्तु चैवानयितुं दूताः शीघ्रजवैर्हयैः}
{भरतं मातुलकुलादद्यैव नृपशासनात्} %||2-16-36||

\twolineshloka
{दण्डकारण्यमेषोऽहमितो गच्छामि सत्वरः}
{अविचार्य पितुर्वाक्यं समावस्तुं चतुर्दश} %||2-16-37||

\twolineshloka
{सा हृष्टा तस्य तद्वाक्यं श्रुत्वा रामस्य कैकयी}
{प्रस्थानं श्रद्दधाना हि त्वरयामास राघवम्} %||2-16-38||

\twolineshloka
{एवं भवतु यास्यन्ति दूताः शीघ्रजवैर्हयैः}
{भरतं मातुलकुलादुपावर्तयितुं नराः} %||2-16-39||

\twolineshloka
{तव त्वहं क्षमं मन्ये नोत्सुकस्य विलम्बनम्}
{राम तस्मादितः शीघ्रं वनं त्वं गन्तुमर्हसि} %||2-16-40||

\twolineshloka
{व्रीडान्वितः स्वयं यच्च नृपस्त्वां नाभिभाषते}
{नैतत्किञ्चिन्नरश्रेष्ठ मन्युरेषोऽपनीयताम्} %||2-16-41||

\twolineshloka
{यावत्त्वं न वनं यातः पुरादस्मादभित्वरन्}
{पिता तावन्न ते राम स्नास्यते भोक्ष्यतेऽपि वा} %||2-16-42||

\twolineshloka
{धिक्कष्टमिति निःश्वस्य राजा शोकपरिप्लुतः}
{मूर्छितो न्यपतत्तस्मिन्पर्यङ्के हेमभूषिते} %||2-16-43||

\twolineshloka
{रामोऽप्युत्थाप्य राजानं कैकेय्याभिप्रचोदितः}
{कशयेवाहतो वाजी वनं गन्तुं कृतत्वरः} %||2-16-44||

\twolineshloka
{तदप्रियमनार्याया वचनं दारुणोदरम्}
{श्रुत्वा गतव्यथो रामः कैकेयीं वाक्यमब्रवीत्} %||2-16-45||

\twolineshloka
{नाहमर्थपरो देवि लोकमावस्तुमुत्सहे}
{विद्धि मामृषिभिस्तुल्यं केवलं धर्ममास्थितम्} %||2-16-46||

\twolineshloka
{यदत्रभवतः किञ्चिच्छक्यं कर्तुं प्रियं मया}
{प्राणानपि परित्यज्य सर्वथा कृतमेव तत्} %||2-16-47||

\twolineshloka
{न ह्यतो धर्मचरणं किञ्चिदस्ति महत्तरम्}
{यथा पितरि शुश्रूषा तस्य वा वचनक्रिया} %||2-16-48||

\twolineshloka
{अनुक्तोऽप्यत्रभवता भवत्या वचनादहम्}
{वने वत्स्यामि विजने वर्षाणीह चतुर्दश} %||2-16-49||

\twolineshloka
{न नूनं मयि कैकेयि किञ्चिदाशंससे गुणम्}
{यद्राजानमवोचस्त्वं ममेश्वरतरा सती} %||2-16-50||

\twolineshloka
{यावन्मातरमापृच्छे सीतां चानुनयाम्यहम्}
{ततोऽद्यैव गमिष्यामि दण्डकानां महद्वनम्} %||2-16-51||

\twolineshloka
{भरतः पालयेद्राज्यं शुश्रूषेच्च पितुर्यथा}
{तथा भवत्या कर्तव्यं स हि धर्मः सनातनः} %||2-16-52||

\twolineshloka
{स रामस्य वचः श्रुत्वा भृशं दुःखहतः पिता}
{शोकादशक्नुवन्बाष्पं प्ररुरोद महास्वनम्} %||2-16-53||

\twolineshloka
{वन्दित्वा चरणौ रामो विसंज्ञस्य पितुस्तदा}
{कैकेय्याश्चाप्यनार्याया निष्पपात महाद्युतिः} %||2-16-54||

\twolineshloka
{स रामः पितरं कृत्वा कैकेयीं च प्रदक्षिणम्}
{निष्क्रम्यान्तःपुरात्तस्मात्स्वं ददर्श सुहृज्जनम्} %||2-16-55||

\twolineshloka
{तं बाष्पपरिपूर्णाक्षः पृष्ठतोऽनुजगाम ह}
{लक्ष्मणः परमक्रुद्धः सुमित्रानन्दवर्धनः} %||2-16-56||

\twolineshloka
{आभिषेचनिकं भाण्डं कृत्वा रामः प्रदक्षिणम्}
{शनैर्जगाम सापेक्षो दृष्टिं तत्राविचालयन्} %||2-16-57||

\twolineshloka
{न चास्य महतीं लक्ष्मीं राज्यनाशोऽपकर्षति}
{लोककान्तस्य कान्तत्वं शीतरश्मेरिव क्षपा} %||2-16-58||

\twolineshloka
{न वनं गन्तुकामस्य त्यजतश्च वसुन्धराम्}
{सर्वलोकातिगस्येव लक्ष्यते चित्तविक्रिया} %||2-16-59||

\twolineshloka
{धारयन्मनसा दुःखमिन्द्रियाणि निगृह्य च}
{प्रविवेशात्मवान्वेश्म मातुरप्रियशंसिवान्} %||2-16-60||

\fourlineindentedshloka
{प्रविश्य वेश्मातिभृशं मुदान्वितम्}
{ समीक्ष्य तां चार्थविपत्तिमागताम्}
{न चैव रामोऽत्र जगाम विक्रियाम्}
{ सुहृज्जनस्यात्मविपत्तिशङ्कया} %||2-17-61||


{॥इत्यार्षे श्रीमद्रामायणे वाल्मीकीये आदिकाव्ये अयोध्याकाण्डे षोडशः सर्गः॥}

\sect{सप्तदशः सर्गः}

\twolineshloka
{रामस्तु भृशमायस्तो निःश्वसन्निव कुञ्जरः}
{जगाम सहितो भ्रात्रा मातुरन्तःपुरं वशी} %||2-17-1||

\twolineshloka
{सोऽपश्यत्पुरुषं तत्र वृद्धं परमपूजितम्}
{उपविष्टं गृहद्वारि तिष्ठतश्चापरान्बहून्} %||2-17-2||

\twolineshloka
{प्रविश्य प्रथमां कक्ष्यां द्वितीयायां ददर्श सः}
{ब्राह्मणान्वेदसम्पन्नान्वृद्धान्राज्ञाभिसत्कृतान्} %||2-17-3||

\twolineshloka
{प्रणम्य रामस्तान्वृद्धांस्तृतीयायां ददर्श सः}
{स्त्रियो वृद्धाश्च बालाश्च द्वाररक्षणतत्पराः} %||2-17-4||

\twolineshloka
{वर्धयित्वा प्रहृष्टास्ताः प्रविश्य च गृहं स्त्रियः}
{न्यवेदयन्त त्वरिता राम मातुः प्रियं तदा} %||2-17-5||

\twolineshloka
{कौसल्यापि तदा देवी रात्रिं स्थित्वा समाहिता}
{प्रभाते त्वकरोत्पूजां विष्णोः पुत्रहितैषिणी} %||2-17-6||

\twolineshloka
{सा क्षौमवसना हृष्टा नित्यं व्रतपरायणा}
{अग्निं जुहोति स्म तदा मन्त्रवत्कृतमङ्गला} %||2-17-7||

\twolineshloka
{प्रविश्य च तदा रामो मातुरन्तःपुरं शुभम्}
{ददर्श मातरं तत्र हावयन्तीं हुताशनम्} %||2-17-8||

\twolineshloka
{सा चिरस्यात्मजं दृष्ट्वा मातृनन्दनमागतम्}
{अभिचक्राम संहृष्टा किशोरं वडवा यथा} %||2-17-9||

\twolineshloka
{तमुवाच दुराधर्षं राघवं सुतमात्मनः}
{कौसल्या पुत्रवात्सल्यादिदं प्रियहितं वचः} %||2-17-10||

\twolineshloka
{वृद्धानां धर्मशीलानां राजर्षीणां महात्मनाम्}
{प्राप्नुह्यायुश्च कीर्तिं च धर्मं चोपहितं कुले} %||2-17-11||

\twolineshloka
{सत्यप्रतिज्ञं पितरं राजानं पश्य राघव}
{अद्यैव हि त्वां धर्मात्मा यौवराज्येऽभिषेक्ष्यति} %||2-17-12||

\twolineshloka
{मातरं राघवः किञ्चित्प्रसार्याञ्जलिमब्रवीत्}
{स स्वभावविनीतश्च गौरवाच्च तदानतः} %||2-17-13||

\twolineshloka
{देवि नूनं न जानीषे महद्भयमुपस्थितम्}
{इदं तव च दुःखाय वैदेह्या लक्ष्मणस्य च} %||2-17-14||

\twolineshloka
{चतुर्दश हि वर्षाणि वत्स्यामि विजने वने}
{मधुमूलफलैर्जीवन्हित्वा मुनिवदामिषम्} %||2-17-15||

\twolineshloka
{भरताय महाराजो यौवराज्यं प्रयच्छति}
{मां पुनर्दण्डकारण्यं विवासयति तापसम्} %||2-17-16||

\twolineshloka
{तामदुःखोचितां दृष्ट्वा पतितां कदलीमिव}
{रामस्तूत्थापयामास मातरं गतचेतसम्} %||2-17-17||

\twolineshloka
{उपावृत्योत्थितां दीनां वडवामिव वाहिताम्}
{पांशुगुण्ठितसर्वाग्नीं विममर्श च पाणिना} %||2-17-18||

\twolineshloka
{सा राघवमुपासीनमसुखार्ता सुखोचिता}
{उवाच पुरुषव्याघ्रमुपशृण्वति लक्ष्मणे} %||2-17-19||

\twolineshloka
{यदि पुत्र न जायेथा मम शोकाय राघव}
{न स्म दुःखमतो भूयः पश्येयमहमप्रजा} %||2-17-20||

\twolineshloka
{एक एव हि वन्ध्यायाः शोको भवति मानवः}
{अप्रजास्मीति सन्तापो न ह्यन्यः पुत्र विद्यते} %||2-17-21||

\twolineshloka
{न दृष्टपूर्वं कल्याणं सुखं वा पतिपौरुषे}
{अपि पुत्रे विपश्येयमिति रामास्थितं मया} %||2-17-22||

\threelineshloka
{सा बहून्यमनोज्ञानि वाक्यानि हृदयच्छिदाम्}
{अहं श्रोष्ये सपत्नीनामवराणां वरा सती}
{अतो दुःखतरं किं नु प्रमदानां भविष्यति} %||2-17-23||

\twolineshloka
{त्वयि संनिहितेऽप्येवमहमासं निराकृता}
{किं पुनः प्रोषिते तात ध्रुवं मरणमेव मे} %||2-17-24||

\twolineshloka
{यो हि मां सेवते कश्चिदथ वाप्यनुवर्तते}
{कैकेय्याः पुत्रमन्वीक्ष्य स जनो नाभिभाषते} %||2-17-25||

\twolineshloka
{दश सप्त च वर्षाणि तव जातस्य राघव}
{अतीतानि प्रकाङ्क्षन्त्या मया दुःखपरिक्षयम्} %||2-17-26||

\twolineshloka
{उपवासैश्च योगैश्च बहुभिश्च परिश्रमैः}
{दुःखं संवर्धितो मोघं त्वं हि दुर्गतया मया} %||2-17-27||

\twolineshloka
{स्थिरं तु हृदयं मन्ये ममेदं यन्न दीर्यते}
{प्रावृषीव महानद्याः स्पृष्टं कूलं नवाम्भसा} %||2-17-28||

\fourlineindentedshloka
{ममैव नूनं मरणं न विद्यते}
{ न चावकाशोऽस्ति यमक्षये मम}
{यदन्तकोऽद्यैव न मां जिहीर्षति}
{ प्रसह्य सिंहो रुदतीं मृगीमिव} %||2-17-29||

\fourlineindentedshloka
{स्थिरं हि नूनं हृदयं ममायसम्}
{ न भिद्यते यद्भुवि नावदीर्यते}
{अनेन दुःखेन च देहमर्पितम्}
{ ध्रुवं ह्यकाले मरणं न विद्यते} %||2-17-30||

\fourlineindentedshloka
{इदं तु दुःखं यदनर्थकानि मे}
{ व्रतानि दानानि च संयमाश्च हि}
{तपश्च तप्तं यदपत्यकारणा}
{त्सुनिष्फलं बीजमिवोप्तमूषरे} %||2-17-31||

\fourlineindentedshloka
{यदि ह्यकाले मरणं स्वयेच्छया}
{ लभेत कश्चिद्गुरु दुःख कर्शितः}
{गताहमद्यैव परेत संसदम्}
{ विना त्वया धेनुरिवात्मजेन वै} %||2-17-32||

\fourlineindentedshloka
{भृशमसुखममर्षिता तदा}
{ बहु विललाप समीक्ष्य राघवम्}
{व्यसनमुपनिशाम्य सा मह}
{त्सुतमिव बद्धमवेक्ष्य किंनरी} %||2-18-33||


{॥इत्यार्षे श्रीमद्रामायणे वाल्मीकीये आदिकाव्ये अयोध्याकाण्डे सप्तदशः सर्गः॥}

\sect{अष्टादशः सर्गः}

\twolineshloka
{तथा तु विलपन्तीं तां कौसल्यां राममातरम्}
{उवाच लक्ष्मणो दीनस्तत्कालसदृशं वचः} %||2-18-1||

\twolineshloka
{न रोचते ममाप्येतदार्ये यद्राघवो वनम्}
{त्यक्त्वा राज्यश्रियं गच्छेत्स्त्रिया वाक्यवशं गतः} %||2-18-2||

\twolineshloka
{विपरीतश्च वृद्धश्च विषयैश्च प्रधर्षितः}
{नृपः किमिव न ब्रूयाच्चोद्यमानः समन्मथः} %||2-18-3||

\twolineshloka
{नास्यापराधं पश्यामि नापि दोषं तथा विधम्}
{येन निर्वास्यते राष्ट्राद्वनवासाय राघवः} %||2-18-4||

\twolineshloka
{न तं पश्याम्यहं लोके परोक्षमपि यो नरः}
{अमित्रोऽपि निरस्तोऽपि योऽस्य दोषमुदाहरेत्} %||2-18-5||

\twolineshloka
{देवकल्पमृजुं दान्तं रिपूणामपि वत्सलम्}
{अवेक्षमाणः को धर्मं त्यजेत्पुत्रमकारणात्} %||2-18-6||

\twolineshloka
{तदिदं वचनं राज्ञः पुनर्बाल्यमुपेयुषः}
{पुत्रः को हृदये कुर्याद्राजवृत्तमनुस्मरन्} %||2-18-7||

\twolineshloka
{यावदेव न जानाति कश्चिदर्थमिमं नरः}
{तावदेव मया सार्धमात्मस्थं कुरु शासनम्} %||2-18-8||

\twolineshloka
{मया पार्श्वे सधनुषा तव गुप्तस्य राघव}
{कः समर्थोऽधिकं कर्तुं कृतान्तस्येव तिष्ठतः} %||2-18-9||

\twolineshloka
{निर्मनुष्यामिमां सर्वामयोध्यां मनुजर्षभ}
{करिष्यामि शरैस्तीक्ष्णैर्यदि स्थास्यति विप्रिये} %||2-18-10||

\twolineshloka
{भरतस्याथ पक्ष्यो वा यो वास्य हितमिच्छति}
{सर्वानेतान्वधिष्यामि मृदुर्हि परिभूयते} %||2-18-11||

\twolineshloka
{त्वया चैव मया चैव कृत्वा वैरमनुत्तमम्}
{कस्य शक्तिः श्रियं दातुं भरतायारिशासन} %||2-18-12||

\twolineshloka
{अनुरक्तोऽस्मि भावेन भ्रातरं देवि तत्त्वतः}
{सत्येन धनुषा चैव दत्तेनेष्टेन ते शपे} %||2-18-13||

\twolineshloka
{दीप्तमग्निमरण्यं वा यदि रामः प्रवेक्ष्यते}
{प्रविष्टं तत्र मां देवि त्वं पूर्वमवधारय} %||2-18-14||

\twolineshloka
{हरामि वीर्याद्दुःखं ते तमः सूर्य इवोदितः}
{देवी पश्यतु मे वीर्यं राघवश्चैव पश्यतु} %||2-18-15||

\twolineshloka
{एतत्तु वचनं श्रुत्वा लक्ष्मणस्य महात्मनः}
{उवाच रामं कौसल्या रुदन्ती शोकलालसा} %||2-18-16||

\twolineshloka
{भ्रातुस्ते वदतः पुत्र लक्ष्मणस्य श्रुतं त्वया}
{यदत्रानन्तरं तत्त्वं कुरुष्व यदि रोचते} %||2-18-17||

\twolineshloka
{न चाधर्म्यं वचः श्रुत्वा सपत्न्या मम भाषितम्}
{विहाय शोकसन्तप्तां गन्तुमर्हसि मामितः} %||2-18-18||

\twolineshloka
{धर्मज्ञ यदि धर्मिष्ठो धर्मं चरितुमिच्छसि}
{शुश्रूष मामिहस्थस्त्वं चर धर्ममनुत्तमम्} %||2-18-19||

\twolineshloka
{शुश्रूषुर्जननीं पुत्र स्वगृहे नियतो वसन्}
{परेण तपसा युक्तः काश्यपस्त्रिदिवं गतः} %||2-18-20||

\twolineshloka
{यथैव राजा पूज्यस्ते गौरवेण तथा ह्यहम्}
{त्वां नाहमनुजानामि न गन्तव्यमितो वनम्} %||2-18-21||

\twolineshloka
{त्वद्वियोगान्न मे कार्यं जीवितेन सुखेन वा}
{त्वया सह मम श्रेयस्तृणानामपि भक्षणम्} %||2-18-22||

\twolineshloka
{यदि त्वं यास्यसि वनं त्यक्त्वा मां शोकलालसाम्}
{अहं प्रायमिहासिष्ये न हि शक्ष्यामि जीवितुम्} %||2-18-23||

\twolineshloka
{ततस्त्वं प्राप्स्यसे पुत्र निरयं लोकविश्रुतम्}
{ब्रह्महत्यामिवाधर्मात्समुद्रः सरितां पतिः} %||2-18-24||

\twolineshloka
{विलपन्तीं तथा दीनां कौसल्यां जननीं ततः}
{उवाच रामो धर्मात्मा वचनं धर्मसंहितम्} %||2-18-25||

\twolineshloka
{नास्ति शक्तिः पितुर्वाक्यं समतिक्रमितुं मम}
{प्रसादये त्वां शिरसा गन्तुमिच्छाम्यहं वनम्} %||2-18-26||

\twolineshloka
{ऋषिणा च पितुर्वाक्यं कुर्वता व्रतचारिणा}
{गौर्हता जानता धर्मं कण्डुनापि विपश्चिता} %||2-18-27||

\twolineshloka
{अस्माकं च कुले पूर्वं सगरस्याज्ञया पितुः}
{खनद्भिः सागरैर्भूतिमवाप्तः सुमहान्वधः} %||2-18-28||

\twolineshloka
{जामदग्न्येन रामेण रेणुका जननी स्वयम्}
{कृत्ता परशुनारण्ये पितुर्वचनकारिणा} %||2-18-29||

\twolineshloka
{न खल्वेतन्मयैकेन क्रियते पितृशासनम्}
{पूर्वैरयमभिप्रेतो गतो मार्गोऽनुगम्यते} %||2-18-30||

\twolineshloka
{तदेतत्तु मया कार्यं क्रियते भुवि नान्यथा}
{पितुर्हि वचनं कुर्वन्न कश्चिन्नाम हीयते} %||2-18-31||

\threelineshloka
{तामेवमुक्त्वा जननीं लक्ष्मणं पुनरब्रवीत्}
{तव लक्ष्मण जानामि मयि स्नेहमनुत्तमम्}
{अभिप्रायमविज्ञाय सत्यस्य च शमस्य च} %||2-18-32||

\twolineshloka
{धर्मो हि परमो लोके धर्मे सत्यं प्रतिष्ठितम्}
{धर्मसंश्रितमेतच्च पितुर्वचनमुत्तमम्} %||2-18-33||

\twolineshloka
{संश्रुत्य च पितुर्वाक्यं मातुर्वा ब्राह्मणस्य वा}
{न कर्तव्यं वृथा वीर धर्ममाश्रित्य तिष्ठता} %||2-18-34||

\twolineshloka
{सोऽहं न शक्ष्यामि पितुर्नियोगमतिवर्तितुम्}
{पितुर्हि वचनाद्वीर कैकेय्याहं प्रचोदितः} %||2-18-35||

\twolineshloka
{तदेनां विसृजानार्यां क्षत्रधर्माश्रितां मतिम्}
{धर्ममाश्रय मा तैक्ष्ण्यं मद्बुद्धिरनुगम्यताम्} %||2-18-36||

\twolineshloka
{तमेवमुक्त्वा सौहार्दाद्भ्रातरं लक्ष्मणाग्रजः}
{उवाच भूयः कौसल्यां प्राञ्जलिः शिरसानतः} %||2-18-37||

\threelineshloka
{अनुमन्यस्व मां देवि गमिष्यन्तमितो वनम्}
{शापितासि मम प्राणैः कुरु स्वस्त्ययनानि मे}
{तीर्णप्रतिज्ञश्च वनात्पुनरेष्याम्यहं पुरीम्} %||2-18-38||

\fourlineindentedshloka
{यशो ह्यहं केवलराज्यकारणा}
{न्न पृष्ठतः कर्तुमलं महोदयम्}
{अदीर्घकाले न तु देवि जीविते}
{ वृणेऽवरामद्य महीमधर्मतः} %||2-18-39||

\fourlineindentedshloka
{प्रसादयन्नरवृषभः स मातरम्}
{ पराक्रमाज्जिगमिषुरेव दण्डकान्}
{अथानुजं भृशमनुशास्य दर्शनम्}
{ चकार तां हृदि जननीं प्रदक्षिणम्} %||2-19-40||


{॥इत्यार्षे श्रीमद्रामायणे वाल्मीकीये आदिकाव्ये अयोध्याकाण्डे अष्टादशः सर्गः॥}

\sect{एकोनविंशः सर्गः}

\twolineshloka
{अथ तं व्यथया दीनं सविशेषममर्षितम्}
{श्वसन्तमिव नागेन्द्रं रोषविस्फारितेक्षणम्} %||2-19-1||

\twolineshloka
{आसाद्य रामः सौमित्रिं सुहृदं भ्रातरं प्रियम्}
{उवाचेदं स धैर्येण धारयन्सत्त्वमात्मवान्} %||2-19-2||

\twolineshloka
{सौमित्रे योऽभिषेकार्थे मम सम्भारसम्भ्रमः}
{अभिषेकनिवृत्त्यर्थे सोऽस्तु सम्भारसम्भ्रमः} %||2-19-3||

\twolineshloka
{यस्या मदभिषेकार्थं मानसं परितप्यते}
{माता नः सा यथा न स्यात्सविशङ्का तथा कुरु} %||2-19-4||

\twolineshloka
{तस्याः शङ्कामयं दुःखं मुहूर्तमपि नोत्सहे}
{मनसि प्रतिसंजातं सौमित्रेऽहमुपेक्षितुम्} %||2-19-5||

\twolineshloka
{न बुद्धिपूर्वं नाबुद्धं स्मरामीह कदाचन}
{मातॄणां वा पितुर्वाहं कृतमल्पं च विप्रियम्} %||2-19-6||

\twolineshloka
{सत्यः सत्याभिसन्धश्च नित्यं सत्यपराक्रमः}
{परलोकभयाद्भीतो निर्भयोऽस्तु पिता मम} %||2-19-7||

\twolineshloka
{तस्यापि हि भवेदस्मिन्कर्मण्यप्रतिसंहृते}
{सत्यं नेति मनस्तापस्तस्य तापस्तपेच्च माम्} %||2-19-8||

\twolineshloka
{अभिषेकविधानं तु तस्मात्संहृत्य लक्ष्मण}
{अन्वगेवाहमिच्छामि वनं गन्तुमितः पुनः} %||2-19-9||

\twolineshloka
{मम प्रव्राजनादद्य कृतकृत्या नृपात्मजा}
{सुतं भरतमव्यग्रमभिषेचयिता ततः} %||2-19-10||

\twolineshloka
{मयि चीराजिनधरे जटामण्डलधारिणि}
{गतेऽरण्यं च कैकेय्या भविष्यति मनःसुखम्} %||2-19-11||

\twolineshloka
{बुद्धिः प्रणीता येनेयं मनश्च सुसमाहितम्}
{तत्तु नार्हामि सङ्क्लेष्टुं प्रव्रजिष्यामि माचिरम्} %||2-19-12||

\twolineshloka
{कृतान्तस्त्वेव सौमित्रे द्रष्टव्यो मत्प्रवासने}
{राज्यस्य च वितीर्णस्य पुनरेव निवर्तने} %||2-19-13||

\twolineshloka
{कैकेय्याः प्रतिपत्तिर्हि कथं स्यान्मम पीडने}
{यदि भावो न दैवोऽयं कृतान्तविहितो भवेत्} %||2-19-14||

\twolineshloka
{जानासि हि यथा सौम्य न मातृषु ममान्तरम्}
{भूतपूर्वं विशेषो वा तस्या मयि सुतेऽपि वा} %||2-19-15||

\twolineshloka
{सोऽभिषेकनिवृत्त्यर्थैः प्रवासार्थैश्च दुर्वचैः}
{उग्रैर्वाक्यैरहं तस्या नान्यद्दैवात्समर्थये} %||2-19-16||

\twolineshloka
{कथं प्रकृतिसम्पन्ना राजपुत्री तथागुणा}
{ब्रूयात्सा प्राकृतेव स्त्री मत्पीडां भर्तृसंनिधौ} %||2-19-17||

\twolineshloka
{यदचिन्त्यं तु तद्दैवं भूतेष्वपि न हन्यते}
{व्यक्तं मयि च तस्यां च पतितो हि विपर्ययः} %||2-19-18||

\twolineshloka
{कश्चिद्दैवेन सौमित्रे योद्धुमुत्सहते पुमान्}
{यस्य न ग्रहणं किञ्चित्कर्मणोऽन्यत्र दृश्यते} %||2-19-19||

\twolineshloka
{सुखदुःखे भयक्रोधौ लाभालाभौ भवाभवौ}
{यस्य किञ्चित्तथाभूतं ननु दैवस्य कर्म तत्} %||2-19-20||

\threelineshloka
{व्याहतेऽप्यभिषेके मे परितापो न विद्यते}
{तस्मादपरितापः संस्त्वमप्यनुविधाय माम्}
{प्रतिसंहारय क्षिप्रमाभिषेचनिकीं क्रियाम्} %||2-19-21||

\fourlineindentedshloka
{न लक्ष्मणास्मिन्मम राज्यविघ्ने}
{ माता यवीयस्यतिशङ्कनीया}
{दैवाभिपन्ना हि वदन्त्यनिष्टम्}
{ जानासि दैवं च तथाप्रभावम्} %||2-20-22||


{॥इत्यार्षे श्रीमद्रामायणे वाल्मीकीये आदिकाव्ये अयोध्याकाण्डे एकोनविंशः सर्गः॥}

\sect{विंशः सर्गः}

\twolineshloka
{इति ब्रुवति रामे तु लक्ष्मणोऽधःशिरा मुहुः}
{श्रुत्वा मध्यं जगामेव मनसा दुःखहर्षयोः} %||2-20-1||

\twolineshloka
{तदा तु बद्ध्वा भ्रुकुटीं भ्रुवोर्मध्ये नरर्षभ}
{निशश्वास महासर्पो बिलस्थ इव रोषितः} %||2-20-2||

\twolineshloka
{तस्य दुष्प्रतिवीक्ष्यं तद्भ्रुकुटीसहितं तदा}
{बभौ क्रुद्धस्य सिंहस्य मुखस्य सदृशं मुखम्} %||2-20-3||

\twolineshloka
{अग्रहस्तं विधुन्वंस्तु हस्ती हस्तमिवात्मनः}
{तिर्यगूर्ध्वं शरीरे च पातयित्वा शिरोधराम्} %||2-20-4||

\twolineshloka
{अग्राक्ष्णा वीक्षमाणस्तु तिर्यग्भ्रातरमब्रवीत्}
{अस्थाने सम्भ्रमो यस्य जातो वै सुमहानयम्} %||2-20-5||

\twolineshloka
{धर्मदोषप्रसङ्गेन लोकस्यानतिशङ्कया}
{कथं ह्येतदसम्भ्रान्तस्त्वद्विधो वक्तुमर्हति} %||2-20-6||

\twolineshloka
{यथा दैवमशौण्डीरं शौण्डीरः क्षत्रियर्षभः}
{किं नाम कृपणं दैवमशक्तमभिशंससि} %||2-20-7||

\twolineshloka
{पापयोस्ते कथं नाम तयोः शङ्का न विद्यते}
{सन्ति धर्मोपधाः श्लक्ष्णा धर्मात्मन्किं न बुध्यसे} %||2-20-8||

\threelineshloka
{लोकविद्विष्टमारब्धं त्वदन्यस्याभिषेचनम्}
{येनेयमागता द्वैधं तव बुद्धिर्महीपते}
{स हि धर्मो मम द्वेष्यः प्रसङ्गाद्यस्य मुह्यसि} %||2-20-9||

\twolineshloka
{यद्यपि प्रतिपत्तिस्ते दैवी चापि तयोर्मतम्}
{तथाप्युपेक्षणीयं ते न मे तदपि रोचते} %||2-20-10||

\twolineshloka
{विक्लवो वीर्यहीनो यः स दैवमनुवर्तते}
{वीराः सम्भावितात्मानो न दैवं पर्युपासते} %||2-20-11||

\twolineshloka
{दैवं पुरुषकारेण यः समर्थः प्रबाधितुम्}
{न दैवेन विपन्नार्थः पुरुषः सोऽवसीदति} %||2-20-12||

\twolineshloka
{द्रक्ष्यन्ति त्वद्य दैवस्य पौरुषं पुरुषस्य च}
{दैवमानुषयोरद्य व्यक्ता व्यक्तिर्भविष्यति} %||2-20-13||

\twolineshloka
{अद्य मत्पौरुषहतं दैवं द्रक्ष्यन्ति वै जनाः}
{यद्दैवादाहतं तेऽद्य दृष्टं राज्याभिषेचनम्} %||2-20-14||

\twolineshloka
{अत्यङ्कुशमिवोद्दामं गजं मदबलोद्धतम्}
{प्रधावितमहं दैवं पौरुषेण निवर्तये} %||2-20-15||

\twolineshloka
{लोकपालाः समस्तास्ते नाद्य रामाभिषेचनम्}
{न च कृत्स्नास्त्रयो लोका विहन्युः किं पुनः पिता} %||2-20-16||

\twolineshloka
{यैर्विवासस्तवारण्ये मिथो राजन्समर्थितः}
{अरण्ये ते विवत्स्यन्ति चतुर्दश समास्तथा} %||2-20-17||

\twolineshloka
{अहं तदाशां छेत्स्यामि पितुस्तस्याश्च या तव}
{अभिषेकविघातेन पुत्रराज्याय वर्तते} %||2-20-18||

\twolineshloka
{मद्बलेन विरुद्धाय न स्याद्दैवबलं तथा}
{प्रभविष्यति दुःखाय यथोग्रं पौरुषं मम} %||2-20-19||

\twolineshloka
{ऊर्ध्वं वर्षसहस्रान्ते प्रजापाल्यमनन्तरम्}
{आर्यपुत्राः करिष्यन्ति वनवासं गते त्वयि} %||2-20-20||

\twolineshloka
{पूर्वराजर्षिवृत्त्या हि वनवासो विधीयते}
{प्रजा निक्षिप्य पुत्रेषु पुत्रवत्परिपालने} %||2-20-21||

\twolineshloka
{स चेद्राजन्यनेकाग्रे राज्यविभ्रमशङ्कया}
{नैवमिच्छसि धर्मात्मन्राज्यं राम त्वमात्मनि} %||2-20-22||

\twolineshloka
{प्रतिजाने च ते वीर मा भूवं वीरलोकभाक्}
{राज्यं च तव रक्षेयमहं वेलेव सागरम्} %||2-20-23||

\twolineshloka
{मङ्गलैरभिषिञ्चस्व तत्र त्वं व्यापृतो भव}
{अहमेको महीपालानलं वारयितुं बलात्} %||2-20-24||

\twolineshloka
{न शोभार्थाविमौ बाहू न धनुर्भूषणाय मे}
{नासिराबन्धनार्थाय न शराः स्तम्भहेतवः} %||2-20-25||

\twolineshloka
{अमित्रदमनार्थं मे सर्वमेतच्चतुष्टयम्}
{न चाहं कामयेऽत्यर्थं यः स्याच्छत्रुर्मतो मम} %||2-20-26||

\twolineshloka
{असिना तीक्ष्णधारेण विद्युच्चलितवर्चसा}
{प्रगृहीतेन वै शत्रुं वज्रिणं वा न कल्पये} %||2-20-27||

\twolineshloka
{खड्गनिष्पेषनिष्पिष्टैर्गहना दुश्चरा च मे}
{हस्त्यश्वनरहस्तोरुशिरोभिर्भविता मही} %||2-20-28||

\twolineshloka
{खड्गधाराहता मेऽद्य दीप्यमाना इवाद्रयः}
{पतिष्यन्ति द्विपा भूमौ मेघा इव सविद्युतः} %||2-20-29||

\twolineshloka
{बद्धगोधाङ्गुलित्राणे प्रगृहीतशरासने}
{कथं पुरुषमानी स्यात्पुरुषाणां मयि स्थिते} %||2-20-30||

\twolineshloka
{बहुभिश्चैकमत्यस्यन्नेकेन च बहूञ्जनान्}
{विनियोक्ष्याम्यहं बाणान्नृवाजिगजमर्मसु} %||2-20-31||

\twolineshloka
{अद्य मेऽस्त्रप्रभावस्य प्रभावः प्रभविष्यति}
{राज्ञश्चाप्रभुतां कर्तुं प्रभुत्वं च तव प्रभो} %||2-20-32||

\twolineshloka
{अद्य चन्दनसारस्य केयूरामोक्षणस्य च}
{वसूनां च विमोक्षस्य सुहृदां पालनस्य च} %||2-20-33||

\twolineshloka
{अनुरूपाविमौ बाहू राम कर्म करिष्यतः}
{अभिषेचनविघ्नस्य कर्तॄणां ते निवारणे} %||2-20-34||

\fourlineindentedshloka
{ब्रवीहि कोऽद्यैव मया वियुज्यताम्}
{ तवासुहृत्प्राणयशः सुहृज्जनैः}
{यथा तवेयं वसुधा वशे भवे}
{त्तथैव मां शाधि तवास्मि किङ्करः} %||2-20-35||

\fourlineindentedshloka
{विमृज्य बाष्पं परिसान्त्व्य चासकृ}
{त्स लक्ष्मणं राघववंशवर्धनः}
{उवाच पित्र्ये वचने व्यवस्थितम्}
{ निबोध मामेष हि सौम्य सत्पथः} %||2-21-36||


{॥इत्यार्षे श्रीमद्रामायणे वाल्मीकीये आदिकाव्ये अयोध्याकाण्डे विंशः सर्गः॥}

\sect{एकविंशः सर्गः}

\twolineshloka
{तं समीक्ष्य त्ववहितं पितुर्निर्देशपालने}
{कौसल्या बाष्पसंरुद्धा वचो धर्मिष्ठमब्रवीत्} %||2-21-1||

\twolineshloka
{अदृष्टदुःखो धर्मात्मा सर्वभूतप्रियंवदः}
{मयि जातो दशरथात्कथमुञ्छेन वर्तयेत्} %||2-21-2||

\twolineshloka
{यस्य भृत्याश्च दासाश्च मृष्टान्यन्नानि भुञ्जते}
{कथं स भोक्ष्यते नाथो वने मूलफलान्ययम्} %||2-21-3||

\twolineshloka
{क एतच्छ्रद्दधेच्छ्रुत्वा कस्य वा न भवेद्भयम्}
{गुणवान्दयितो राज्ञो राघवो यद्विवास्यते} %||2-21-4||

\twolineshloka
{त्वया विहीनामिह मां शोकाग्निरतुलो महान्}
{प्रधक्ष्यति यथा कक्षं चित्रभानुर्हिमात्यये} %||2-21-5||

\twolineshloka
{कथं हि धेनुः स्वं वत्सं गच्छन्तं नानुगच्छति}
{अहं त्वानुगमिष्यामि यत्र पुत्र गमिष्यसि} %||2-21-6||

\twolineshloka
{तथा निगदितं मात्रा तद्वाक्यं पुरुषर्षभः}
{श्रुत्वा रामोऽब्रवीद्वाक्यं मातरं भृशदुःखिताम्} %||2-21-7||

\twolineshloka
{कैकेय्या वञ्चितो राजा मयि चारण्यमाश्रिते}
{भवत्या च परित्यक्तो न नूनं वर्तयिष्यति} %||2-21-8||

\twolineshloka
{भर्तुः किल परित्यागो नृशंसः केवलं स्त्रियाः}
{स भवत्या न कर्तव्यो मनसापि विगर्हितः} %||2-21-9||

\twolineshloka
{यावज्जीवति काकुत्स्थः पिता मे जगतीपतिः}
{शुश्रूषा क्रियतां तावत्स हि धर्मः सनातनः} %||2-21-10||

\twolineshloka
{एवमुक्ता तु रामेण कौसल्या शुभ दर्शना}
{तथेत्युवाच सुप्रीता राममक्लिष्टकारिणम्} %||2-21-11||

\twolineshloka
{एवमुक्तस्तु वचनं रामो धर्मभृतां वरः}
{भूयस्तामब्रवीद्वाक्यं मातरं भृशदुःखिताम्} %||2-21-12||

\twolineshloka
{मया चैव भवत्या च कर्तव्यं वचनं पितुः}
{राजा भर्ता गुरुः श्रेष्ठः सर्वेषामीश्वरः प्रभुः} %||2-21-13||

\twolineshloka
{इमानि तु महारण्ये विहृत्य नव पञ्च च}
{वर्षाणि परमप्रीतः स्थास्यामि वचने तव} %||2-21-14||

\twolineshloka
{एवमुक्ता प्रियं पुत्रं बाष्पपूर्णानना तदा}
{उवाच परमार्ता तु कौसल्या पुत्रवत्सला} %||2-21-15||

\threelineshloka
{आसां राम सपत्नीनां वस्तुं मध्ये न मे क्षमम्}
{नय मामपि काकुत्स्थ वनं वन्यं मृगीं यथा}
{यदि ते गमने बुद्धिः कृता पितुरपेक्षया} %||2-21-16||

\threelineshloka
{तां तथा रुदतीं रामो रुदन्वचनमब्रवीत्}
{जीवन्त्या हि स्त्रिया भर्ता दैवतं प्रभुरेव च}
{भवत्या मम चैवाद्य राजा प्रभवति प्रभुः} %||2-21-17||

\twolineshloka
{भरतश्चापि धर्मात्मा सर्वभूतप्रियंवदः}
{भवतीमनुवर्तेत स हि धर्मरतः सदा} %||2-21-18||

\twolineshloka
{यथा मयि तु निष्क्रान्ते पुत्रशोकेन पार्थिवः}
{श्रमं नावाप्नुयात्किञ्चिदप्रमत्ता तथा कुरु} %||2-21-19||

\twolineshloka
{व्रतोपवासनिरता या नारी परमोत्तमा}
{भर्तारं नानुवर्तेत सा च पापगतिर्भवेत्} %||2-21-20||

\twolineshloka
{शुश्रूषमेव कुर्वीत भर्तुः प्रियहिते रता}
{एष धर्मः पुरा दृष्टो लोके वेदे श्रुतः स्मृतः} %||2-21-21||

\twolineshloka
{पूज्यास्ते मत्कृते देवि ब्राह्मणाश्चैव सुव्रताः}
{एवं कालं प्रतीक्षस्व ममागमनकाङ्क्षिणी} %||2-21-22||

\twolineshloka
{प्राप्स्यसे परमं कामं मयि प्रत्यागते सति}
{यदि धर्मभृतां श्रेष्ठो धारयिष्यति जीवितम्} %||2-21-23||

\threelineshloka
{एवमुक्ता तु रामेण बाष्पपर्याकुलेक्षणा}
{कौसल्या पुत्रशोकार्ता रामं वचनमब्रवीत्}
{गच्छ पुत्र त्वमेकाग्रो भद्रं तेऽस्तु सदा विभो} %||2-21-24||

\fourlineindentedshloka
{तथा हि रामं वनवासनिश्चितम्}
{ समीक्ष्य देवी परमेण चेतसा}
{उवाच रामं शुभलक्षणं वचो}
{ बभूव च स्वस्त्ययनाभिकाङ्क्षिणी} %||2-22-25||


{॥इत्यार्षे श्रीमद्रामायणे वाल्मीकीये आदिकाव्ये अयोध्याकाण्डे एकविंशः सर्गः॥}

\sect{द्वाविंशः सर्गः}

\twolineshloka
{सापनीय तमायासमुपस्पृश्य जलं शुचि}
{चकार माता रामस्य मङ्गलानि मनस्विनी} %||2-22-1||

\twolineshloka
{स्वस्ति साध्याश्च विश्वे च मरुतश्च महर्षयः}
{स्वस्ति धाता विधाता च स्वस्ति पूषा भगोऽर्यमा} %||2-22-2||

\twolineshloka
{ऋतवश्चैव पक्षाश्च मासाः संवत्सराः क्षपाः}
{दिनानि च मुहूर्ताश्च स्वस्ति कुर्वन्तु ते सदा} %||2-22-3||

\twolineshloka
{स्मृतिर्धृतिश्च धर्मश्च पान्तु त्वां पुत्र सर्वतः}
{स्कन्दश्च भगवान्देवः सोमश्च सबृहस्पतिः} %||2-22-4||

\threelineshloka
{सप्तर्षयो नारदश्च ते त्वां रक्षन्तु सर्वतः}
{नक्षत्राणि च सर्वाणि ग्रहाश्च सहदेवताः}
{महावनानि चरतो मुनिवेषस्य धीमतः} %||2-22-5||

\twolineshloka
{प्लवगा वृश्चिका दंशा मशकाश्चैव कानने}
{सरीसृपाश्च कीटाश्च मा भूवन्गहने तव} %||2-22-6||

\twolineshloka
{महाद्विपाश्च सिंहाश्च व्याघ्रा ऋक्षाश्च दंष्ट्रिणः}
{महिषाः शृङ्गिणो रौद्रा न ते द्रुह्यन्तु पुत्रक} %||2-22-7||

\twolineshloka
{नृमांसभोजना रौद्रा ये चान्ये सत्त्वजातयः}
{मा च त्वां हिंसिषुः पुत्र मया सम्पूजितास्त्विह} %||2-22-8||

\twolineshloka
{आगमास्ते शिवाः सन्तु सिध्यन्तु च पराक्रमाः}
{सर्वसम्पत्तयो राम स्वस्तिमान्गच्छ पुत्रक} %||2-22-9||

\twolineshloka
{स्वस्ति तेऽस्त्वान्तरिक्षेभ्यः पार्थिवेभ्यः पुनः पुनः}
{सर्वेभ्यश्चैव देवेभ्यो ये च ते परिपन्थिनः} %||2-22-10||

\twolineshloka
{सर्वलोकप्रभुर्ब्रह्मा भूतभर्ता तथर्षयः}
{ये च शेषाः सुरास्ते त्वां रक्षन्तु वनवासिनम्} %||2-22-11||

\twolineshloka
{इति माल्यैः सुरगणान्गन्धैश्चापि यशस्विनी}
{स्तुतिभिश्चानुरूपाभिरानर्चायतलोचना} %||2-22-12||

\twolineshloka
{यन्मङ्गलं सहस्राक्षे सर्वदेवनमस्कृते}
{वृत्रनाशे समभवत्तत्ते भवतु मङ्गलम्} %||2-22-13||

\twolineshloka
{यन्मङ्गलं सुपर्णस्य विनताकल्पयत्पुरा}
{अमृतं प्रार्थयानस्य तत्ते भवतु मङ्गलम्} %||2-22-14||

\twolineshloka
{ओषधीं चापि सिद्धार्थां विशल्यकरणीं शुभाम्}
{चकार रक्षां कौसल्या मन्त्रैरभिजजाप च} %||2-22-15||

\twolineshloka
{आनम्य मूर्ध्नि चाघ्राय परिष्वज्य यशस्विनी}
{अवदत्पुत्र सिद्धार्थो गच्छ राम यथासुखम्} %||2-22-16||

\twolineshloka
{अरोगं सर्वसिद्धार्थमयोध्यां पुनरागतम्}
{पश्यामि त्वां सुखं वत्स सुस्थितं राजवेश्मनि} %||2-22-17||

\fourlineindentedshloka
{मयार्चिता देवगणाः शिवादयो}
{ महर्षयो भूतमहासुरोरगाः}
{अभिप्रयातस्य वनं चिराय ते}
{ हितानि काङ्क्षन्तु दिशश्च राघव} %||2-22-18||

\fourlineindentedshloka
{इतीव चाश्रुप्रतिपूर्णलोचना}
{ समाप्य च स्वस्त्ययनं यथाविधि}
{प्रदक्षिणं चैव चकार राघवम्}
{ पुनः पुनश्चापि निपीड्य सस्वजे} %||2-22-19||

\fourlineindentedshloka
{तथा तु देव्या स कृतप्रदक्षिणो}
{ निपीड्य मातुश्चरणौ पुनः पुनः}
{जगाम सीतानिलयं महायशाः}
{ स राघवः प्रज्वलितः स्वया श्रिया} %||2-23-20||


{॥इत्यार्षे श्रीमद्रामायणे वाल्मीकीये आदिकाव्ये अयोध्याकाण्डे द्वाविंशः सर्गः॥}

\sect{त्रयोविंशः सर्गः}

\twolineshloka
{अभिवाद्य तु कौसल्यां रामः सम्प्रस्थितो वनम्}
{कृतस्वस्त्ययनो मात्रा धर्मिष्ठे वर्त्मनि स्थितः} %||2-23-1||

\twolineshloka
{विराजयन्राजसुतो राजमार्गं नरैर्वृतम्}
{हृदयान्याममन्थेव जनस्य गुणवत्तया} %||2-23-2||

\twolineshloka
{वैदेही चापि तत्सर्वं न शुश्राव तपस्विनी}
{तदेव हृदि तस्याश्च यौवराज्याभिषेचनम्} %||2-23-3||

\twolineshloka
{देवकार्यं स्म सा कृत्वा कृतज्ञा हृष्टचेतना}
{अभिज्ञा राजधर्माणां राजपुत्रं प्रतीक्षते} %||2-23-4||

\twolineshloka
{प्रविवेशाथ रामस्तु स्ववेश्म सुविभूषितम्}
{प्रहृष्टजनसम्पूर्णं ह्रिया किञ्चिदवाङ्मुखः} %||2-23-5||

\twolineshloka
{अथ सीता समुत्पत्य वेपमाना च तं पतिम्}
{अपश्यच्छोकसन्तप्तं चिन्ताव्याकुलिलेन्द्रियम्} %||2-23-6||

\twolineshloka
{विवर्णवदनं दृष्ट्वा तं प्रस्विन्नममर्षणम्}
{आह दुःखाभिसन्तप्ता किमिदानीमिदं प्रभो} %||2-23-7||

\twolineshloka
{अद्य बार्हस्पतः श्रीमान्युक्तः पुष्यो न राघव}
{प्रोच्यते ब्राह्मणैः प्राज्ञैः केन त्वमसि दुर्मनाः} %||2-23-8||

\twolineshloka
{न ते शतशलाकेन जलफेननिभेन च}
{आवृतं वदनं वल्गु छत्रेणाभिविराजते} %||2-23-9||

\twolineshloka
{व्यजनाभ्यां च मुख्याभ्यां शतपत्रनिभेक्षणम्}
{चन्द्रहंसप्रकाशाभ्यां वीज्यते न तवाननम्} %||2-23-10||

\twolineshloka
{वाग्मिनो बन्दिनश्चापि प्रहृष्टास्त्वं नरर्षभ}
{स्तुवन्तो नाद्य दृश्यन्ते मङ्गलैः सूतमागधाः} %||2-23-11||

\twolineshloka
{न ते क्षौद्रं च दधि च ब्राह्मणा वेदपारगाः}
{मूर्ध्नि मूर्धावसिक्तस्य दधति स्म विधानतः} %||2-23-12||

\twolineshloka
{न त्वां प्रकृतयः सर्वा श्रेणीमुख्याश्च भूषिताः}
{अनुव्रजितुमिच्छन्ति पौरजापपदास्तथा} %||2-23-13||

\twolineshloka
{चतुर्भिर्वेगसम्पन्नैर्हयैः काञ्चनभूषणैः}
{मुख्यः पुष्यरथो युक्तः किं न गच्छति तेऽग्रतः} %||2-23-14||

\twolineshloka
{न हस्ती चाग्रतः श्रीमांस्तव लक्षणपूजितः}
{प्रयाणे लक्ष्यते वीर कृष्णमेघगिरि प्रभः} %||2-23-15||

\twolineshloka
{न च काञ्चनचित्रं ते पश्यामि प्रियदर्शन}
{भद्रासनं पुरस्कृत्य यान्तं वीरपुरःसरम्} %||2-23-16||

\twolineshloka
{अभिषेको यदा सज्जः किमिदानीमिदं तव}
{अपूर्वो मुखवर्णश्च न प्रहर्षश्च लक्ष्यते} %||2-23-17||

\twolineshloka
{इतीव विलपन्तीं तां प्रोवाच रघुनन्दनः}
{सीते तत्रभवांस्तात प्रव्राजयति मां वनम्} %||2-23-18||

\twolineshloka
{कुले महति सम्भूते धर्मज्ञे धर्मचारिणि}
{शृणु जानकि येनेदं क्रमेणाभ्यागतं मम} %||2-23-19||

\twolineshloka
{राज्ञा सत्यप्रतिज्ञेन पित्रा दशरथेन मे}
{कैकेय्यै प्रीतमनसा पुरा दत्तौ महावरौ} %||2-23-20||

\twolineshloka
{तयाद्य मम सज्जेऽस्मिन्नभिषेके नृपोद्यते}
{प्रचोदितः स समयो धर्मेण प्रतिनिर्जितः} %||2-23-21||

\threelineshloka
{चतुर्दश हि वर्षाणि वस्तव्यं दण्डके मया}
{पित्रा मे भरतश्चापि यौवराज्ये नियोजितः}
{सोऽहं त्वामागतो द्रष्टुं प्रस्थितो विजनं वनम्} %||2-23-22||

\threelineshloka
{भरतस्य समीपे ते नाहं कथ्यः कदाचन}
{ऋद्धियुक्ता हि पुरुषा न सहन्ते परस्तवम्}
{तस्मान्न ते गुणाः कथ्या भरतस्याग्रतो मम} %||2-23-23||

\twolineshloka
{नापि त्वं तेन भर्तव्या विशेषेण कदाचन}
{अनुकूलतया शक्यं समीपे तस्य वर्तितुम्} %||2-23-24||

\twolineshloka
{अहं चापि प्रतिज्ञां तां गुरोः समनुपालयन्}
{वनमद्यैव यास्यामि स्थिरा भव मनस्विनि} %||2-23-25||

\twolineshloka
{याते च मयि कल्याणि वनं मुनिनिषेवितम्}
{व्रतोपवासरतया भवितव्यं त्वयानघे} %||2-23-26||

\twolineshloka
{काल्यमुत्थाय देवानां कृत्वा पूजां यथाविधि}
{वन्दितव्यो दशरथः पिता मम नरेश्वरः} %||2-23-27||

\twolineshloka
{माता च मम कौसल्या वृद्धा सन्तापकर्शिता}
{धर्ममेवाग्रतः कृत्वा त्वत्तः सम्मानमर्हति} %||2-23-28||

\twolineshloka
{वन्दितव्याश्च ते नित्यं याः शेषा मम मातरः}
{स्नेहप्रणयसम्भोगैः समा हि मम मातरः} %||2-23-29||

\twolineshloka
{भ्रातृपुत्रसमौ चापि द्रष्टव्यौ च विशेषतः}
{त्वया लक्ष्मणशत्रुघ्नौ प्राणैः प्रियतरौ मम} %||2-23-30||

\twolineshloka
{विप्रियं न च कर्तव्यं भरतस्य कदाचन}
{स हि राजा प्रभुश्चैव देशस्य च कुलस्य च} %||2-23-31||

\twolineshloka
{आराधिता हि शीलेन प्रयत्नैश्चोपसेविताः}
{राजानः सम्प्रसीदन्ति प्रकुप्यन्ति विपर्यये} %||2-23-32||

\twolineshloka
{औरसानपि पुत्रान्हि त्यजन्त्यहितकारिणः}
{समर्थान्सम्प्रगृह्णन्ति जनानपि नराधिपाः} %||2-23-33||

\fourlineindentedshloka
{अहं गमिष्यामि महावनं प्रिये}
{ त्वया हि वस्तव्यमिहैव भामिनि}
{यथा व्यलीकं कुरुषे न कस्य चि}
{त्तथा त्वया कार्यमिदं वचो मम} %||2-24-34||


{॥इत्यार्षे श्रीमद्रामायणे वाल्मीकीये आदिकाव्ये अयोध्याकाण्डे त्रयोविंशः सर्गः॥}

\sect{चतुर्विंशः सर्गः}

\twolineshloka
{एवमुक्ता तु वैदेही प्रियार्हा प्रियवादिनी}
{प्रणयादेव सङ्क्रुद्धा भर्तारमिदमब्रवीत्} %||2-24-1||

\twolineshloka
{आर्यपुत्र पिता माता भ्राता पुत्रस्तथा स्नुषा}
{स्वानि पुण्यानि भुञ्जानाः स्वं स्वं भाग्यमुपासते} %||2-24-2||

\twolineshloka
{भर्तुर्भाग्यं तु भार्यैका प्राप्नोति पुरुषर्षभ}
{अतश्चैवाहमादिष्टा वने वस्तव्यमित्यपि} %||2-24-3||

\twolineshloka
{न पिता नात्मजो नात्मा न माता न सखीजनः}
{इह प्रेत्य च नारीणां पतिरेको गतिः सदा} %||2-24-4||

\twolineshloka
{यदि त्वं प्रस्थितो दुर्गं वनमद्यैव राघव}
{अग्रतस्ते गमिष्यामि मृद्नन्ती कुशकण्टकान्} %||2-24-5||

\twolineshloka
{ईर्ष्या रोषौ बहिष्कृत्य भुक्तशेषमिवोदकम्}
{नय मां वीर विश्रब्धः पापं मयि न विद्यते} %||2-24-6||

\twolineshloka
{प्रासादाग्रैर्विमानैर्वा वैहायसगतेन वा}
{सर्वावस्थागता भर्तुः पादच्छाया विशिष्यते} %||2-24-7||

\twolineshloka
{अनुशिष्टास्मि मात्रा च पित्रा च विविधाश्रयम्}
{नास्मि सम्प्रति वक्तव्या वर्तितव्यं यथा मया} %||2-24-8||

\twolineshloka
{सुखं वने निवत्स्यामि यथैव भवने पितुः}
{अचिन्तयन्ती त्रीँल्लोकांश्चिन्तयन्ती पतिव्रतम्} %||2-24-9||

\twolineshloka
{शुश्रूषमाणा ते नित्यं नियता ब्रह्मचारिणी}
{सह रंस्ये त्वया वीर वनेषु मधुगन्धिषु} %||2-24-10||

\twolineshloka
{त्वं हि कर्तुं वने शक्तो राम सम्परिपालनम्}
{अन्यस्य पै जनस्येह किं पुनर्मम मानद} %||2-24-11||

\twolineshloka
{फलमूलाशना नित्यं भविष्यामि न संशयः}
{न ते दुःखं करिष्यामि निवसन्ती सह त्वया} %||2-24-12||

\twolineshloka
{इच्छामि सरितः शैलान्पल्वलानि वनानि च}
{द्रष्टुं सर्वत्र निर्भीता त्वया नाथेन धीमता} %||2-24-13||

\twolineshloka
{हंसकारण्डवाकीर्णाः पद्मिनीः साधुपुष्पिताः}
{इच्छेयं सुखिनी द्रष्टुं त्वया वीरेण सङ्गता} %||2-24-14||

\twolineshloka
{सह त्वया विशालाक्ष रंस्ये परमनन्दिनी}
{एवं वर्षसहस्राणां शतं वाहं त्वया सह} %||2-24-15||

\twolineshloka
{स्वर्गेऽपि च विना वासो भविता यदि राघव}
{त्वया मम नरव्याघ्र नाहं तमपि रोचये} %||2-24-16||

\fourlineindentedshloka
{अहं गमिष्यामि वनं सुदुर्गमम्}
{ मृगायुतं वानरवारणैर्युतम्}
{वने निवत्स्यामि यथा पितुर्गृहे}
{ तवैव पादावुपगृह्य सम्मता} %||2-24-17||

\fourlineindentedshloka
{अनन्यभावामनुरक्तचेतसम्}
{ त्वया वियुक्तां मरणाय निश्चिताम्}
{नयस्व मां साधु कुरुष्व याचनाम्}
{ न ते मयातो गुरुता भविष्यति} %||2-24-18||

\fourlineindentedshloka
{तथा ब्रुवाणामपि धर्मवत्सलो}
{ न च स्म सीतां नृवरो निनीषति}
{उवाच चैनां बहु संनिवर्तने}
{ वने निवासस्य च दुःखितां प्रति} %||2-25-19||


{॥इत्यार्षे श्रीमद्रामायणे वाल्मीकीये आदिकाव्ये अयोध्याकाण्डे चतुर्विंशः सर्गः॥}

\sect{पञ्चविंशः सर्गः}

\twolineshloka
{स एवं ब्रुवतीं सीतां धर्मज्ञो धर्मवत्सलः}
{निवर्तनार्थे धर्मात्मा वाक्यमेतदुवाच ह} %||2-25-1||

\twolineshloka
{सीते महाकुलीनासि धर्मे च निरता सदा}
{इहाचर स्वधर्मं त्वं मा यथा मनसः सुखम्} %||2-25-2||

\twolineshloka
{सीते यथा त्वां वक्ष्यामि तथा कार्यं त्वयाबले}
{वने दोषा हि बहवो वदतस्तान्निबोध मे} %||2-25-3||

\twolineshloka
{सीते विमुच्यतामेषा वनवासकृता मतिः}
{बहुदोषं हि कान्तारं वनमित्यभिधीयते} %||2-25-4||

\twolineshloka
{हितबुद्ध्या खलु वचो मयैतदभिधीयते}
{सदा सुखं न जानामि दुःखमेव सदा वनम्} %||2-25-5||

\twolineshloka
{गिरिनिर्झरसम्भूता गिरिकन्दरवासिनाम्}
{सिंहानां निनदा दुःखाः श्रोतुं दुःखमतो वनम्} %||2-25-6||

\twolineshloka
{सुप्यते पर्णशय्यासु स्वयं भग्नासु भूतले}
{रात्रिषु श्रमखिन्नेन तस्माद्दुःखतरं वनम्} %||2-25-7||

\twolineshloka
{उपवासश्च कर्तव्या यथाप्राणेन मैथिलि}
{जटाभारश्च कर्तव्यो वल्कलाम्बरधारिणा} %||2-25-8||

\twolineshloka
{अतीव वातस्तिमिरं बुभुक्षा चात्र नित्यशः}
{भयानि च महान्त्यत्र ततो दुःखतरं वनम्} %||2-25-9||

\twolineshloka
{सरीसृपाश्च बहवो बहुरूपाश्च भामिनि}
{चरन्ति पृथिवीं दर्पादतो दुखतरं वनम्} %||2-25-10||

\twolineshloka
{नदीनिलयनाः सर्पा नदीकुटिलगामिनः}
{तिष्ठन्त्यावृत्य पन्थानमतो दुःखतरं वनम्} %||2-25-11||

\twolineshloka
{पतङ्गा वृश्चिकाः कीटा दंशाश्च मशकैः सह}
{बाधन्ते नित्यमबले सर्वं दुःखमतो वनम्} %||2-25-12||

\twolineshloka
{द्रुमाः कण्टकिनश्चैव कुशकाशाश्च भामिनि}
{वने व्याकुलशाखाग्रास्तेन दुःखतरं वनम्} %||2-25-13||

\twolineshloka
{तदलं ते वनं गत्वा क्षमं न हि वनं तव}
{विमृशन्निह पश्यामि बहुदोषतरं वनम्} %||2-25-14||

\fourlineindentedshloka
{वनं तु नेतुं न कृता मतिस्तदा}
{ बभूव रामेण यदा महात्मना}
{न तस्य सीता वचनं चकार त}
{त्ततोऽब्रवीद्राममिदं सुदुःखिता} %||2-26-15||


{॥इत्यार्षे श्रीमद्रामायणे वाल्मीकीये आदिकाव्ये अयोध्याकाण्डे पञ्चविंशः सर्गः॥}

\sect{षड्विंशः सर्गः}

\twolineshloka
{एतत्तु वचनं श्रुत्वा सीता रामस्य दुःखिता}
{प्रसक्ताश्रुमुखी मन्दमिदं वचनमब्रवीत्} %||2-26-1||

\twolineshloka
{ये त्वया कीर्तिता दोषा वने वस्तव्यतां प्रति}
{गुणानित्येव तान्विद्धि तव स्नेहपुरस्कृतान्} %||2-26-2||

\twolineshloka
{त्वया च सह गन्तव्यं मया गुरुजनाज्ञया}
{त्वद्वियोगेन मे राम त्यक्तव्यमिह जीवितम्} %||2-26-3||

\twolineshloka
{न च मां त्वत्समीपस्थमपि शक्नोति राघव}
{सुराणामीश्वरः शक्रः प्रधर्षयितुमोजसा} %||2-26-4||

\twolineshloka
{पतिहीना तु या नारी न सा शक्ष्यति जीवितुम्}
{काममेवंविधं राम त्वया मम विदर्शितम्} %||2-26-5||

\twolineshloka
{अथ चापि महाप्राज्ञ ब्राह्मणानां मया श्रुतम्}
{पुरा पितृगृहे सत्यं वस्तव्यं किल मे वने} %||2-26-6||

\twolineshloka
{लक्षणिभ्यो द्विजातिभ्यः श्रुत्वाहं वचनं गृहे}
{वनवासकृतोत्साहा नित्यमेव महाबल} %||2-26-7||

\twolineshloka
{आदेशो वनवासस्य प्राप्तव्यः स मया किल}
{सा त्वया सह तत्राहं यास्यामि प्रिय नान्यथा} %||2-26-8||

\twolineshloka
{कृतादेशा भविष्यामि गमिष्यामि सह त्वया}
{कालश्चायं समुत्पन्नः सत्यवाग्भवतु द्विजः} %||2-26-9||

\twolineshloka
{वनवासे हि जानामि दुःखानि बहुधा किल}
{प्राप्यन्ते नियतं वीर पुरुषैरकृतात्मभिः} %||2-26-10||

\twolineshloka
{कन्यया च पितुर्गेहे वनवासः श्रुतो मया}
{भिक्षिण्याः साधुवृत्ताया मम मातुरिहाग्रतः} %||2-26-11||

\twolineshloka
{प्रसादितश्च वै पूर्वं त्वं वै बहुविधं प्रभो}
{गमनं वनवासस्य काङ्क्षितं हि सह त्वया} %||2-26-12||

\twolineshloka
{कृतक्षणाहं भद्रं ते गमनं प्रति राघव}
{वनवासस्य शूरस्य चर्या हि मम रोचते} %||2-26-13||

\twolineshloka
{शुद्धात्मन्प्रेमभावाद्धि भविष्यामि विकल्मषा}
{भर्तारमनुगच्छन्ती भर्ता हि मम दैवतम्} %||2-26-14||

\twolineshloka
{प्रेत्यभावेऽपि कल्याणः सङ्गमो मे सह त्वया}
{श्रुतिर्हि श्रूयते पुण्या ब्राह्मणानां यशस्विनाम्} %||2-26-15||

\twolineshloka
{इह लोके च पितृभिर्या स्त्री यस्य महामते}
{अद्भिर्दत्ता स्वधर्मेण प्रेत्यभावेऽपि तस्य सा} %||2-26-16||

\twolineshloka
{एवमस्मात्स्वकां नारीं सुवृत्तां हि पतिव्रताम्}
{नाभिरोचयसे नेतुं त्वं मां केनेह हेतुना} %||2-26-17||

\twolineshloka
{भक्तां पतिव्रतां दीनां मां समां सुखदुःखयोः}
{नेतुमर्हसि काकुत्स्थ समानसुखदुःखिनीम्} %||2-26-18||

\twolineshloka
{यदि मां दुःखितामेवं वनं नेतुं न चेच्छसि}
{विषमग्निं जलं वाहमास्थास्ये मृत्युकारणात्} %||2-26-19||

\twolineshloka
{एवं बहुविधं तं सा याचते गमनं प्रति}
{नानुमेने महाबाहुस्तां नेतुं विजनं वनम्} %||2-26-20||

\twolineshloka
{एवमुक्ता तु सा चिन्तां मैथिली समुपागता}
{स्नापयन्तीव गामुष्णैरश्रुभिर्नयनच्युतैः} %||2-26-21||

\twolineshloka
{चिन्तयन्तीं तथा तां तु निवर्तयितुमात्मवान्}
{क्रोधाविष्टां तु वैदेहीं काकुत्स्थो बह्वसान्त्वयत्} %||2-27-22||


{॥इत्यार्षे श्रीमद्रामायणे वाल्मीकीये आदिकाव्ये अयोध्याकाण्डे षड्विंशः सर्गः॥}

\sect{सप्तविंशः सर्गः}

\twolineshloka
{सान्त्व्यमाना तु रामेण मैथिली जनकात्मजा}
{वनवासनिमित्ताय भर्तारमिदमब्रवीत्} %||2-27-1||

\twolineshloka
{सा तमुत्तमसंविग्ना सीता विपुलवक्षसम्}
{प्रणयाच्चाभिमानाच्च परिचिक्षेप राघवम्} %||2-27-2||

\twolineshloka
{किं त्वामन्यत वैदेहः पिता मे मिथिलाधिपः}
{राम जामातरं प्राप्य स्त्रियं पुरुषविग्रहम्} %||2-27-3||

\twolineshloka
{अनृतं बललोकोऽयमज्ञानाद्यद्धि वक्ष्यति}
{तेजो नास्ति परं रामे तपतीव दिवाकरे} %||2-27-4||

\twolineshloka
{किं हि कृत्वा विषण्णस्त्वं कुतो वा भयमस्ति ते}
{यत्परित्यक्तुकामस्त्वं मामनन्यपरायणाम्} %||2-27-5||

\twolineshloka
{द्युमत्सेनसुतं वीर सत्यवन्तमनुव्रताम्}
{सावित्रीमिव मां विद्धि त्वमात्मवशवर्तिनीम्} %||2-27-6||

\twolineshloka
{न त्वहं मनसाप्यन्यं द्रष्टास्मि त्वदृतेऽनघ}
{त्वया राघव गच्छेयं यथान्या कुलपांसनी} %||2-27-7||

\twolineshloka
{स्वयं तु भार्यां कौमारीं चिरमध्युषितां सतीम्}
{शैलूष इव मां राम परेभ्यो दातुमिच्छसि} %||2-27-8||

\twolineshloka
{स मामनादाय वनं न त्वं प्रस्थातुमर्हसि}
{तपो वा यदि वारण्यं स्वर्गो वा स्यात्सह त्वया} %||2-27-9||

\twolineshloka
{न च मे भविता तत्र कश्चित्पथि परिश्रमः}
{पृष्ठतस्तव गच्छन्त्या विहारशयनेष्वपि} %||2-27-10||

\twolineshloka
{कुशकाशशरेषीका ये च कण्टकिनो द्रुमाः}
{तूलाजिनसमस्पर्शा मार्गे मम सह त्वया} %||2-27-11||

\twolineshloka
{महावात समुद्धूतं यन्मामवकरिष्यति}
{रजो रमण तन्मन्ये परार्ध्यमिव चन्दनम्} %||2-27-12||

\twolineshloka
{शाद्वलेषु यदासिष्ये वनान्ते वनगोरचा}
{कुथास्तरणतल्पेषु किं स्यात्सुखतरं ततः} %||2-27-13||

\twolineshloka
{पत्रं मूलं फलं यत्त्वमल्पं वा यदि वा बहु}
{दास्यसि स्वयमाहृत्य तन्मेऽमृतरसोपमम्} %||2-27-14||

\twolineshloka
{न मातुर्न पितुस्तत्र स्मरिष्यामि न वेश्मनः}
{आर्तवान्युपभुञ्जाना पुष्पाणि च फलानि च} %||2-27-15||

\twolineshloka
{न च तत्र गतः किञ्चिद्द्रष्टुमर्हसि विप्रियम्}
{मत्कृते न च ते शोको न भविष्यामि दुर्भरा} %||2-27-16||

\twolineshloka
{यस्त्वया सह स स्वर्गो निरयो यस्त्वया विना}
{इति जानन्परां प्रीतिं गच्छ राम मया सह} %||2-27-17||

\twolineshloka
{अथ मामेवमव्यग्रां वनं नैव नयिष्यसि}
{विषमद्यैव पास्यामि मा विशं द्विषतां वशम्} %||2-27-18||

\twolineshloka
{पश्चादपि हि दुःखेन मम नैवास्ति जीवितम्}
{उज्झितायास्त्वया नाथ तदैव मरणं वरम्} %||2-27-19||

\twolineshloka
{इदं हि सहितुं शोकं मुहूर्तमपि नोत्सहे}
{किं पुनर्दशवर्षाणि त्रीणि चैकं च दुःखिता} %||2-27-20||

\twolineshloka
{इति सा शोकसन्तप्ता विलप्य करुणं बहु}
{चुक्रोश पतिमायस्ता भृशमालिङ्ग्य सस्वरम्} %||2-27-21||

\twolineshloka
{सा विद्धा बहुभिर्वाक्यैर्दिग्धैरिव गजाङ्गना}
{चिर संनियतं बाष्पं मुमोचाग्निमिवारणिः} %||2-27-22||

\twolineshloka
{तस्याः स्फटिकसङ्काशं वारि सन्तापसम्भवम्}
{नेत्राभ्यां परिसुस्राव पङ्कजाभ्यामिवोदकम्} %||2-27-23||

\twolineshloka
{तां परिष्वज्य बाहुभ्यां विसंज्ञामिव दुःखिताम्}
{उवाच वचनं रामः परिविश्वासयंस्तदा} %||2-27-24||

\twolineshloka
{न देवि तव दुःखेन स्वर्गमप्यभिरोचये}
{न हि मेऽस्ति भयं किञ्चित्स्वयम्भोरिव सर्वतः} %||2-27-25||

\twolineshloka
{तव सर्वमभिप्रायमविज्ञाय शुभानने}
{वासं न रोचयेऽरण्ये शक्तिमानपि रक्षणे} %||2-27-26||

\twolineshloka
{यत्सृष्टासि मया सार्धं वनवासाय मैथिलि}
{न विहातुं मया शक्या कीर्तिरात्मवता यथा} %||2-27-27||

\twolineshloka
{धर्मस्तु गजनासोरु सद्भिराचरितः पुरा}
{तं चाहमनुवर्तेऽद्य यथा सूर्यं सुवर्चला} %||2-27-28||

\twolineshloka
{एष धर्मस्तु सुश्रोणि पितुर्मातुश्च वश्यता}
{अतश्चाज्ञां व्यतिक्रम्य नाहं जीवितुमुत्सहे} %||2-27-29||

\threelineshloka
{स मां पिता यथा शास्ति सत्यधर्मपथे स्थितः}
{तथा वर्तितुमिच्छामि स हि धर्मः सनातनः}
{अनुगच्छस्व मां भीरु सहधर्मचरी भव} %||2-27-30||

\twolineshloka
{ब्राह्मणेभ्यश्च रत्नानि भिक्षुकेभ्यश्च भोजनम्}
{देहि चाशंसमानेभ्यः सन्त्वरस्व च माचिरम्} %||2-27-31||

\twolineshloka
{अनुकूलं तु सा भर्तुर्ज्ञात्वा गमनमात्मनः}
{क्षिप्रं प्रमुदिता देवी दातुमेवोपचक्रमे} %||2-27-32||

\fourlineindentedshloka
{ततः प्रहृष्टा परिपूर्णमानसा}
{ यशस्विनी भर्तुरवेक्ष्य भाषितम्}
{धनानि रत्नानि च दातुमङ्गना}
{ प्रचक्रमे धर्मभृतां मनस्विनी} %||2-28-33||


{॥इत्यार्षे श्रीमद्रामायणे वाल्मीकीये आदिकाव्ये अयोध्याकाण्डे सप्तविंशः सर्गः॥}

\sect{अष्टाविंशः सर्गः}

\twolineshloka
{ततोऽब्रवीन्महातेजा रामो लक्ष्मणमग्रतः}
{स्थितं प्राग्गामिनं वीरं याचमानं कृताञ्जलिम्} %||2-28-1||

\twolineshloka
{मयाद्य सह सौमित्रे त्वयि गच्छति तद्वनम्}
{को भरिष्यति कौसल्यां सुमित्रां वा यशस्विनीम्} %||2-28-2||

\twolineshloka
{अभिवर्षति कामैर्यः पर्जन्यः पृथिवीमिव}
{स कामपाशपर्यस्तो महातेजा महीपतिः} %||2-28-3||

\twolineshloka
{सा हि राज्यमिदं प्राप्य नृपस्याश्वपतेः सुता}
{दुःखितानां सपत्नीनां न करिष्यति शोभनम्} %||2-28-4||

\twolineshloka
{एवमुक्तस्तु रामेण लक्ष्मणः श्लक्ष्णया गिरा}
{प्रत्युवाच तदा रामं वाक्यज्ञो वाक्यकोविदम्} %||2-28-5||

\twolineshloka
{तवैव तेजसा वीर भरतः पूजयिष्यति}
{कौसल्यां च सुमित्रां च प्रयतो नात्र संशयः} %||2-28-6||

\twolineshloka
{कौसल्या बिभृयादार्या सहस्रमपि मद्विधान्}
{यस्याः सहस्रं ग्रामाणां सम्प्राप्तमुपजीवनम्} %||2-28-7||

\twolineshloka
{धनुरादाय सशरं खनित्रपिटकाधरः}
{अग्रतस्ते गमिष्यामि पन्थानमनुदर्शयन्} %||2-28-8||

\twolineshloka
{आहरिष्यामि ते नित्यं मूलानि च फलानि च}
{वन्यानि यानि चान्यानि स्वाहाराणि तपस्विनाम्} %||2-28-9||

\twolineshloka
{भवांस्तु सह वैदेह्या गिरिसानुषु रंस्यते}
{अहं सर्वं करिष्यामि जाग्रतः स्वपतश्च ते} %||2-28-10||

\twolineshloka
{रामस्त्वनेन वाक्येन सुप्रीतः प्रत्युवाच तम्}
{व्रजापृच्छस्व सौमित्रे सर्वमेव सुहृज्जनम्} %||2-28-11||

\twolineshloka
{ये च राज्ञो ददौ दिव्ये महात्मा वरुणः स्वयम्}
{जनकस्य महायज्ञे धनुषी रौद्रदर्शने} %||2-28-12||

\twolineshloka
{अभेद्यकवचे दिव्ये तूणी चाक्षयसायकौ}
{आदित्यविमलौ चोभौ खड्गौ हेमपरिष्कृतौ} %||2-28-13||

\twolineshloka
{सत्कृत्य निहितं सर्वमेतदाचार्यसद्मनि}
{स त्वमायुधमादाय क्षिप्रमाव्रज लक्ष्मण} %||2-28-14||

\twolineshloka
{स सुहृज्जनमामन्त्र्य वनवासाय निश्चितः}
{इक्ष्वाकुगुरुमामन्त्र्य जग्राहायुधमुत्तमम्} %||2-28-15||

\twolineshloka
{तद्दिव्यं राजशार्दूलः सत्कृतं माल्यभूषितम्}
{रामाय दर्शयामास सौमित्रिः सर्वमायुधम्} %||2-28-16||

\twolineshloka
{तमुवाचात्मवान्रामः प्रीत्या लक्ष्मणमागतम्}
{काले त्वमागतः सौम्य काङ्क्षिते मम लक्ष्मण} %||2-28-17||

\twolineshloka
{अहं प्रदातुमिच्छामि यदिदं मामकं धनम्}
{ब्राह्मणेभ्यस्तपस्विभ्यस्त्वया सह परन्तप} %||2-28-18||

\twolineshloka
{वसन्तीह दृढं भक्त्या गुरुषु द्विजसत्तमाः}
{तेषामपि च मे भूयः सर्वेषां चोपजीविनाम्} %||2-28-19||

\fourlineindentedshloka
{वसिष्ठपुत्रं तु सुयज्ञमार्यम्}
{ त्वमानयाशु प्रवरं द्विजानाम्}
{अभिप्रयास्यामि वनं समस्ता}
{नभ्यर्च्य शिष्टानपरान्द्विजातीन्} %||2-29-20||


{॥इत्यार्षे श्रीमद्रामायणे वाल्मीकीये आदिकाव्ये अयोध्याकाण्डे अष्टाविंशः सर्गः॥}

\sect{एकोनत्रिंशः सर्गः}

\twolineshloka
{ततः शासनमाज्ञाय भ्रातुः शुभतरं प्रियम्}
{गत्वा स प्रविवेशाशु सुयज्ञस्य निवेशनम्} %||2-29-1||

\twolineshloka
{तं विप्रमग्न्यगारस्थं वन्दित्वा लक्ष्मणोऽब्रवीत्}
{सखेऽभ्यागच्छ पश्य त्वं वेश्म दुष्करकारिणः} %||2-29-2||

\twolineshloka
{ततः सन्ध्यामुपास्याशु गत्वा सौमित्रिणा सह}
{जुष्टं तत्प्राविशल्लक्ष्म्या रम्यं रामनिवेशनम्} %||2-29-3||

\twolineshloka
{तमागतं वेदविदं प्राञ्जलिः सीतया सह}
{सुयज्ञमभिचक्राम राघवोऽग्निमिवार्चितम्} %||2-29-4||

\twolineshloka
{जातरूपमयैर्मुख्यैरङ्गदैः कुण्डलैः शुभैः}
{सहेम सूत्रैर्मणिभिः केयूरैर्वलयैरपि} %||2-29-5||

\twolineshloka
{अन्यैश्च रत्नैर्बहुभिः काकुत्स्थः प्रत्यपूजयत्}
{सुयज्ञं स तदोवाच रामः सीताप्रचोदितः} %||2-29-6||

\twolineshloka
{हारं च हेमसूत्रं च भार्यायै सौम्य हारय}
{रशनां चाधुना सीता दातुमिच्छति ते सखे} %||2-29-7||

\twolineshloka
{पर्यङ्कमग्र्यास्तरणं नानारत्नविभूषितम्}
{तमपीच्छति वैदेही प्रतिष्ठापयितुं त्वयि} %||2-29-8||

\twolineshloka
{नागः शत्रुं जयो नाम मातुलो यं ददौ मम}
{तं ते गजसहस्रेण ददामि द्विजपुङ्गव} %||2-29-9||

\twolineshloka
{इत्युक्तः स हि रामेण सुयज्ञः प्रतिगृह्य तत्}
{रामलक्ष्मणसीतानां प्रयुयोजाशिषः शिवाः} %||2-29-10||

\twolineshloka
{अथ भ्रातरमव्यग्रं प्रियं रामः प्रियंवदः}
{सौमित्रिं तमुवाचेदं ब्रह्मेव त्रिदशेश्वरम्} %||2-29-11||

\twolineshloka
{अगस्त्यं कौशिकं चैव तावुभौ ब्राह्मणोत्तमौ}
{अर्चयाहूय सौमित्रे रत्नैः सस्यमिवाम्बुभिः} %||2-29-12||

\twolineshloka
{कौसल्यां च य आशीर्भिर्भक्तः पर्युपतिष्ठति}
{आचार्यस्तैत्तिरीयाणामभिरूपश्च वेदवित्} %||2-29-13||

\twolineshloka
{तस्य यानं च दासीश्च सौमित्रे सम्प्रदापय}
{कौशेयानि च वस्त्राणि यावत्तुष्यति स द्विजः} %||2-29-14||

\twolineshloka
{सूतश्चित्ररथश्चार्यः सचिवः सुचिरोषितः}
{तोषयैनं महार्हैश्च रत्नैर्वस्त्रैर्धनैस्तथा} %||2-29-15||

\twolineshloka
{शालिवाहसहस्रं च द्वे शते भद्रकांस्तथा}
{व्यञ्जनार्थं च सौमित्रे गोसहस्रमुपाकुरु} %||2-29-16||

\twolineshloka
{ततः स पुरुषव्याघ्रस्तद्धनं लक्ष्मणः स्वयम्}
{यथोक्तं ब्राह्मणेन्द्राणामददाद्धनदो यथा} %||2-29-17||

\twolineshloka
{अथाब्रवीद्बाष्पकलांस्तिष्ठतश्चोपजीविनः}
{सम्प्रदाय बहु द्रव्यमेकैकस्योपजीविनः} %||2-29-18||

\twolineshloka
{लक्ष्मणस्य च यद्वेश्म गृहं च यदिदं मम}
{अशून्यं कार्यमेकैकं यावदागमनं मम} %||2-29-19||

\threelineshloka
{इत्युक्त्वा दुःखितं सर्वं जनं तमुपजीविनम्}
{उवाचेदं धनध्यक्षं धनमानीयतामिति}
{ततोऽस्य धनमाजह्रुः सर्वमेवोपजीविनः} %||2-29-20||

\twolineshloka
{ततः स पुरुषव्याघ्रस्तद्धनं सहलक्ष्मणः}
{द्विजेभ्यो बालवृद्धेभ्यः कृपणेभ्योऽभ्यदापयत्} %||2-29-21||

\twolineshloka
{तत्रासीत्पिङ्गलो गार्ग्यस्त्रिजटो नाम वै द्विजः}
{आ पञ्चमायाः कक्ष्याया नैनं कश्चिदवारयत्} %||2-29-22||

\threelineshloka
{स राजपुत्रमासाद्य त्रिजटो वाक्यमब्रवीत्}
{निर्धनो बहुपुत्रोऽस्मि राजपुत्र महायशः}
{उञ्छवृत्तिर्वने नित्यं प्रत्यवेक्षस्व मामिति} %||2-29-23||

\threelineshloka
{तमुवाच ततो रामः परिहाससमन्वितम्}
{गवां सहस्रमप्येकं न तु विश्राणितं मया}
{परिक्षिपसि दण्डेन यावत्तावदवाप्स्यसि} %||2-29-24||

\twolineshloka
{स शाटीं त्वरितः कट्यां सम्भ्रान्तः परिवेष्ट्य ताम्}
{आविध्य दण्डं चिक्षेप सर्वप्राणेन वेगितः} %||2-29-25||

\twolineshloka
{उवाच च ततो रामस्तं गार्ग्यमभिसान्त्वयन्}
{मन्युर्न खलु कर्तव्यः परिहासो ह्ययं मम} %||2-29-26||

\fourlineindentedshloka
{ततः सभार्यस्त्रिजटो महामुनि}
{र्गवामनीकं प्रतिगृह्य मोदितः}
{यशोबलप्रीतिसुखोपबृंहिणी}
{स्तदाशिषः प्रत्यवदन्महात्मनः} %||2-30-27||


{॥इत्यार्षे श्रीमद्रामायणे वाल्मीकीये आदिकाव्ये अयोध्याकाण्डे एकोनत्रिंशः सर्गः॥}

\sect{त्रिंशः सर्गः}

\twolineshloka
{दत्त्वा तु सह वैदेह्या ब्राह्मणेभ्यो धनं बहु}
{जग्मतुः पितरं द्रष्टुं सीतया सह राघवौ} %||2-30-1||

\twolineshloka
{ततो गृहीते दुष्प्रेक्ष्ये अशोभेतां तदायुधे}
{मालादामभिरासक्ते सीतया समलङ्कृते} %||2-30-2||

\twolineshloka
{ततः प्रासादहर्म्याणि विमानशिखराणि च}
{अधिरुह्य जनः श्रीमानुदासीनो व्यलोकयत्} %||2-30-3||

\twolineshloka
{न हि रथ्याः स्म शक्यन्ते गन्तुं बहुजनाकुलाः}
{आरुह्य तस्मात्प्रासादान्दीनाः पश्यन्ति राघवम्} %||2-30-4||

\twolineshloka
{पदातिं वर्जितच्छत्रं रामं दृष्ट्वा तदा जनाः}
{ऊचुर्बहुविधा वाचः शोकोपहतचेतसः} %||2-30-5||

\twolineshloka
{यं यान्तमनुयाति स्म चतुरङ्गबलं महत्}
{तमेकं सीतया सार्धमनुयाति स्म लक्ष्मणः} %||2-30-6||

\twolineshloka
{ऐश्वर्यस्य रसज्ञः सन्कामिनां चैव कामदः}
{नेच्छत्येवानृतं कर्तुं पितरं धर्मगौरवात्} %||2-30-7||

\twolineshloka
{या न शक्या पुरा द्रष्टुं भूतैराकाशगैरपि}
{तामद्य सीतां पश्यन्ति राजमार्गगता जनाः} %||2-30-8||

\twolineshloka
{अङ्गरागोचितां सीतां रक्तचन्दन सेविनीम्}
{वर्षमुष्णं च शीतं च नेष्यत्याशु विवर्णताम्} %||2-30-9||

\twolineshloka
{अद्य नूनं दशरथः सत्त्वमाविश्य भाषते}
{न हि राजा प्रियं पुत्रं विवासयितुमर्हति} %||2-30-10||

\twolineshloka
{निर्गुणस्यापि पुत्रस्या काथं स्याद्विप्रवासनम्}
{किं पुनर्यस्य लोकोऽयं जितो वृत्तेन केवलम्} %||2-30-11||

\twolineshloka
{आनृशंस्यमनुक्रोशः श्रुतं शीलं दमः शमः}
{राघवं शोभयन्त्येते षड्गुणाः पुरुषोत्तमम्} %||2-30-12||

\twolineshloka
{तस्मात्तस्योपघातेन प्रजाः परमपीडिताः}
{औदकानीव सत्त्वानि ग्रीष्मे सलिलसङ्क्षयात्} %||2-30-13||

\twolineshloka
{पीडया पीडितं सर्वं जगदस्य जगत्पतेः}
{मूलस्येवोपघातेन वृक्षः पुष्पफलोपगः} %||2-30-14||

\twolineshloka
{ते लक्ष्मण इव क्षिप्रं सपत्न्यः सहबान्धवाः}
{गच्छन्तमनुगच्छामो येन गच्छति राघवः} %||2-30-15||

\twolineshloka
{उद्यानानि परित्यज्य क्षेत्राणि च गृहाणि च}
{एकदुःखसुखा राममनुगच्छाम धार्मिकम्} %||2-30-16||

\twolineshloka
{समुद्धृतनिधानानि परिध्वस्ताजिराणि च}
{उपात्तधनधान्यानि हृतसाराणि सर्वशः} %||2-30-17||

\twolineshloka
{रजसाभ्यवकीर्णानि परित्यक्तानि दैवतैः}
{अस्मत्त्यक्तानि वेश्मानि कैकेयी प्रतिपद्यताम्} %||2-30-18||

\twolineshloka
{वनं नगरमेवास्तु येन गच्छति राघवः}
{अस्माभिश्च परित्यक्तं पुरं सम्पद्यतां वनम्} %||2-30-19||

\twolineshloka
{बिलानि दंष्ट्रिणः सर्वे सानूनि मृगपक्षिणः}
{अस्मत्त्यक्तं प्रपद्यन्तां सेव्यमानं त्यजन्तु च} %||2-30-20||

\twolineshloka
{इत्येवं विविधा वाचो नानाजनसमीरिताः}
{शुश्राव रामः श्रुत्वा च न विचक्रेऽस्य मानसम्} %||2-30-21||

\fourlineindentedshloka
{प्रतीक्षमाणोऽभिजनं तदार्त}
{मनार्तरूपः प्रहसन्निवाथ}
{जगाम रामः पितरं दिदृक्षुः}
{ पितुर्निदेशं विधिवच्चिकीर्षुः} %||2-30-22||

\fourlineindentedshloka
{तत्पूर्वमैक्ष्वाकसुतो महात्मा}
{ रामो गमिष्यन्वनमार्तरूपम्}
{व्यतिष्ठत प्रेक्ष्य तदा सुमन्त्रम्}
{ पितुर्महात्मा प्रतिहारणार्थम्} %||2-30-23||

\fourlineindentedshloka
{पितुर्निदेशेन तु धर्मवत्सलो}
{ वनप्रवेशे कृतबुद्धिनिश्चयः}
{स राघवः प्रेक्ष्य सुमन्त्रमब्रवी}
{न्निवेदयस्वागमनं नृपाय मे} %||2-31-24||


{॥इत्यार्षे श्रीमद्रामायणे वाल्मीकीये आदिकाव्ये अयोध्याकाण्डे त्रिंशः सर्गः॥}

\sect{एकत्रिंशः सर्गः}

\twolineshloka
{स रामप्रेषितः क्षिप्रं सन्तापकलुषेन्द्रियः}
{प्रविश्य नृपतिं सूतो निःश्वसन्तं ददर्श ह} %||2-31-1||

\twolineshloka
{आलोक्य तु महाप्राज्ञः परमाकुल चेतसम्}
{राममेवानुशोचन्तं सूतः प्राञ्जलिरासदत्} %||2-31-2||

\twolineshloka
{अयं स पुरुषव्याघ्र द्वारि तिष्ठति ते सुतः}
{ब्राह्मणेभ्यो धनं दत्त्वा सर्वं चैवोपजीविनाम्} %||2-31-3||

\twolineshloka
{स त्वा पश्यतु भद्रं ते रामः सत्यपराक्रमः}
{सर्वान्सुहृद आपृच्छ्य त्वामिदानीं दिदृक्षते} %||2-31-4||

\twolineshloka
{गमिष्यति महारण्यं तं पश्य जगतीपते}
{वृतं राजगुणैः सर्वैरादित्यमिव रश्मिभिः} %||2-31-5||

\twolineshloka
{स सत्यवादी धर्मात्मा गाम्भीर्यात्सागरोपमः}
{आकाश इव निष्पङ्को नरेन्द्रः प्रत्युवाच तम्} %||2-31-6||

\twolineshloka
{सुमन्त्रानय मे दारान्ये केचिदिह मामकाः}
{दारैः परिवृतः सर्वैर्द्रष्टुमिच्छामि राघवम्} %||2-31-7||

\twolineshloka
{सोऽन्तःपुरमतीत्यैव स्त्रियस्ता वाक्यमब्रवीत्}
{आर्यो ह्वयति वो राजा गम्यतां तत्र माचिरम्} %||2-31-8||

\twolineshloka
{एवमुक्ताः स्त्रियः सर्वाः सुमन्त्रेण नृपाज्ञया}
{प्रचक्रमुस्तद्भवनं भर्तुराज्ञाय शासनम्} %||2-31-9||

\twolineshloka
{अर्धसप्तशतास्तास्तु प्रमदास्ताम्रलोचनाः}
{कौसल्यां परिवार्याथ शनैर्जग्मुर्धृतव्रताः} %||2-31-10||

\twolineshloka
{आगतेषु च दारेषु समवेक्ष्य महीपतिः}
{उवाच राजा तं सूतं सुमन्त्रानय मे सुतम्} %||2-31-11||

\twolineshloka
{स सूतो राममादाय लक्ष्मणं मैथिलीं तदा}
{जगामाभिमुखस्तूर्णं सकाशं जगतीपतेः} %||2-31-12||

\twolineshloka
{स राजा पुत्रमायान्तं दृष्ट्वा दूरात्कृताञ्जलिम्}
{उत्पपातासनात्तूर्णमार्तः स्त्रीजनसंवृतः} %||2-31-13||

\twolineshloka
{सोऽभिदुद्राव वेगेन रामं दृष्ट्वा विशां पतिः}
{तमसम्प्राप्य दुःखार्तः पपात भुवि मूर्छितः} %||2-31-14||

\twolineshloka
{तं रामोऽभ्यपातत्क्षिप्रं लक्ष्मणश्च महारथः}
{विसंज्ञमिव दुःखेन सशोकं नृपतिं तदा} %||2-31-15||

\twolineshloka
{स्त्रीसहस्रनिनादश्च संजज्ञे राजवेश्मनि}
{हाहा रामेति सहसा भूषणध्वनिमूर्छितः} %||2-31-16||

\twolineshloka
{तं परिष्वज्य बाहुभ्यां तावुभौ रामलक्ष्मणौ}
{पर्यङ्के सीतया सार्धं रुदन्तः समवेशयन्} %||2-31-17||

\twolineshloka
{अथ रामो मुहूर्तेन लब्धसंज्ञं महीपतिम्}
{उवाच प्राञ्जलिर्भूत्वा शोकार्णवपरिप्लुतम्} %||2-31-18||

\twolineshloka
{आपृच्छे त्वां महाराज सर्वेषामीश्वरोऽसि नः}
{प्रस्थितं दण्डकारण्यं पश्य त्वं कुशलेन माम्} %||2-31-19||

\twolineshloka
{लक्ष्मणं चानुजानीहि सीता चान्वेति मां वनम्}
{कारणैर्बहुभिस्तथ्यैर्वार्यमाणौ न चेच्छतः} %||2-31-20||

\twolineshloka
{अनुजानीहि सर्वान्नः शोकमुत्सृज्य मानद}
{लक्ष्मणं मां च सीतां च प्रजापतिरिव प्रजाः} %||2-31-21||

\twolineshloka
{प्रतीक्षमाणमव्यग्रमनुज्ञां जगतीपतेः}
{उवाच रर्जा सम्प्रेक्ष्य वनवासाय राघवम्} %||2-31-22||

\twolineshloka
{अहं राघव कैकेय्या वरदानेन मोहितः}
{अयोध्यायास्त्वमेवाद्य भव राजा निगृह्य माम्} %||2-31-23||

\twolineshloka
{एवमुक्तो नृपतिना रामो धर्मभृतां वरः}
{प्रत्युवाचाञ्जलिं कृत्वा पितरं वाक्यकोविदः} %||2-31-24||

\twolineshloka
{भवान्वर्षसहस्राय पृथिव्या नृपते पतिः}
{अहं त्वरण्ये वत्स्यामि न मे कार्यं त्वयानृतम्} %||2-31-25||

\twolineshloka
{श्रेयसे वृद्धये तात पुनरागमनाय च}
{गच्छस्वारिष्टमव्यग्रः पन्थानमकुतोभयम्} %||2-31-26||

\threelineshloka
{अद्य त्विदानीं रजनीं पुत्र मा गच्छ सर्वथा}
{मातरं मां च सम्पश्यन्वसेमामद्य शर्वरीम्}
{तर्पितः सर्वकामैस्त्वं श्वःकाले साधयिष्यसि} %||2-31-27||

\twolineshloka
{अथ रामस्तथा श्रुत्वा पितुरार्तस्य भाषितम्}
{लक्ष्मणेन सह भ्रात्रा दीनो वचनमब्रवीत्} %||2-31-28||

\twolineshloka
{प्राप्स्यामि यानद्य गुणान्को मे श्वस्तान्प्रदास्यति}
{अपक्रमणमेवातः सर्वकामैरहं वृणे} %||2-31-29||

\twolineshloka
{इयं सराष्ट्रा सजना धनधान्यसमाकुला}
{मया विसृष्टा वसुधा भरताय प्रदीयताम्} %||2-31-30||

\twolineshloka
{अपगच्छतु ते दुःखं मा भूर्बाष्पपरिप्लुतः}
{न हि क्षुभ्यति दुर्धर्षः समुद्रः सरितां पतिः} %||2-31-31||

\twolineshloka
{नैवाहं राज्यमिच्छामि न सुखं न च मैथिलीम्}
{त्वामहं सत्यमिच्छामि नानृतं पुरुषर्षभ} %||2-31-32||

\fourlineindentedshloka
{पुरं च राष्ट्रं च मही च केवला}
{ मया निसृष्टा भरताय दीयताम्}
{अहं निदेशं भवतोऽनुपालय}
{न्वनं गमिष्यामि चिराय सेवितुम्} %||2-31-33||

\fourlineindentedshloka
{मया निसृष्टां भरतो महीमिमाम्}
{ सशैलखण्डां सपुरां सकाननाम्}
{शिवां सुसीमामनुशास्तु केवलम्}
{ त्वया यदुक्तं नृपते यथास्तु तत्} %||2-31-34||

\fourlineindentedshloka
{न मे तथा पार्थिव धीयते मनो}
{ महत्सु कामेषु न चात्मनः प्रिये}
{यथा निदेशे तव शिष्टसम्मते}
{ व्यपैतु दुःखं तव मत्कृतेऽनघ} %||2-31-35||

\fourlineindentedshloka
{तदद्य नैवानघ राज्यमव्ययम्}
{ न सर्वकामान्न सुखं न मैथिलीम्}
{न जीवितं त्वामनृतेन योजय}
{न्वृणीय सत्यं व्रतमस्तु ते तथा} %||2-31-36||

\fourlineindentedshloka
{फलानि मूलानि च भक्षयन्वने}
{ गिरींश्च पश्यन्सरितः सरांसि च}
{वनं प्रविश्यैव विचित्रपादपम्}
{ सुखी भविष्यामि तवास्तु निर्वृतिः} %||2-32-37||


{॥इत्यार्षे श्रीमद्रामायणे वाल्मीकीये आदिकाव्ये अयोध्याकाण्डे एकत्रिंशः सर्गः॥}

\sect{द्वात्रिंशः सर्गः}

\twolineshloka
{ततः सुमन्त्रमैक्ष्वाकः पीडितोऽत्र प्रतिज्ञया}
{सबाष्पमतिनिःश्वस्य जगादेदं पुनः पुनः} %||2-32-1||

\twolineshloka
{सूत रत्नसुसम्पूर्णा चतुर्विधबला चमूः}
{रागवस्यानुयात्रार्थं क्षिप्रं प्रतिविधीयताम्} %||2-32-2||

\twolineshloka
{रूपाजीवा च शालिन्यो वणिजश्च महाधनाः}
{शोभयन्तु कुमारस्य वाहिनीं सुप्रसारिताः} %||2-32-3||

\twolineshloka
{ये चैनमुपजीवन्ति रमते यैश्च वीर्यतः}
{तेषां बहुविधं दत्त्वा तानप्यत्र नियोजय} %||2-32-4||

\twolineshloka
{निघ्नन्मृगान्कुञ्जरांश्च पिबंश्चारण्यकं मधु}
{नदीश्च विविधाः पश्यन्न राज्यं संस्मरिष्यति} %||2-32-5||

\twolineshloka
{धान्यकोशश्च यः कश्चिद्धनकोशश्च मामकः}
{तौ राममनुगच्छेतां वसन्तं निर्जने वने} %||2-32-6||

\twolineshloka
{यजन्पुण्येषु देशेषु विसृजंश्चाप्तदक्षिणाः}
{ऋषिभिश्च समागम्य प्रवत्स्यति सुखं वने} %||2-32-7||

\twolineshloka
{भरतश्च महाबाहुरयोध्यां पालयिष्यति}
{सर्वकामैः पुनः श्रीमान्रामः संसाध्यतामिति} %||2-32-8||

\twolineshloka
{एवं ब्रुवति काकुत्स्थे कैकेय्या भयमागतम्}
{मुखं चाप्यगमाच्छेषं स्वरश्चापि न्यरुध्यत} %||2-32-9||

\threelineshloka
{सा विषण्णा च सन्त्रस्ता कैकेयी वाक्यमब्रवीत्}
{राज्यं गतजनं साधो पीतमण्डां सुरामिव}
{निरास्वाद्यतमं शून्यं भरतो नाभिपत्स्यते} %||2-32-10||

\threelineshloka
{कैकेय्यां मुक्तलज्जायां वदन्त्यामतिदारुणम्}
{राजा दशरथो वाक्यमुवाचायतलोचनाम्}
{वहन्तं किं तुदसि मां नियुज्य धुरि माहिते} %||2-32-11||

\threelineshloka
{कैकेयी द्विगुणं क्रुद्धा राजानमिदमब्रवीत्}
{तवैव वंशे सगरो ज्येष्ठं पुत्रमुपारुधत्}
{असमञ्ज इति ख्यातं तथायं गन्तुमर्हति} %||2-32-12||

\twolineshloka
{एवमुक्तो धिगित्येव राजा दशरथोऽब्रवीत्}
{व्रीडितश्च जनः सर्वः सा च तन्नावबुध्यत} %||2-32-13||

\twolineshloka
{तत्र वृद्धो महामात्रः सिद्धार्थो नाम नामतः}
{शुचिर्बहुमतो राज्ञः कैकेयीमिदमब्रवीत्} %||2-32-14||

\twolineshloka
{असमञ्जो गृहीत्वा तु क्रीडितः पथि दारकान्}
{सरय्वाः प्रक्षिपन्नप्सु रमते तेन दुर्मतिः} %||2-32-15||

\twolineshloka
{तं दृष्ट्वा नागरः सर्वे क्रुद्धा राजानमब्रुवन्}
{असमञ्जं वृषीण्वैकमस्मान्वा राष्ट्रवर्धन} %||2-32-16||

\twolineshloka
{तानुवाच ततो राजा किंनिमित्तमिदं भयम्}
{ताश्चापि राज्ञा सम्पृष्टा वाक्यं प्रकृतयोऽब्रुवन्} %||2-32-17||

\twolineshloka
{क्रीडितस्त्वेष नः पुत्रान्बालानुद्भ्रान्तचेतनः}
{सरय्वां प्रक्षिपन्मौर्ख्यादतुलां प्रीतिमश्नुते} %||2-32-18||

\twolineshloka
{स तासां वचनं श्रुत्वा प्रकृतीनां नराधिप}
{तं तत्याजाहितं पुत्रं तासां प्रियचिकीर्षया} %||2-32-19||

\twolineshloka
{इत्येवमत्यजद्राजा सगरो वै सुधार्मिकः}
{रामः किमकरोत्पापं येनैवमुपरुध्यते} %||2-32-20||

\twolineshloka
{श्रुत्वा तु सिद्धार्थवचो राजा श्रान्ततरस्वनः}
{शोकोपहतया वाचा कैकेयीमिदमब्रवीत्} %||2-32-21||

\fourlineindentedshloka
{अनुव्रजिष्याम्यहमद्य रामम्}
{ राज्यं परित्यज्य सुखं धनं च}
{सहैव राज्ञा भरतेन च त्वम्}
{ यथा सुखं भुङ्क्ष्व चिराय राज्यम्} %||2-33-22||


{॥इत्यार्षे श्रीमद्रामायणे वाल्मीकीये आदिकाव्ये अयोध्याकाण्डे द्वात्रिंशः सर्गः॥}

\sect{त्रयस्त्रिंशः सर्गः}

\twolineshloka
{महामात्रवचः श्रुत्वा रामो दशरथं तदा}
{अन्वभाषत वाक्यं तु विनयज्ञो विनीतवत्} %||2-33-1||

\twolineshloka
{त्यक्तभोगस्य मे राजन्वने वन्येन जीवतः}
{किं कार्यमनुयात्रेण त्यक्तसङ्गस्य सर्वतः} %||2-33-2||

\twolineshloka
{यो हि दत्त्वा द्विपश्रेष्ठं कक्ष्यायां कुरुते मनः}
{रज्जुस्नेहेन किं तस्य त्यजतः कुञ्जरोत्तमम्} %||2-33-3||

\twolineshloka
{तथा मम सतां श्रेष्ठ किं ध्वजिन्या जगत्पते}
{सर्वाण्येवानुजानामि चीराण्येवानयन्तु मे} %||2-33-4||

\twolineshloka
{खनित्रपिटके चोभे ममानयत गच्छतः}
{चतुर्दश वने वासं वर्षाणि वसतो मम} %||2-33-5||

\twolineshloka
{अथ चीराणि कैकेयी स्वयमाहृत्य राघवम्}
{उवाच परिधत्स्वेति जनौघे निरपत्रपा} %||2-33-6||

\twolineshloka
{स चीरे पुरुषव्याघ्रः कैकेय्याः प्रतिगृह्य ते}
{सूक्ष्मवस्त्रमवक्षिप्य मुनिवस्त्राण्यवस्त ह} %||2-33-7||

\twolineshloka
{लक्ष्मणश्चापि तत्रैव विहाय वसने शुभे}
{तापसाच्छादने चैव जग्राह पितुरग्रतः} %||2-33-8||

\twolineshloka
{अथात्मपरिधानार्थं सीता कौशेयवासिनी}
{समीक्ष्य चीरं सन्त्रस्ता पृषती वागुरामिव} %||2-33-9||

\threelineshloka
{सा व्यपत्रपमाणेव प्रतिगृह्य च दुर्मनाः}
{गन्धर्वराजप्रतिमं भर्तारमिदमब्रवीत्}
{कथं नु चीरं बध्नन्ति मुनयो वनवासिनः} %||2-33-10||

\twolineshloka
{कृत्वा कण्ठे च सा चीरमेकमादाय पाणिना}
{तस्थौ ह्यकुषला तत्र व्रीडिता जनकात्मज} %||2-33-11||

\twolineshloka
{तस्यास्तत्क्षिप्रमागम्य रामो धर्मभृतां वरः}
{चीरं बबन्ध सीतायाः कौशेयस्योपरि स्वयम्} %||2-33-12||

\twolineshloka
{तस्यां चीरं वसानायां नाथवत्यामनाथवत्}
{प्रचुक्रोश जनः सर्वो धिक्त्वां दशरथं त्विति} %||2-33-13||

\twolineshloka
{स निःश्वस्योष्णमैक्ष्वाकस्तां भार्यामिदमब्रवीत्}
{कैकेयि कुशचीरेण न सीता गन्तुमर्हति} %||2-33-14||

\twolineshloka
{ननु पर्याप्तमेतत्ते पापे रामविवासनम्}
{किमेभिः कृपणैर्भूयः पातकैरपि ते कृतैः} %||2-33-15||

\twolineshloka
{एवं ब्रुवन्तं पितरं रामः सम्प्रस्थितो वनम्}
{अवाक्शिरसमासीनमिदं वचनमब्रवीत्} %||2-33-16||

\twolineshloka
{इयं धार्मिक कौसल्या मम माता यशस्विनी}
{वृद्धा चाक्षुद्रशीला च न च त्वां देवगर्हिते} %||2-33-17||

\twolineshloka
{मया विहीनां वरद प्रपन्नां शोकसागरम्}
{अदृष्टपूर्वव्यसनां भूयः सम्मन्तुमर्हसि} %||2-33-18||

\fourlineindentedshloka
{इमां महेन्द्रोपमजातगर्भिणीम्}
{ तथा विधातुं जनमीं ममार्हसि}
{यथा वनस्थे मयि शोककर्शिता}
{ न जीवितं न्यस्य यमक्षयं व्रजेत्} %||2-34-19||


{॥इत्यार्षे श्रीमद्रामायणे वाल्मीकीये आदिकाव्ये अयोध्याकाण्डे त्रयस्त्रिंशः सर्गः॥}

\sect{चतुस्त्रिंशः सर्गः}

\twolineshloka
{रामस्य तु वचः श्रुत्वा मुनिवेषधरं च तम्}
{समीक्ष्य सह भार्याभी राजा विगतचेतनः} %||2-34-1||

\twolineshloka
{नैनं दुःखेन सन्तप्तः प्रत्यवैक्षत राघवम्}
{न चैनमभिसम्प्रेक्ष्य प्रत्यभाषत दुर्मनाः} %||2-34-2||

\twolineshloka
{स मुहूर्तमिवासंज्ञो दुःखितश्च महीपतिः}
{विललाप महाबाहू राममेवानुचिन्तयन्} %||2-34-3||

\twolineshloka
{मन्ये खलु मया पूर्वं विवत्सा बहवः कृताः}
{प्राणिनो हिंसिता वापि तस्मादिदमुपस्थितम्} %||2-34-4||

\twolineshloka
{न त्वेवानागते काले देहाच्च्यवति जीवितम्}
{कैकेय्या क्लिश्यमानस्य मृत्युर्मम न विद्यते} %||2-34-5||

\twolineshloka
{योऽहं पावकसङ्काशं पश्यामि पुरतः स्थितम्}
{विहाय वसने सूक्ष्मे तापसाच्छादमात्मजम्} %||2-34-6||

\twolineshloka
{एकस्याः खलु कैकेय्याः कृतेऽयं क्लिश्यते जनः}
{स्वार्थे प्रयतमानायाः संश्रित्य निकृतिं त्विमाम्} %||2-34-7||

\twolineshloka
{एवमुक्त्वा तु वचनं बाष्पेण पिहितेक्ष्णह}
{रामेति सकृदेवोक्त्वा व्याहर्तुं न शशाक ह} %||2-34-8||

\twolineshloka
{संज्ञां तु प्रतिलभ्यैव मुहूर्तात्स महीपतिः}
{नेत्राभ्यामश्रुपूर्णाभ्यां सुमन्त्रमिदमब्रवीत्} %||2-34-9||

\twolineshloka
{औपवाह्यं रथं युक्त्वा त्वमायाहि हयोत्तमैः}
{प्रापयैनं महाभागमितो जनपदात्परम्} %||2-34-10||

\twolineshloka
{एवं मन्ये गुणवतां गुणानां फलमुच्यते}
{पित्रा मात्रा च यत्साधुर्वीरो निर्वास्यते वनम्} %||2-34-11||

\twolineshloka
{राज्ञो वचनमाज्ञाय सुमन्त्रः शीघ्रविक्रमः}
{योजयित्वाययौ तत्र रथमश्वैरलङ्कृतम्} %||2-34-12||

\twolineshloka
{तं रथं राजपुत्राय सूतः कनकभूषितम्}
{आचचक्षेऽञ्जलिं कृत्वा युक्तं परमवाजिभिः} %||2-34-13||

\twolineshloka
{राजा सत्वरमाहूय व्यापृतं वित्तसञ्चये}
{उवाच देशकालज्ञो निश्चितं सर्वतः शुचि} %||2-34-14||

\twolineshloka
{वासांसि च महार्हाणि भूषणानि वराणि च}
{वर्षाण्येतानि सङ्ख्याय वैदेह्याः क्षिप्रमानय} %||2-34-15||

\twolineshloka
{नरेन्द्रेणैवमुक्तस्तु गत्वा कोशगृहं ततः}
{प्रायच्छत्सर्वमाहृत्य सीतायै क्षिप्रमेव तत्} %||2-34-16||

\twolineshloka
{सा सुजाता सुजातानि वैदेही प्रस्थिता वनम्}
{भूषयामास गात्राणि तैर्विचित्रैर्विभूषणैः} %||2-34-17||

\twolineshloka
{व्यराजयत वैदेही वेश्म तत्सुविभूषिता}
{उद्यतोंऽशुमतः काले खं प्रभेव विवस्वतः} %||2-34-18||

\twolineshloka
{तां भुजाभ्यां परिष्वज्य श्वश्रूर्वचनमब्रवीत्}
{अनाचरन्तीं कृपणं मूध्न्युपाघ्राय मैथिलीम्} %||2-34-19||

\twolineshloka
{असत्यः सर्वलोकेऽस्मिन्सततं सत्कृताः प्रियैः}
{भर्तारं नानुमन्यन्ते विनिपातगतं स्त्रियः} %||2-34-20||

\twolineshloka
{स त्वया नावमन्तव्यः पुत्रः प्रव्राजितो मम}
{तव दैवतमस्त्वेष निर्धनः सधनोऽपि वा} %||2-34-21||

\twolineshloka
{विज्ञाय वचनं सीता तस्या धर्मार्थसंहितम्}
{कृताञ्जलिरुवाचेदं श्वश्रूमभिमुखे स्थिता} %||2-34-22||

\twolineshloka
{करिष्ये सर्वमेवाहमार्या यदनुशास्ति माम्}
{अभिज्ञास्मि यथा भर्तुर्वर्तितव्यं श्रुतं च मे} %||2-34-23||

\twolineshloka
{न मामसज्जनेनार्या समानयितुमर्हति}
{धर्माद्विचलितुं नाहमलं चन्द्रादिव प्रभा} %||2-34-24||

\twolineshloka
{नातन्त्री वाद्यते वीणा नाचक्रो वर्तते रथः}
{नापतिः सुखमेधते या स्यादपि शतात्मजा} %||2-34-25||

\twolineshloka
{मितं ददाति हि पिता मितं माता मितं सुतः}
{अमितस्य हि दातारं भर्तारं का न पूजयेत्} %||2-34-26||

\twolineshloka
{साहमेवङ्गता श्रेष्ठा श्रुतधर्मपरावरा}
{आर्ये किमवमन्येयं स्त्रीणां भर्ता हि दैवतम्} %||2-34-27||

\twolineshloka
{सीताया वचनं श्रुत्वा कौसल्या हृदयङ्गमम्}
{शुद्धसत्त्वा मुमोचाश्रु सहसा दुःखहर्षजम्} %||2-34-28||

\twolineshloka
{तां प्राञ्जलिरभिक्रम्य मातृमध्येऽतिसत्कृताम्}
{रामः परमधर्मज्ञो मातरं वाक्यमब्रवीत्} %||2-34-29||

\twolineshloka
{अम्ब मा दुःखिता भूस्त्वं पश्य त्वं पितरं मम}
{क्षयो हि वनवासस्य क्षिप्रमेव भविष्यति} %||2-34-30||

\twolineshloka
{सुप्तायास्ते गमिष्यन्ति नववर्षाणि पञ्च च}
{सा समग्रमिह प्राप्तं मां द्रक्ष्यसि सुहृद्वृतम्} %||2-34-31||

\twolineshloka
{एतावदभिनीतार्थमुक्त्वा स जननीं वचः}
{त्रयः शतशतार्धा हि ददर्शावेक्ष्य मातरः} %||2-34-32||

\twolineshloka
{ताश्चापि स तथैवार्ता मातॄर्दशरथात्मजः}
{धर्मयुक्तमिदं वाक्यं निजगाद कृताञ्जलिः} %||2-34-33||

\twolineshloka
{संवासात्परुषं किञ्चिदज्ञानाद्वापि यत्कृतम्}
{तन्मे समनुजानीत सर्वाश्चामन्त्रयामि वः} %||2-34-34||

\twolineshloka
{जज्ञेऽथ तासां संनादः क्रौञ्चीनामिव निःस्वनः}
{मानवेन्द्रस्य भार्याणामेवं वदति राघवे} %||2-34-35||

\fourlineindentedshloka
{मुरजपणवमेघघोषव}
{द्दशरथवेश्म बभूव यत्पुरा}
{विलपित परिदेवनाकुलम्}
{ व्यसनगतं तदभूत्सुदुःखितम्} %||2-35-36||


{॥इत्यार्षे श्रीमद्रामायणे वाल्मीकीये आदिकाव्ये अयोध्याकाण्डे चतुस्त्रिंशः सर्गः॥}

\sect{पञ्चत्रिंशः सर्गः}

\twolineshloka
{अथ रामश्च सीता च लक्ष्मणश्च कृताञ्जलिः}
{उपसङ्गृह्य राजानं चक्रुर्दीनाः प्रदक्षिणम्} %||2-35-1||

\twolineshloka
{तं चापि समनुज्ञाप्य धर्मज्ञः सीतया सह}
{राघवः शोकसम्मूढो जननीमभ्यवादयत्} %||2-35-2||

\twolineshloka
{अन्वक्षं लक्ष्मणो भ्रातुः कौसल्यामभ्यवादयत्}
{अथ मातुः सुमित्राया जग्राह चरणौ पुनः} %||2-35-3||

\twolineshloka
{तं वन्दमानं रुदती माता सौमित्रिमब्रवीत्}
{हितकामा महाबाहुं मूर्ध्न्युपाघ्राय लक्ष्मणम्} %||2-35-4||

\twolineshloka
{सृष्टस्त्वं वनवासाय स्वनुरक्तः सुहृज्जने}
{रामे प्रमादं मा कार्षीः पुत्र भ्रातरि गच्छति} %||2-35-5||

\twolineshloka
{व्यसनी वा समृद्धो वा गतिरेष तवानघ}
{एष लोके सतां धर्मो यज्ज्येष्ठवशगो भवेत्} %||2-35-6||

\twolineshloka
{इदं हि वृत्तमुचितं कुलस्यास्य सनातनम्}
{दानं दीक्षा च यज्ञेषु तनुत्यागो मृधेषु च} %||2-35-7||

\twolineshloka
{रामं दशरथं विद्धि मां विद्धि जनकात्मजाम्}
{अयोध्यामटवीं विद्धि गच्छ तात यथासुखम्} %||2-35-8||

\twolineshloka
{ततः सुमन्त्रः काकुत्स्थं प्राञ्जलिर्वाक्यमब्रवीत्}
{विनीतो विनयज्ञश्च मातलिर्वासवं यथा} %||2-35-9||

\twolineshloka
{रथमारोह भद्रं ते राजपुत्र महायशः}
{क्षिप्रं त्वां प्रापयिष्यामि यत्र मां राम वक्ष्यसि} %||2-35-10||

\twolineshloka
{चतुर्दश हि वर्षाणि वस्तव्यानि वने त्वया}
{तान्युपक्रमितव्यानि यानि देव्यासि चोदितः} %||2-35-11||

\twolineshloka
{तं रथं सूर्यसङ्काशं सीता हृष्टेन चेतसा}
{आरुरोह वरारोहा कृत्वालङ्कारमात्मनः} %||2-35-12||

\twolineshloka
{तथैवायुधजातानि भ्रातृभ्यां कवचानि च}
{रथोपस्थे प्रतिन्यस्य सचर्मकठिनं च तत्} %||2-35-13||

\twolineshloka
{सीतातृतीयानारूढान्दृष्ट्वा धृष्टमचोदयत्}
{सुमन्त्रः सम्मतानश्वान्वायुवेगसमाञ्जवे} %||2-35-14||

\twolineshloka
{प्रयाते तु महारण्यं चिररात्राय राघवे}
{बभूव नगरे मूर्च्छा बलमूर्च्छा जनस्य च} %||2-35-15||

\twolineshloka
{तत्समाकुलसम्भ्रान्तं मत्तसङ्कुपित द्विपम्}
{हयशिञ्जितनिर्घोषं पुरमासीन्महास्वनम्} %||2-35-16||

\twolineshloka
{ततः सबालवृद्धा सा पुरी परमपीडिता}
{राममेवाभिदुद्राव घर्मार्तः सलिलं यथा} %||2-35-17||

\twolineshloka
{पार्श्वतः पृष्ठतश्चापि लम्बमानास्तदुन्मुखाः}
{बाष्पपूर्णमुखाः सर्वे तमूचुर्भृशदुःखिताः} %||2-35-18||

\twolineshloka
{संयच्छ वाजिनां रश्मीन्सूत याहि शनैः शनैः}
{मुखं द्रक्ष्यामि रामस्य दुर्दर्शं नो भविष्यति} %||2-35-19||

\twolineshloka
{आयसं हृदयं नूनं राममातुरसंशयम्}
{यद्देवगर्भप्रतिमे वनं याति न भिद्यते} %||2-35-20||

\twolineshloka
{कृतकृत्या हि वैदेही छायेवानुगता पतिम्}
{न जहाति रता धर्मे मेरुमर्कप्रभा यथा} %||2-35-21||

\twolineshloka
{अहो लक्ष्मण सिद्धार्थः सततां प्रियवादिनम्}
{भ्रातरं देवसङ्काशं यस्त्वं परिचरिष्यसि} %||2-35-22||

\threelineshloka
{महत्येषा हि ते सिद्धिरेष चाभ्युदयो महान्}
{एष स्वर्गस्य मार्गश्च यदेनमनुगच्छसि}
{एवं वदन्तस्ते सोढुं न शेकुर्बाष्पमागतम्} %||2-35-23||

\twolineshloka
{अथ राजा वृतः स्त्रीभिर्दीनाभिर्दीनचेतनः}
{निर्जगाम प्रियं पुत्रं द्रक्ष्यामीति ब्रुवन्गृहात्} %||2-35-24||

\twolineshloka
{शुश्रुवे चाग्रतः स्त्रीणां रुदन्तीनां महास्वनः}
{यथा नादः करेणूनां बद्धे महति कुञ्जरे} %||2-35-25||

\twolineshloka
{पिता च राजा काकुत्स्थः श्रीमान्सन्नस्तदा बभौ}
{परिपूर्णः शशी काले ग्रहेणोपप्लुतो यथा} %||2-35-26||

\twolineshloka
{ततो हलहलाशब्दो जज्ञे रामस्य पृष्ठतः}
{नराणां प्रेक्ष्य राजानं सीदन्तं भृशदुःखितम्} %||2-35-27||

\twolineshloka
{हा रामेति जनाः केचिद्राममातेति चापरे}
{अन्तःपुरं समृद्धं च क्रोशन्तं पर्यदेवयन्} %||2-35-28||

\threelineshloka
{अन्वीक्षमाणो रामस्तु विषण्णं भ्रान्तचेतसम्}
{राजानं मातरं चैव ददर्शानुगतौ पथि}
{धर्मपाशेन सङ्क्षिप्तः प्रकाशं नाभ्युदैक्षत} %||2-35-29||

\twolineshloka
{पदातिनौ च यानार्हावदुःखार्हौ सुखोचितौ}
{दृष्ट्वा सञ्चोदयामास शीघ्रं याहीति सारथिम्} %||2-35-30||

\twolineshloka
{न हि तत्पुरुषव्याघ्रो दुःखदं दर्शनं पितुः}
{मातुश्च सहितुं शक्तस्तोत्रार्दित इव द्विपः} %||2-35-31||

\threelineshloka
{तथा रुदन्तीं कौसल्यां रथं तमनुधावतीम्}
{क्रोशन्तीं राम रामेति हा सीते लक्ष्मणेति च}
{असकृत्प्रैक्षत तदा नृत्यन्तीमिव मातरम्} %||2-35-32||

\twolineshloka
{तिष्ठेति राजा चुक्रोष याहि याहीति राघवः}
{सुमन्त्रस्य बभूवात्मा चक्रयोरिव चान्तरा} %||2-35-33||

\twolineshloka
{नाश्रौषमिति राजानमुपालब्धोऽपि वक्ष्यसि}
{चिरं दुःखस्य पापिष्ठमिति रामस्तमब्रवीत्} %||2-35-34||

\twolineshloka
{रामस्य स वचः कुर्वन्ननुज्ञाप्य च तं जनम्}
{व्रजतोऽपि हयाञ्शीघ्रं चोदयामास सारथिः} %||2-35-35||

\twolineshloka
{न्यवर्तत जनो राज्ञो रामं कृत्वा प्रदक्षिणम्}
{मनसाप्यश्रुवेगैश्च न न्यवर्तत मानुषम्} %||2-35-36||

\twolineshloka
{यमिच्छेत्पुनरायान्तं नैनं दूरमनुव्रजेत्}
{इत्यमात्या महाराजमूचुर्दशरथं वचः} %||2-35-37||

\fourlineindentedshloka
{तेषां वचः सर्वगुणोपपन्नम्}
{ प्रस्विन्नगात्रः प्रविषण्णरूपः}
{निशम्य राजा कृपणः सभार्यो}
{ व्यवस्थितस्तं सुतमीक्षमाणः} %||2-36-38||


{॥इत्यार्षे श्रीमद्रामायणे वाल्मीकीये आदिकाव्ये अयोध्याकाण्डे पञ्चत्रिंशः सर्गः॥}

\sect{षट्त्रिंशः सर्गः}

\twolineshloka
{तस्मिंस्तु पुरुषव्याघ्रे निष्क्रामति कृताञ्जलौ}
{आर्तशब्दो हि संजज्ञे स्त्रीणामन्तःपुरे महान्} %||2-36-1||

\twolineshloka
{अनाथस्य जनस्यास्य दुर्बलस्य तपस्विनः}
{यो गतिं शरणं चासीत्स नाथः क्व नु गच्छति} %||2-36-2||

\twolineshloka
{न क्रुध्यत्यभिशस्तोऽपि क्रोधनीयानि वर्जयन्}
{क्रुद्धान्प्रसादयन्सर्वान्समदुःखः क्व गच्छति} %||2-36-3||

\twolineshloka
{कौसल्यायां महातेजा यथा मातरि वर्तते}
{तथा यो वर्ततेऽस्मासु महात्मा क्व नु गच्छति} %||2-36-4||

\twolineshloka
{कैकेय्या क्लिश्यमानेन राज्ञा सञ्चोदितो वनम्}
{परित्राता जनस्यास्य जगतः क्व नु गच्छति} %||2-36-5||

\twolineshloka
{अहो निश्चेतनो राजा जीवलोकस्य सम्प्रियम्}
{धर्म्यं सत्यव्रतं रामं वनवासो प्रवत्स्यति} %||2-36-6||

\twolineshloka
{इति सर्वा महिष्यस्ता विवत्सा इव धेनवः}
{रुरुदुश्चैव दुःखार्ताः सस्वरं च विचुक्रुशुः} %||2-36-7||

\twolineshloka
{स तमन्तःपुरे घोरमार्तशब्दं महीपतिः}
{पुत्रशोकाभिसन्तप्तः श्रुत्वा चासीत्सुदुःखितः} %||2-36-8||

\twolineshloka
{नाग्निहोत्राण्यहूयन्त सूर्यश्चान्तरधीयत}
{व्यसृजन्कवलान्नागा गावो वत्सान्न पाययन्} %||2-36-9||

\twolineshloka
{त्रिशङ्कुर्लोहिताङ्गश्च बृहस्पतिबुधावपि}
{दारुणाः सोममभ्येत्य ग्रहाः सर्वे व्यवस्थिताः} %||2-36-10||

\twolineshloka
{नक्षत्राणि गतार्चींषि ग्रहाश्च गततेजसः}
{विशाखाश्च सधूमाश्च नभसि प्रचकाशिरे} %||2-36-11||

\twolineshloka
{अकस्मान्नागरः सर्वो जनो दैन्यमुपागमत्}
{आहारे वा विहारे वा न कश्चिदकरोन्मनः} %||2-36-12||

\twolineshloka
{बाष्पपर्याकुलमुखो राजमार्गगतो जनः}
{न हृष्टो लक्ष्यते कश्चित्सर्वः शोकपरायणः} %||2-36-13||

\twolineshloka
{न वाति पवनः शीतो न शशी सौम्यदर्शनः}
{न सूर्यस्तपते लोकं सर्वं पर्याकुलं जगत्} %||2-36-14||

\twolineshloka
{अनर्थिनः सुताः स्त्रीणां भर्तारो भ्रातरस्तथा}
{सर्वे सर्वं परित्यज्य राममेवान्वचिन्तयन्} %||2-36-15||

\twolineshloka
{ये तु रामस्य सुहृदः सर्वे ते मूढचेतसः}
{शोकभारेण चाक्रान्ताः शयनं न जुहुस्तदा} %||2-36-16||

\fourlineindentedshloka
{ततस्त्वयोध्या रहिता महात्मना}
{ पुरन्दरेणेव मही सपर्वता}
{चचाल घोरं भयभारपीडिता}
{ सनागयोधाश्वगणा ननाद च} %||2-37-17||


{॥इत्यार्षे श्रीमद्रामायणे वाल्मीकीये आदिकाव्ये अयोध्याकाण्डे षट्त्रिंशः सर्गः॥}

\sect{सप्तत्रिंशः सर्गः}

\twolineshloka
{यावत्तु निर्यतस्तस्य रजोरूपमदृश्यत}
{नैवेक्ष्वाकुवरस्तावत्संजहारात्मचक्षुषी} %||2-37-1||

\twolineshloka
{यावद्राजा प्रियं पुत्रं पश्यत्यत्यन्तधार्मिकम्}
{तावद्व्यवर्धतेवास्य धरण्यां पुत्रदर्शने} %||2-37-2||

\twolineshloka
{न पश्यति रजोऽप्यस्य यदा रामस्य भूमिपः}
{तदार्तश्च विषण्णश्च पपात धरणीतले} %||2-37-3||

\twolineshloka
{तस्य दक्षिणमन्वगात्कौसल्या बाहुमङ्गना}
{वामं चास्यान्वगात्पार्श्वं कैकेयी भरतप्रिया} %||2-37-4||

\twolineshloka
{तां नयेन च सम्पन्नो धर्मेण निवयेन च}
{उवाच राजा कैकेयीं समीक्ष्य व्यथितेन्द्रियः} %||2-37-5||

\twolineshloka
{कैकेयि मा ममाङ्गानि स्प्राक्षीस्त्वं दुष्टचारिणी}
{न हि त्वां द्रष्टुमिच्छामि न भार्या न च बान्धवी} %||2-37-6||

\twolineshloka
{ये च त्वामुपजीवन्ति नाहं तेषां न ते मम}
{केवलार्थपरां हि त्वां त्यक्तधर्मां त्यजाम्यहम्} %||2-37-7||

\twolineshloka
{अगृह्णां यच्च ते पाणिमग्निं पर्यणयं च यत्}
{अनुजानामि तत्सर्वमस्मिँल्लोके परत्र च} %||2-37-8||

\twolineshloka
{भरतश्चेत्प्रतीतः स्याद्राज्यं प्राप्येदमव्ययम्}
{यन्मे स दद्यात्पित्रर्थं मा मा तद्दत्तमागमत्} %||2-37-9||

\twolineshloka
{अथ रेणुसमुध्वस्तं तमुत्थाप्य नराधिपम्}
{न्यवर्तत तदा देवी कौसल्या शोककर्शिता} %||2-37-10||

\twolineshloka
{हत्वेव ब्राह्मणं कामात्स्पृष्ट्वाग्निमिव पाणिना}
{अन्वतप्यत धर्मात्मा पुत्रं सञ्चिन्त्य तापसम्} %||2-37-11||

\twolineshloka
{निवृत्यैव निवृत्यैव सीदतो रथवर्त्मसु}
{राज्ञो नातिबभौ रूपं ग्रस्तस्यांशुमतो यथा} %||2-37-12||

\twolineshloka
{विललाप च दुःखार्तः प्रियं पुत्रमनुस्मरन्}
{नगरान्तमनुप्राप्तं बुद्ध्वा पुत्रमथाब्रवीत्} %||2-37-13||

\twolineshloka
{वाहनानां च मुख्यानां वहतां तं ममात्मजम्}
{पदानि पथि दृश्यन्ते स महात्मा न दृश्यते} %||2-37-14||

\twolineshloka
{स नूनं क्वचिदेवाद्य वृक्षमूलमुपाश्रितः}
{काष्ठं वा यदि वाश्मानमुपधाय शयिष्यते} %||2-37-15||

\twolineshloka
{उत्थास्यति च मेदिन्याः कृपणः पांशुगुण्ठितः}
{विनिःश्वसन्प्रस्रवणात्करेणूनामिवर्षभः} %||2-37-16||

\twolineshloka
{द्रक्ष्यन्ति नूनं पुरुषा दीर्घबाहुं वनेचराः}
{राममुत्थाय गच्छन्तं लोकनाथमनाथवत्} %||2-37-17||

\twolineshloka
{सकामा भव कैकेयि विधवा राज्यमावस}
{न हि तं पुरुषव्याघ्रं विना जीवितुमुत्सहे} %||2-37-18||

\twolineshloka
{इत्येवं विलपन्राजा जनौघेनाभिसंवृतः}
{अपस्नात इवारिष्टं प्रविवेश पुरोत्तमम्} %||2-37-19||

\twolineshloka
{शून्यचत्वरवेश्मान्तां संवृतापणदेवताम्}
{क्लान्तदुर्बलदुःखार्तां नात्याकीर्णमहापथाम्} %||2-37-20||

\twolineshloka
{तामवेक्ष्य पुरीं सर्वां राममेवानुचिन्तयन्}
{विलपन्प्राविशद्राजा गृहं सूर्य इवाम्बुदम्} %||2-37-21||

\twolineshloka
{महाह्रदमिवाक्षोभ्यं सुपर्णेन हृतोरगम्}
{रामेण रहितं वेश्म वैदेह्या लक्ष्मणेन च} %||2-37-22||

\twolineshloka
{कौसल्याया गृहं शीघ्रं राम मातुर्नयन्तु माम्}
{इति ब्रुवन्तं राजानमनयन्द्वारदर्शितः} %||2-37-23||

\twolineshloka
{ततस्तत्र प्रविष्टस्य कौसल्याया निवेशनम्}
{अधिरुह्यापि शयनं बभूव लुलितं मनः} %||2-37-24||

\twolineshloka
{तच्च दृष्ट्वा महाराजो भुजमुद्यम्य वीर्यवान्}
{उच्चैः स्वरेण चुक्रोश हा राघव जहासि माम्} %||2-37-25||

\twolineshloka
{सुखिता बत तं कालं जीविष्यन्ति नरोत्तमाः}
{परिष्वजन्तो ये रामं द्रक्ष्यन्ति पुनरागतम्} %||2-37-26||

\twolineshloka
{न त्वां पश्यामि कौसल्ये साधु मां पाणिना स्पृश}
{रामं मेऽनुगता दृष्टिरद्यापि न निवर्तते} %||2-37-27||

\fourlineindentedshloka
{तं राममेवानुविचिन्तयन्तम्}
{ समीक्ष्य देवी शयने नरेन्द्रम्}
{उपोपविश्याधिकमार्तरूपा}
{ विनिःश्वसन्ती विललाप कृच्छ्रम्} %||2-38-28||


{॥इत्यार्षे श्रीमद्रामायणे वाल्मीकीये आदिकाव्ये अयोध्याकाण्डे सप्तत्रिंशः सर्गः॥}

\sect{अष्टात्रिंशः सर्गः}

\twolineshloka
{ततः समीक्ष्य शयने सन्नं शोकेन पार्थिवम्}
{कौसल्या पुत्रशोकार्ता तमुवाच महीपतिम्} %||2-38-1||

\twolineshloka
{राघवो नरशार्दूल विषमुप्त्वा द्विजिह्ववत्}
{विचरिष्यति कैकेयी निर्मुक्तेव हि पन्नगी} %||2-38-2||

\twolineshloka
{विवास्य रामं सुभगा लब्धकामा समाहिता}
{त्रासयिष्यति मां भूयो दुष्टाहिरिव वेश्मनि} %||2-38-3||

\twolineshloka
{अथ स्म नगरे रामश्चरन्भैक्षं गृहे वसेत्}
{कामकारो वरं दातुमपि दासं ममात्मजम्} %||2-38-4||

\twolineshloka
{पातयित्वा तु कैकेय्या रामं स्थानाद्यथेष्टतः}
{प्रदिष्टो रक्षसां भागः पर्वणीवाहिताग्निना} %||2-38-5||

\twolineshloka
{गजराजगतिर्वीरो महाबाहुर्धनुर्धरः}
{वनमाविशते नूनं सभार्यः सहलक्ष्मणः} %||2-38-6||

\twolineshloka
{वने त्वदृष्टदुःखानां कैकेय्यानुमते त्वया}
{त्यक्तानां वनवासाय का न्ववस्था भविष्यति} %||2-38-7||

\twolineshloka
{ते रत्नहीनास्तरुणाः फलकाले विवासिताः}
{कथं वत्स्यन्ति कृपणाः फलमूलैः कृताशनाः} %||2-38-8||

\twolineshloka
{अपीदानीं स कालः स्यान्मम शोकक्षयः शिवः}
{सभार्यं यत्सह भ्रात्रा पश्येयमिह राघवम्} %||2-38-9||

\twolineshloka
{श्रुत्वैवोपस्थितौ वीरौ कदायोध्या भविष्यति}
{यशस्विनी हृष्टजना सूच्छ्रितध्वजमालिनी} %||2-38-10||

\twolineshloka
{कदा प्रेक्ष्य नरव्याघ्रावरण्यात्पुनरागतौ}
{नन्दिष्यति पुरी हृष्टा समुद्र इव पर्वणि} %||2-38-11||

\twolineshloka
{कदायोध्यां महाबाहुः पुरीं वीरः प्रवेक्ष्यति}
{पुरस्कृत्य रथे सीतां वृषभो गोवधूमिव} %||2-38-12||

\twolineshloka
{कदा प्राणिसहस्राणि राजमार्गे ममात्मजौ}
{लाजैरवकरिष्यन्ति प्रविशन्तावरिन्दमौ} %||2-38-13||

\twolineshloka
{कदा सुमनसः कन्या द्विजातीनां फलानि च}
{प्रदिशन्त्यः पुरीं हृष्टाः करिष्यन्ति प्रदक्षिणम्} %||2-38-14||

\twolineshloka
{कदा परिणतो बुद्ध्या वयसा चामरप्रभः}
{अभ्युपैष्यति धर्मज्ञस्त्रिवर्ष इव मां ललन्} %||2-38-15||

\twolineshloka
{निःसंशयं मया मन्ये पुरा वीर कदर्यया}
{पातु कामेषु वत्सेषु मातॄणां शातिताः स्तनाः} %||2-38-16||

\twolineshloka
{साहं गौरिव सिंहेन विवत्सा वत्सला कृता}
{कैकेय्या पुरुषव्याघ्र बालवत्सेव गौर्बलात्} %||2-38-17||

\twolineshloka
{न हि तावद्गुणैर्जुष्टं सर्वशास्त्रविशारदम्}
{एकपुत्रा विना पुत्रमहं जीवितुमुत्सहे} %||2-38-18||

\twolineshloka
{न हि मे जीविते किञ्चित्सामर्थमिह कल्प्यते}
{अपश्यन्त्याः प्रियं पुत्रं महाबाहुं महाबलम्} %||2-38-19||

\fourlineindentedshloka
{अयं हि मां दीपयते समुत्थित}
{स्तनूजशोकप्रभवो हुताशनः}
{महीमिमां रश्मिभिरुत्तमप्रभो}
{ यथा निदाघे भगवान्दिवाकरः} %||2-39-20||


{॥इत्यार्षे श्रीमद्रामायणे वाल्मीकीये आदिकाव्ये अयोध्याकाण्डे अष्टात्रिंशः सर्गः॥}

\sect{एकोनचत्वारिंशः सर्गः}

\twolineshloka
{विलपन्तीं तथा तां तु कौसल्यां प्रमदोत्तमाम्}
{इदं धर्मे स्थिता धर्म्यं सुमित्रा वाक्यमब्रवीत्} %||2-39-1||

\twolineshloka
{तवार्ये सद्गुणैर्युक्तः पुत्रः स पुरुषोत्तमः}
{किं ते विलपितेनैवं कृपणं रुदितेन वा} %||2-39-2||

\twolineshloka
{यस्तवार्ये गतः पुत्रस्त्यक्त्वा राज्यं महाबलः}
{साधु कुर्वन्महात्मानं पितरं सत्यवादिनाम्} %||2-39-3||

\twolineshloka
{शिष्टैराचरिते सम्यक्शश्वत्प्रेत्य फलोदये}
{रामो धर्मे स्थितः श्रेष्ठो न स शोच्यः कदाचन} %||2-39-4||

\twolineshloka
{वर्तते चोत्तमां वृत्तिं लक्ष्मणोऽस्मिन्सदानघः}
{दयावान्सर्वभूतेषु लाभस्तस्य महात्मनः} %||2-39-5||

\twolineshloka
{अरण्यवासे यद्दुःखं जानती वै सुखोचिता}
{अनुगच्छति वैदेही धर्मात्मानं तवात्मजम्} %||2-39-6||

\twolineshloka
{कीर्तिभूतां पताकां यो लोके भ्रामयति प्रभुः}
{दमसत्यव्रतपरः किं न प्राप्तस्तवात्मजः} %||2-39-7||

\twolineshloka
{व्यक्तं रामस्य विज्ञाय शौचं माहात्म्यमुत्तमम्}
{न गात्रमंशुभिः सूर्यः सन्तापयितुमर्हति} %||2-39-8||

\twolineshloka
{शिवः सर्वेषु कालेषु काननेभ्यो विनिःसृतः}
{राघवं युक्तशीतोष्णः सेविष्यति सुखोऽनिलः} %||2-39-9||

\twolineshloka
{शयानमनघं रात्रौ पितेवाभिपरिष्वजन्}
{रश्मिभिः संस्पृशञ्शीतैश्चन्द्रमा ह्लादयिष्यति} %||2-39-10||

\twolineshloka
{ददौ चास्त्राणि दिव्यानि यस्मै ब्रह्मा महौजसे}
{दानवेन्द्रं हतं दृष्ट्वा तिमिध्वजसुतं रणे} %||2-39-11||

\twolineshloka
{पृथिव्या सह वैदेह्या श्रिया च पुरुषर्षभः}
{क्षिप्रं तिसृभिरेताभिः सह रामोऽभिषेक्ष्यते} %||2-39-12||

\twolineshloka
{दुःखजं विसृजन्त्यस्रं निष्क्रामन्तमुदीक्ष्य यम्}
{समुत्स्रक्ष्यसि नेत्राभ्यां क्षिप्रमानन्दजं पयः} %||2-39-13||

\twolineshloka
{अभिवादयमानं तं दृष्ट्वा ससुहृदं सुतम्}
{मुदाश्रु मोक्ष्यसे क्षिप्रं मेघलेकेव वार्षिकी} %||2-39-14||

\twolineshloka
{पुत्रस्ते वरदः क्षिप्रमयोध्यां पुनरागतः}
{कराभ्यां मृदुपीनाभ्यां चरणौ पीडयिष्यति} %||2-39-15||

\fourlineindentedshloka
{निशम्य तल्लक्ष्मणमातृवाक्यम्}
{ रामस्य मातुर्नरदेवपत्न्याः}
{सद्यः शरीरे विननाश शोकः}
{ शरद्गतो मेघ इवाल्पतोयः} %||2-40-16||


{॥इत्यार्षे श्रीमद्रामायणे वाल्मीकीये आदिकाव्ये अयोध्याकाण्डे एकोनचत्वारिंशः सर्गः॥}

\sect{चत्वारिंशः सर्गः}

\twolineshloka
{अनुरक्ता महात्मानं रामं सत्यपरक्रमम्}
{अनुजग्मुः प्रयान्तं तं वनवासाय मानवाः} %||2-40-1||

\twolineshloka
{निवर्तितेऽपि च बलात्सुहृद्वर्गे च राजिनि}
{नैव ते संन्यवर्तन्त रामस्यानुगता रथम्} %||2-40-2||

\twolineshloka
{अयोध्यानिलयानां हि पुरुषाणां महायशाः}
{बभूव गुणसम्पन्नः पूर्णचन्द्र इव प्रियः} %||2-40-3||

\twolineshloka
{स याच्यमानः काकुत्स्थः स्वाभिः प्रकृतिभिस्तदा}
{कुर्वाणः पितरं सत्यं वनमेवान्वपद्यत} %||2-40-4||

\twolineshloka
{अवेक्षमाणः सस्नेहं चक्षुषा प्रपिबन्निव}
{उवाच रामः स्नेहेन ताः प्रजाः स्वाः प्रजा इव} %||2-40-5||

\twolineshloka
{या प्रीतिर्बहुमानश्च मय्ययोध्यानिवासिनाम्}
{मत्प्रियार्थं विशेषेण भरते सा निवेश्यताम्} %||2-40-6||

\twolineshloka
{स हि कल्याण चारित्रः कैकेय्यानन्दवर्धनः}
{करिष्यति यथावद्वः प्रियाणि च हितानि च} %||2-40-7||

\twolineshloka
{ज्ञानवृद्धो वयोबालो मृदुर्वीर्यगुणान्वितः}
{अनुरूपः स वो भर्ता भविष्यति भयापहः} %||2-40-8||

\twolineshloka
{स हि राजगुणैर्युक्तो युवराजः समीक्षितः}
{अपि चापि मया शिष्टैः कार्यं वो भर्तृशासनम्} %||2-40-9||

\twolineshloka
{न च तप्येद्यथा चासौ वनवासं गते मयि}
{महाराजस्तथा कार्यो मम प्रियचिकीर्षया} %||2-40-10||

\twolineshloka
{यथा यथा दाशरथिर्धर्ममेवास्थितोऽभवत्}
{तथा तथा प्रकृतयो रामं पतिमकामयन्} %||2-40-11||

\twolineshloka
{बाष्पेण पिहितं दीनं रामः सौमित्रिणा सह}
{चकर्षेव गुणैर्बद्ध्वा जनं पुनरिवासनम्} %||2-40-12||

\twolineshloka
{ते द्विजास्त्रिविधं वृद्धा ज्ञानेन वयसौजसा}
{वयःप्रकम्पशिरसो दूरादूचुरिदं वचः} %||2-40-13||

\threelineshloka
{वहन्तो जवना रामं भो भो जात्यास्तुरङ्गमाः}
{निवर्तध्वं न गन्तव्यं हिता भवत भर्तरि}
{उपवाह्यस्तु वो भर्ता नापवाह्यः पुराद्वनम्} %||2-40-14||

\twolineshloka
{एवमार्तप्रलापांस्तान्वृद्धान्प्रलपतो द्विजान्}
{अवेक्ष्य सहसा रामो रथादवततार ह} %||2-40-15||

\twolineshloka
{पद्भ्यामेव जगामाथ ससीतः सहलक्ष्मणः}
{संनिकृष्टपदन्यासो रामो वनपरायणः} %||2-40-16||

\twolineshloka
{द्विजातींस्तु पदातींस्तान्रामश्चारित्रवत्सलः}
{न शशाक घृणाचक्षुः परिमोक्तुं रथेन सः} %||2-40-17||

\twolineshloka
{गच्छन्तमेव तं दृष्ट्वा रामं सम्भ्रान्तमानसाः}
{ऊचुः परमसन्तप्ता रामं वाक्यमिदं द्विजाः} %||2-40-18||

\twolineshloka
{ब्राह्मण्यं कृत्स्नमेतत्त्वां ब्रह्मण्यमनुगच्छति}
{द्विजस्कन्धाधिरूढास्त्वामग्नयोऽप्यनुयान्त्यमी} %||2-40-19||

\twolineshloka
{वाजपेयसमुत्थानि छत्राण्येतानि पश्य नः}
{पृष्ठतोऽनुप्रयातानि हंसानिव जलात्यये} %||2-40-20||

\twolineshloka
{अनवाप्तातपत्रस्य रश्मिसन्तापितस्य ते}
{एभिश्छायां करिष्यामः स्वैश्छत्रैर्वाजपेयिकैः} %||2-40-21||

\twolineshloka
{या हि नः सततं बुद्धिर्वेदमन्त्रानुसारिणी}
{त्वत्कृते सा कृता वत्स वनवासानुसारिणी} %||2-40-22||

\twolineshloka
{हृदयेष्ववतिष्ठन्ते वेदा ये नः परं धनम्}
{वत्स्यन्त्यपि गृहेष्वेव दाराश्चारित्ररक्षिताः} %||2-40-23||

\twolineshloka
{न पुनर्निश्चयः कार्यस्त्वद्गतौ सुकृता मतिः}
{त्वयि धर्मव्यपेक्षे तु किं स्याद्धर्ममवेक्षितुम्} %||2-40-24||

\twolineshloka
{याचितो नो निवर्तस्व हंसशुक्लशिरोरुहैः}
{शिरोभिर्निभृताचार महीपतनपांशुलैः} %||2-40-25||

\twolineshloka
{बहूनां वितता यज्ञा द्विजानां य इहागताः}
{तेषां समाप्तिरायत्ता तव वत्स निवर्तने} %||2-40-26||

\twolineshloka
{भक्तिमन्ति हि भूतानि जङ्गमाजङ्गमानि च}
{याचमानेषु तेषु त्वं भक्तिं भक्तेषु दर्शय} %||2-40-27||

\twolineshloka
{अनुगन्तुमशक्तास्त्वां मूलैरुद्धृतवेगिभिः}
{उन्नता वायुवेगेन विक्रोशन्तीव पादपाः} %||2-40-28||

\twolineshloka
{निश्चेष्टाहारसञ्चारा वृक्षैकस्थानविष्ठिताः}
{पक्षिणोऽपि प्रयाचन्ते सर्वभूतानुकम्पिनम्} %||2-40-29||

\twolineshloka
{एवं विक्रोशतां तेषां द्विजातीनां निवर्तने}
{ददृशे तमसा तत्र वारयन्तीव राघवम्} %||2-41-30||


{॥इत्यार्षे श्रीमद्रामायणे वाल्मीकीये आदिकाव्ये अयोध्याकाण्डे चत्वारिंशः सर्गः॥}

\sect{एकचत्वारिंशः सर्गः}

\twolineshloka
{ततस्तु तमसा तीरं रम्यमाश्रित्य राघवः}
{सीतामुद्वीक्ष्य सौमित्रिमिदं वचनमब्रवीत्} %||2-41-1||

\twolineshloka
{इयमद्य निशा पूर्वा सौमित्रे प्रस्थिता वनम्}
{वनवासस्य भद्रं ते स नोत्कण्ठितुमर्हसि} %||2-41-2||

\twolineshloka
{पश्य शून्यान्यरण्यानि रुदन्तीव समन्ततः}
{यथानिलयमायद्भिर्निलीनानि मृगद्विजैः} %||2-41-3||

\twolineshloka
{अद्यायोध्या तु नगरी राजधानी पितुर्मम}
{सस्त्रीपुंसा गतानस्माञ्शोचिष्यति न संशयः} %||2-41-4||

\twolineshloka
{भरतः खलु धर्मात्मा पितरं मातरं च मे}
{धर्मार्थकामसहितैर्वाक्यैराश्वासयिष्यति} %||2-41-5||

\twolineshloka
{भरतस्यानृशंसत्वं सञ्चिन्त्याहं पुनः पुनः}
{नानुशोचामि पितरं मातरं चापि लक्ष्मण} %||2-41-6||

\twolineshloka
{त्वया कार्यं नरव्याघ्र मामनुव्रजता कृतम्}
{अन्वेष्टव्या हि वैदेह्या रक्षणार्थे सहायता} %||2-41-7||

\twolineshloka
{अद्भिरेव तु सौमित्रे वत्स्याम्यद्य निशामिमाम्}
{एतद्धि रोचते मह्यं वन्येऽपि विविधे सति} %||2-41-8||

\twolineshloka
{एवमुक्त्वा तु सौमित्रं सुमन्त्रमपि राघवः}
{अप्रमत्तस्त्वमश्वेषु भव सौम्येत्युवाच ह} %||2-41-9||

\twolineshloka
{सोऽश्वान्सुमन्त्रः संयम्य सूर्येऽस्तं समुपागते}
{प्रभूतयवसान्कृत्वा बभूव प्रत्यनन्तरः} %||2-41-10||

\twolineshloka
{उपास्यतु शिवां सन्ध्यां दृष्ट्वा रात्रिमुपस्थिताम्}
{रामस्य शयनं चक्रे सूतः सौमित्रिणा सह} %||2-41-11||

\twolineshloka
{तां शय्यां तमसातीरे वीक्ष्य वृक्षदलैः कृताम्}
{रामः सौमित्रिणां सार्धं सभार्यः संविवेश ह} %||2-41-12||

\twolineshloka
{सभार्यं सम्प्रसुप्तं तं भ्रातरं वीक्ष्य लक्ष्मणः}
{कथयामास सूताय रामस्य विविधान्गुणान्} %||2-41-13||

\twolineshloka
{जाग्रतो ह्येव तां रात्रिं सौमित्रेरुदितो रविः}
{सूतस्य तमसातीरे रामस्य ब्रुवतो गुणान्} %||2-41-14||

\twolineshloka
{गोकुलाकुलतीरायास्तमसाया विदूरतः}
{अवसत्तत्र तां रात्रिं रामः प्रकृतिभिः सह} %||2-41-15||

\twolineshloka
{उत्थाय तु महातेजाः प्रकृतीस्ता निशाम्य च}
{अब्रवीद्भ्रातरं रामो लक्ष्मणं पुण्यलक्षणम्} %||2-41-16||

\twolineshloka
{अस्मद्व्यपेक्षान्सौमित्रे निरपेक्षान्गृहेष्वपि}
{वृक्षमूलेषु संसुप्तान्पश्य लक्ष्मण साम्प्रतम्} %||2-41-17||

\twolineshloka
{यथैते नियमं पौराः कुर्वन्त्यस्मन्निवर्तने}
{अपि प्राणानसिष्यन्ति न तु त्यक्ष्यन्ति निश्चयम्} %||2-41-18||

\twolineshloka
{यावदेव तु संसुप्तास्तावदेव वयं लघु}
{रथमारुह्य गच्छामः पन्थानमकुतोभयम्} %||2-41-19||

\twolineshloka
{अतो भूयोऽपि नेदानीमिक्ष्वाकुपुरवासिनः}
{स्वपेयुरनुरक्ता मां वृक्षमूलानि संश्रिताः} %||2-41-20||

\twolineshloka
{पौरा ह्यात्मकृताद्दुःखाद्विप्रमोच्या नृपात्मजैः}
{न तु खल्वात्मना योज्या दुःखेन पुरवासिनः} %||2-41-21||

\twolineshloka
{अब्रवील्लक्ष्मणो रामं साक्षाद्धर्ममिव स्थितम्}
{रोचते मे महाप्राज्ञ क्षिप्रमारुह्यतामिति} %||2-41-22||

\twolineshloka
{सूतस्ततः सन्त्वरितः स्यन्दनं तैर्हयोत्तमैः}
{योजयित्वाथ रामाय प्राञ्जलिः प्रत्यवेदयत्} %||2-41-23||

\twolineshloka
{मोहनार्थं तु पौराणां सूतं रामोऽब्रवीद्वचः}
{उदङ्मुखः प्रयाहि त्वं रथमास्थाय सारथे} %||2-41-24||

\twolineshloka
{मुहूर्तं त्वरितं गत्वा निर्गतय रथं पुनः}
{यथा न विद्युः पौरा मां तथा कुरु समाहितः} %||2-41-25||

\twolineshloka
{रामस्य वचनं श्रुत्वा तथा चक्रे स सारथिः}
{प्रत्यागम्य च रामस्य स्यन्दनं प्रत्यवेदयत्} %||2-41-26||

\twolineshloka
{तं स्यन्दनमधिष्ठाय राघवः सपरिच्छदः}
{शीघ्रगामाकुलावर्तां तमसामतरन्नदीम्} %||2-41-27||

\twolineshloka
{स सन्तीर्य महाबाहुः श्रीमाञ्शिवमकण्टकम्}
{प्रापद्यत महामार्गमभयं भयदर्शिनाम्} %||2-41-28||

\twolineshloka
{प्रभातायां तु शर्वर्यां पौरास्ते राघवो विना}
{शोकोपहतनिश्चेष्टा बभूवुर्हतचेतसः} %||2-41-29||

\twolineshloka
{शोकजाश्रुपरिद्यूना वीक्षमाणास्ततस्ततः}
{आलोकमपि रामस्य न पश्यन्ति स्म दुःखिताः} %||2-41-30||

\twolineshloka
{ततो मार्गानुसारेण गत्वा किञ्चित्क्षणं पुनः}
{मार्गनाशाद्विषादेन महता समभिप्लुतः} %||2-41-31||

\twolineshloka
{रथस्य मार्गनाशेन न्यवर्तन्त मनस्विनः}
{किमिदं किं करिष्यामो दैवेनोपहता इति} %||2-41-32||

\twolineshloka
{ततो यथागतेनैव मार्गेण क्लान्तचेतसः}
{अयोध्यामगमन्सर्वे पुरीं व्यथितसज्जनाम्} %||2-42-33||


{॥इत्यार्षे श्रीमद्रामायणे वाल्मीकीये आदिकाव्ये अयोध्याकाण्डे एकचत्वारिंशः सर्गः॥}

\sect{द्विचत्वारिंशः सर्गः}

\twolineshloka
{अनुगम्य निवृत्तानां रामं नगरवासिनाम्}
{उद्गतानीव सत्त्वानि बभूवुरमनस्विनाम्} %||2-42-1||

\twolineshloka
{स्वं स्वं निलयमागम्य पुत्रदारैः समावृताः}
{अश्रूणि मुमुचुः सर्वे बाष्पेण पिहिताननाः} %||2-42-2||

\twolineshloka
{न चाहृष्यन्न चामोदन्वणिजो न प्रसारयन्}
{न चाशोभन्त पण्यानि नापचन्गृहमेधिनः} %||2-42-3||

\twolineshloka
{नष्टं दृष्ट्वा नाभ्यनन्दन्विपुलं वा धनागमम्}
{पुत्रं प्रथमजं लब्ध्वा जननी नाभ्यनन्दत} %||2-42-4||

\twolineshloka
{गृहे गृहे रुदन्त्यश्च भर्तारं गृहमागतम्}
{व्यगर्हयन्तो दुःखार्ता वाग्भिस्तोत्रैरिव द्विपान्} %||2-42-5||

\twolineshloka
{किं नु तेषां गृहैः कार्यं किं दारैः किं धनेन वा}
{पुत्रैर्वा किं सुखैर्वापि ये न पश्यन्ति राघवम्} %||2-42-6||

\twolineshloka
{एकः सत्पुरुषो लोके लक्ष्मणः सह सीतया}
{योऽनुगच्छति काकुत्स्थं रामं परिचरन्वने} %||2-42-7||

\twolineshloka
{आपगाः कृतपुण्यास्ताः पद्मिन्यश्च सरांसि च}
{येषु स्नास्यति काकुत्स्थो विगाह्य सलिलं शुचि} %||2-42-8||

\twolineshloka
{शोभयिष्यन्ति काकुत्स्थमटव्यो रम्यकाननाः}
{आपगाश्च महानूपाः सानुमन्तश्च पर्वताः} %||2-42-9||

\twolineshloka
{काननं वापि शैलं वा यं रामोऽभिगमिष्यति}
{प्रियातिथिमिव प्राप्तं नैनं शक्ष्यन्त्यनर्चितुम्} %||2-42-10||

\threelineshloka
{विचित्रकुसुमापीडा बहुमञ्जरिधारिणः}
{अकाले चापि मुख्यानि पुष्पाणि च फलानि च}
{दर्शयिष्यन्त्यनुक्रोशाद्गिरयो राममागतम्} %||2-42-11||

\twolineshloka
{विदर्शयन्तो विविधान्भूयश्चित्रांश्च निर्झरान्}
{पादपाः पर्वताग्रेषु रमयिष्यन्ति राघवम्} %||2-42-12||

\twolineshloka
{यत्र रामो भयं नात्र नास्ति तत्र पराभवः}
{स हि शूरो महाबाहुः पुत्रो दशरथस्य च} %||2-42-13||

\threelineshloka
{पुरा भवति नो दूरादनुगच्छाम राघवम्}
{पादच्छाया सुखा भर्तुस्तादृशस्य महात्मनः}
{स हि नाथो जनस्यास्य स गतिः स परायणम्} %||2-42-14||

\twolineshloka
{वयं परिचरिष्यामः सीतां यूयं तु राघवम्}
{इति पौरस्त्रियो भर्तॄन्दुःखार्तास्तत्तदब्रुवन्} %||2-42-15||

\twolineshloka
{युष्माकं राघवोऽरण्ये योगक्षेमं विधास्यति}
{सीता नारीजनस्यास्य योगक्षेमं करिष्यति} %||2-42-16||

\twolineshloka
{को न्वनेनाप्रतीतेन सोत्कण्ठितजनेन च}
{सम्प्रीयेतामनोज्ञेन वासेन हृतचेतसा} %||2-42-17||

\twolineshloka
{कैकेय्या यदि चेद्राज्यं स्यादधर्म्यमनाथवत्}
{न हि नो जीवितेनार्थः कुतः पुत्रैः कुतो धनैः} %||2-42-18||

\twolineshloka
{यया पुत्रश्च भर्ता च त्यक्तावैश्वर्यकारणात्}
{कं सा परिहरेदन्यं कैकेयी कुलपांसनी} %||2-42-19||

\twolineshloka
{कैकेय्या न वयं राज्ये भृतका निवसेमहि}
{जीवन्त्या जातु जीवन्त्यः पुत्रैरपि शपामहे} %||2-42-20||

\twolineshloka
{या पुत्रं पार्थिवेन्द्रस्य प्रवासयति निर्घृणा}
{कस्तां प्राप्य सुखं जीवेदधर्म्यां दुष्टचारिणीम्} %||2-42-21||

\twolineshloka
{न हि प्रव्रजिते रामे जीविष्यति महीपतिः}
{मृते दशरथे व्यक्तं विलोपस्तदनन्तरम्} %||2-42-22||

\twolineshloka
{ते विषं पिबतालोड्य क्षीणपुण्याः सुदुर्गताः}
{राघवं वानुगच्छध्वमश्रुतिं वापि गच्छत} %||2-42-23||

\twolineshloka
{मिथ्या प्रव्राजितो रामः सभार्यः सहलक्ष्मणः}
{भरते संनिषृष्टाः स्मः सौनिके पशवो यथा} %||2-42-24||

\twolineshloka
{तास्तथा विलपन्त्यस्तु नगरे नागरस्त्रियः}
{चुक्रुशुर्भृशसन्तप्ता मृत्योरिव भयागमे} %||2-42-25||

\fourlineindentedshloka
{तथा स्त्रियो रामनिमित्तमातुरा}
{ यथा सुते भ्रातरि वा विवासिते}
{विलप्य दीना रुरुदुर्विचेतसः}
{ सुतैर्हि तासामधिको हि सोऽभवत्} %||2-43-26||


{॥इत्यार्षे श्रीमद्रामायणे वाल्मीकीये आदिकाव्ये अयोध्याकाण्डे द्विचत्वारिंशः सर्गः॥}

\sect{त्रिचत्वारिंशः सर्गः}

\twolineshloka
{रामोऽपि रात्रिशेषेण तेनैव महदन्तरम्}
{जगाम पुरुषव्याघ्रः पितुराज्ञामनुस्मरन्} %||2-43-1||

\twolineshloka
{तथैव गच्छतस्तस्य व्यपायाद्रजनी शिवा}
{उपास्य स शिवां सन्ध्यां विषयान्तं व्यगाहत} %||2-43-2||

\twolineshloka
{ग्रामान्विकृष्टसीमांस्तान्पुष्पितानि वनानि च}
{पश्यन्नतिययौ शीघ्रं शरैरिव हयोत्तमैः} %||2-43-3||

\twolineshloka
{शृण्वन्वाचो मनुष्याणां ग्रामसंवासवासिनाम्}
{राजानं धिग्दशरथं कामस्य वशमागतम्} %||2-43-4||

\twolineshloka
{हा नृशंसाद्य कैकेयी पापा पापानुबन्धिनी}
{तीक्ष्णा सम्भिन्नमर्यादा तीक्ष्णे कर्मणि वर्तते} %||2-43-5||

\twolineshloka
{या पुत्रमीदृशं राज्ञः प्रवासयति धार्मिकम्}
{वन वासे महाप्राज्ञं सानुक्रोशमतन्द्रितम्} %||2-43-6||

\twolineshloka
{एता वाचो मनुष्याणां ग्रामसंवासवासिनाम्}
{शृण्वन्नतिययौ वीरः कोसलान्कोसलेश्वरः} %||2-43-7||

\twolineshloka
{ततो वेदश्रुतिं नाम शिववारिवहां नदीम्}
{उत्तीर्याभिमुखः प्रायादगस्त्याध्युषितां दिशम्} %||2-43-8||

\twolineshloka
{गत्वा तु सुचिरं कालं ततः शीतजलां नदीम्}
{गोमतीं गोयुतानूपामतरत्सागरङ्गमाम्} %||2-43-9||

\twolineshloka
{गोमतीं चाप्यतिक्रम्य राघवः शीघ्रगैर्हयैः}
{मयूरहंसाभिरुतां ततार स्यन्दिकां नदीम्} %||2-43-10||

\twolineshloka
{स महीं मनुना राज्ञा दत्तामिक्ष्वाकवे पुरा}
{स्फीतां राष्ट्रावृतां रामो वैदेहीमन्वदर्शयत्} %||2-43-11||

\twolineshloka
{सूत इत्येव चाभाष्य सारथिं तमभीक्ष्णशः}
{हंसमत्तस्वरः श्रीमानुवाच पुरुषर्षभः} %||2-43-12||

\twolineshloka
{कदाहं पुनरागम्य सरय्वाः पुष्पिते वने}
{मृगयां पर्याटष्यामि मात्रा पित्रा च सङ्गतः} %||2-43-13||

\twolineshloka
{नात्यर्थमभिकाङ्क्षामि मृगयां सरयूवने}
{रतिर्ह्येषातुला लोके राजर्षिगणसम्मता} %||2-43-14||

\twolineshloka
{स तमध्वानमैक्ष्वाकः सूतं मधुरया गिरा}
{तं तमर्थमभिप्रेत्य ययौवाक्यमुदीरयन्} %||2-44-15||


{॥इत्यार्षे श्रीमद्रामायणे वाल्मीकीये आदिकाव्ये अयोध्याकाण्डे त्रिचत्वारिंशः सर्गः॥}

\sect{चतुश्चत्वारिंशः सर्गः}

\twolineshloka
{विशालान्कोसलान्रम्यान्यात्वा लक्ष्मणपूर्वजः}
{आससाद महाबाहुः शृङ्गवेरपुरं प्रति} %||2-44-1||

\twolineshloka
{तत्र त्रिपथगां दिव्यां शिवतोयामशैवलाम्}
{ददर्श राघवो गङ्गां पुण्यामृषिनिसेविताम्} %||2-44-2||

\twolineshloka
{हंससारससङ्घुष्टां चक्रवाकोपकूजिताम्}
{शिंशुमरैश्च नक्रैश्च भुजङ्गैश्च निषेविताम्} %||2-44-3||

\twolineshloka
{तामूर्मिकलिलावर्तामन्ववेक्ष्य महारथः}
{सुमन्त्रमब्रवीत्सूतमिहैवाद्य वसामहे} %||2-44-4||

\twolineshloka
{अविदूरादयं नद्या बहुपुष्पप्रवालवान्}
{सुमहानिङ्गुदीवृक्षो वसामोऽत्रैव सारथे} %||2-44-5||

\twolineshloka
{लक्षणश्च सुमन्त्रश्च बाढमित्येव राघवम्}
{उक्त्वा तमिङ्गुदीवृक्षं तदोपययतुर्हयैः} %||2-44-6||

\twolineshloka
{रामोऽभियाय तं रम्यं वृक्षमिक्ष्वाकुनन्दनः}
{रथादवातरत्तस्मात्सभार्यः सहलक्ष्मणः} %||2-44-7||

\twolineshloka
{सुमन्त्रोऽप्यवतीर्यैव मोचयित्वा हयोत्तमान्}
{वृक्षमूलगतं राममुपतस्थे कृताञ्जलिः} %||2-44-8||

\twolineshloka
{तत्र राजा गुहो नाम रामस्यात्मसमः सखा}
{निषादजात्यो बलवान्स्थपतिश्चेति विश्रुतः} %||2-44-9||

\twolineshloka
{स श्रुत्वा पुरुषव्याघ्रं रामं विषयमागतम्}
{वृद्धैः परिवृतोऽमात्यैर्ज्ञातिभिश्चाप्युपागतः} %||2-44-10||

\twolineshloka
{ततो निषादाधिपतिं दृष्ट्वा दूरादवस्थितम्}
{सह सौमित्रिणा रामः समागच्छद्गुहेन सः} %||2-44-11||

\twolineshloka
{तमार्तः सम्परिष्वज्य गुहो राघवमब्रवीत्}
{यथायोध्या तथेदं ते राम किं करवाणि ते} %||2-44-12||

\twolineshloka
{ततो गुणवदन्नाद्यमुपादाय पृथग्विधम्}
{अर्घ्यं चोपानयत्क्षिप्रं वाक्यं चेदमुवाच ह} %||2-44-13||

\twolineshloka
{स्वागतं ते महाबाहो तवेयमखिला मही}
{वयं प्रेष्या भवान्भर्ता साधु राज्यं प्रशाधि नः} %||2-44-14||

\twolineshloka
{भक्ष्यं भोज्यं च पेयं च लेह्यं चेदमुपस्थितम्}
{शयनानि च मुख्यानि वाजिनां खादनं च ते} %||2-44-15||

\twolineshloka
{गुहमेव ब्रुवाणं तं राघवः प्रत्युवाच ह}
{अर्चिताश्चैव हृष्टाश्च भवता सर्वथा वयम्} %||2-44-16||

\twolineshloka
{पद्भ्यामभिगमाच्चैव स्नेहसन्दर्शनेन च}
{भुजाभ्यां साधुवृत्ताभ्यां पीडयन्वाक्यमब्रवीत्} %||2-44-17||

\twolineshloka
{दिष्ट्या त्वां गुह पश्यामि अरोगं सह बान्धवैः}
{अपि ते कूशलं राष्ट्रे मित्रेषु च धनेषु च} %||2-44-18||

\twolineshloka
{यत्त्विदं भवता किञ्चित्प्रीत्या समुपकल्पितम्}
{सर्वं तदनुजानामि न हि वर्ते प्रतिग्रहे} %||2-44-19||

\twolineshloka
{कुशचीराजिनधरं फलमूलाशनं च माम्}
{विद्धि प्रणिहितं धर्मे तापसं वनगोचरम्} %||2-44-20||

\twolineshloka
{अश्वानां खादनेनाहमर्थी नान्येन केनचित्}
{एतावतात्रभवता भविष्यामि सुपूजितः} %||2-44-21||

\twolineshloka
{एते हि दयिता राज्ञः पितुर्दशरथस्य मे}
{एतैः सुविहितैरश्वैर्भविष्याम्यहमर्चितः} %||2-44-22||

\twolineshloka
{अश्वानां प्रतिपानं च खादनं चैव सोऽन्वशात्}
{गुहस्तत्रैव पुरुषांस्त्वरितं दीयतामिति} %||2-44-23||

\twolineshloka
{ततश्चीरोत्तरासङ्गः सन्ध्यामन्वास्य पश्चिमाम्}
{जलमेवाददे भोज्यं लक्ष्मणेनाहृतं स्वयम्} %||2-44-24||

\twolineshloka
{तस्य भूमौ शयानस्य पादौ प्रक्षाल्य लक्ष्मणः}
{सभार्यस्य ततोऽभ्येत्य तस्थौ वृक्षमुपाश्रितः} %||2-44-25||

\twolineshloka
{गुहोऽपि सह सूतेन सौमित्रिमनुभाषयन्}
{अन्वजाग्रत्ततो राममप्रमत्तो धनुर्धरः} %||2-44-26||

\fourlineindentedshloka
{तथा शयानस्य ततोऽस्य धीमतो}
{ यशस्विनो दाशरथेर्महात्मनः}
{अदृष्टदुःखस्य सुखोचितस्य सा}
{ तदा व्यतीयाय चिरेण शर्वरी} %||2-45-27||


{॥इत्यार्षे श्रीमद्रामायणे वाल्मीकीये आदिकाव्ये अयोध्याकाण्डे चतुश्चत्वारिंशः सर्गः॥}

\sect{पञ्चचत्वारिंशः सर्गः}

\twolineshloka
{तं जाग्रतमदम्भेन भ्रातुरर्थाय लक्ष्मणम्}
{गुहः सन्तापसन्तप्तो राघवं वाक्यमब्रवीत्} %||2-45-1||

\twolineshloka
{इयं तात सुखा शय्या त्वदर्थमुपकल्पिता}
{प्रत्याश्वसिहि साध्वस्यां राजपुत्र यथासुखम्} %||2-45-2||

\twolineshloka
{उचितोऽयं जनः सर्वः क्लेशानां त्वं सुखोचितः}
{गुप्त्यर्थं जागरिष्यामः काकुत्स्थस्य वयं निशाम्} %||2-45-3||

\twolineshloka
{न हि रामात्प्रियतरो ममास्ति भुवि कश्चन}
{ब्रवीम्येतदहं सत्यं सत्येनैव च ते शपे} %||2-45-4||

\twolineshloka
{अस्य प्रसादादाशंसे लोकेऽस्मिन्सुमहद्यशः}
{धर्मावाप्तिं च विपुलामर्थावाप्तिं च केवलाम्} %||2-45-5||

\twolineshloka
{सोऽहं प्रियसखं रामं शयानं सह सीतया}
{रक्षिष्यामि धनुष्पाणिः सर्वतो ज्ञातिभिः सह} %||2-45-6||

\twolineshloka
{न हि मेऽविदितं किञ्चिद्वनेऽस्मिंश्चरतः सदा}
{चतुरङ्गं ह्यपि बलं सुमहत्प्रसहेमहि} %||2-45-7||

\twolineshloka
{लक्ष्मणस्तं तदोवाच रक्ष्यमाणास्त्वयानघ}
{नात्र भीता वयं सर्वे धर्ममेवानुपश्यता} %||2-45-8||

\twolineshloka
{कथं दाशरथौ भूमौ शयाने सह सीतया}
{शक्या निद्रा मया लब्धुं जीवितं वा सुखानि वा} %||2-45-9||

\twolineshloka
{यो न देवासुरैः सर्वैः शक्यः प्रसहितुं युधि}
{तं पश्य सुखसंविष्टं तृणेषु सह सीतया} %||2-45-10||

\twolineshloka
{यो मन्त्र तपसा लब्धो विविधैश्च परिश्रमैः}
{एको दशरथस्यैष पुत्रः सदृशलक्षणः} %||2-45-11||

\twolineshloka
{अस्मिन्प्रव्रजितो राजा न चिरं वर्तयिष्यति}
{विधवा मेदिनी नूनं क्षिप्रमेव भविष्यति} %||2-45-12||

\twolineshloka
{विनद्य सुमहानादं श्रमेणोपरताः स्त्रियः}
{निर्घोषोपरतं तात मन्ये राजनिवेशनम्} %||2-45-13||

\twolineshloka
{कौसल्या चैव राजा च तथैव जननी मम}
{नाशंसे यदि जीवन्ति सर्वे ते शर्वरीमिमाम्} %||2-45-14||

\twolineshloka
{जीवेदपि हि मे माता शत्रुघ्नस्यान्ववेक्षया}
{तद्दुःखं यत्तु कौसल्या वीरसूर्विनशिष्यति} %||2-45-15||

\twolineshloka
{अनुरक्तजनाकीर्णा सुखालोकप्रियावहा}
{राजव्यसनसंसृष्टा सा पुरी विनशिष्यति} %||2-45-16||

\twolineshloka
{अतिक्रान्तमतिक्रान्तमनवाप्य मनोरथम्}
{राज्ये राममनिक्षिप्य पिता मे विनशिष्यति} %||2-45-17||

\twolineshloka
{सिद्धार्थाः पितरं वृत्तं तस्मिन्काले ह्युपस्थिते}
{प्रेतकार्येषु सर्वेषु संस्करिष्यन्ति भूमिपम्} %||2-45-18||

\twolineshloka
{रम्यचत्वरसंस्थानां सुविभक्तमहापथाम्}
{हर्म्यप्रासादसम्पन्नां गणिकावरशोभिताम्} %||2-45-19||

\twolineshloka
{रथाश्वगजसम्बाधां तूर्यनादविनादिताम्}
{सर्वकल्याणसम्पूर्णां हृष्टपुष्टजनाकुलाम्} %||2-45-20||

\twolineshloka
{आरामोद्यानसम्पन्नां समाजोत्सवशालिनीम्}
{सुखिता विचरिष्यन्ति राजधानीं पितुर्मम} %||2-45-21||

\twolineshloka
{अपि सत्यप्रतिज्ञेन सार्धं कुशलिना वयम्}
{निवृत्ते वनवासेऽस्मिन्नयोध्यां प्रविशेमहि} %||2-45-22||

\twolineshloka
{परिदेवयमानस्य दुःखार्तस्य महात्मनः}
{तिष्ठतो राजपुत्रस्य शर्वरी सात्यवर्तत} %||2-45-23||

\fourlineindentedshloka
{तथा हि सत्यं ब्रुवति प्रजाहिते}
{ नरेन्द्रपुत्रे गुरुसौहृदाद्गुहः}
{मुमोच बाष्पं व्यसनाभिपीडितो}
{ ज्वरातुरो नाग इव व्यथातुरः} %||2-46-24||


{॥इत्यार्षे श्रीमद्रामायणे वाल्मीकीये आदिकाव्ये अयोध्याकाण्डे पञ्चचत्वारिंशः सर्गः॥}

\sect{षट्चत्वारिंशः सर्गः}

\twolineshloka
{प्रभातायां तु शर्वर्यां पृथु वृक्षा महायशाः}
{उवाच रामः सौमित्रिं लक्ष्मणं शुभलक्षणम्} %||2-46-1||

\twolineshloka
{भास्करोदयकालोऽयं गता भगवती निशा}
{असौ सुकृष्णो विहगः कोकिलस्तात कूजति} %||2-46-2||

\twolineshloka
{बर्हिणानां च निर्घोषः श्रूयते नदतां वने}
{तराम जाह्नवीं सौम्य शीघ्रगां सागरङ्गमाम्} %||2-46-3||

\twolineshloka
{विज्ञाय रामस्य वचः सौमित्रिर्मित्रनन्दनः}
{गुहमामन्त्र्य सूतं च सोऽतिष्ठद्भ्रातुरग्रतः} %||2-46-4||

\twolineshloka
{ततः कलापान्संनह्य खड्गौ बद्ध्वा च धन्विनौ}
{जग्मतुर्येन तौ गङ्गां सीतया सह राघवौ} %||2-46-5||

\twolineshloka
{राममेव तु धर्मज्ञमुपगम्य विनीतवत्}
{किमहं करवाणीति सूतः प्राञ्जलिरब्रवीत्} %||2-46-6||

\twolineshloka
{निवर्तस्वेत्युवाचैनमेतावद्धि कृतं मम}
{यानं विहाय पद्भ्यां तु गमिष्यामो महावनम्} %||2-46-7||

\twolineshloka
{आत्मानं त्वभ्यनुज्ञातमवेक्ष्यार्तः स सारथिः}
{सुमन्त्रः पुरुषव्याघ्रमैक्ष्वाकमिदमब्रवीत्} %||2-46-8||

\twolineshloka
{नातिक्रान्तमिदं लोके पुरुषेणेह केनचित्}
{तव सभ्रातृभार्यस्य वासः प्राकृतवद्वने} %||2-46-9||

\twolineshloka
{न मन्ये ब्रह्मचर्येऽस्ति स्वधीते वा फलोदयः}
{मार्दवार्जवयोर्वापि त्वां चेद्व्यसनमागतम्} %||2-46-10||

\twolineshloka
{सह राघव वैदेह्या भ्रात्रा चैव वने वसन्}
{त्वं गतिं प्राप्स्यसे वीर त्रीँल्लोकांस्तु जयन्निव} %||2-46-11||

\twolineshloka
{वयं खलु हता राम ये तयाप्युपवञ्चिताः}
{कैकेय्या वशमेष्यामः पापाया दुःखभागिनः} %||2-46-12||

\twolineshloka
{इति ब्रुवन्नात्म समं सुमन्त्रः सारथिस्तदा}
{दृष्ट्वा दुर गतं रामं दुःखार्तो रुरुदे चिरम्} %||2-46-13||

\twolineshloka
{ततस्तु विगते बाष्पे सूतं स्पृष्टोदकं शुचिम्}
{रामस्तु मधुरं वाक्यं पुनः पुनरुवाच तम्} %||2-46-14||

\twolineshloka
{इक्ष्वाकूणां त्वया तुल्यं सुहृदं नोपलक्षये}
{यथा दशरथो राजा मां न शोचेत्तथा कुरु} %||2-46-15||

\twolineshloka
{शोकोपहत चेताश्च वृद्धश्च जगतीपतिः}
{काम भारावसन्नश्च तस्मादेतद्ब्रवीमि ते} %||2-46-16||

\twolineshloka
{यद्यदाज्ञापयेत्किञ्चित्स महात्मा महीपतिः}
{कैकेय्याः प्रियकामार्थं कार्यं तदविकाङ्क्षया} %||2-46-17||

\twolineshloka
{एतदर्थं हि राज्यानि प्रशासति नरेश्वराः}
{यदेषां सर्वकृत्येषु मनो न प्रतिहन्यते} %||2-46-18||

\twolineshloka
{तद्यथा स महाराजो नालीकमधिगच्छति}
{न च ताम्यति दुःखेन सुमन्त्र कुरु तत्तथा} %||2-46-19||

\twolineshloka
{अदृष्टदुःखं राजानं वृद्धमार्यं जितेन्द्रियम्}
{ब्रूयास्त्वमभिवाद्यैव मम हेतोरिदं वचः} %||2-46-20||

\twolineshloka
{नैवाहमनुशोचामि लक्ष्मणो न च मैथिली}
{अयोध्यायाश्च्युताश्चेति वने वत्स्यामहेति वा} %||2-46-21||

\twolineshloka
{चतुर्दशसु वर्षेषु निवृत्तेषु पुनः पुनः}
{लक्ष्मणं मां च सीतां च द्रक्ष्यसि क्षिप्रमागतान्} %||2-46-22||

\twolineshloka
{एवमुक्त्वा तु राजानं मातरं च सुमन्त्र मे}
{अन्याश्च देवीः सहिताः कैकेयीं च पुनः पुनः} %||2-46-23||

\twolineshloka
{आरोग्यं ब्रूहि कौसल्यामथ पादाभिवन्दनम्}
{सीताया मम चार्यस्य वचनाल्लक्ष्मणस्य च} %||2-46-24||

\twolineshloka
{ब्रूयाश्च हि महाराजं भरतं क्षिप्रमानय}
{आगतश्चापि भरतः स्थाप्यो नृपमते पदे} %||2-46-25||

\twolineshloka
{भरतं च परिष्वज्य यौवराज्येऽभिषिच्य च}
{अस्मत्सन्तापजं दुःखं न त्वामभिभविष्यति} %||2-46-26||

\twolineshloka
{भरतश्चापि वक्तव्यो यथा राजनि वर्तसे}
{तथा मातृषु वर्तेथाः सर्वास्वेवाविशेषतः} %||2-46-27||

\twolineshloka
{यथा च तव कैकेयी सुमित्रा चाविशेषतः}
{तथैव देवी कौसल्या मम माता विशेषतः} %||2-46-28||

\twolineshloka
{निवर्त्यमानो रामेण सुमन्त्रः शोककर्शितः}
{तत्सर्वं वचनं श्रुत्वा स्नेहात्काकुत्स्थमब्रवीत्} %||2-46-29||

\twolineshloka
{यदहं नोपचारेण ब्रूयां स्नेहादविक्लवः}
{भक्तिमानिति तत्तावद्वाक्यं त्वं क्षन्तुमर्हसि} %||2-46-30||

\twolineshloka
{कथं हि त्वद्विहीनोऽहं प्रतियास्यामि तां पुरीम्}
{तव तात वियोगेन पुत्रशोकाकुलामिव} %||2-46-31||

\twolineshloka
{सराममपि तावन्मे रथं दृष्ट्वा तदा जनः}
{विना रामं रथं दृष्ट्वा विदीर्येतापि सा पुरी} %||2-46-32||

\twolineshloka
{दैन्यं हि नगरी गच्छेद्दृष्ट्वा शून्यमिमं रथम्}
{सूतावशेषं स्वं सैन्यं हतवीरमिवाहवे} %||2-46-33||

\twolineshloka
{दूरेऽपि निवसन्तं त्वां मानसेनाग्रतः स्थितम्}
{चिन्तयन्त्योऽद्य नूनं त्वां निराहाराः कृताः प्रजाः} %||2-46-34||

\twolineshloka
{आर्तनादो हि यः पौरैर्मुक्तस्तद्विप्रवासने}
{रथस्थं मां निशाम्यैव कुर्युः शतगुणं ततः} %||2-46-35||

\twolineshloka
{अहं किं चापि वक्ष्यामि देवीं तव सुतो मया}
{नीतोऽसौ मातुलकुलं सन्तापं मा कृथा इति} %||2-46-36||

\twolineshloka
{असत्यमपि नैवाहं ब्रूयां वचनमीदृशम्}
{कथमप्रियमेवाहं ब्रूयां सत्यमिदं वचः} %||2-46-37||

\twolineshloka
{मम तावन्नियोगस्थास्त्वद्बन्धुजनवाहिनः}
{कथं रथं त्वया हीनं प्रवक्ष्यन्ति हयोत्तमाः} %||2-46-38||

\twolineshloka
{यदि मे याचमानस्य त्यागमेव करिष्यसि}
{सरथोऽग्निं प्रवेक्ष्यामि त्यक्त मात्र इह त्वया} %||2-46-39||

\twolineshloka
{भविष्यन्ति वने यानि तपोविघ्नकराणि ते}
{रथेन प्रतिबाधिष्ये तानि सत्त्वानि राघव} %||2-46-40||

\twolineshloka
{तत्कृतेन मया प्राप्तं रथ चर्या कृतं सुखम्}
{आशंसे त्वत्कृतेनाहं वनवासकृतं सुखम्} %||2-46-41||

\twolineshloka
{प्रसीदेच्छामि तेऽरण्ये भवितुं प्रत्यनन्तरः}
{प्रीत्याभिहितमिच्छामि भव मे पत्यनन्तरः} %||2-46-42||

\twolineshloka
{तव शुश्रूषणं मूर्ध्ना करिष्यामि वने वसन्}
{अयोध्यां देवलोकं वा सर्वथा प्रजहाम्यहम्} %||2-46-43||

\twolineshloka
{न हि शक्या प्रवेष्टुं सा मयायोध्या त्वया विना}
{राजधानी महेन्द्रस्य यथा दुष्कृतकर्मणा} %||2-46-44||

\twolineshloka
{इमे चापि हया वीर यदि ते वनवासिनः}
{परिचर्यां करिष्यन्ति प्राप्स्यन्ति परमां गतिम्} %||2-46-45||

\twolineshloka
{वनवासे क्षयं प्राप्ते ममैष हि मनोरथः}
{यदनेन रथेनैव त्वां वहेयं पुरीं पुनः} %||2-46-46||

\twolineshloka
{चतुर्दश हि वर्षाणि सहितस्य त्वया वने}
{क्षणभूतानि यास्यन्ति शतशस्तु ततोऽन्यथा} %||2-46-47||

\twolineshloka
{भृत्यवत्सल तिष्ठन्तं भर्तृपुत्रगते पथि}
{भक्तं भृत्यं स्थितं स्थित्यां त्वं न मां हातुमर्हसि} %||2-46-48||

\twolineshloka
{एवं बहुविधं दीनं याचमानं पुनः पुनः}
{रामो भृत्यानुकम्पी तु सुमन्त्रमिदमब्रवीत्} %||2-46-49||

\twolineshloka
{जानामि परमां भक्तिं मयि ते भर्तृवत्सल}
{शृणु चापि यदर्थं त्वां प्रेषयामि पुरीमितः} %||2-46-50||

\twolineshloka
{नगरीं त्वां गतं दृष्ट्वा जननी मे यवीयसी}
{कैकेयी प्रत्ययं गच्छेदिति रामो वनं गतः} %||2-46-51||

\twolineshloka
{परितुष्टा हि सा देवि वनवासं गते मयि}
{राजानं नातिशङ्केत मिथ्यावादीति धार्मिकम्} %||2-46-52||

\twolineshloka
{एष मे प्रथमः कल्पो यदम्बा मे यवीयसी}
{भरतारक्षितं स्फीतं पुत्रराज्यमवाप्नुयात्} %||2-46-53||

\twolineshloka
{मम प्रियार्थं राज्ञश्च सरथस्त्वं पुरीं व्रज}
{सन्दिष्टश्चासि यानर्थांस्तांस्तान्ब्रूयास्तथातथा} %||2-46-54||

\threelineshloka
{इत्युक्त्वा वचनं सूतं सान्त्वयित्वा पुनः पुनः}
{गुहं वचनमक्लीबं रामो हेतुमदब्रवीत्}
{जटाः कृत्वा गमिष्यामि न्यग्रोधक्षीरमानय} %||2-46-55||

\twolineshloka
{तत्क्षीरं राजपुत्राय गुहः क्षिप्रमुपाहरत्}
{लक्ष्मणस्यात्मनश्चैव रामस्तेनाकरोज्जटाः} %||2-46-56||

\twolineshloka
{तौ तदा चीरवसनौ जटामण्डलधारिणौ}
{अशोभेतामृषिसमौ भ्रातरौ रामलक्ष्मणौ} %||2-46-57||

\twolineshloka
{ततो वैखानसं मार्गमास्थितः सहलक्ष्मणः}
{व्रतमादिष्टवान्रामः सहायं गुहमब्रवीत्} %||2-46-58||

\twolineshloka
{अप्रमत्तो बले कोशे दुर्गे जनपदे तथा}
{भवेथा गुह राज्यं हि दुरारक्षतमं मतम्} %||2-46-59||

\twolineshloka
{ततस्तं समनुज्ञाय गुहमिक्ष्वाकुनन्दनः}
{जगाम तूर्णमव्यग्रः सभार्यः सहलक्ष्मणः} %||2-46-60||

\twolineshloka
{स तु दृष्ट्वा नदीतीरे नावमिक्ष्वाकुनन्दनः}
{तितीर्षुः शीघ्रगां गङ्गामिदं लक्ष्मणमब्रवीत्} %||2-46-61||

\twolineshloka
{आरोह त्वं नर व्याघ्र स्थितां नावमिमां शनैः}
{सीतां चारोपयान्वक्षं परिगृह्य मनस्विनीम्} %||2-46-62||

\twolineshloka
{स भ्रातुः शासनं श्रुत्वा सर्वमप्रतिकूलयन्}
{आरोप्य मैथिलीं पूर्वमारुरोहात्मवांस्ततः} %||2-46-63||

\twolineshloka
{अथारुरोह तेजस्वी स्वयं लक्ष्मणपूर्वजः}
{ततो निषादाधिपतिर्गुहो ज्ञातीनचोदयत्} %||2-46-64||

\twolineshloka
{अनुज्ञाय सुमन्त्रं च सबलं चैव तं गुहम्}
{आस्थाय नावं रामस्तु चोदयामास नाविकान्} %||2-46-65||

\twolineshloka
{ततस्तैश्चोदिता सा नौः कर्णधारसमाहिता}
{शुभस्फ्यवेगाभिहता शीघ्रं सलिलमत्यगात्} %||2-46-66||

\twolineshloka
{मध्यं तु समनुप्राप्य भागीरथ्यास्त्वनिन्दिता}
{वैदेही प्राञ्जलिर्भूत्वा तां नदीमिदमब्रवीत्} %||2-46-67||

\twolineshloka
{पुत्रो दशरथस्यायं महाराजस्य धीमतः}
{निदेशं पालयत्वेनं गङ्गे त्वदभिरक्षितः} %||2-46-68||

\twolineshloka
{चतुर्दश हि वर्षाणि समग्राण्युष्य कानने}
{भ्रात्रा सह मया चैव पुनः प्रत्यागमिष्यति} %||2-46-69||

\twolineshloka
{ततस्त्वां देवि सुभगे क्षेमेण पुनरागता}
{यक्ष्ये प्रमुदिता गङ्गे सर्वकामसमृद्धये} %||2-46-70||

\twolineshloka
{त्वं हि त्रिपथगा देवि ब्रह्म लोकं समीक्षसे}
{भार्या चोदधिराजस्य लोकेऽस्मिन्सम्प्रदृश्यसे} %||2-46-71||

\twolineshloka
{सा त्वां देवि नमस्यामि प्रशंसामि च शोभने}
{प्राप्त राज्ये नरव्याघ्र शिवेन पुनरागते} %||2-46-72||

\twolineshloka
{गवां शतसहस्राणि वस्त्राण्यन्नं च पेशलम्}
{ब्राह्मणेभ्यः प्रदास्यामि तव प्रियचिकीर्षया} %||2-46-73||

\twolineshloka
{तथा सम्भाषमाणा सा सीता गङ्गामनिन्दिता}
{दक्षिणा दक्षिणं तीरं क्षिप्रमेवाभ्युपागमत्} %||2-46-74||

\twolineshloka
{तीरं तु समनुप्राप्य नावं हित्वा नरर्षभः}
{प्रातिष्ठत सह भ्रात्रा वैदेह्या च परन्तपः} %||2-46-75||

\twolineshloka
{अथाब्रवीन्महाबाहुः सुमित्रानन्दवर्धनम्}
{अग्रतो गच्छ सौमित्रे सीता त्वामनुगच्छतु} %||2-46-76||

\twolineshloka
{पृष्ठतोऽहं गमिष्यामि त्वां च सीतां च पालयन्}
{अद्य दुःखं तु वैदेही वनवासस्य वेत्स्यति} %||2-46-77||

\fourlineindentedshloka
{गतं तु गङ्गापरपारमाशु}
{ रामं सुमन्त्रः प्रततं निरीक्ष्य}
{अध्वप्रकर्षाद्विनिवृत्तदृष्टि}
{र्मुमोच बाष्पं व्यथितस्तपस्वी} %||2-46-78||

\fourlineindentedshloka
{तौ तत्र हत्वा चतुरो महामृगा}
{न्वराहमृश्यं पृषतं महारुरुम्}
{आदाय मेध्यं त्वरितं बुभुक्षितौ}
{ वासाय काले ययतुर्वनस्पतिम्} %||2-47-79||


{॥इत्यार्षे श्रीमद्रामायणे वाल्मीकीये आदिकाव्ये अयोध्याकाण्डे षट्चत्वारिंशः सर्गः॥}

\sect{सप्तचत्वारिंशः सर्गः}

\twolineshloka
{स तं वृक्षं समासाद्य सन्ध्यामन्वास्य पश्चिमाम्}
{रामो रमयतां श्रेष्ठ इति होवाच लक्ष्मणम्} %||2-47-1||

\twolineshloka
{अद्येयं प्रथमा रात्रिर्याता जनपदाद्बहिः}
{या सुमन्त्रेण रहिता तां नोत्कण्ठितुमर्हसि} %||2-47-2||

\twolineshloka
{जागर्तव्यमतन्द्रिभ्यामद्य प्रभृति रात्रिषु}
{योगक्षेमो हि सीताया वर्तते लक्ष्मणावयोः} %||2-47-3||

\twolineshloka
{रात्रिं कथञ्चिदेवेमां सौमित्रे वर्तयामहे}
{उपावर्तामहे भूमावास्तीर्य स्वयमार्जितैः} %||2-47-4||

\twolineshloka
{स तु संविश्य मेदिन्यां महार्हशयनोचितः}
{इमाः सौमित्रये रामो व्याजहार कथाः शुभाः} %||2-47-5||

\twolineshloka
{ध्रुवमद्य महाराजो दुःखं स्वपिति लक्ष्मण}
{कृतकामा तु कैकेयी तुष्टा भवितुमर्हति} %||2-47-6||

\twolineshloka
{सा हि देवी महाराजं कैकेयी राज्यकारणात्}
{अपि न च्यावयेत्प्राणान्दृष्ट्वा भरतमागतम्} %||2-47-7||

\twolineshloka
{अनाथश्चैव वृद्धश्च मया चैव विनाकृतः}
{किं करिष्यति कामात्मा कैकेय्या वशमागतः} %||2-47-8||

\twolineshloka
{इदं व्यसनमालोक्य राज्ञश्च मतिविभ्रमम्}
{काम एवार्धधर्माभ्यां गरीयानिति मे मतिः} %||2-47-9||

\twolineshloka
{को ह्यविद्वानपि पुमान्प्रमदायाः कृते त्यजेत्}
{छन्दानुवर्तिनं पुत्रं तातो मामिव लक्ष्मण} %||2-47-10||

\twolineshloka
{सुखी बत सभार्यश्च भरतः केकयीसुतः}
{मुदितान्कोसलानेको यो भोक्ष्यत्यधिराजवत्} %||2-47-11||

\twolineshloka
{स हि सर्वस्य राज्यस्य मुखमेकं भविष्यति}
{ताते च वयसातीते मयि चारण्यमाश्रिते} %||2-47-12||

\twolineshloka
{अर्थधर्मौ परित्यज्य यः काममनुवर्तते}
{एवमापद्यते क्षिप्रं राजा दशरथो यथा} %||2-47-13||

\twolineshloka
{मन्ये दशरथान्ताय मम प्रव्राजनाय च}
{कैकेयी सौम्य सम्प्राप्ता राज्याय भरतस्य च} %||2-47-14||

\twolineshloka
{अपीदानीं न कैकेयी सौभाग्यमदमोहिता}
{कौसल्यां च सुमित्रां च सम्प्रबाधेत मत्कृते} %||2-47-15||

\twolineshloka
{मा स्म मत्कारणाद्देवी सुमित्रा दुःखमावसेत्}
{अयोध्यामित एव त्वं काले प्रविश लक्ष्मण} %||2-47-16||

\twolineshloka
{अहमेको गमिष्यामि सीतया सह दण्डकान्}
{अनाथाया हि नाथस्त्वं कौसल्याया भविष्यसि} %||2-47-17||

\twolineshloka
{क्षुद्रकर्मा हि कैकेयी द्वेषादन्याय्यमाचरेत्}
{परिदद्या हि धर्मज्ञे भरते मम मातरम्} %||2-47-18||

\twolineshloka
{नूनं जात्यन्तरे कस्मिं स्त्रियः पुत्रैर्वियोजिताः}
{जनन्या मम सौमित्रे तदप्येतदुपस्थितम्} %||2-47-19||

\twolineshloka
{मया हि चिरपुष्टेन दुःखसंवर्धितेन च}
{विप्रायुज्यत कौसल्या फलकाले धिगस्तु माम्} %||2-47-20||

\twolineshloka
{मा स्म सीमन्तिनी काचिज्जनयेत्पुत्रमीदृशम्}
{सौमित्रे योऽहमम्बाया दद्मि शोकमनन्तकम्} %||2-47-21||

\twolineshloka
{मन्ये प्रीतिविशिष्टा सा मत्तो लक्ष्मणसारिका}
{यस्यास्तच्छ्रूयते वाक्यं शुक पादमरेर्दश} %||2-47-22||

\twolineshloka
{शोचन्त्याश्चाल्पभाग्याया न किञ्चिदुपकुर्वता}
{पुर्त्रेण किमपुत्राया मया कार्यमरिन्दम} %||2-47-23||

\twolineshloka
{अल्पभाग्या हि मे माता कौसल्या रहिता मया}
{शेते परमदुःखार्ता पतिता शोकसागरे} %||2-47-24||

\twolineshloka
{एको ह्यहमयोध्यां च पृथिवीं चापि लक्ष्मण}
{तरेयमिषुभिः क्रुद्धो ननु वीर्यमकारणम्} %||2-47-25||

\twolineshloka
{अधर्मभय भीतश्च परलोकस्य चानघ}
{तेन लक्ष्मण नाद्याहमात्मानमभिषेचये} %||2-47-26||

\twolineshloka
{एतदन्यच्च करुणं विलप्य विजने बहु}
{अश्रुपूर्णमुखो रामो निशि तूष्णीमुपाविशत्} %||2-47-27||

\twolineshloka
{विलप्योपरतं रामं गतार्चिषमिवानलम्}
{समुद्रमिव निर्वेगमाश्वासयत लक्ष्मणः} %||2-47-28||

\twolineshloka
{ध्रुवमद्य पुरी राम अयोध्या युधिनां वर}
{निष्प्रभा त्वयि निष्क्रान्ते गतचन्द्रेव शर्वरी} %||2-47-29||

\twolineshloka
{नैतदौपयिकं राम यदिदं परितप्यसे}
{विषादयसि सीतां च मां चैव पुरुषर्षभ} %||2-47-30||

\twolineshloka
{न च सीता त्वया हीना न चाहमपि राघव}
{मुहूर्तमपि जीवावो जलान्मत्स्याविवोद्धृतौ} %||2-47-31||

\twolineshloka
{न हि तातं न शत्रुघ्नं न सुमित्रां परन्तप}
{द्रष्टुमिच्छेयमद्याहं स्वर्गं वापि त्वया विना} %||2-47-32||

\fourlineindentedshloka
{स लक्ष्मणस्योत्तम पुष्कलं वचो}
{ निशम्य चैवं वनवासमादरात्}
{समाः समस्ता विदधे परन्तपः}
{ प्रपद्य धर्मं सुचिराय राघवः} %||2-48-33||


{॥इत्यार्षे श्रीमद्रामायणे वाल्मीकीये आदिकाव्ये अयोध्याकाण्डे सप्तचत्वारिंशः सर्गः॥}

\sect{अष्टचत्वारिंशः सर्गः}

\twolineshloka
{ते तु तस्मिन्महावृक्ष उषित्वा रजनीं शिवाम्}
{विमलेऽभ्युदिते सूर्ये तस्माद्देशात्प्रतस्थिरे} %||2-48-1||

\twolineshloka
{यत्र भागीरथी गङ्गा यमुनामभिवर्तते}
{जग्मुस्तं देशमुद्दिश्य विगाह्य सुमहद्वनम्} %||2-48-2||

\twolineshloka
{ते भूमिमागान्विविधान्देशांश्चापि मनोरमान्}
{अदृष्टपूर्वान्पश्यन्तस्तत्र तत्र यशस्विनः} %||2-48-3||

\twolineshloka
{यथाक्षेमेण गच्छन्स पश्यंश्च विविधान्द्रुमान्}
{निवृत्तमात्रे दिवसे रामः सौमित्रिमब्रवीत्} %||2-48-4||

\twolineshloka
{प्रयागमभितः पश्य सौमित्रे धूममुन्नतम्}
{अग्नेर्भगवतः केतुं मन्ये संनिहितो मुनिः} %||2-48-5||

\twolineshloka
{नूनं प्राप्ताः स्म सम्भेदं गङ्गायमुनयोर्वयम्}
{तथा हि श्रूयते शम्ब्दो वारिणा वारिघट्टितः} %||2-48-6||

\twolineshloka
{दारूणि परिभिन्नानि वनजैरुपजीविभिः}
{भरद्वाजाश्रमे चैते दृश्यन्ते विविधा द्रुमाः} %||2-48-7||

\twolineshloka
{धन्विनौ तौ सुखं गत्वा लम्बमाने दिवाकरे}
{गङ्गायमुनयोः सन्धौ प्रापतुर्निलयं मुनेः} %||2-48-8||

\twolineshloka
{रामस्त्वाश्रममासाद्य त्रासयन्मृगपक्षिणः}
{गत्वा मुहूर्तमध्वानं भरद्वाजमुपागमत्} %||2-48-9||

\twolineshloka
{ततस्त्वाश्रममासाद्य मुनेर्दर्शनकाङ्क्षिणौ}
{सीतयानुगतौ वीरौ दूरादेवावतस्थतुः} %||2-48-10||

\twolineshloka
{हुताग्निहोत्रं दृष्ट्वैव महाभागं कृताञ्जलिः}
{रामः सौमित्रिणा सार्धं सीतया चाभ्यवादयत्} %||2-48-11||

\twolineshloka
{न्यवेदयत चात्मानं तस्मै लक्ष्मणपूर्वजः}
{पुत्रौ दशरथस्यावां भगवन्रामलक्ष्मणौ} %||2-48-12||

\twolineshloka
{भार्या ममेयं वैदेही कल्याणी जनकात्मजा}
{मां चानुयाता विजनं तपोवनमनिन्दिता} %||2-48-13||

\twolineshloka
{पित्रा प्रव्राज्यमानं मां सौमित्रिरनुजः प्रियः}
{अयमन्वगमद्भ्राता वनमेव दृढव्रतः} %||2-48-14||

\twolineshloka
{पित्रा नियुक्ता भगवन्प्रवेष्यामस्तपोवनम्}
{धर्ममेवाचरिष्यामस्तत्र मूलफलाशनाः} %||2-48-15||

\twolineshloka
{तस्य तद्वचनं श्रुत्वा राजपुत्रस्य धीमतः}
{उपानयत धर्मात्मा गामर्घ्यमुदकं ततः} %||2-48-16||

\twolineshloka
{मृगपक्षिभिरासीनो मुनिभिश्च समन्ततः}
{राममागतमभ्यर्च्य स्वागतेनाह तं मुनिः} %||2-48-17||

\twolineshloka
{प्रतिगृह्य च तामर्चामुपविष्टं सराघवम्}
{भरद्वाजोऽब्रवीद्वाक्यं धर्मयुक्तमिदं तदा} %||2-48-18||

\twolineshloka
{चिरस्य खलु काकुत्स्थ पश्यामि त्वामिहागतम्}
{श्रुतं तव मया चेदं विवासनमकारणम्} %||2-48-19||

\twolineshloka
{अवकाशो विविक्तोऽयं महानद्योः समागमे}
{पुण्यश्च रमणीयश्च वसत्विह भगान्सुखम्} %||2-48-20||

\twolineshloka
{एवमुक्तस्तु वचनं भरद्वाजेन राघवः}
{प्रत्युवाच शुभं वाक्यं रामः सर्वहिते रतः} %||2-48-21||

\threelineshloka
{भगवन्नित आसन्नः पौरजानपदो जनः}
{आगमिष्यति वैदेहीं मां चापि प्रेक्षको जनः}
{अनेन कारणेनाहमिह वासं न रोचये} %||2-48-22||

\twolineshloka
{एकान्ते पश्य भगवन्नाश्रमस्थानमुत्तमम्}
{रमते यत्र वैदेही सुखार्हा जनकात्मजा} %||2-48-23||

\twolineshloka
{एतच्छ्रुत्वा शुभं वाक्यं भरद्वाजो महामुनिः}
{राघवस्य ततो वाक्यमर्थ ग्राहकमब्रवीत्} %||2-48-24||

\twolineshloka
{दशक्रोश इतस्तात गिरिर्यस्मिन्निवत्स्यसि}
{महर्षिसेवितः पुण्यः सर्वतः सुख दर्शनः} %||2-48-25||

\twolineshloka
{गोलाङ्गूलानुचरितो वानरर्क्षनिषेवितः}
{चित्रकूट इति ख्यातो गन्धमादनसंनिभः} %||2-48-26||

\twolineshloka
{यावता चित्र कूटस्य नरः शृङ्गाण्यवेक्षते}
{कल्याणानि समाधत्ते न पापे कुरुते मनः} %||2-48-27||

\twolineshloka
{ऋषयस्तत्र बहवो विहृत्य शरदां शतम्}
{तपसा दिवमारूढाः कपालशिरसा सह} %||2-48-28||

\twolineshloka
{प्रविविक्तमहं मन्ये तं वासं भवतः सुखम्}
{इह वा वनवासाय वस राम मया सह} %||2-48-29||

\twolineshloka
{स रामं सर्वकामैस्तं भरद्वाजः प्रियातिथिम्}
{सभार्यं सह च भ्रात्रा प्रतिजग्राह धर्मवित्} %||2-48-30||

\twolineshloka
{तस्य प्रयागे रामस्य तं महर्षिमुपेयुषः}
{प्रपन्ना रजनी पुण्या चित्राः कथयतः कथाः} %||2-48-31||

\twolineshloka
{प्रभातायां रजन्यां तु भरद्वाजमुपागमत्}
{उवाच नरशार्दूलो मुनिं ज्वलिततेजसम्} %||2-48-32||

\twolineshloka
{शर्वरीं भवनन्नद्य सत्यशील तवाश्रमे}
{उषिताः स्मेह वसतिमनुजानातु नो भवान्} %||2-48-33||

\twolineshloka
{रात्र्यां तु तस्यां व्युष्टायां भरद्वाजोऽब्रवीदिदम्}
{मधुमूलफलोपेतं चित्रकूटं व्रजेति ह} %||2-48-34||

\twolineshloka
{तत्र कुञ्जरयूथानि मृगयूथानि चाभितः}
{विचरन्ति वनान्तेषु तानि द्रक्ष्यसि राघव} %||2-48-35||

\fourlineindentedshloka
{प्रहृष्टकोयष्टिककोकिलस्वनै}
{र्विनादितं तं वसुधाधरं शिवम्}
{मृगैश्च मत्तैर्बहुभिश्च कुञ्जरैः}
{ सुरम्यमासाद्य समावसाश्रमम्} %||2-49-36||


{॥इत्यार्षे श्रीमद्रामायणे वाल्मीकीये आदिकाव्ये अयोध्याकाण्डे अष्टचत्वारिंशः सर्गः॥}

\sect{एकोनपञ्चाशत्तमः सर्गः}

\twolineshloka
{उषित्वा रजनीं तत्र राजपुत्रावरिन्दमौ}
{महर्षिमभिवाद्याथ जग्मतुस्तं गिरिं प्रति} %||2-49-1||

\twolineshloka
{प्रस्थितांश्चैव तान्प्रेक्ष्य पिता पुत्रानिवान्वगात्}
{ततः प्रचक्रमे वक्तुं वचनं स महामुनिः} %||2-49-2||

\twolineshloka
{अथासाद्य तु कालिन्दीं शीघ्रस्रोतसमापगाम्}
{तत्र यूयं प्लवं कृत्वा तरतांशुमतीं नदीम्} %||2-49-3||

\twolineshloka
{ततो न्यग्रोधमासाद्य महान्तं हरितच्छदम्}
{विवृद्धं बहुभिर्वृक्षैः श्यामं सिद्धोपसेवितम्} %||2-49-4||

\twolineshloka
{क्रोशमात्रं ततो गत्वा नीलं द्रक्ष्यथ काननम्}
{पलाशबदरीमिश्रं राम वंशैश्च यामुनैः} %||2-49-5||

\threelineshloka
{स पन्थाश्चित्रकूटस्य गतः सुबहुशो मया}
{रम्यो मार्दवयुक्तश्च वनदावैर्विवर्जितः}
{इति पन्थानमावेद्य महर्षिः स न्यवर्तत} %||2-49-6||

\twolineshloka
{उपावृत्ते मुनौ तस्मिन्रामो लक्ष्मणमब्रवीत्}
{कृतपुण्याः स्म सौमित्रे मुनिर्यन्नोऽनुकम्पते} %||2-49-7||

\twolineshloka
{इति तौ पुरुषव्याघ्रौ मन्त्रयित्वा मनस्विनौ}
{सीतामेवाग्रतः कृत्वा कालिन्दीं जग्मतुर्नदीम्} %||2-49-8||

\twolineshloka
{तौ काष्ठसङ्घाटमथो चक्रतुः सुमहाप्लवम्}
{चकार लक्ष्मणश्छित्त्वा सीतायाः सुखमानसम्} %||2-49-9||

\twolineshloka
{तत्र श्रियमिवाचिन्त्यां रामो दाशरथिः प्रियाम्}
{ईषत्संलज्जमानां तामध्यारोपयत प्लवम्} %||2-49-10||

\twolineshloka
{ततः प्लवेनांशुमतीं शीघ्रगामूर्मिमालिनीम्}
{तीरजैर्बहुभिर्वृक्षैः सन्तेरुर्यमुनां नदीम्} %||2-49-11||

\twolineshloka
{ते तीर्णाः प्लवमुत्सृज्य प्रस्थाय यमुनावनात्}
{श्यामं न्यग्रोधमासेदुः शीतलं हरितच्छदम्} %||2-49-12||

\twolineshloka
{कौसल्यां चैव पश्येयं सुमित्रां च यशस्विनीम्}
{इति सीताञ्जलिं कृत्वा पर्यगछद्वनस्पतिम्} %||2-49-13||

\twolineshloka
{क्रोशमात्रं ततो गत्वा भ्रातरौ रामलक्ष्मणौ}
{बहून्मेध्यान्मृगान्हत्वा चेरतुर्यमुनावने} %||2-49-14||

\fourlineindentedshloka
{विहृत्य ते बर्हिणपूगनादिते}
{ शुभे वने वारणवानरायुते}
{समं नदीवप्रमुपेत्य सम्मतम्}
{ निवासमाजग्मुरदीनदर्शनः} %||2-50-15||


{॥इत्यार्षे श्रीमद्रामायणे वाल्मीकीये आदिकाव्ये अयोध्याकाण्डे एकोनपञ्चाशत्तमः सर्गः॥}

\sect{पञ्चाशत्तमः सर्गः}

\twolineshloka
{अथ रात्र्यां व्यतीतायामवसुप्तमनन्तरम्}
{प्रबोधयामास शनैर्लक्ष्मणं रघुनन्दनः} %||2-50-1||

\twolineshloka
{सौमित्रे शृणु वन्यानां वल्गु व्याहरतां स्वनम्}
{सम्प्रतिष्ठामहे कालः प्रस्थानस्य परन्तप} %||2-50-2||

\twolineshloka
{स सुप्तः समये भ्रात्रा लक्ष्मणः प्रतिबोधितः}
{जहौ निद्रां च तन्द्रीं च प्रसक्तं च पथि श्रमम्} %||2-50-3||

\twolineshloka
{तत उत्थाय ते सर्वे स्पृष्ट्वा नद्याः शिवं जलम्}
{पन्थानमृषिणोद्दिष्टं चित्रकूटस्य तं ययुः} %||2-50-4||

\twolineshloka
{ततः सम्प्रस्थितः काले रामः सौमित्रिणा सह}
{सीतां कमलपत्राक्षीमिदं वचनमब्रवीत्} %||2-50-5||

\twolineshloka
{आदीप्तानिव वैदेहि सर्वतः पुष्पितान्नगान्}
{स्वैः पुष्पैः किंशुकान्पश्य मालिनः शिशिरात्यये} %||2-50-6||

\twolineshloka
{पश्य भल्लातकान्फुल्लान्नरैरनुपसेवितान्}
{फलपत्रैरवनतान्नूनं शक्ष्यामि जीवितुम्} %||2-50-7||

\twolineshloka
{पश्य द्रोणप्रमाणानि लम्बमानानि लक्ष्मण}
{मधूनि मधुकारीभिः सम्भृतानि नगे नगे} %||2-50-8||

\twolineshloka
{एष क्रोशति नत्यूहस्तं शिखी प्रतिकूजति}
{रमणीये वनोद्देशे पुष्पसंस्तरसङ्कटे} %||2-50-9||

\twolineshloka
{मातङ्गयूथानुसृतं पक्षिसङ्घानुनादितम्}
{चित्रकूटमिमं पश्य प्रवृद्धशिखरं गिरिम्} %||2-50-10||

\twolineshloka
{ततस्तौ पादचारेण गच्छन्तौ सह सीतया}
{रम्यमासेदतुः शैलं चित्रकूटं मनोरमम्} %||2-50-11||

\twolineshloka
{तं तु पर्वतमासाद्य नानापक्षिगणायुतम्}
{अयं वासो भवेत्तावदत्र सौम्य रमेमहि} %||2-50-12||

\twolineshloka
{लक्ष्मणानय दारूणि दृढानि च वराणि च}
{कुरुष्वावसथं सौम्य वासे मेऽभिरतं मनः} %||2-50-13||

\twolineshloka
{तस्य तद्वचनं श्रुत्वा सौमित्रिर्विविधान्द्रुमान्}
{आजहार ततश्चक्रे पर्ण शालामरिं दम} %||2-50-14||

\twolineshloka
{शुश्रूषमाणमेकाग्रमिदं वचनमब्रवीत्}
{ऐणेयं मांसमाहृत्य शालां यक्ष्यामहे वयम्} %||2-50-15||

\twolineshloka
{स लक्ष्मणः कृष्णमृगं हत्वा मेध्यं पतापवान्}
{अथ चिक्षेप सौमित्रिः समिद्धे जातवेदसि} %||2-50-16||

\twolineshloka
{तं तु पक्वं समाज्ञाय निष्टप्तं छिन्नशोणितम्}
{लक्ष्मणः पुरुषव्याघ्रमथ राघवमब्रवीत्} %||2-50-17||

\twolineshloka
{अयं कृष्णः समाप्ताङ्गः शृतः कृष्ण मृगो यथा}
{देवता देवसङ्काश यजस्व कुशलो ह्यसि} %||2-50-18||

\twolineshloka
{रामः स्नात्वा तु नियतो गुणवाञ्जप्यकोविदः}
{पापसंशमनं रामश्चकार बलिमुत्तमम्} %||2-50-19||

\fourlineindentedshloka
{तां वृक्षपर्णच्छदनां मनोज्ञाम्}
{ यथाप्रदेशं सुकृतां निवाताम्}
{वासाय सर्वे विविशुः समेताः}
{ सभां यथा देव गणाः सुधर्माम्} %||2-50-20||

\fourlineindentedshloka
{अनेकनानामृगपक्षिसङ्कुले}
{ विचित्रपुष्पस्तबलैर्द्रुमैर्युते}
{वनोत्तमे व्यालमृगानुनादिते}
{ तथा विजह्रुः सुसुखं जितेन्द्रियाः} %||2-50-21||

\fourlineindentedshloka
{सुरम्यमासाद्य तु चित्रकूटम्}
{ नदीं च तां माल्यवतीं सुतीर्थाम्}
{ननन्द हृष्टो मृगपक्षिजुष्टाम्}
{ जहौ च दुःखं पुरविप्रवासात्} %||2-51-22||


{॥इत्यार्षे श्रीमद्रामायणे वाल्मीकीये आदिकाव्ये अयोध्याकाण्डे पञ्चाशत्तमः सर्गः॥}

\sect{एकपञ्चाशत्तमः सर्गः}

\twolineshloka
{कथयित्वा सुदुःखार्तः सुमन्त्रेण चिरं सह}
{रामे दक्षिण कूलस्थे जगाम स्वगृहं गुहः} %||2-51-1||

\twolineshloka
{अनुज्ञातः सुमन्त्रोऽथ योजयित्वा हयोत्तमान्}
{अयोध्यामेव नगरीं प्रययौ गाढदुर्मनाः} %||2-51-2||

\twolineshloka
{स वनानि सुगन्धीनि सरितश्च सरांसि च}
{पश्यन्नतिययौ शीघ्रं ग्रामाणि नगराणि च} %||2-51-3||

\twolineshloka
{ततः सायाह्नसमये तृतीयेऽहनि सारथिः}
{अयोध्यां समनुप्राप्य निरानन्दां ददर्श ह} %||2-51-4||

\twolineshloka
{स शून्यामिव निःशब्दां दृष्ट्वा परमदुर्मनाः}
{सुमन्त्रश्चिन्तयामास शोकवेगसमाहतः} %||2-51-5||

\threelineshloka
{कच्चिन्न सगजा साश्वा सजना सजनाधिपा}
{राम सन्तापदुःखेन दग्धा शोकाग्निना पुरी}
{इति चिन्तापरः सूतस्त्वरितः प्रविवेश ह} %||2-51-6||

\twolineshloka
{सुमन्त्रमभियान्तं तं शतशोऽथ सहस्रशः}
{क्व राम इति पृच्छन्तः सूतमभ्यद्रवन्नराः} %||2-51-7||

\twolineshloka
{तेषां शशंस गङ्गायामहमापृच्छ्य राघवम्}
{अनुज्ञातो निवृत्तोऽस्मि धार्मिकेण महात्मना} %||2-51-8||

\twolineshloka
{ते तीर्णा इति विज्ञाय बाष्पपूर्णमुखा जनाः}
{अहो धिगिति निःश्वस्य हा रामेति च चुक्रुशुः} %||2-51-9||

\twolineshloka
{शुश्राव च वचस्तेषां वृन्दं वृन्दं च तिष्ठताम्}
{हताः स्म खलु ये नेह पश्याम इति राघवम्} %||2-51-10||

\twolineshloka
{दानयज्ञविवाहेषु समाजेषु महत्सु च}
{न द्रक्ष्यामः पुनर्जातु धार्मिकं राममन्तरा} %||2-51-11||

\twolineshloka
{किं समर्थं जनस्यास्य किं प्रियं किं सुखावहम्}
{इति रामेण नगरं पितृवत्परिपालितम्} %||2-51-12||

\twolineshloka
{वातायनगतानां च स्त्रीणामन्वन्तरापणम्}
{रामशोकाभितप्तानां शुश्राव परिदेवनम्} %||2-51-13||

\twolineshloka
{स राजमार्गमध्येन सुमन्त्रः पिहिताननः}
{यत्र राजा दशरथस्तदेवोपययौ गृहम्} %||2-51-14||

\twolineshloka
{सोऽवतीर्य रथाच्छीघ्रं राजवेश्म प्रविश्य च}
{कक्ष्याः सप्ताभिचक्राम महाजनसमाकुलाः} %||2-51-15||

\twolineshloka
{ततो दशरथस्त्रीणां प्रासादेभ्यस्ततस्ततः}
{रामशोकाभितप्तानां मन्दं शुश्राव जल्पितम्} %||2-51-16||

\twolineshloka
{सह रामेण निर्यातो विना राममिहागतः}
{सूतः किं नाम कौसल्यां शोचन्तीं प्रतिवक्ष्यति} %||2-51-17||

\twolineshloka
{यथा च मन्ये दुर्जीवमेवं न सुकरं ध्रुवम्}
{आच्छिद्य पुत्रे निर्याते कौसल्या यत्र जीवति} %||2-51-18||

\twolineshloka
{सत्य रूपं तु तद्वाक्यं राज्ञः स्त्रीणां निशामयन्}
{प्रदीप्तमिव शोकेन विवेश सहसा गृहम्} %||2-51-19||

\twolineshloka
{स प्रविश्याष्टमीं कक्ष्यां राजानं दीनमातुलम्}
{पुत्रशोकपरिद्यूनमपश्यत्पाण्डरे गृहे} %||2-51-20||

\twolineshloka
{अभिगम्य तमासीनं नरेन्द्रमभिवाद्य च}
{सुमन्त्रो रामवचनं यथोक्तं प्रत्यवेदयत्} %||2-51-21||

\twolineshloka
{स तूष्णीमेव तच्छ्रुत्वा राजा विभ्रान्त चेतनः}
{मूर्छितो न्यपतद्भूमौ रामशोकाभिपीडितः} %||2-51-22||

\twolineshloka
{ततोऽन्तःपुरमाविद्धं मूर्छिते पृथिवीपतौ}
{उद्धृत्य बाहू चुक्रोश नृपतौ पतिते क्षितौ} %||2-51-23||

\twolineshloka
{सुमित्रया तु सहिता कौसल्या पतितं पतिम्}
{उत्थापयामास तदा वचनं चेदमब्रवीत्} %||2-51-24||

\twolineshloka
{इमं तस्य महाभाग दूतं दुष्करकारिणः}
{वनवासादनुप्राप्तं कस्मान्न प्रतिभाषसे} %||2-51-25||

\twolineshloka
{अद्येममनयं कृत्वा व्यपत्रपसि राघव}
{उत्तिष्ठ सुकृतं तेऽस्तु शोके न स्यात्सहायता} %||2-51-26||

\twolineshloka
{देव यस्या भयाद्रामं नानुपृच्छसि सारथिम्}
{नेह तिष्ठति कैकेयी विश्रब्धं प्रतिभाष्यताम्} %||2-51-27||

\twolineshloka
{सा तथोक्त्वा महाराजं कौसल्या शोकलालसा}
{धरण्यां निपपाताशु बाष्पविप्लुतभाषिणी} %||2-51-28||

\twolineshloka
{एवं विलपतीं दृष्ट्वा कौसल्यां पतितां भुवि}
{पतिं चावेक्ष्य ताः सर्वाः सस्वरं रुरुदुः स्त्रियः} %||2-51-29||

\fourlineindentedshloka
{ततस्तमन्तःपुरनादमुत्थितम्}
{ समीक्ष्य वृद्धास्तरुणाश्च मानवाः}
{स्त्रियश्च सर्वा रुरुदुः समन्ततः}
{ पुरं तदासीत्पुनरेव सङ्कुलम्} %||2-52-30||


{॥इत्यार्षे श्रीमद्रामायणे वाल्मीकीये आदिकाव्ये अयोध्याकाण्डे एकपञ्चाशत्तमः सर्गः॥}

\sect{द्विपञ्चाशत्तमः सर्गः}

\twolineshloka
{प्रत्याश्वस्तो यदा राजा मोहात्प्रत्यागतः पुनः}
{अथाजुहाव तं सूतं रामवृत्तान्तकारणात्} %||2-52-1||

\twolineshloka
{वृद्धं परमसन्तप्तं नवग्रहमिव द्विपम्}
{विनिःश्वसन्तं ध्यायन्तमस्वस्थमिव कुञ्जरम्} %||2-52-2||

\twolineshloka
{राजा तु रजसा सूतं ध्वस्ताङ्गं समुपस्थितम्}
{अश्रु पूर्णमुखं दीनमुवाच परमार्तवत्} %||2-52-3||

\threelineshloka
{क्व नु वत्स्यति धर्मात्मा वृक्षमूलमुपाश्रितः}
{सोऽत्यन्तसुखितः सूत किमशिष्यति राघवः}
{भूमिपालात्मजो भूमौ शेते कथमनाथवत्} %||2-52-4||

\twolineshloka
{यं यान्तमनुयान्ति स्म पदाति रथकुञ्जराः}
{स वत्स्यति कथं रामो विजनं वनमाश्रितः} %||2-52-5||

\twolineshloka
{व्यालैर्मृगैराचरितं कृष्णसर्पनिषेवितम्}
{कथं कुमारौ वैदेह्या सार्धं वनमुपस्थितौ} %||2-52-6||

\twolineshloka
{सुकुमार्या तपस्विन्या सुमन्त्र सह सीतया}
{राजपुत्रौ कथं पादैरवरुह्य रथाद्गतौ} %||2-52-7||

\twolineshloka
{सिद्धार्थः खलु सूत त्वं येन दृष्टौ ममात्मजौ}
{वनान्तं प्रविशन्तौ तावश्विनाविव मन्दरम्} %||2-52-8||

\threelineshloka
{किमुवाच वचो रामः किमुवाच च लक्ष्मणः}
{सुमन्त्र वनमासाद्य किमुवाच च मैथिली}
{आसितं शयितं भुक्तं सूत रामस्य कीर्तय} %||2-52-9||

\twolineshloka
{इति सूतो नरेन्द्रेण चोदितः सज्जमानया}
{उवाच वाचा राजानं सबाष्पपरिरब्धया} %||2-52-10||

\twolineshloka
{अब्रवीन्मां महाराज धर्ममेवानुपालयन्}
{अञ्जलिं राघवः कृत्वा शिरसाभिप्रणम्य च} %||2-52-11||

\twolineshloka
{सूत मद्वचनात्तस्य तातस्य विदितात्मनः}
{शिरसा वन्दनीयस्य वन्द्यौ पादौ महात्मनः} %||2-52-12||

\twolineshloka
{सर्वमन्तःपुरं वाच्यं सूत मद्वचनात्त्वया}
{आरोग्यमविशेषेण यथार्हं चाभिवादनम्} %||2-52-13||

\twolineshloka
{माता च मम कौसल्या कुशलं चाभिवादनम्}
{देवि देवस्य पादौ च देववत्परिपालय} %||2-52-14||

\twolineshloka
{भरतः कुशलं वाच्यो वाच्यो मद्वचनेन च}
{सर्वास्वेव यथान्यायं वृत्तिं वर्तस्व मातृषु} %||2-52-15||

\twolineshloka
{वक्तव्यश्च महाबाहुरिक्ष्वाकुकुलनन्दनः}
{पितरं यौवराज्यस्थो राज्यस्थमनुपालय} %||2-52-16||

\twolineshloka
{इत्येवं मां महाराज ब्रुवन्नेव महायशाः}
{रामो राजीवताम्राक्षो भृशमश्रूण्यवर्तयत्} %||2-52-17||

\twolineshloka
{लक्ष्मणस्तु सुसङ्क्रुद्धो निःश्वसन्वाक्यमब्रवीत्}
{केनायमपराधेन राजपुत्रो विवासितः} %||2-52-18||

\threelineshloka
{यदि प्रव्राजितो रामो लोभकारणकारितम्}
{वरदाननिमित्तं वा सर्वथा दुष्कृतं कृतम्}
{रामस्य तु परित्यागे न हेतुमुपलक्षये} %||2-52-19||

\twolineshloka
{असमीक्ष्य समारब्धं विरुद्धं बुद्धिलाघवात्}
{जनयिष्यति सङ्क्रोशं राघवस्य विवासनम्} %||2-52-20||

\twolineshloka
{अहं तावन्महाराजे पितृत्वं नोपलक्षये}
{भ्राता भर्ता च बन्धुश्च पिता च मम राघवः} %||2-52-21||

\twolineshloka
{सर्वलोकप्रियं त्यक्त्वा सर्वलोकहिते रतम्}
{सर्वलोकोऽनुरज्येत कथं त्वानेन कर्मणा} %||2-52-22||

\twolineshloka
{जानकी तु महाराज निःश्वसन्ती तपस्विनी}
{भूतोपहतचित्तेव विष्ठिता वृष्मृता स्थिता} %||2-52-23||

\twolineshloka
{अदृष्टपूर्वव्यसना राजपुत्री यशस्विनी}
{तेन दुःखेन रुदती नैव मां किञ्चिदब्रवीत्} %||2-52-24||

\twolineshloka
{उद्वीक्षमाणा भर्तारं मुखेन परिशुष्यता}
{मुमोच सहसा बाष्पं मां प्रयान्तमुदीक्ष्य सा} %||2-52-25||

\fourlineindentedshloka
{तथैव रामोऽश्रुमुखः कृताञ्जलिः}
{ स्थितोऽभवल्लक्ष्मणबाहुपालितः}
{तथैव सीता रुदती तपस्विनी}
{ निरीक्षते राजरथं तथैव माम्} %||2-53-26||


{॥इत्यार्षे श्रीमद्रामायणे वाल्मीकीये आदिकाव्ये अयोध्याकाण्डे द्विपञ्चाशत्तमः सर्गः॥}

\sect{त्रिपञ्चाशत्तमः सर्गः}

\twolineshloka
{मम त्वश्वा निवृत्तस्य न प्रावर्तन्त वर्त्मनि}
{उष्णमश्रु विमुञ्चन्तो रामे सम्प्रस्थिते वनम्} %||2-53-1||

\twolineshloka
{उभाभ्यां राजपुत्राभ्यामथ कृत्वाहमज्ञलिम्}
{प्रस्थितो रथमास्थाय तद्दुःखमपि धारयन्} %||2-53-2||

\twolineshloka
{गुहेव सार्धं तत्रैव स्थितोऽस्मि दिवसान्बहून्}
{आशया यदि मां रामः पुनः शब्दापयेदिति} %||2-53-3||

\twolineshloka
{विषये ते महाराज रामव्यसनकर्शिताः}
{अपि वृक्षाः परिम्लानः सपुष्पाङ्कुरकोरकाः} %||2-53-4||

\twolineshloka
{न च सर्पन्ति सत्त्वानि व्याला न प्रसरन्ति च}
{रामशोकाभिभूतं तन्निष्कूजमभवद्वनम्} %||2-53-5||

\twolineshloka
{लीनपुष्करपत्राश्च नरेन्द्र कलुषोदकाः}
{सन्तप्तपद्माः पद्मिन्यो लीनमीनविहङ्गमाः} %||2-53-6||

\twolineshloka
{जलजानि च पुष्पाणि माल्यानि स्थलजानि च}
{नाद्य भान्त्यल्पगन्धीनि फलानि च यथा पुरम्} %||2-53-7||

\twolineshloka
{प्रविशन्तमयोध्यां मां न कश्चिदभिनन्दति}
{नरा राममपश्यन्तो निःश्वसन्ति मुहुर्मुहुः} %||2-53-8||

\twolineshloka
{हर्म्यैर्विमानैः प्रासादैरवेक्ष्य रथमागतम्}
{हाहाकारकृता नार्यो रामादर्शनकर्शिताः} %||2-53-9||

\twolineshloka
{आयतैर्विमलैर्नेत्रैरश्रुवेगपरिप्लुतैः}
{अन्योन्यमभिवीक्षन्ते व्यक्तमार्ततराः स्त्रियः} %||2-53-10||

\twolineshloka
{नामित्राणां न मित्राणामुदासीनजनस्य च}
{अहमार्ततया कञ्चिद्विशेषं नोपलक्षये} %||2-53-11||

\twolineshloka
{अप्रहृष्टमनुष्या च दीननागतुरङ्गमा}
{आर्तस्वरपरिम्लाना विनिःश्वसितनिःस्वना} %||2-53-12||

\twolineshloka
{निरानन्दा महाराज रामप्रव्राजनातुला}
{कौसल्या पुत्र हीनेव अयोध्या प्रतिभाति मा} %||2-53-13||

\twolineshloka
{सूतस्य वचनं श्रुत्वा वाचा परमदीनया}
{बाष्पोपहतया राजा तं सूतमिदमब्रवीत्} %||2-53-14||

\twolineshloka
{कैकेय्या विनियुक्तेन पापाभिजनभावया}
{मया न मन्त्रकुशलैर्वृद्धैः सह समर्थितम्} %||2-53-15||

\twolineshloka
{न सुहृद्भिर्न चामात्यैर्मन्त्रयित्वा न नैगमैः}
{मयायमर्थः सम्मोहात्स्त्रीहेतोः सहसा कृतः} %||2-53-16||

\twolineshloka
{भवितव्यतया नूनमिदं वा व्यसनं महत्}
{कुलस्यास्य विनाशाय प्राप्तं सूत यदृच्छया} %||2-53-17||

\twolineshloka
{सूत यद्यस्ति ते किञ्चिन्मयापि सुकृतं कृतम्}
{त्वं प्रापयाशु मां रामं प्राणाः सन्त्वरयन्ति माम्} %||2-53-18||

\twolineshloka
{यद्यद्यापि ममैवाज्ञा निवर्तयतु राघवम्}
{न शक्ष्यामि विना राम मुहूर्तमपि जीवितुम्} %||2-53-19||

\twolineshloka
{अथ वापि महाबाहुर्गतो दूरं भविष्यति}
{मामेव रथमारोप्य शीघ्रं रामाय दर्शय} %||2-53-20||

\twolineshloka
{वृत्तदंष्ट्रो महेष्वासः क्वासौ लक्ष्मणपूर्वजः}
{यदि जीवामि साध्वेनं पश्येयं सह सीतया} %||2-53-21||

\twolineshloka
{लोहिताक्षं महाबाहुमामुक्तमणिकुण्डलम्}
{रामं यदि न पश्यामि गमिष्यामि यमक्षयम्} %||2-53-22||

\twolineshloka
{अतो नु किं दुःखतरं योऽहमिक्ष्वाकुनन्दनम्}
{इमामवस्थामापन्नो नेह पश्यामि राघवम्} %||2-53-23||

\threelineshloka
{हा राम रामानुज हा हा वैदेहि तपस्विनी}
{न मां जानीत दुःखेन म्रियमाणमनाथवत्}
{दुस्तरो जीवता देवि मयायं शोकसागरः} %||2-53-24||

\fourlineindentedshloka
{अशोभनं योऽहमिहाद्य राघवम्}
{ दिदृक्षमाणो न लभे सलक्ष्मणम्}
{इतीव राजा विलपन्महायशाः}
{ पपात तूर्णं शयने स मूर्छितः} %||2-53-25||

\fourlineindentedshloka
{इति विलपति पार्थिवे प्रनष्टे}
{ करुणतरं द्विगुणं च रामहेतोः}
{वचनमनुनिशम्य तस्य देवी}
{ भयमगमत्पुनरेव राममाता} %||2-54-26||


{॥इत्यार्षे श्रीमद्रामायणे वाल्मीकीये आदिकाव्ये अयोध्याकाण्डे त्रिपञ्चाशत्तमः सर्गः॥}

\sect{चतुःपञ्चाशत्तमः सर्गः}

\twolineshloka
{ततो भूतोपसृष्टेव वेपमाना पुनः पुनः}
{धरण्यां गतसत्त्वेव कौसल्या सूतमब्रवीत्} %||2-54-1||

\twolineshloka
{नय मां यत्र काकुत्स्थः सीता यत्र च लक्ष्मणः}
{तान्विना क्षणमप्यत्र जीवितुं नोत्सहे ह्यहम्} %||2-54-2||

\twolineshloka
{निवर्तय रथं शीघ्रं दण्डकान्नय मामपि}
{अथ तान्नानुगच्छामि गमिष्यामि यमक्षयम्} %||2-54-3||

\twolineshloka
{बाष्पवेगौपहतया स वाचा सज्जमानया}
{इदमाश्वासयन्देवीं सूतः प्राञ्जलिरब्रवीत्} %||2-54-4||

\twolineshloka
{त्यज शोकं च मोहं च सम्भ्रमं दुःखजं तथा}
{व्यवधूय च सन्तापं वने वत्स्यति राघवः} %||2-54-5||

\twolineshloka
{लक्ष्मणश्चापि रामस्य पादौ परिचरन्वने}
{आराधयति धर्मज्ञः परलोकं जितेन्द्रियः} %||2-54-6||

\twolineshloka
{विजनेऽपि वने सीता वासं प्राप्य गृहेष्विव}
{विस्रम्भं लभतेऽभीता रामे संन्यस्त मानसा} %||2-54-7||

\twolineshloka
{नास्या दैन्यं कृतं किञ्चित्सुसूक्ष्ममपि लक्षये}
{उचितेव प्रवासानां वैदेही प्रतिभाति मा} %||2-54-8||

\twolineshloka
{नगरोपवनं गत्वा यथा स्म रमते पुरा}
{तथैव रमते सीता निर्जनेषु वनेष्वपि} %||2-54-9||

\twolineshloka
{बालेव रमते सीता बालचन्द्रनिभानना}
{रामा रामे ह्यदीनात्मा विजनेऽपि वने सती} %||2-54-10||

\twolineshloka
{तद्गतं हृदयं ह्यस्यास्तदधीनं च जीवितम्}
{अयोध्यापि भवेत्तस्या राम हीना तथा वनम्} %||2-54-11||

\twolineshloka
{पथि पृच्छति वैदेही ग्रामांश्च नगराणि च}
{गतिं दृष्ट्वा नदीनां च पादपान्विविधानपि} %||2-54-12||

\twolineshloka
{अध्वना वात वेगेन सम्भ्रमेणातपेन च}
{न हि गच्छति वैदेह्याश्चन्द्रांशुसदृशी प्रभा} %||2-54-13||

\twolineshloka
{सदृशं शतपत्रस्य पूर्णचन्द्रोपमप्रभम्}
{वदनं तद्वदान्याया वैदेह्या न विकम्पते} %||2-54-14||

\twolineshloka
{अलक्तरसरक्ताभावलक्तरसवर्जितौ}
{अद्यापि चरणौ तस्याः पद्मकोशसमप्रभौ} %||2-54-15||

\twolineshloka
{नूपुरोद्घुष्टहेलेव खेलं गच्छति भामिनी}
{इदानीमपि वैदेही तद्रागा न्यस्तभूषणा} %||2-54-16||

\twolineshloka
{गजं वा वीक्ष्य सिंहं वा व्याघ्रं वा वनमाश्रिता}
{नाहारयति सन्त्रासं बाहू रामस्य संश्रिता} %||2-54-17||

\twolineshloka
{न शोच्यास्ते न चात्मा ते शोच्यो नापि जनाधिपः}
{इदं हि चरितं लोके प्रतिष्ठास्यति शाश्वतम्} %||2-54-18||

\fourlineindentedshloka
{विधूय शोकं परिहृष्टमानसा}
{ महर्षियाते पथि सुव्यवस्थिताः}
{वने रता वन्यफलाशनाः पितुः}
{ शुभां प्रतिज्ञां परिपालयन्ति ते} %||2-54-19||

\fourlineindentedshloka
{तथापि सूतेन सुयुक्तवादिना}
{ निवार्यमाणा सुतशोककर्शिता}
{न चैव देवी विरराम कूजिता}
{त्प्रियेति पुत्रेति च राघवेति च} %||2-55-20||


{॥इत्यार्षे श्रीमद्रामायणे वाल्मीकीये आदिकाव्ये अयोध्याकाण्डे चतुःपञ्चाशत्तमः सर्गः॥}

\sect{पञ्चपञ्चाशत्तमः सर्गः}

\twolineshloka
{वनं गते धर्मपरे रामे रमयतां वरे}
{कौसल्या रुदती स्वार्ता भर्तारमिदमब्रवीत्} %||2-55-1||

\twolineshloka
{यद्यपित्रिषु लोकेषु प्रथितं ते मयद्यशः}
{सानुक्रोशो वदान्यश्च प्रियवादी च राघवः} %||2-55-2||

\twolineshloka
{कथं नरवरश्रेष्ठ पुत्रौ तौ सह सीतया}
{दुःखितौ सुखसंवृद्धौ वने दुःखं सहिष्यतः} %||2-55-3||

\twolineshloka
{सा नूनं तरुणी श्यामा सुकुमारी सुखोचिता}
{कथमुष्णं च शीतं च मैथिली प्रसहिष्यते} %||2-55-4||

\twolineshloka
{भुक्त्वाशनं विशालाक्षी सूपदंशान्वितं शुभम्}
{वन्यं नैवारमाहारं कथं सीतोपभोक्ष्यते} %||2-55-5||

\twolineshloka
{गीतवादित्रनिर्घोषं श्रुत्वा शुभमनिन्दिता}
{कथं क्रव्यादसिंहानां शब्दं श्रोष्यत्यशोभनम्} %||2-55-6||

\twolineshloka
{महेन्द्रध्वजसङ्काशः क्व नु शेते महाभुजः}
{भुजं परिघसङ्काशमुपधाय महाबलः} %||2-55-7||

\twolineshloka
{पद्मवर्णं सुकेशान्तं पद्मनिःश्वासमुत्तमम्}
{कदा द्रक्ष्यामि रामस्य वदनं पुष्करेक्षणम्} %||2-55-8||

\twolineshloka
{वज्रसारमयं नूनं हृदयं मे न संशयः}
{अपश्यन्त्या न तं यद्वै फलतीदं सहस्रधा} %||2-55-9||

\twolineshloka
{यदि पञ्चदशे वर्षे राघवः पुनरेष्यति}
{जह्याद्राज्यं च कोशं च भरतेनोपभोक्ष्यते} %||2-55-10||

\twolineshloka
{एवं कनीयसा भ्रात्रा भुक्तं राज्यं विशां पते}
{भ्राता ज्येष्ठा वरिष्ठाश्च किमर्थं नावमंस्यते} %||2-55-11||

\twolineshloka
{न परेणाहृतं भक्ष्यं व्याघ्रः खादितुमिच्छति}
{एवमेव नरव्याघ्रः परलीढं न मंस्यते} %||2-55-12||

\twolineshloka
{हविराज्यं पुरोडाशाः कुशा यूपाश्च खादिराः}
{नैतानि यातयामानि कुर्वन्ति पुनरध्वरे} %||2-55-13||

\twolineshloka
{तथा ह्यात्तमिदं राज्यं हृतसारां सुरामिव}
{नाभिमन्तुमलं रामो नष्टसोममिवाध्वरम्} %||2-55-14||

\twolineshloka
{नैवंविधमसत्कारं राघवो मर्षयिष्यति}
{बलवानिव शार्दूलो बालधेरभिमर्शनम्} %||2-55-15||

\twolineshloka
{स तादृशः सिंहबलो वृषभाक्षो नरर्षभः}
{स्वयमेव हतः पित्रा जलजेनात्मजो यथा} %||2-55-16||

\twolineshloka
{द्विजाति चरितो धर्मः शास्त्रदृष्टः सनातनः}
{यदि ते धर्मनिरते त्वया पुत्रे विवासिते} %||2-55-17||

\twolineshloka
{गतिरेवाक्पतिर्नार्या द्वितीया गतिरात्मजः}
{तृतीया ज्ञातयो राजंश्चतुर्थी नेह विद्यते} %||2-55-18||

\twolineshloka
{तत्र त्वं चैव मे नास्ति रामश्च वनमाश्रितः}
{न वनं गन्तुमिच्छामि सर्वथा हि हता त्वया} %||2-55-19||

\fourlineindentedshloka
{हतं त्वया राज्यमिदं सराष्ट्रम्}
{ हतस्तथात्मा सह मन्त्रिभिश्च}
{हता सपुत्रास्मि हताश्च पौराः}
{ सुतश्च भार्या च तव प्रहृष्टौ} %||2-55-20||

\fourlineindentedshloka
{इमां गिरं दारुणशब्दसंश्रिताम्}
{ निशम्य राजापि मुमोह दुःखितः}
{ततः स शोकं प्रविवेश पार्थिवः}
{ स्वदुष्कृतं चापि पुनस्तदास्मरत्} %||2-56-21||


{॥इत्यार्षे श्रीमद्रामायणे वाल्मीकीये आदिकाव्ये अयोध्याकाण्डे पञ्चपञ्चाशत्तमः सर्गः॥}

\sect{षट्पञ्चाशत्तमः सर्गः}

\twolineshloka
{एवं तु क्रुद्धया राजा राममात्रा सशोकया}
{श्रावितः परुषं वाक्यं चिन्तयामास दुःखितः} %||2-56-1||

\twolineshloka
{तस्य चिन्तयमानस्य प्रत्यभात्कर्म दुष्कृतम्}
{यदनेन कृतं पूर्वमज्ञानाच्छब्दवेधिना} %||2-56-2||

\twolineshloka
{अमनास्तेन शोकेन रामशोकेन च प्रभुः}
{दह्यमानस्तु शोकाभ्यां कौसल्यामाह भूपतिः} %||2-56-3||

\twolineshloka
{प्रसादये त्वां कौसल्ये रचितोऽयं मयाञ्जलिः}
{वत्सला चानृशंसा च त्वं हि नित्यं परेष्वपि} %||2-56-4||

\twolineshloka
{भर्ता तु खलु नारीणां गुणवान्निर्गुणोऽपि वा}
{धर्मं विमृशमानानां प्रत्यक्षं देवि दैवतम्} %||2-56-5||

\twolineshloka
{सा त्वं धर्मपरा नित्यं दृष्टलोकपरावर}
{नार्हसे विप्रियं वक्तुं दुःखितापि सुदुःखितम्} %||2-56-6||

\twolineshloka
{तद्वाक्यं करुणं राज्ञः श्रुत्वा दीनस्य भाषितम्}
{कौसल्या व्यसृजद्बाष्पं प्रणालीव नवोदकम्} %||2-56-7||

\twolineshloka
{स मूद्र्ह्णि बद्ध्वा रुदती राज्ञः पद्ममिवाञ्जलिम्}
{सम्भ्रमादब्रवीत्त्रस्ता त्वरमाणाक्षरं वचः} %||2-56-8||

\twolineshloka
{प्रसीद शिरसा याचे भूमौ निततितास्मि ते}
{याचितास्मि हता देव हन्तव्याहं न हि त्वया} %||2-56-9||

\twolineshloka
{नैषा हि सा स्त्री भवति श्लाघनीयेन धीमता}
{उभयोर्लोकयोर्वीर पत्या या सम्प्रसाद्यते} %||2-56-10||

\twolineshloka
{जानामि धर्मं धर्मज्ञ त्वां जाने सत्यवादिनम्}
{पुत्रशोकार्तया तत्तु मया किमपि भाषितम्} %||2-56-11||

\twolineshloka
{शोको नाशयते धैर्यं शोको नाशयते श्रुतम्}
{शोको नाशयते सर्वं नास्ति शोकसमो रिपुः} %||2-56-12||

\twolineshloka
{शयमापतितः सोढुं प्रहरो रिपुहस्ततः}
{सोढुमापतितः शोकः सुसूक्ष्मोऽपि न शक्यते} %||2-56-13||

\twolineshloka
{वनवासाय रामस्य पञ्चरात्रोऽद्य गण्यते}
{यः शोकहतहर्षायाः पञ्चवर्षोपमो मम} %||2-56-14||

\twolineshloka
{तं हि चिन्तयमानायाः शोकोऽयं हृदि वर्धते}
{अदीनामिव वेगेन समुद्रसलिलं महत्} %||2-56-15||

\twolineshloka
{एवं हि कथयन्त्यास्तु कौसल्यायाः शुभं वचः}
{मन्दरश्मिरभूत्सुर्यो रजनी चाभ्यवर्तत} %||2-56-16||

\twolineshloka
{अथ प्रह्लादितो वाक्यैर्देव्या कौसल्यया नृपः}
{शोकेन च समाक्रान्तो निद्राया वशमेयिवान्} %||2-57-17||


{॥इत्यार्षे श्रीमद्रामायणे वाल्मीकीये आदिकाव्ये अयोध्याकाण्डे षट्पञ्चाशत्तमः सर्गः॥}

\sect{सप्तपञ्चाशत्तमः सर्गः}

\twolineshloka
{प्रतिबुद्धो मुहुर्तेन शोकोपहतचेतनः}
{अथ राजा दशरथः स चिन्तामभ्यपद्यत} %||2-57-1||

\twolineshloka
{रामलक्ष्मणयोश्चैव विवासाद्वासवोपमम्}
{आविवेशोपसर्गस्तं तमः सूर्यमिवासुरम्} %||2-57-2||

\threelineshloka
{स राजा रजनीं षष्ठीं रामे प्रव्रजिते वनम्}
{अर्धरात्रे दशरथः संस्मरन्दुष्कृतं कृतम्}
{कौसल्यां पुत्रशोकार्तामिदं वचनमब्रवीत्} %||2-57-3||

\twolineshloka
{यदाचरति कल्याणि शुभं वा यदि वाशुभम्}
{तदेव लभते भद्रे कर्ता कर्मजमात्मनः} %||2-57-4||

\twolineshloka
{गुरु लाघवमर्थानामारम्भे कर्मणां फलम्}
{दोषं वा यो न जानाति स बाल इति होच्यते} %||2-57-5||

\twolineshloka
{कश्चिदाम्रवणं छित्त्वा पलाशांश्च निषिञ्चति}
{पुष्पं दृष्ट्वा फले गृध्नुः स शोचति फलागमे} %||2-57-6||

\twolineshloka
{सोऽहमाम्रवणं छित्त्वा पलाशांश्च न्यषेचयम्}
{रामं फलागमे त्यक्त्वा पश्चाच्छोचामि दुर्मतिः} %||2-57-7||

\threelineshloka
{लब्धशब्देन कौसल्ये कुमारेण धनुष्मता}
{कुमारः शब्दवेधीति मया पापमिदं कृतम्}
{तदिदं मेऽनुसम्प्राप्तं देवि दुःखं स्वयं कृतम्} %||2-57-8||

\twolineshloka
{सम्मोहादिह बालेन यथा स्याद्भक्षितं विषम्}
{एवं ममाप्यविज्ञातं शब्दवेध्यमयं फलम्} %||2-57-9||

\twolineshloka
{देव्यनूढा त्वमभवो युवराजो भवाम्यहम्}
{ततः प्रावृडनुप्राप्ता मदकामविवर्धिनी} %||2-57-10||

\twolineshloka
{उपास्यहि रसान्भौमांस्तप्त्वा च जगदंशुभिः}
{परेताचरितां भीमां रविराविशते दिशम्} %||2-57-11||

\twolineshloka
{उष्णमन्तर्दधे सद्यः स्निग्धा ददृशिरे घनाः}
{ततो जहृषिरे सर्वे भेकसारङ्गबर्हिणः} %||2-57-12||

\twolineshloka
{पतितेनाम्भसा छन्नः पतमानेन चासकृत्}
{आबभौ मत्तसारङ्गस्तोयराशिरिवाचलः} %||2-57-13||

\twolineshloka
{तस्मिन्नतिसुखे काले धनुष्मानिषुमान्रथी}
{व्यायाम कृतसङ्कल्पः सरयूमन्वगां नदीम्} %||2-57-14||

\twolineshloka
{निपाने महिषं रात्रौ गजं वाभ्यागतं नदीम्}
{अन्यं वा श्वापदं कञ्चिज्जिघांसुरजितेन्द्रियः} %||2-57-15||

\twolineshloka
{अथान्धकारे त्वश्रौषं जले कुम्भस्य पर्यतः}
{अचक्षुर्विषये घोषं वारणस्येव नर्दतः} %||2-57-16||

\twolineshloka
{ततोऽहं शरमुद्धृत्य दीप्तमाशीविषोपमम्}
{अमुञ्चं निशितं बाणमहमाशीविषोपमम्} %||2-57-17||

\threelineshloka
{तत्र वागुषसि व्यक्ता प्रादुरासीद्वनौकसः}
{हा हेति पततस्तोये वागभूत्तत्र मानुषी}
{कथमस्मद्विधे शस्त्रं निपतेत्तु तपस्विनि} %||2-57-18||

\twolineshloka
{प्रविविक्तां नदीं रात्रावुदाहारोऽहमागतः}
{इषुणाभिहतः केन कस्य वा किं कृतं मया} %||2-57-19||

\twolineshloka
{ऋषेर्हि न्यस्त दण्डस्य वने वन्येन जीवतः}
{कथं नु शस्त्रेण वधो मद्विधस्य विधीयते} %||2-57-20||

\twolineshloka
{जटाभारधरस्यैव वल्कलाजिनवाससः}
{को वधेन ममार्थी स्यात्किं वास्यापकृतं मया} %||2-57-21||

\twolineshloka
{एवं निष्फलमारब्धं केवलानर्थसंहितम्}
{न कश्चित्साधु मन्येत यथैव गुरुतल्पगम्} %||2-57-22||

\twolineshloka
{नेमं तथानुशोचामि जीवितक्षयमात्मनः}
{मातरं पितरं चोभावनुशोचामि मद्विधे} %||2-57-23||

\twolineshloka
{तदेतान्मिथुनं वृद्धं चिरकालभृतं मया}
{मयि पञ्चत्वमापन्ने कां वृत्तिं वर्तयिष्यति} %||2-57-24||

\twolineshloka
{वृद्धौ च मातापितरावहं चैकेषुणा हतः}
{केन स्म निहताः सर्वे सुबालेनाकृतात्मना} %||2-57-25||

\twolineshloka
{तं गिरं करुणां श्रुत्वा मम धर्मानुकाङ्क्षिणः}
{कराभ्यां सशरं चापं व्यथितस्यापतद्भुवि} %||2-57-26||

\twolineshloka
{तं देशमहमागम्य दीनसत्त्वः सुदुर्मनाः}
{अपश्यमिषुणा तीरे सरय्वास्तापसं हतम्} %||2-57-27||

\twolineshloka
{स मामुद्वीक्ष्य नेत्राभ्यां त्रस्तमस्वस्थचेतसम्}
{इत्युवाच वचः क्रूरं दिधक्षन्निव तेजसा} %||2-57-28||

\twolineshloka
{किं तवापकृतं राजन्वने निवसता मया}
{जिहीर्षुरम्भो गुर्वर्थं यदहं ताडितस्त्वया} %||2-57-29||

\twolineshloka
{एकेन खलु बाणेन मर्मण्यभिहते मयि}
{द्वावन्धौ निहतौ वृद्धौ माता जनयिता च मे} %||2-57-30||

\twolineshloka
{तौ नूनं दुर्बलावन्धौ मत्प्रतीक्षौ पिपासितौ}
{चिरमाशाकृतां तृष्णां कष्टां सन्धारयिष्यतः} %||2-57-31||

\twolineshloka
{न नूनं तपसो वास्ति फलयोगः श्रुतस्य वा}
{पिता यन्मां न जानाति शयानं पतितं भुवि} %||2-57-32||

\twolineshloka
{जानन्नपि च किं कुर्यादशक्तिरपरिक्रमः}
{भिद्यमानमिवाशक्तस्त्रातुमन्यो नगो नगम्} %||2-57-33||

\twolineshloka
{पितुस्त्वमेव मे गत्वा शीघ्रमाचक्ष्व राघव}
{न त्वामनुदहेत्क्रुद्धो वनं वह्निरिवैधितः} %||2-57-34||

\twolineshloka
{इयमेकपदी राजन्यतो मे पितुराश्रमः}
{तं प्रसादय गत्वा त्वं न त्वां स कुपितः शपेत्} %||2-57-35||

\twolineshloka
{विशल्यं कुरु मां राजन्मर्म मे निशितः शरः}
{रुणद्धि मृदु सोत्सेधं तीरमम्बुरयो यथा} %||2-57-36||

\twolineshloka
{न द्विजातिरहं राजन्मा भूत्ते मनसो व्यथा}
{शूद्रायामस्मि वैश्येन जातो जनपदाधिप} %||2-57-37||

\twolineshloka
{इतीव वदतः कृच्छ्राद्बाणाभिहतमर्मणः}
{तस्य त्वानम्यमानस्य तं बाणमहमुद्धरम्} %||2-57-38||

\fourlineindentedshloka
{जलार्द्रगात्रं तु विलप्य कृच्छा}
{न्मर्मव्रणं सन्ततमुच्छसन्तम्}
{ततः सरय्वां तमहं शयानम्}
{ समीक्ष्य भद्रे सुभृशं विषण्णः} %||2-58-39||


{॥इत्यार्षे श्रीमद्रामायणे वाल्मीकीये आदिकाव्ये अयोध्याकाण्डे सप्तपञ्चाशत्तमः सर्गः॥}

\sect{अष्टपञ्चाशत्तमः सर्गः}

\twolineshloka
{तदज्ञानान्महत्पापं कृत्वा सङ्कुलितेन्द्रियः}
{एकस्त्वचिन्तयं बुद्ध्या कथं नु सुकृतं भवेत्} %||2-58-1||

\twolineshloka
{ततस्तं घटमादय पूर्णं परमवारिणा}
{आश्रमं तमहं प्राप्य यथाख्यातपथं गतः} %||2-58-2||

\twolineshloka
{तत्राहं दुर्बलावन्धौ वृद्धावपरिणायकौ}
{अपश्यं तस्य पितरौ लूनपक्षाविव द्विजौ} %||2-58-3||

\twolineshloka
{तन्निमित्ताभिरासीनौ कथाभिरपरिक्रमौ}
{तामाशां मत्कृते हीनावुदासीनावनाथवत्} %||2-58-4||

\twolineshloka
{पदशब्दं तु मे श्रुत्वा मुनिर्वाक्यमभाषत}
{किं चिरायसि मे पुत्र पानीयं क्षिप्रमानय} %||2-58-5||

\twolineshloka
{यन्निमित्तमिदं तात सलिले क्रीडितं त्वया}
{उत्कण्ठिता ते मातेयं प्रविश क्षिप्रमाश्रमम्} %||2-58-6||

\twolineshloka
{यद्व्यलीकं कृतं पुत्र मात्रा ते यदि वा मया}
{न तन्मनसि कर्तव्यं त्वया तात तपस्विना} %||2-58-7||

\twolineshloka
{त्वं गतिस्त्वगतीनां च चक्षुस्त्वं हीनचक्षुषाम्}
{समासक्तास्त्वयि प्राणाः किञ्चिन्नौ नाभिभाषसे} %||2-58-8||

\twolineshloka
{मुनिमव्यक्तया वाचा तमहं सज्जमानया}
{हीनव्यञ्जनया प्रेक्ष्य भीतो भीत इवाब्रुवम्} %||2-58-9||

\twolineshloka
{मनसः कर्म चेष्टाभिरभिसंस्तभ्य वाग्बलम्}
{आचचक्षे त्वहं तस्मै पुत्रव्यसनजं भयम्} %||2-58-10||

\twolineshloka
{क्षत्रियोऽहं दशरथो नाहं पुत्रो महात्मनः}
{सज्जनावमतं दुःखमिदं प्राप्तं स्वकर्मजम्} %||2-58-11||

\twolineshloka
{भगवंश्चापहस्तोऽहं सरयूतीरमागतः}
{जिघांसुः श्वापदं किञ्चिन्निपाने वागतं गजम्} %||2-58-12||

\twolineshloka
{तत्र श्रुतो मया शब्दो जले कुम्भस्य पूर्यतः}
{द्विपोऽयमिति मत्वा हि बाणेनाभिहतो मया} %||2-58-13||

\twolineshloka
{गत्वा नद्यास्ततस्तीरमपश्यमिषुणा हृदि}
{विनिर्भिन्नं गतप्राणं शयानं भुवि तापसम्} %||2-58-14||

\twolineshloka
{भगवञ्शब्दमालक्ष्य मया गजजिघांसुना}
{विसृष्टोऽम्भसि नाराचस्तेन ते निहतः सुतः} %||2-58-15||

\twolineshloka
{स चोद्धृतेन बाणेन तत्रैव स्वर्गमास्थितः}
{भगवन्तावुभौ शोचन्नन्धाविति विलप्य च} %||2-58-16||

\twolineshloka
{अज्ञानाद्भवतः पुत्रः सहसाभिहतो मया}
{शेषमेवङ्गते यत्स्यात्तत्प्रसीदतु मे मुनिः} %||2-58-17||

\twolineshloka
{स तच्छ्रुत्वा वचः क्रूरं निःश्वसञ्शोककर्शितः}
{मामुवाच महातेजाः कृताञ्जलिमुपस्थितम्} %||2-58-18||

\twolineshloka
{यद्येतदशुभं कर्म न स्म मे कथयेः स्वयम्}
{फलेन्मूर्धा स्म ते राजन्सद्यः शतसहस्रधा} %||2-58-19||

\twolineshloka
{क्षत्रियेण वधो राजन्वानप्रस्थे विशेषतः}
{ज्ञानपूर्वं कृतः स्थानाच्च्यावयेदपि वज्रिणम्} %||2-58-20||

\twolineshloka
{अज्ञानाद्धि कृतं यस्मादिदं तेनैव जीवसि}
{अपि ह्यद्य कुलं नस्याद्राघवाणां कुतो भवान्} %||2-58-21||

\twolineshloka
{नय नौ नृप तं देशमिति मां चाभ्यभाषत}
{अद्य तं द्रष्टुमिच्छावः पुत्रं पश्चिमदर्शनम्} %||2-58-22||

\twolineshloka
{रुधिरेणावसिताङ्गं प्रकीर्णाजिनवाससम्}
{शयानं भुवि निःसंज्ञं धर्मराजवशं गतम्} %||2-58-23||

\twolineshloka
{अथाहमेकस्तं देशं नीत्वा तौ भृशदुःखितौ}
{अस्पर्शयमहं पुत्रं तं मुनिं सह भार्यया} %||2-58-24||

\twolineshloka
{तौ पुत्रमात्मनः स्पृष्ट्वा तमासाद्य तपस्विनौ}
{निपेततुः शरीरेऽस्य पिता चास्येदमब्रवीत्} %||2-58-25||

\twolineshloka
{न न्वहं ते प्रियः पुत्र मातरं पश्य धार्मिक}
{किं नु नालिङ्गसे पुत्र सुकुमार वचो वद} %||2-58-26||

\twolineshloka
{कस्य वापररात्रेऽहं श्रोष्यामि हृदयङ्गमम्}
{अधीयानस्य मधुरं शास्त्रं वान्यद्विशेषतः} %||2-58-27||

\twolineshloka
{को मां सन्ध्यामुपास्यैव स्नात्वा हुतहुताशनः}
{श्लाघयिष्यत्युपासीनः पुत्रशोकभयार्दितम्} %||2-58-28||

\twolineshloka
{कन्दमूलफलं हृत्वा को मां प्रियमिवातिथिम्}
{भोजयिष्यत्यकर्मण्यमप्रग्रहमनायकम्} %||2-58-29||

\twolineshloka
{इमामन्धां च वृद्धां च मातरं ते तपस्विनीम्}
{कथं पुत्र भरिष्यामि कृपणां पुत्रगर्धिनीम्} %||2-58-30||

\twolineshloka
{तिष्ठ मा मा गमः पुत्र यमस्य सदनं प्रति}
{श्वो मया सह गन्तासि जनन्या च समेधितः} %||2-58-31||

\twolineshloka
{उभावपि च शोकार्तावनाथौ कृपणौ वने}
{क्षिप्रमेव गमिष्यावस्त्वया हीनौ यमक्षयम्} %||2-58-32||

\twolineshloka
{ततो वैवस्वतं दृष्ट्वा तं प्रवक्ष्यामि भारतीम्}
{क्षमतां धर्मराजो मे बिभृयात्पितरावयम्} %||2-58-33||

\twolineshloka
{अपापोऽसि यथा पुत्र निहतः पापकर्मणा}
{तेन सत्येन गच्छाशु ये लोकाः शस्त्रयोधिनाम्} %||2-58-34||

\twolineshloka
{यान्ति शूरा गतिं यां च सङ्ग्रामेष्वनिवर्तिनः}
{हतास्त्वभिमुखाः पुत्र गतिं तां परमां व्रज} %||2-58-35||

\twolineshloka
{यां गतिं सगरः शैब्यो दिलीपो जनमेजयः}
{नहुषो धुन्धुमारश्च प्राप्तास्तां गच्छ पुत्रक} %||2-58-36||

\twolineshloka
{या गतिः सर्वसाधूनां स्वाध्यायात्पतसश्च या}
{भूमिदस्याहिताग्नेश्च एकपत्नीव्रतस्य च} %||2-58-37||

\threelineshloka
{गोसहस्रप्रदातॄणां या या गुरुभृतामपि}
{देहन्यासकृतां या च तां गतिं गच्छ पुत्रक}
{न हि त्वस्मिन्कुले जातो गच्छत्यकुशलां गतिम्} %||2-58-38||

\twolineshloka
{एवं स कृपणं तत्र पर्यदेवयतासकृत्}
{ततोऽस्मै कर्तुमुदकं प्रवृत्तः सह भार्यया} %||2-58-39||

\twolineshloka
{स तु दिव्येन रूपेण मुनिपुत्रः स्वकर्मभिः}
{आश्वास्य च मुहूर्तं तु पितरौ वाक्यमब्रवीत्} %||2-58-40||

\twolineshloka
{स्थानमस्मि महत्प्राप्तो भवतोः परिचारणात्}
{भवन्तावपि च क्षिप्रं मम मूलमुपैष्यतः} %||2-58-41||

\twolineshloka
{एवमुक्त्वा तु दिव्येन विमानेन वपुष्मता}
{आरुरोह दिवं क्षिप्रं मुनिपुत्रो जितेन्द्रियः} %||2-58-42||

\twolineshloka
{स कृत्वा तूदकं तूर्णं तापसः सह भार्यया}
{मामुवाच महातेजाः कृताञ्जलिमुपस्थितम्} %||2-58-43||

\twolineshloka
{अद्यैव जहि मां राजन्मरणे नास्ति मे व्यथा}
{यच्छरेणैकपुत्रं मां त्वमकार्षीरपुत्रकम्} %||2-58-44||

\twolineshloka
{त्वया तु यदविज्ञानान्निहतो मे सुतः शुचिः}
{तेन त्वामभिशप्स्यामि सुदुःखमतिदारुणम्} %||2-58-45||

\twolineshloka
{पुत्रव्यसनजं दुःखं यदेतन्मम साम्प्रतम्}
{एवं त्वं पुत्रशोकेन राजन्कालं करिष्यसि} %||2-58-46||

\twolineshloka
{तस्मान्मामागतं भद्रे तस्योदारस्य तद्वचः}
{यदहं पुत्रशोकेन सन्त्यक्ष्याम्यद्य जीवितम्} %||2-58-47||

\twolineshloka
{यदि मां संस्पृशेद्रामः सकृदद्यालभेत वा}
{न तन्मे सदृशं देवि यन्मया राघवे कृतम्} %||2-58-48||

\twolineshloka
{चक्षुषा त्वां न पश्यामि स्मृतिर्मम विलुप्यते}
{दूता वैवस्वतस्यैते कौसल्ये त्वरयन्ति माम्} %||2-58-49||

\twolineshloka
{अतस्तु किं दुःखतरं यदहं जीवितक्षये}
{न हि पश्यामि धर्मज्ञं रामं सत्यपराक्यमम्} %||2-58-50||

\twolineshloka
{न ते मनुष्या देवास्ते ये चारुशुभकुण्डलम्}
{मुखं द्रक्ष्यन्ति रामस्य वर्षे पञ्चदशे पुनः} %||2-58-51||

\twolineshloka
{पद्मपत्रेक्षणं सुभ्रु सुदंष्ट्रं चारुनासिकम्}
{धन्या द्रक्ष्यन्ति रामस्य ताराधिपनिभं मुखम्} %||2-58-52||

\twolineshloka
{सदृशं शारदस्येन्दोः फुल्लस्य कमलस्य च}
{सुगन्धि मम नाथस्य धन्या द्रक्ष्यन्ति तन्मुखम्} %||2-58-53||

\twolineshloka
{निवृत्तवनवासं तमयोध्यां पुनरागतम्}
{द्रक्ष्यन्ति सुखिनो रामं शुक्रं मार्गगतं यथा} %||2-58-54||

\twolineshloka
{अयमात्मभवः शोको मामनाथमचेतनम्}
{संसादयति वेगेन यथा कूलं नदीरयः} %||2-58-55||

\twolineshloka
{हा राघव महाबाहो हा ममायास नाशन}
{राजा दशरथः शोचञ्जीवितान्तमुपागमत्} %||2-58-56||

\fourlineindentedshloka
{तथा तु दीनं कथयन्नराधिपः}
{ प्रियस्य पुत्रस्य विवासनातुरः}
{गतेऽर्धरात्रे भृशदुःखपीडित}
{स्तदा जहौ प्राणमुदारदर्शनः} %||2-59-57||


{॥इत्यार्षे श्रीमद्रामायणे वाल्मीकीये आदिकाव्ये अयोध्याकाण्डे अष्टपञ्चाशत्तमः सर्गः॥}

\sect{एकोनषष्टितमः सर्गः}

\twolineshloka
{अथ रात्र्यां व्यतीतायां प्रातरेवापरेऽहनि}
{बन्दिनः पर्युपातिष्ठंस्तत्पार्थिवनिवेशनम्} %||2-59-1||

\twolineshloka
{ततः शुचिसमाचाराः पर्युपस्थान कोविदः}
{स्त्रीवर्षवरभूयिष्ठा उपतस्थुर्यथापुरम्} %||2-59-2||

\twolineshloka
{हरिचन्दनसम्पृक्तमुदकं काञ्चनैर्घटैः}
{आनिन्युः स्नानशिक्षाज्ञा यथाकालं यथाविधि} %||2-59-3||

\twolineshloka
{मङ्गलालम्भनीयानि प्राशनीयानुपस्करान्}
{उपनिन्युस्तथाप्यन्याः कुमारी बहुलाः स्त्रियः} %||2-59-4||

\twolineshloka
{अथ याः कोसलेन्द्रस्य शयनं प्रत्यनन्तराः}
{ताः स्त्रियस्तु समागम्य भर्तारं प्रत्यबोधयन्} %||2-59-5||

\twolineshloka
{ता वेपथुपरीताश्च राज्ञः प्राणेषु शङ्किताः}
{प्रतिस्रोतस्तृणाग्राणां सदृशं सञ्चकम्पिरे} %||2-59-6||

\twolineshloka
{अथ संवेपमनानां स्त्रीणां दृष्ट्वा च पार्थिवम्}
{यत्तदाशङ्कितं पापं तस्य जज्ञे विनिश्चयः} %||2-59-7||

\twolineshloka
{ततः प्रचुक्रुशुर्दीनाः सस्वरं ता वराङ्गनाः}
{करेणव इवारण्ये स्थानप्रच्युतयूथपाः} %||2-59-8||

\twolineshloka
{तासामाक्रन्द शब्देन सहसोद्गतचेतने}
{कौसल्या च सुमित्राच त्यक्तनिद्रे बभूवतुः} %||2-59-9||

\twolineshloka
{कौसल्या च सुमित्रा च दृष्ट्वा स्पृष्ट्वा च पार्थिवम्}
{हा नाथेति परिक्रुश्य पेततुर्धरणीतले} %||2-59-10||

\twolineshloka
{सा कोसलेन्द्रदुहिता वेष्टमाना महीतले}
{न बभ्राज रजोध्वस्ता तारेव गगनच्युता} %||2-59-11||

\twolineshloka
{तत्समुत्त्रस्तसम्भ्रान्तं पर्युत्सुकजनाकुलम्}
{सर्वतस्तुमुलाक्रन्दं परितापार्तबान्धवम्} %||2-59-12||

\twolineshloka
{सद्यो निपतितानन्दं दीनविक्लवदर्शनम्}
{बभूव नरदेवस्य सद्म दिष्टान्तमीयुषः} %||2-59-13||

\fourlineindentedshloka
{अतीतमाज्ञाय तु पार्थिवर्षभम्}
{ यशस्विनं सम्परिवार्य पत्नयः}
{भृशं रुदन्त्यः करुणं सुदुःखिताः}
{ प्रगृह्य बाहू व्यलपन्ननाथवत्} %||2-60-14||


{॥इत्यार्षे श्रीमद्रामायणे वाल्मीकीये आदिकाव्ये अयोध्याकाण्डे एकोनषष्टितमः सर्गः॥}

\sect{षष्टितमः सर्गः}

\twolineshloka
{तमग्निमिव संशान्तमम्बुहीनमिवार्णवम्}
{हतप्रभमिवादित्यं स्वर्गथं प्रेक्ष्य भूमिपम्} %||2-60-1||

\twolineshloka
{कौसल्या बाष्पपूर्णाक्षी विविधं शोककर्शिता}
{उपगृह्य शिरो राज्ञः कैकेयीं प्रत्यभाषत} %||2-60-2||

\twolineshloka
{सकामा भव कैकेयि भुङ्क्ष्व राज्यमकण्टकम्}
{त्यक्त्वा राजानमेकाग्रा नृशंसे दुष्टचारिणि} %||2-60-3||

\twolineshloka
{विहाय मां गतो रामो भर्ता च स्वर्गतो मम}
{विपथे सार्थहीनेव नाहं जीवितुमुत्सहे} %||2-60-4||

\twolineshloka
{भर्तारं तं परित्यज्य का स्त्री दैवतमात्मनः}
{इच्छेज्जीवितुमन्यत्र कैकेय्यास्त्यक्तधर्मणः} %||2-60-5||

\twolineshloka
{न लुब्धो बुध्यते दोषान्किं पाकमिव भक्षयन्}
{कुब्जानिमित्तं कैकेय्या राघवाणान्कुलं हतम्} %||2-60-6||

\twolineshloka
{अनियोगे नियुक्तेन राज्ञा रामं विवासितम्}
{सभार्यं जनकः श्रुत्वा पतितप्स्यत्यहं यथा} %||2-60-7||

\threelineshloka
{रामः कमलपत्राक्षो जीवनाशमितो गतः}
{विदेहराजस्य सुता तहा सीता तपस्विनी}
{दुःखस्यानुचिता दुःखं वने पर्युद्विजिष्यति} %||2-60-8||

\twolineshloka
{नदतां भीमघोषाणां निशासु मृगपक्षिणाम्}
{निशम्य नूनं संस्त्रस्ता राघवं संश्रयिष्यति} %||2-60-9||

\twolineshloka
{वृद्धश्चैवाल्पपुत्रश्च वैदेहीमनिचिन्तयन्}
{सोऽपि शोकसमाविष्टो ननु त्यक्ष्यति जीवितम्} %||2-60-10||

\twolineshloka
{तां ततः सम्परिष्वज्य विलपन्तीं तपस्विनीम्}
{व्यपनिन्युः सुदुःखार्तां कौसल्यां व्यावहारिकाः} %||2-60-11||

\twolineshloka
{तैलद्रोण्यामथामात्याः संवेश्य जगतीपतिम्}
{राज्ञः सर्वाण्यथादिष्टाश्चक्रुः कर्माण्यनन्तरम्} %||2-60-12||

\twolineshloka
{न तु सङ्कलनं राज्ञो विना पुत्रेण मन्त्रिणः}
{सर्वज्ञाः कर्तुमीषुस्ते ततो रक्षन्ति भूमिपम्} %||2-60-13||

\twolineshloka
{तैलद्रोण्यां तु सचिवैः शायितं तं नराधिपम्}
{हा मृतोऽयमिति ज्ञात्वा स्त्रियस्ताः पर्यदेवयन्} %||2-60-14||

\twolineshloka
{बाहूनुद्यम्य कृपणा नेत्रप्रस्रवणैर्मुखैः}
{रुदन्त्यः शोकसन्तप्ताः कृपणं पर्यदेवयन्} %||2-60-15||

\twolineshloka
{निशानक्षत्रहीनेव स्त्रीव भर्तृविवर्जिता}
{पुरी नाराजतायोध्या हीना राज्ञा महात्मना} %||2-60-16||

\twolineshloka
{बाष्पपर्याकुलजना हाहाभूतकुलाङ्गना}
{शून्यचत्वरवेश्मान्ता न बभ्राज यथापुरम्} %||2-60-17||

\fourlineindentedshloka
{गतप्रभा द्यौरिव भास्करं विना}
{ व्यपेतनक्षत्रगणेव शर्वरी}
{पुरी बभासे रहिता महात्मना}
{ न चास्रकण्ठाकुलमार्गचत्वरा} %||2-60-18||

\fourlineindentedshloka
{नराश्च नार्यश्च समेत्य सङ्घशो}
{ विगर्हमाणा भरतस्य मातरम्}
{तदा नगर्यां नरदेवसङ्क्षये}
{ बभूवुरार्ता न च शर्म लेभिरे} %||2-61-19||


{॥इत्यार्षे श्रीमद्रामायणे वाल्मीकीये आदिकाव्ये अयोध्याकाण्डे षष्टितमः सर्गः॥}

\sect{एकषष्टितमः सर्गः}

\twolineshloka
{व्यतीतायां तु शर्वर्यामादित्यस्योदये ततः}
{समेत्य राजकर्तारः सभामीयुर्द्विजातयः} %||2-61-1||

\twolineshloka
{मार्कण्डेयोऽथ मौद्गल्यो वामदेवश्च काश्यपः}
{कात्ययनो गौतमश्च जाबालिश्च महायशाः} %||2-61-2||

\twolineshloka
{एते द्विजाः सहामात्यैः पृथग्वाचमुदीरयन्}
{वसिष्ठमेवाभिमुखाः श्रेष्ठो राजपुरोहितम्} %||2-61-3||

\twolineshloka
{अतीता शर्वरी दुःखं या नो वर्षशतोपमा}
{अस्मिन्पञ्चत्वमापन्ने पुत्रशोकेन पार्थिवे} %||2-61-4||

\twolineshloka
{स्वर्गतश्च महाराजो रामश्चारण्यमाश्रितः}
{लक्ष्मणश्चापि तेजस्वी रामेणैव गतः सह} %||2-61-5||

\twolineshloka
{उभौ भरतशत्रुघ्नौ क्केकयेषु परन्तपौ}
{पुरे राजगृहे रम्ये मातामहनिवेशने} %||2-61-6||

\twolineshloka
{इक्ष्वाकूणामिहाद्यैव कश्चिद्राजा विधीयताम्}
{अराजकं हि नो राष्ट्रं न विनाशमवाप्नुयात्} %||2-61-7||

\twolineshloka
{नाराजले जनपदे विद्युन्माली महास्वनः}
{अभिवर्षति पर्जन्यो महीं दिव्येन वारिणा} %||2-61-8||

\twolineshloka
{नाराजके जनपदे बीजमुष्टिः प्रकीर्यते}
{नाराकके पितुः पुत्रो भार्या वा वर्तते वशे} %||2-61-9||

\twolineshloka
{अराजके धनं नास्ति नास्ति भार्याप्यराजके}
{इदमत्याहितं चान्यत्कुतः सत्यमराजके} %||2-61-10||

\twolineshloka
{नाराजके जनपदे कारयन्ति सभां नराः}
{उद्यानानि च रम्याणि हृष्टाः पुण्यगृहाणि च} %||2-61-11||

\twolineshloka
{नाराजके जनपदे यज्ञशीला द्विजातयः}
{सत्राण्यन्वासते दान्ता ब्राह्मणाः संशितव्रताः} %||2-61-12||

\twolineshloka
{नाराजके जनपदे प्रभूतनटनर्तकाः}
{उत्सवाश्च समाजाश्च वर्धन्ते राष्ट्रवर्धनाः} %||2-61-13||

\twolineshloka
{नारजके जनपदे सिद्धार्था व्यवहारिणः}
{कथाभिरनुरज्यन्ते कथाशीलाः कथाप्रियैः} %||2-61-14||

\twolineshloka
{नाराजके जनपदे वाहनैः शीघ्रगामिभिः}
{नरा निर्यान्त्यरण्यानि नारीभिः सह कामिनः} %||2-61-15||

\twolineshloka
{नाराकजे जनपदे धनवन्तः सुरक्षिताः}
{शेरते विवृत द्वाराः कृषिगोरक्षजीविनः} %||2-61-16||

\twolineshloka
{नाराजके जनपदे वणिजो दूरगामिनः}
{गच्छन्ति क्षेममध्वानं बहुपुण्यसमाचिताः} %||2-61-17||

\twolineshloka
{नाराजके जनपदे चरत्येकचरो वशी}
{भावयन्नात्मनात्मानं यत्रसायङ्गृहो मुनिः} %||2-61-18||

\twolineshloka
{नाराजके जनपदे योगक्षेमं प्रवर्तते}
{न चाप्यराजके सेना शत्रून्विषहते युधि} %||2-61-19||

\twolineshloka
{यथा ह्यनुदका नद्यो यथा वाप्यतृणं वनम्}
{अगोपाला यथा गावस्तथा राष्ट्रमराजकम्} %||2-61-20||

\twolineshloka
{नाराजके जनपदे स्वकं भवति कस्यचित्}
{मत्स्या इव नरा नित्यं भक्षयन्ति परस्परम्} %||2-61-21||

\twolineshloka
{येहि सम्भिन्नमर्यादा नास्तिकाश्छिन्नसंशयाः}
{तेऽपि भावाय कल्पन्ते राजदण्डनिपीडिताः} %||2-61-22||

\twolineshloka
{अहो तम इवेदं स्यान्न प्रज्ञायेत किञ्चन}
{राजा चेन्न भवेँल्लोके विभजन्साध्वसाधुनी} %||2-61-23||

\twolineshloka
{जीवत्यपि महाराजे तवैव वचनं वयम्}
{नातिक्रमामहे सर्वे वेलां प्राप्येव सागरः} %||2-61-24||

\fourlineindentedshloka
{स नः समीक्ष्य द्विजवर्यवृत्तम्}
{ नृपं विना राज्यमरण्यभूतम्}
{कुमारमिक्ष्वाकुसुतं वदान्यम्}
{ त्वमेव राजानमिहाभिषिञ्चय} %||2-62-25||


{॥इत्यार्षे श्रीमद्रामायणे वाल्मीकीये आदिकाव्ये अयोध्याकाण्डे एकषष्टितमः सर्गः॥}

\sect{द्विषष्टितमः सर्गः}

\twolineshloka
{तेषां तद्वचनं श्रुत्वा वसिष्ठः प्रत्युवाच ह}
{मित्रामात्यगणान्सर्वान्ब्राह्मणांस्तानिदं वचः} %||2-62-1||

\twolineshloka
{यदसौ मातुलकुले पुरे राजगृहे सुखी}
{भरतो वसति भ्रात्रा शत्रुघ्नेन समन्वितः} %||2-62-2||

\twolineshloka
{तच्छीघ्रं जवना दूता गच्छन्तु त्वरितैर्हयैः}
{आनेतुं भ्रातरौ वीरौ किं समीक्षामहे वयम्} %||2-62-3||

\twolineshloka
{गच्छन्त्विति ततः सर्वे वसिष्ठं वाक्यमब्रुवन्}
{तेषां तद्वचनं श्रुत्वा वसिष्ठो वाक्यमब्रवीत्} %||2-62-4||

\twolineshloka
{एहि सिद्धार्थ विजय जयन्ताशोकनन्दन}
{श्रूयतामितिकर्तव्यं सर्वानेव ब्रवीमि वः} %||2-62-5||

\twolineshloka
{पुरं राजगृहं गत्वा शीघ्रं शीघ्रजवैर्हयैः}
{त्यक्तशोकैरिदं वाच्यः शासनाद्भरतो मम} %||2-62-6||

\twolineshloka
{पुरोहितस्त्वां कुशलं प्राह सर्वे च मन्त्रिणः}
{त्वरमाणश्च निर्याहि कृत्यमात्ययिकं त्वया} %||2-62-7||

\twolineshloka
{मा चास्मै प्रोषितं रामं मा चास्मै पितरं मृतम्}
{भवन्तः शंसिषुर्गत्वा राघवाणामिमं क्षयम्} %||2-62-8||

\threelineshloka
{कौशेयानि च वस्त्राणि भूषणानि वराणि च}
{क्षिप्रमादाय राज्ञश्च भरतस्य च गच्छत}
{वसिष्ठेनाभ्यनुज्ञाता दूताः सन्त्वरिता ययुः} %||2-62-9||

\twolineshloka
{ते हस्तिन पुरे गङ्गां तीर्त्वा प्रत्यङ्मुखा ययुः}
{पाञ्चालदेशमासाद्य मध्येन कुरुजाङ्गलम्} %||2-62-10||

\twolineshloka
{ते प्रसन्नोदकां दिव्यां नानाविहगसेविताम्}
{उपातिजग्मुर्वेगेन शरदण्डां जनाकुलाम्} %||2-62-11||

\twolineshloka
{निकूलवृक्षमासाद्य दिव्यं सत्योपयाचनम्}
{अभिगम्याभिवाद्यं तं कुलिङ्गां प्राविशन्पुरीम्} %||2-62-12||

\threelineshloka
{अभिकालं ततः प्राप्य तेजोऽभिभवनाच्च्युताः}
{ययुर्मध्येन बाह्लीकान्सुदामानं च पर्वतम्}
{विष्णोः पदं प्रेक्षमाणा विपाशां चापि शाल्मलीम्} %||2-62-13||

\twolineshloka
{ते श्रान्तवाहना दूता विकृष्टेन सता पथा}
{गिरि व्रजं पुर वरं शीघ्रमासेदुरञ्जसा} %||2-62-14||

\fourlineindentedshloka
{भर्तुः प्रियार्थं कुलरक्षणार्थम्}
{ भर्तुश्च वंशस्य परिग्रहार्थम्}
{अहेडमानास्त्वरया स्म दूता}
{ रात्र्यां तु ते तत्पुरमेव याताः} %||2-63-15||


{॥इत्यार्षे श्रीमद्रामायणे वाल्मीकीये आदिकाव्ये अयोध्याकाण्डे द्विषष्टितमः सर्गः॥}

\sect{त्रिषष्टितमः सर्गः}

\twolineshloka
{यामेव रात्रिं ते दूताः प्रविशन्ति स्म तां पुरीम्}
{भरतेनापि तां रात्रिं स्वप्नो दृष्टोऽयमप्रियः} %||2-63-1||

\twolineshloka
{व्युष्टामेव तु तां रात्रिं दृष्ट्वा तं स्वप्नमप्रियम्}
{पुत्रो राजाधिराजस्य सुभृशं पर्यतप्यत} %||2-63-2||

\twolineshloka
{तप्यमानं समाज्ञाय वयस्याः प्रियवादिनः}
{आयासं हि विनेष्यन्तः सभायां चक्रिरे कथाः} %||2-63-3||

\twolineshloka
{वादयन्ति तथा शान्तिं लासयन्त्यपि चापरे}
{नाटकान्यपरे प्राहुर्हास्यानि विविधानि च} %||2-63-4||

\twolineshloka
{स तैर्महात्मा भरतः सखिभिः प्रिय वादिभिः}
{गोष्ठीहास्यानि कुर्वद्भिर्न प्राहृष्यत राघवः} %||2-63-5||

\twolineshloka
{तमब्रवीत्प्रियसखो भरतं सखिभिर्वृतम्}
{सुहृद्भिः पर्युपासीनः किं सखे नानुमोदसे} %||2-63-6||

\twolineshloka
{एवं ब्रुवाणं सुहृदं भरतः प्रत्युवाच ह}
{शृणु त्वं यन्निमित्तम्मे दैन्यमेतदुपागतम्} %||2-63-7||

\twolineshloka
{स्वप्ने पितरमद्राक्षं मलिनं मुक्तमूर्धजम्}
{पतन्तमद्रिशिखरात्कलुषे गोमये ह्रदे} %||2-63-8||

\twolineshloka
{प्लवमानश्च मे दृष्टः स तस्मिन्गोमयह्रदे}
{पिबन्नञ्जलिना तैलं हसन्निव मुहुर्मुहुः} %||2-63-9||

\twolineshloka
{ततस्तिलोदनं भुक्त्वा पुनः पुनरधःशिराः}
{तैलेनाभ्यक्तसर्वाङ्गस्तैलमेवावगाहत} %||2-63-10||

\twolineshloka
{स्वप्नेऽपि सागरं शुष्कं चन्द्रं च पतितं भुवि}
{सहसा चापि संशन्तं ज्वलितं जातवेदसम्} %||2-63-11||

\twolineshloka
{अवदीर्णां च पृथिवीं शुष्कांश्च विविधान्द्रुमान्}
{अहं पश्यामि विध्वस्तान्सधूमांश्चैव पार्वतान्} %||2-63-12||

\twolineshloka
{पीठे कार्ष्णायसे चैनं निषण्णं कृष्णवाससम्}
{प्रहसन्ति स्म राजानं प्रमदाः कृष्णपिङ्गलाः} %||2-63-13||

\twolineshloka
{त्वरमाणश्च धर्मात्मा रक्तमाल्यानुलेपनः}
{रथेन खरयुक्तेन प्रयातो दक्षिणामुखः} %||2-63-14||

\twolineshloka
{एवमेतन्मया दृष्टमिमां रात्रिं भयावहाम्}
{अहं रामोऽथ वा राजा लक्ष्मणो वा मरिष्यति} %||2-63-15||

\threelineshloka
{नरो यानेन यः स्वप्ने खरयुक्तेन याति हि}
{अचिरात्तस्य धूमाग्रं चितायां सम्प्रदृश्यते}
{एतन्निमित्तं दीनोऽहं तन्न वः प्रतिपूजये} %||2-63-16||

\twolineshloka
{शुष्यतीव च मे कण्ठो न स्वस्थमिव मे मनः}
{जुगुप्सन्निव चात्मानं न च पश्यामि कारणम्} %||2-63-17||

\fourlineindentedshloka
{इमां हि दुःस्वप्नगतिं निशाम्य ता}
{मनेकरूपामवितर्कितां पुरा}
{भयं महत्तद्धृदयान्न याति मे}
{ विचिन्त्य राजानमचिन्त्यदर्शनम्} %||2-64-18||


{॥इत्यार्षे श्रीमद्रामायणे वाल्मीकीये आदिकाव्ये अयोध्याकाण्डे त्रिषष्टितमः सर्गः॥}

\sect{चतुःषष्टितमः सर्गः}

\twolineshloka
{भरते ब्रुवति स्वप्नं दूतास्ते क्लान्तवाहनाः}
{प्रविश्यासह्यपरिखं रम्यं राजगृहं पुरम्} %||2-64-1||

\twolineshloka
{समागम्य तु राज्ञा च राजपुत्रेण चार्चिताः}
{राज्ञः पादौ गृहीत्वा तु तमूचुर्भरतं वचः} %||2-64-2||

\twolineshloka
{पुरोहितस्त्वा कुशलं प्राह सर्वे च मन्त्रिणः}
{त्वरमाणश्च निर्याहि कृत्यमात्ययिकं त्वया} %||2-64-3||

\twolineshloka
{अत्र विंशतिकोट्यस्तु नृपतेर्मातुलस्य ते}
{दशकोट्यस्तु सम्पूर्णास्तथैव च नृपात्मज} %||2-64-4||

\twolineshloka
{प्रतिगृह्य च तत्सर्वं स्वनुरक्तः सुहृज्जने}
{दूतानुवाच भरतः कामैः सम्प्रतिपूज्य तान्} %||2-64-5||

\twolineshloka
{कच्चित्सुकुशली राजा पिता दशरथो मम}
{कच्चिच्चारागता रामे लक्ष्मणे वा महात्मनि} %||2-64-6||

\twolineshloka
{आर्या च धर्मनिरता धर्मज्ञा धर्मदर्शिनी}
{अरोगा चापि कौसल्या माता रामस्य धीमतः} %||2-64-7||

\twolineshloka
{कच्चित्सुमित्रा धर्मज्ञा जननी लक्ष्मणस्य या}
{शत्रुघ्नस्य च वीरस्य सारोगा चापि मध्यमा} %||2-64-8||

\twolineshloka
{आत्मकामा सदा चण्डी क्रोधना प्राज्ञमानिनी}
{अरोगा चापि कैकेयी माता मे किमुवाच ह} %||2-64-9||

\threelineshloka
{एवमुक्तास्तु ते दूता भरतेन महात्मना}
{ऊचुः सम्प्रश्रितं वाक्यमिदं तं भरतं तदा}
{कुशलास्ते नरव्याघ्र येषां कुशलमिच्छसि} %||2-64-10||

\twolineshloka
{भरतश्चापि तान्दूतानेवमुक्तोऽभ्यभाषत}
{आपृच्छेऽहं महाराजं दूताः सन्त्वरयन्ति माम्} %||2-64-11||

\twolineshloka
{एवमुक्त्वा तु तान्दूतान्भरतः पार्थिवात्मजः}
{दूतैः सञ्चोदितो वाक्यं मातामहमुवाच ह} %||2-64-12||

\twolineshloka
{राजन्पितुर्गमिष्यामि सकाशं दूतचोदितः}
{पुनरप्यहमेष्यामि यदा मे त्वं स्मरिष्यसि} %||2-64-13||

\twolineshloka
{भरतेनैवमुक्तस्तु नृपो मातामहस्तदा}
{तमुवाच शुभं वाक्यं शिरस्याघ्राय राघवम्} %||2-64-14||

\twolineshloka
{गच्छ तातानुजाने त्वां कैकेयी सुप्रजास्त्वया}
{मातरं कुशलं ब्रूयाः पितरं च परन्तप} %||2-64-15||

\twolineshloka
{पुरोहितं च कुशलं ये चान्ये द्विजसत्तमाः}
{तौ च तात महेष्वासौ भ्रातरु रामलक्ष्मणौ} %||2-64-16||

\twolineshloka
{तस्मै हस्त्युत्तमांश्चित्रान्कम्बलानजिनानि च}
{अभिसत्कृत्य कैकेयो भरताय धनं ददौ} %||2-64-17||

\twolineshloka
{रुक्म निष्कसहस्रे द्वे षोडशाश्वशतानि च}
{सत्कृत्य कैकेयी पुत्रं केकयो धनमादिशत्} %||2-64-18||

\twolineshloka
{तथामात्यानभिप्रेतान्विश्वास्यांश्च गुणान्वितान्}
{ददावश्वपतिः शीघ्रं भरतायानुयायिनः} %||2-64-19||

\twolineshloka
{ऐरावतानैन्द्रशिरान्नागान्वै प्रियदर्शनान्}
{खराञ्शीघ्रान्सुसंयुक्तान्मातुलोऽस्मै धनं ददौ} %||2-64-20||

\twolineshloka
{अन्तःपुरेऽतिसंवृद्धान्व्याघ्रवीर्यबलान्वितान्}
{दंष्ट्रायुधान्महाकायाञ्शुनश्चोपायनं ददौ} %||2-64-21||

\twolineshloka
{स मातामहमापृच्छ्य मातुलं च युधाजितम्}
{रथमारुह्य भरतः शत्रुघ्नसहितो ययौ} %||2-64-22||

\twolineshloka
{रथान्मण्डलचक्रांश्च योजयित्वा परःशतम्}
{उष्ट्रगोऽश्वखरैर्भृत्या भरतं यान्तमन्वयुः} %||2-64-23||

\fourlineindentedshloka
{बलेन गुप्तो भरतो महात्मा}
{ सहार्यकस्यात्मसमैरमात्यैः}
{आदाय शत्रुघ्नमपेतशत्रु}
{र्गृहाद्ययौ सिद्ध इवेन्द्रलोकात्} %||2-65-24||


{॥इत्यार्षे श्रीमद्रामायणे वाल्मीकीये आदिकाव्ये अयोध्याकाण्डे चतुःषष्टितमः सर्गः॥}

\sect{पञ्चषष्टितमः सर्गः}

\threelineshloka
{स प्राङ्मुखो राजगृहादभिनिर्याय वीर्यवान्}
{ह्रादिनीं दूरपारां च प्रत्यक्स्रोतस्तरङ्गिणीम्}
{शतद्रूमतरच्छ्रीमान्नदीमिक्ष्वाकुनन्दनः} %||2-65-1||

\twolineshloka
{एलधाने नदीं तीर्त्वा प्राप्य चापरपर्पटान्}
{शिलामाकुर्वतीं तीर्त्वा आग्नेयं शल्यकर्तनम्} %||2-65-2||

\twolineshloka
{सत्यसन्धः शुचिः श्रीमान्प्रेक्षमाणः शिलावहाम्}
{अत्ययात्स महाशैलान्वनं चैत्ररथं प्रति} %||2-65-3||

\twolineshloka
{वेगिनीं च कुलिङ्गाख्यां ह्रादिनीं पर्वतावृताम्}
{यमुनां प्राप्य सन्तीर्णो बलमाश्वासयत्तदा} %||2-65-4||

\twolineshloka
{शीतीकृत्य तु गात्राणि क्लान्तानाश्वास्य वाजिनः}
{तत्र स्नात्वा च पीत्वा च प्रायादादाय चोदकम्} %||2-65-5||

\twolineshloka
{राजपुत्रो महारण्यमनभीक्ष्णोपसेवितम्}
{भद्रो भद्रेण यानेन मारुतः खमिवात्ययात्} %||2-65-6||

\twolineshloka
{तोरणं दक्षिणार्धेन जम्बूप्रस्थमुपागमत्}
{वरूथं च ययौ रम्यं ग्रामं दशरथात्मजः} %||2-65-7||

\twolineshloka
{तत्र रम्ये वने वासं कृत्वासौ प्राङ्मुखो ययौ}
{उद्यानमुज्जिहानायाः प्रियका यत्र पादपाः} %||2-65-8||

\twolineshloka
{सालांस्तु प्रियकान्प्राप्य शीघ्रानास्थाय वाजिनः}
{अनुज्ञाप्याथ भरतो वाहिनीं त्वरितो ययौ} %||2-65-9||

\twolineshloka
{वासं कृत्वा सर्वतीर्थे तीर्त्वा चोत्तानकां नदीम्}
{अन्या नदीश्च विविधाः पार्वतीयैस्तुरङ्गमैः} %||2-65-10||

\threelineshloka
{हस्तिपृष्ठकमासाद्य कुटिकामत्यवर्तत}
{ततार च नरव्याघ्रो लौहित्ये स कपीवतीम्}
{एकसाले स्थाणुमतीं विनते गोमतीं नदीम्} %||2-65-11||

\twolineshloka
{कलिङ्ग नगरे चापि प्राप्य सालवनं तदा}
{भरतः क्षिप्रमागच्छत्सुपरिश्रान्तवाहनः} %||2-65-12||

\twolineshloka
{वनं च समतीत्याशु शर्वर्यामरुणोदये}
{अयोध्यां मनुना राज्ञा निर्मितां स ददर्श ह} %||2-65-13||

\twolineshloka
{तां पुरीं पुरुषव्याघ्रः सप्तरात्रोषिटः पथि}
{अयोध्यामग्रतो दृष्ट्वा रथे सारथिमब्रवीत्} %||2-65-14||

\twolineshloka
{एषा नातिप्रतीता मे पुण्योद्याना यशस्विनी}
{अयोध्या दृश्यते दूरात्सारथे पाण्डुमृत्तिका} %||2-65-15||

\twolineshloka
{यज्वभिर्गुणसम्पन्नैर्ब्राह्मणैर्वेदपारगैः}
{भूयिष्ठमृष्हैराकीर्णा राजर्षिवरपालिता} %||2-65-16||

\twolineshloka
{अयोध्यायां पुराशब्दः श्रूयते तुमुलो महान्}
{समन्तान्नरनारीणां तमद्य न शृणोम्यहम्} %||2-65-17||

\twolineshloka
{उद्यानानि हि सायाह्ने क्रीडित्वोपरतैर्नरैः}
{समन्ताद्विप्रधावद्भिः प्रकाशन्ते ममान्यदा} %||2-65-18||

\twolineshloka
{तान्यद्यानुरुदन्तीव परित्यक्तानि कामिभिः}
{अरण्यभूतेव पुरी सारथे प्रतिभाति मे} %||2-65-19||

\twolineshloka
{न ह्यत्र यानैर्दृश्यन्ते न गजैर्न च वाजिभिः}
{निर्यान्तो वाभियान्तो वा नरमुख्या यथापुरम्} %||2-65-20||

\twolineshloka
{अनिष्टानि च पापानि पश्यामि विविधानि च}
{निमित्तान्यमनोज्ञानि तेन सीदति ते मनः} %||2-65-21||

\twolineshloka
{द्वारेण वैजयन्तेन प्राविशच्छ्रान्तवाहनः}
{द्वाःस्थैरुत्थाय विजयं पृष्टस्तैः सहितो ययौ} %||2-65-22||

\twolineshloka
{स त्वनेकाग्रहृदयो द्वाःस्थं प्रत्यर्च्य तं जनम्}
{सूतमश्वपतेः क्लान्तमब्रवीत्तत्र राघवः} %||2-65-23||

\twolineshloka
{श्रुता नो यादृशाः पूर्वं नृपतीनां विनाशने}
{आकारास्तानहं सर्वानिह पश्यामि सारथे} %||2-65-24||

\twolineshloka
{मलिनं चाश्रुपूर्णाक्षं दीनं ध्यानपरं कृशम्}
{सस्त्री पुंसं च पश्यामि जनमुत्कण्ठितं पुरे} %||2-65-25||

\twolineshloka
{इत्येवमुक्त्वा भरतः सूतं तं दीनमानसः}
{तान्यनिष्टान्ययोध्यायां प्रेक्ष्य राजगृहं ययौ} %||2-65-26||

\fourlineindentedshloka
{तां शून्यशृङ्गाटकवेश्मरथ्याम्}
{ रजोऽरुणद्वारकपाटयन्त्राम्}
{दृष्ट्वा पुरीमिन्द्रपुरी प्रकाशाम्}
{ दुःखेन सम्पूर्णतरो बभूव} %||2-65-27||

\fourlineindentedshloka
{बहूनि पश्यन्मनसोऽप्रियाणि}
{ यान्यन्न्यदा नास्य पुरे बभूवुः}
{अवाक्शिरा दीनमना नहृष्टः}
{ पितुर्महात्मा प्रविवेश वेश्म} %||2-66-28||


{॥इत्यार्षे श्रीमद्रामायणे वाल्मीकीये आदिकाव्ये अयोध्याकाण्डे पञ्चषष्टितमः सर्गः॥}

\sect{षट्षष्टितमः सर्गः}

\twolineshloka
{अपश्यंस्तु ततस्तत्र पितरं पितुरालये}
{जगाम भरतो द्रष्टुं मातरं मातुरालये} %||2-66-1||

\twolineshloka
{अनुप्राप्तं तु तं दृष्ट्वा कैकेयी प्रोषितं सुतम्}
{उत्पपात तदा हृष्टा त्यक्त्वा सौवर्णमानसम्} %||2-66-2||

\twolineshloka
{स प्रविश्यैव धर्मात्मा स्वगृहं श्रीविवर्जितम्}
{भरतः प्रेक्ष्य जग्राह जनन्याश्चरणौ शुभौ} %||2-66-3||

\twolineshloka
{तं मूर्ध्नि समुपाघ्राय परिष्वज्य यशस्विनम्}
{अङ्के भरतमारोप्य प्रष्टुं समुपचक्रमे} %||2-66-4||

\twolineshloka
{अद्य ते कतिचिद्रात्र्यश्च्युतस्यार्यकवेश्मनः}
{अपि नाध्वश्रमः शीघ्रं रथेनापततस्तव} %||2-66-5||

\twolineshloka
{आर्यकस्ते सुकुशलो युधाजिन्मातुलस्तव}
{प्रवासाच्च सुखं पुत्र सर्वं मे वक्तुमर्हसि} %||2-66-6||

\twolineshloka
{एवं पृष्ठस्तु कैकेय्या प्रियं पार्थिवनन्दनः}
{आचष्ट भरतः सर्वं मात्रे राजीवलोचनः} %||2-66-7||

\twolineshloka
{अद्य मे सप्तमी रात्रिश्च्युतस्यार्यकवेश्मनः}
{अम्बायाः कुशली तातो युधाजिन्मातुलश्च मे} %||2-66-8||

\twolineshloka
{यन्मे धनं च रत्नं च ददौ राजा परन्तपः}
{परिश्रान्तं पथ्यभवत्ततोऽहं पूर्वमागतः} %||2-66-9||

\twolineshloka
{राजवाक्यहरैर्दूतैस्त्वर्यमाणोऽहमागतः}
{यदहं प्रष्टुमिच्छामि तदम्बा वक्तुमर्हसि} %||2-66-10||

\twolineshloka
{शून्योऽयं शयनीयस्ते पर्यङ्को हेमभूषितः}
{न चायमिक्ष्वाकुजनः प्रहृष्टः प्रतिभाति मे} %||2-66-11||

\twolineshloka
{राजा भवति भूयिष्ठ्गमिहाम्बाया निवेशने}
{तमहं नाद्य पश्यामि द्रष्टुमिच्छन्निहागतः} %||2-66-12||

\twolineshloka
{पितुर्ग्रहीष्ये चरणौ तं ममाख्याहि पृच्छतः}
{आहोस्विदम्ब ज्येष्ठायाः कौसल्याया निवेशने} %||2-66-13||

\threelineshloka
{तं प्रत्युवाच कैकेयी प्रियवद्घोरमप्रियम्}
{अजानन्तं प्रजानन्ती राज्यलोभेन मोहिता}
{या गतिः सर्वभूतानां तां गतिं ते पिता गतः} %||2-66-14||

\twolineshloka
{तच्छ्रुत्वा भरतो वाक्यं धर्माभिजनवाञ्शुचिः}
{पपात सहसा भूमौ पितृशोकबलार्दितः} %||2-66-15||

\twolineshloka
{ततः शोकेन संवीतः पितुर्मरणदुःखितः}
{विललाप महातेजा भ्रान्ताकुलितचेतनः} %||2-66-16||

\twolineshloka
{एतत्सुरुचिरं भाति पितुर्मे शयनं पुरा}
{तदिदं न विभात्यद्य विहीनं तेन धीमता} %||2-66-17||

\twolineshloka
{तमार्तं देवसङ्काशं समीक्ष्य पतितं भुवि}
{उत्थापयित्वा शोकार्तं वचनं चेदमब्रवीत्} %||2-66-18||

\twolineshloka
{उत्तिष्ठोत्तिष्ठ किं शेषे राजपुत्र महायशः}
{त्वद्विधा न हि शोचन्ति सन्तः सदसि सम्मताः} %||2-66-19||

\twolineshloka
{स रुदत्या चिरं कालं भूमौ विपरिवृत्य च}
{जननीं प्रत्युवाचेदं शोकैर्बहुभिरावृतः} %||2-66-20||

\twolineshloka
{अभिषेक्ष्यति रामं तु राजा यज्ञं नु यक्ष्यति}
{इत्यहं कृतसङ्कल्पो हृष्टो यात्रामयासिषम्} %||2-66-21||

\twolineshloka
{तदिदं ह्यन्यथा भूतं व्यवदीर्णं मनो मम}
{पितरं यो न पश्यामि नित्यं प्रियहिते रतम्} %||2-66-22||

\twolineshloka
{अम्ब केनात्यगाद्राजा व्याधिना मय्यनागते}
{धन्या रामादयः सर्वे यैः पिता संस्कृतः स्वयम्} %||2-66-23||

\twolineshloka
{न नूनं मां महाराजः प्राप्तं जानाति कीर्तिमान्}
{उपजिघ्रेद्धि मां मूर्ध्नि तातः संनम्य सत्वरम्} %||2-66-24||

\twolineshloka
{क्व स पाणिः सुखस्पर्शस्तातस्याक्लिष्टकर्मणः}
{येन मां रजसा ध्वस्तमभीक्ष्णं परिमार्जति} %||2-66-25||

\twolineshloka
{यो मे भ्राता पिता बन्धुर्यस्य दासोऽस्मि धीमतः}
{तस्य मां शीघ्रमाख्याहि रामस्याक्लिष्ट कर्मणः} %||2-66-26||

\twolineshloka
{पिता हि भवति ज्येष्ठो धर्ममार्यस्य जानतः}
{तस्य पादौ ग्रहीष्यामि स हीदानीं गतिर्मम} %||2-66-27||

\twolineshloka
{आर्ये किमब्रवीद्राजा पिता मे सत्यविक्रमः}
{पश्चिमं साधुसन्देशमिच्छामि श्रोतुमात्मनः} %||2-66-28||

\threelineshloka
{इति पृष्टा यथातत्त्वं कैकेयी वाक्यमब्रवीत्}
{रामेति राजा विलपन्हा सीते लक्ष्मणेति च}
{स महात्मा परं लोकं गतो गतिमतां वरः} %||2-66-29||

\twolineshloka
{इमां तु पश्चिमां वाचं व्याजहार पिता तव}
{काल धर्मपरिक्षिप्तः पाशैरिव महागजः} %||2-66-30||

\twolineshloka
{सिद्धार्थास्तु नरा राममागतं सीतया सह}
{लक्ष्मणं च महाबाहुं द्रक्ष्यन्ति पुनरागतम्} %||2-66-31||

\twolineshloka
{तच्छ्रुत्वा विषसादैव द्वितीया प्रियशंसनात्}
{विषण्णवदनो भूत्वा भूयः पप्रच्छ मातरम्} %||2-66-32||

\twolineshloka
{क्व चेदानीं स धर्मात्मा कौसल्यानन्दवर्धनः}
{लक्ष्मणेन सह भ्रात्रा सीतया च समं गतः} %||2-66-33||

\twolineshloka
{तथा पृष्टा यथातत्त्वमाख्यातुमुपचक्रमे}
{मातास्य युगपद्वाक्यं विप्रियं प्रियशङ्कया} %||2-66-34||

\twolineshloka
{स हि राजसुतः पुत्र चीरवासा महावनम्}
{दण्डकान्सह वैदेह्या लक्ष्मणानुचरो गतः} %||2-66-35||

\twolineshloka
{तच्छ्रुत्वा भरतस्त्रस्तो भ्रातुश्चारित्रशङ्कया}
{स्वस्य वंशस्य माहात्म्यात्प्रष्टुं समुपचक्रमे} %||2-66-36||

\twolineshloka
{कच्चिन्न ब्राह्मणवधं हृतं रामेण कस्यचित्}
{कच्चिन्नाढ्यो दरिद्रो वा तेनापापो विहिंसितः} %||2-66-37||

\twolineshloka
{कच्चिन्न परदारान्वा राजपुत्रोऽभिमन्यते}
{कस्मात्स दण्डकारण्ये भ्रूणहेव विवासितः} %||2-66-38||

\twolineshloka
{अथास्य चपला माता तत्स्वकर्म यथातथम्}
{तेनैव स्त्रीस्वभावेन व्याहर्तुमुपचक्रमे} %||2-66-39||

\threelineshloka
{न ब्राह्मण धनं किञ्चिद्धृतं रामेण कस्यचित्}
{कश्चिन्नाढ्यो दरिद्रो वा तेनापापो विहिंसितः}
{न रामः परदारांश्च चक्षुर्भ्यामपि पश्यति} %||2-66-40||

\twolineshloka
{मया तु पुत्र श्रुत्वैव रामस्यैवाभिषेचनम्}
{याचितस्ते पिता राज्यं रामस्य च विवासनम्} %||2-66-41||

\twolineshloka
{स स्ववृत्तिं समास्थाय पिता ते तत्तथाकरोत्}
{रामश्च सहसौमित्रिः प्रेषितः सह सीतया} %||2-66-42||

\twolineshloka
{तमपश्यन्प्रियं पुत्रं महीपालो महायशाः}
{पुत्रशोकपरिद्यूनः पञ्चत्वमुपपेदिवान्} %||2-66-43||

\twolineshloka
{त्वया त्विदानीं धर्मज्ञ राजत्वमवलम्ब्यताम्}
{त्वत्कृते हि मया सर्वमिदमेवंविधं कृतम्} %||2-66-44||

\fourlineindentedshloka
{तत्पुत्र शीघ्रं विधिना विधिज्ञै}
{र्वसिष्ठमुख्यैः सहितो द्विजेन्द्रैः}
{सङ्काल्य राजानमदीनसत्त्व}
{मात्मानमुर्व्यामभिषेचयस्व} %||2-67-45||


{॥इत्यार्षे श्रीमद्रामायणे वाल्मीकीये आदिकाव्ये अयोध्याकाण्डे षट्षष्टितमः सर्गः॥}

\sect{सप्तषष्टितमः सर्गः}

\twolineshloka
{श्रुत्वा तु पितरं वृत्तं भ्रातरु च विवासितौ}
{भरतो दुःखसन्तप्त इदं वचनमब्रवीत्} %||2-67-1||

\twolineshloka
{किं नुण्कार्यं हतस्येह मम राज्येन शोचतः}
{विहीनस्याथ पित्रा च भ्रात्रा पितृसमेन च} %||2-67-2||

\twolineshloka
{दुःखे मे दुःखमकरोर्व्रणे क्षारमिवादधाः}
{राजानं प्रेतभावस्थं कृत्वा रामं च तापसम्} %||2-67-3||

\twolineshloka
{कुलस्य त्वमभावाय कालरात्रिरिवागता}
{अङ्गारमुपगूह्य स्म पिता मे नावबुद्धवान्} %||2-67-4||

\twolineshloka
{कौसल्या च सुमित्रा च पुत्रशोकाभिपीडिते}
{दुष्करं यदि जीवेतां प्राप्य त्वां जननीं मम} %||2-67-5||

\twolineshloka
{ननु त्वार्योऽपि धर्मात्मा त्वयि वृत्तिमनुत्तमाम्}
{वर्तते गुरुवृत्तिज्ञो यथा मातरि वर्तते} %||2-67-6||

\twolineshloka
{तथा ज्येष्ठा हि मे माता कौसल्या दीर्घदर्शिनी}
{त्वयि धर्मं समास्थाय भगिन्यामिव वर्तते} %||2-67-7||

\twolineshloka
{तस्याः पुत्रं कृतात्मानं चीरवल्कलवाससम्}
{प्रस्थाप्य वनवासाय कथं पापे न शोचसि} %||2-67-8||

\twolineshloka
{अपापदर्शिनं शूरं कृतात्मानं यशस्विनम्}
{प्रव्राज्य चीरवसनं किं नु पश्यसि कारणम्} %||2-67-9||

\twolineshloka
{लुब्धाया विदितो मन्ये न तेऽहं राघवं प्रति}
{तथा ह्यनर्थो राज्यार्थं त्वया नीतो महानयम्} %||2-67-10||

\twolineshloka
{अहं हि पुरुषव्याघ्रावपश्यन्रामलक्ष्मणौ}
{केन शक्तिप्रभावेन राज्यं रक्षितुमुत्सहे} %||2-67-11||

\twolineshloka
{तं हि नित्यं महाराजो बलवन्तं महाबलः}
{अपाश्रितोऽभूद्धर्मात्मा मेरुर्मेरुवनं यथा} %||2-67-12||

\twolineshloka
{सोऽहं कथमिमं भारं महाधुर्यसमुद्यतम्}
{दम्यो धुरमिवासाद्य सहेयं केन चौजसा} %||2-67-13||

\threelineshloka
{अथ वा मे भवेच्छक्तिर्योगैर्बुद्धिबलेन वा}
{सकामां न करिष्यामि त्वामहं पुत्रगर्धिनीम्}
{निवर्तयिष्यामि वनाद्भ्रातरं स्वजनप्रियम्} %||2-67-14||

\fourlineindentedshloka
{इत्येवमुक्त्वा भरतो महात्मा}
{ प्रियेतरैर्वाक्यगणैस्तुदंस्ताम्}
{शोकातुरश्चापि ननाद भूयः}
{ सिंहो यथा पर्वतगह्वरस्थः} %||2-68-15||


{॥इत्यार्षे श्रीमद्रामायणे वाल्मीकीये आदिकाव्ये अयोध्याकाण्डे सप्तषष्टितमः सर्गः॥}

\sect{अष्टषष्टितमः सर्गः}

\twolineshloka
{तां तथा गर्हयित्वा तु मातरं भरतस्तदा}
{रोषेण महताविष्टः पुनरेवाब्रवीद्वचः} %||2-68-1||

\twolineshloka
{राज्याद्भ्रंशस्व कैकेयि नृशंसे दुष्टचारिणि}
{परित्यक्ता च धर्मेण मा मृतं रुदती भव} %||2-68-2||

\twolineshloka
{किं नु तेऽदूषयद्राजा रामो वा भृशधार्मिकः}
{ययोर्मृत्युर्विवासश्च त्वत्कृते तुल्यमागतौ} %||2-68-3||

\twolineshloka
{भ्रूणहत्यामसि प्राप्ता कुलस्यास्य विनाशनात्}
{कैकेयि नरकं गच्छ मा च भर्तुः सलोकताम्} %||2-68-4||

\twolineshloka
{यत्त्वया हीदृशं पापं कृतं घोरेण कर्मणा}
{सर्वलोकप्रियं हित्वा ममाप्यापादितं भयम्} %||2-68-5||

\twolineshloka
{त्वत्कृते मे पिता वृत्तो रामश्चारण्यमाश्रितः}
{अयशो जीवलोके च त्वयाहं प्रतिपादितः} %||2-68-6||

\twolineshloka
{मातृरूपे ममामित्रे नृशंसे राज्यकामुके}
{न तेऽहमभिभाष्योऽस्मि दुर्वृत्ते पतिघातिनि} %||2-68-7||

\twolineshloka
{कौसल्या च सुमित्रा च याश्चान्या मम मातरः}
{दुःखेन महताविष्टास्त्वां प्राप्य कुलदूषिणीम्} %||2-68-8||

\twolineshloka
{न त्वमश्वपतेः कन्या धर्मराजस्य धीमतः}
{राक्षसी तत्र जातासि कुलप्रध्वंसिनी पितुः} %||2-68-9||

\twolineshloka
{यत्त्वया धार्मिको रामो नित्यं सत्यपरायणः}
{वनं प्रस्थापितो दुःखात्पिता च त्रिदिवं गतः} %||2-68-10||

\twolineshloka
{यत्प्रधानासि तत्पापं मयि पित्रा विनाकृते}
{भ्रातृभ्यां च परित्यक्ते सर्वलोकस्य चाप्रिये} %||2-68-11||

\twolineshloka
{कौसल्यां धर्मसंयुक्तां वियुक्तां पापनिश्चये}
{कृत्वा कं प्राप्स्यसे त्वद्य लोकं निरयगामिनी} %||2-68-12||

\twolineshloka
{किं नावबुध्यसे क्रूरे नियतं बन्धुसंश्रयम्}
{ज्येष्ठं पितृसमं रामं कौसल्यायात्मसम्भवम्} %||2-68-13||

\twolineshloka
{अङ्गप्रत्यङ्गजः पुत्रो हृदयाच्चापि जायते}
{तस्मात्प्रियतरो मातुः प्रियत्वान्न तु बान्धवः} %||2-68-14||

\twolineshloka
{अन्यदा किल धर्मज्ञा सुरभिः सुरसम्मता}
{वहमानौ ददर्शोर्व्यां पुत्रौ विगतचेतसौ} %||2-68-15||

\twolineshloka
{तावर्धदिवसे श्रान्तौ दृष्ट्वा पुत्रौ महीतले}
{रुरोद पुत्र शोकेन बाष्पपर्याकुलेक्षणा} %||2-68-16||

\twolineshloka
{अधस्ताद्व्रजतस्तस्याः सुरराज्ञो महात्मनः}
{बिन्दवः पतिता गात्रे सूक्ष्माः सुरभिगन्धिनः} %||2-68-17||

\twolineshloka
{तां दृष्ट्वा शोकसन्तप्तां वज्रपाणिर्यशस्विनीम्}
{इन्द्रः प्राञ्जलिरुद्विग्नः सुरराजोऽब्रवीद्वचः} %||2-68-18||

\twolineshloka
{भयं कच्चिन्न चास्मासु कुतश्चिद्विद्यते महत्}
{कुतो निमित्तः शोकस्ते ब्रूहि सर्वहितैषिणि} %||2-68-19||

\twolineshloka
{एवमुक्ता तु सुरभिः सुरराजेन धीमता}
{पत्युवाच ततो धीरा वाक्यं वाक्यविशारदा} %||2-68-20||

\twolineshloka
{शान्तं पातं न वः किञ्चित्कुतश्चिदमराधिप}
{अहं तु मग्नौ शोचामि स्वपुत्रौ विषमे स्थितौ} %||2-68-21||

\twolineshloka
{एतौ दृष्ट्वा कृषौ दीनौ सूर्यरश्मिप्रतापिनौ}
{वध्यमानौ बलीवर्दौ कर्षकेण सुराधिप} %||2-68-22||

\twolineshloka
{मम कायात्प्रसूतौ हि दुःखितौ भार पीडितौ}
{यौ दृष्ट्वा परितप्येऽहं नास्ति पुत्रसमः प्रियः} %||2-68-23||

\twolineshloka
{यस्याः पुत्र सहस्राणि सापि शोचति कामधुक्}
{किं पुनर्या विना रामं कौसल्या वर्तयिष्यति} %||2-68-24||

\twolineshloka
{एकपुत्रा च साध्वी च विवत्सेयं त्वया कृता}
{तस्मात्त्वं सततं दुःखं प्रेत्य चेह च लप्स्यसे} %||2-68-25||

\twolineshloka
{अहं ह्यपचितिं भ्रातुः पितुश्च सकलामिमाम्}
{वर्धनं यशसश्चापि करिष्यामि न संशयः} %||2-68-26||

\twolineshloka
{आनाययित्वा तनयं कौसल्याया महाद्युतिम्}
{स्वयमेव प्रवेक्ष्यामि वनं मुनिनिषेवितम्} %||2-68-27||

\twolineshloka
{इति नाग इवारण्ये तोमराङ्कुशचोदितः}
{पपात भुवि सङ्क्रुद्धो निःश्वसन्निव पन्नगः} %||2-68-28||

\fourlineindentedshloka
{संरक्तनेत्रः शिथिलाम्बरस्तदा}
{ विधूतसर्वाभरणः परन्तपः}
{बभूव भूमौ पतितो नृपात्मजः}
{ शचीपतेः केतुरिवोत्सवक्षये} %||2-69-29||


{॥इत्यार्षे श्रीमद्रामायणे वाल्मीकीये आदिकाव्ये अयोध्याकाण्डे अष्टषष्टितमः सर्गः॥}

\sect{एकोनसप्ततितमः सर्गः}

\twolineshloka
{तथैव क्रोशतस्तस्य भरतस्य महात्मनः}
{कौसल्या शब्दमाज्ञाय सुमित्रामिदमब्रवीत्} %||2-69-1||

\twolineshloka
{आगतः क्रूरकार्यायाः कैकेय्या भरतः सुतः}
{तमहं द्रष्टुमिच्छामि भरतं दीर्घदर्शिनम्} %||2-69-2||

\twolineshloka
{एवमुक्त्वा सुमित्रां सा विवर्णा मलिनाम्बरा}
{प्रतस्थे भरतो यत्र वेपमाना विचेतना} %||2-69-3||

\twolineshloka
{स तु रामानुजश्चापि शत्रुघ्नसहितस्तदा}
{प्रतस्थे भरतो यत्र कौसल्याया निवेशनम्} %||2-69-4||

\twolineshloka
{ततः शत्रुघ्न भरतौ कौसल्यां प्रेक्ष्य दुःखितौ}
{पर्यष्वजेतां दुःखार्तां पतितां नष्टचेतनाम्} %||2-69-5||

\threelineshloka
{भरतं प्रत्युवाचेदं कौसल्या भृशदुःखिता}
{इदं ते राज्यकामस्य राज्यं प्राप्तमकण्टकम्}
{सम्प्राप्तं बत कैकेय्या शीघ्रं क्रूरेण कर्मणा} %||2-69-6||

\twolineshloka
{प्रस्थाप्य चीरवसनं पुत्रं मे वनवासिनम्}
{कैकेयी कं गुणं तत्र पश्यति क्रूरदर्शिनी} %||2-69-7||

\twolineshloka
{क्षिप्रं मामपि कैकेयी प्रस्थापयितुमर्हति}
{हिरण्यनाभो यत्रास्ते सुतो मे सुमहायशाः} %||2-69-8||

\twolineshloka
{अथ वा स्वयमेवाहं सुमित्रानुचरा सुखम्}
{अग्निहोत्रं पुरस्कृत्य प्रस्थास्ये यत्र राघवः} %||2-69-9||

\twolineshloka
{कामं वा स्वयमेवाद्य तत्र मां नेतुमर्हसि}
{यत्रासौ पुरुषव्याघ्रस्तप्यते मे तपः सुतः} %||2-69-10||

\twolineshloka
{इदं हि तव विस्तीर्णं धनधान्यसमाचितम्}
{हस्त्यश्वरथसम्पूर्णं राज्यं निर्यातितं तया} %||2-69-11||

\twolineshloka
{एवं विलपमानां तां भरतः प्राञ्जलिस्तदा}
{कौसल्यां प्रत्युवाचेदं शोकैर्बहुभिरावृताम्} %||2-69-12||

\twolineshloka
{आर्ये कस्मादजानन्तं गर्हसे मामकिल्बिषम्}
{विपुलां च मम प्रीतिं स्थिरां जानासि राघवे} %||2-69-13||

\twolineshloka
{कृता शास्त्रानुगा बुद्धिर्मा भूत्तस्य कदाचन}
{सत्यसन्धः सतां श्रेष्ठो यस्यार्योऽनुमते गतः} %||2-69-14||

\twolineshloka
{प्रैष्यं पापीयसां यातु सूर्यं च प्रति मेहतु}
{हन्तु पादेन गां सुप्तां यस्यार्योऽनुमते गतः} %||2-69-15||

\twolineshloka
{कारयित्वा महत्कर्म भर्ता भृत्यमनर्थकम्}
{अधर्मो योऽस्य सोऽस्यास्तु यस्यार्योऽनुमते गतः} %||2-69-16||

\twolineshloka
{परिपालयमानस्य राज्ञो भूतानि पुत्रवत्}
{ततस्तु द्रुह्यतां पापं यस्यार्योऽनुमते गतः} %||2-69-17||

\twolineshloka
{बलिषड्भागमुद्धृत्य नृपस्यारक्षतः प्रजाः}
{अधर्मो योऽस्य सोऽस्यास्तु यस्यार्योऽनुमते गतः} %||2-69-18||

\twolineshloka
{संश्रुत्य च तपस्विभ्यः सत्रे वै यज्ञदक्षिणाम्}
{तां विप्रलपतां पापं यस्यार्योऽनुमते गतः} %||2-69-19||

\twolineshloka
{हस्त्यश्वरथसम्बाधे युद्धे शस्त्रसमाकुले}
{मा स्म कार्षीत्सतां धर्मं यस्यार्योऽनुमते गतः} %||2-69-20||

\twolineshloka
{उपदिष्टं सुसूक्ष्मार्थं शास्त्रं यत्नेन धीमता}
{स नाशयतु दुष्टात्मा यस्यार्योऽनुमते गतः} %||2-69-21||

\twolineshloka
{पायसं कृसरं छागं वृथा सोऽश्नातु निर्घृणः}
{गुरूंश्चाप्यवजानातु यस्यार्योऽनुमते गतः} %||2-69-22||

\twolineshloka
{पुत्रैर्दारैश्च भृत्यैश्च स्वगृहे परिवारितः}
{स एको मृष्टमश्नातु यस्यार्योऽनुमते गतः} %||2-69-23||

\twolineshloka
{राजस्त्रीबालवृद्धानां वधे यत्पापमुच्यते}
{भृत्यत्यागे च यत्पापं तत्पापं प्रतिपद्यताम्} %||2-69-24||

\twolineshloka
{उभे सन्ध्ये शयानस्य यत्पापं परिकल्प्यते}
{तच्च पापं भवेत्तस्य यस्यार्योऽनुमते गतः} %||2-69-25||

\twolineshloka
{यदग्निदायके पापं यत्पापं गुरुतल्पगे}
{मित्रद्रोहे च यत्पापं तत्पापं प्रतिपद्यताम्} %||2-69-26||

\twolineshloka
{देवतानां पितॄणां च माता पित्रोस्तथैव च}
{मा स्म कार्षीत्स शुश्रूषां यस्यार्योऽनुमते गतः} %||2-69-27||

\twolineshloka
{सतां लोकात्सतां कीर्त्याः सज्जुष्टात्कर्मणस्तथा}
{भ्रश्यतु क्षिप्रमद्यैव यस्यार्योऽनुमते गतः} %||2-69-28||

\twolineshloka
{विहीनां पतिपुत्राभ्यां कौसल्यां पार्थिवात्मजः}
{एवमाश्वसयन्नेव दुःखार्तो निपपात ह} %||2-69-29||

\twolineshloka
{तथा तु शपथैः कष्टैः शपमानमचेतनम्}
{भरतं शोकसन्तप्तं कौसल्या वाक्यमब्रवीत्} %||2-69-30||

\twolineshloka
{मम दुःखमिदं पुत्र भूयः समुपजायते}
{शपथैः शपमानो हि प्राणानुपरुणत्सि मे} %||2-69-31||

\twolineshloka
{दिष्ट्या न चलितो धर्मादात्मा ते सहलक्ष्मणः}
{वत्स सत्यप्रतिज्ञो मे सतां लोकानवाप्स्यसि} %||2-69-32||

\twolineshloka
{एवं विलपमानस्य दुःखार्तस्य महात्मनः}
{मोहाच्च शोकसंरोधाद्बभूव लुलितं मनः} %||2-69-33||

\fourlineindentedshloka
{लालप्यमानस्य विचेतनस्य}
{ प्रनष्टबुद्धेः पतितस्य भूमौ}
{मुहुर्मुहुर्निःश्वसतश्च दीर्घम्}
{ सा तस्य शोकेन जगाम रात्रिः} %||2-70-34||


{॥इत्यार्षे श्रीमद्रामायणे वाल्मीकीये आदिकाव्ये अयोध्याकाण्डे एकोनसप्ततितमः सर्गः॥}

\sect{सप्ततितमः सर्गः}

\twolineshloka
{तमेवं शोकसन्तप्तं भरतं केकयीसुतम्}
{उवाच वदतां श्रेष्ठो वसिष्ठः श्रेष्ठवागृषिः} %||2-70-1||

\twolineshloka
{अलं शोकेन भद्रं ते राजपुत्र महायशः}
{प्राप्तकालं नरपतेः कुरु संयानमुत्तरम्} %||2-70-2||

\twolineshloka
{वसिष्ठस्य वचः श्रुत्वा भरतो धारणां गतः}
{प्रेतकार्याणि सर्वाणि कारयामास धर्मवित्} %||2-70-3||

\twolineshloka
{उद्धृतं तैलसङ्क्लेदात्स तु भूमौ निवेशितम्}
{आपीतवर्णवदनं प्रसुप्तमिव भूमिपम्} %||2-70-4||

\twolineshloka
{निवेश्य शयने चाग्र्ये नानारत्नपरिष्कृते}
{ततो दशरथं पुत्रो विललाप सुदुःखितः} %||2-70-5||

\twolineshloka
{किं ते व्यवसितं राजन्प्रोषिते मय्यनागते}
{विवास्य रामं धर्मज्ञं लक्ष्मणं च महाबलम्} %||2-70-6||

\twolineshloka
{क्व यास्यसि महाराज हित्वेमं दुःखितं जनम्}
{हीनं पुरुषसिंहेन रामेणाक्लिष्टकर्मणा} %||2-70-7||

\twolineshloka
{योगक्षेमं तु ते राजन्कोऽस्मिन्कल्पयिता पुरे}
{त्वयि प्रयाते स्वस्तात रामे च वनमाश्रिते} %||2-70-8||

\twolineshloka
{विधवा पृथिवी राजंस्त्वया हीना न राजते}
{हीनचन्द्रेव रजनी नगरी प्रतिभाति माम्} %||2-70-9||

\twolineshloka
{एवं विलपमानं तं भरतं दीनमानसम्}
{अब्रवीद्वचनं भूयो वसिष्ठस्तु महानृषिः} %||2-70-10||

\twolineshloka
{प्रेतकार्याणि यान्यस्य कर्तव्यानि विशाम्पतेः}
{तान्यव्यग्रं महाबाहो क्रियतामविचारितम्} %||2-70-11||

\twolineshloka
{तथेति भरतो वाक्यं वसिष्ठस्याभिपूज्य तत्}
{ऋत्विक्पुरोहिताचार्यांस्त्वरयामास सर्वशः} %||2-70-12||

\twolineshloka
{ये त्वग्रतो नरेन्द्रस्य अग्न्यगाराद्बहिष्कृताः}
{ऋत्विग्भिर्याजकैश्चैव ते ह्रियन्ते यथाविधि} %||2-70-13||

\twolineshloka
{शिबिलायामथारोप्य राजानं गतचेतनम्}
{बाष्पकण्ठा विमनसस्तमूहुः परिचारकाः} %||2-70-14||

\twolineshloka
{हिरण्यं च सुवर्णं च वासांसि विविधानि च}
{प्रकिरन्तो जना मार्गं नृपतेरग्रतो ययुः} %||2-70-15||

\twolineshloka
{चन्दनागुरुनिर्यासान्सरलं पद्मकं तथा}
{देवदारूणि चाहृत्य चितां चक्रुस्तथापरे} %||2-70-16||

\twolineshloka
{गन्धानुच्चावचांश्चान्यांस्तत्र दत्त्वाथ भूमिपम्}
{ततः संवेशयामासुश्चितामध्ये तमृत्विजः} %||2-70-17||

\twolineshloka
{तथा हुताशनं हुत्वा जेपुस्तस्य तदर्त्विजः}
{जगुश्च ते यथाशास्त्रं तत्र सामानि सामगाः} %||2-70-18||

\twolineshloka
{शिबिकाभिश्च यानैश्च यथार्हं तस्य योषितः}
{नगरान्निर्ययुस्तत्र वृद्धैः परिवृतास्तदा} %||2-70-19||

\twolineshloka
{प्रसव्यं चापि तं चक्रुरृत्विजोऽग्निचितं नृपम्}
{स्त्रियश्च शोकसन्तप्ताः कौसल्या प्रमुखास्तदा} %||2-70-20||

\twolineshloka
{क्रौञ्चीनामिव नारीणां निनादस्तत्र शुश्रुवे}
{आर्तानां करुणं काले क्रोशन्तीनां सहस्रशः} %||2-70-21||

\twolineshloka
{ततो रुदन्त्यो विवशा विलप्य च पुनः पुनः}
{यानेभ्यः सरयूतीरमवतेरुर्वराङ्गनाः} %||2-70-22||

\fourlineindentedshloka
{कृतोदकं ते भरतेन सार्धम्}
{ नृपाङ्गना मन्त्रिपुरोहिताश्च}
{पुरं प्रविश्याश्रुपरीतनेत्रा}
{ भूमौ दशाहं व्यनयन्त दुःखम्} %||2-71-23||


{॥इत्यार्षे श्रीमद्रामायणे वाल्मीकीये आदिकाव्ये अयोध्याकाण्डे सप्ततितमः सर्गः॥}

\sect{एकसप्ततितमः सर्गः}

\twolineshloka
{ततो दशाहेऽतिगते कृतशौचो नृपात्मजः}
{द्वादशेऽहनि सम्प्राप्ते श्राद्धकर्माण्यकारयत्} %||2-71-1||

\twolineshloka
{ब्राह्मणेभ्यो ददौ रत्नं धनमन्नं च पुष्कलम्}
{बास्तिकं बहुशुक्लं च गाश्चापि शतशस्तथा} %||2-71-2||

\twolineshloka
{दासीदासं च यानं च वेश्मानि सुमहान्ति च}
{ब्राह्मणेभ्यो ददौ पुत्रो राज्ञस्तस्यौर्ध्वदैहिकम्} %||2-71-3||

\twolineshloka
{ततः प्रभातसमये दिवसेऽथ त्रयोदशे}
{विललाप महाबाहुर्भरतः शोकमूर्छितः} %||2-71-4||

\twolineshloka
{शब्दापिहितकण्ठश्च शोधनार्थमुपागतः}
{चितामूले पितुर्वाक्यमिदमाह सुदुःखितः} %||2-71-5||

\twolineshloka
{तात यस्मिन्निषृष्टोऽहं त्वया भ्रातरि राघवे}
{तस्मिन्वनं प्रव्रजिते शून्ये त्यक्तोऽस्म्यहं त्वया} %||2-71-6||

\twolineshloka
{यथागतिरनाथायाः पुत्रः प्रव्राजितो वनम्}
{तामम्बां तात कौसल्यां त्यक्त्वा त्वं क्व गतो नृप} %||2-71-7||

\twolineshloka
{दृष्ट्वा भस्मारुणं तच्च दग्धास्थिस्थानमण्डलम्}
{पितुः शरीर निर्वाणं निष्टनन्विषसाद ह} %||2-71-8||

\twolineshloka
{स तु दृष्ट्वा रुदन्दीनः पपात धरणीतले}
{उत्थाप्यमानः शक्रस्य यन्त्र ध्वज इव च्युतः} %||2-71-9||

\twolineshloka
{अभिपेतुस्ततः सर्वे तस्यामात्याः शुचिव्रतम्}
{अन्तकाले निपतितं ययातिमृषयो यथा} %||2-71-10||

\twolineshloka
{शत्रुघ्नश्चापि भरतं दृष्ट्वा शोकपरिप्लुतम्}
{विसंज्ञो न्यपतद्भूमौ भूमिपालमनुस्मरन्} %||2-71-11||

\twolineshloka
{उन्मत्त इव निश्चेता विललाप सुदुःखितः}
{स्मृत्वा पितुर्गुणाङ्गानि तानि तानि तदा तदा} %||2-71-12||

\twolineshloka
{मन्थरा प्रभवस्तीव्रः कैकेयीग्राहसङ्कुलः}
{वरदानमयोऽक्षोभ्योऽमज्जयच्छोकसागरः} %||2-71-13||

\twolineshloka
{सुकुमारं च बालं च सततं लालितं त्वया}
{क्व तात भरतं हित्वा विलपन्तं गतो भवान्} %||2-71-14||

\twolineshloka
{ननु भोज्येषु पानेषु वस्त्रेष्वाभरणेषु च}
{प्रवारयसि नः सर्वांस्तन्नः कोऽद्य करिष्यति} %||2-71-15||

\twolineshloka
{अवदारण काले तु पृथिवी नावदीर्यते}
{विहीना या त्वया राज्ञा धर्मज्ञेन महात्मना} %||2-71-16||

\twolineshloka
{पितरि स्वर्गमापन्ने रामे चारण्यमाश्रिते}
{किं मे जीवित सामर्थ्यं प्रवेक्ष्यामि हुताशनम्} %||2-71-17||

\twolineshloka
{हीनो भ्रात्रा च पित्रा च शून्यामिक्ष्वाकुपालिताम्}
{अयोध्यां न प्रवेक्ष्यामि प्रवेक्ष्यामि तपोवनम्} %||2-71-18||

\twolineshloka
{तयोर्विलपितं श्रुत्वा व्यसनं चान्ववेक्ष्य तत्}
{भृशमार्ततरा भूयः सर्व एवानुगामिनः} %||2-71-19||

\twolineshloka
{ततो विषण्णौ श्रान्तौ च शत्रुघ्न भरतावुभौ}
{धरण्यां संव्यचेष्टेतां भग्नशृङ्गाविवर्षभौ} %||2-71-20||

\twolineshloka
{ततः प्रकृतिमान्वैद्यः पितुरेषां पुरोहितः}
{वसिष्ठो भरतं वाक्यमुत्थाप्य तमुवाच ह} %||2-71-21||

\twolineshloka
{त्रीणि द्वन्द्वानि भूतेषु प्रवृत्तान्यविशेषतः}
{तेषु चापरिहार्येषु नैवं भवितुमर्हति} %||2-71-22||

\twolineshloka
{सुमन्त्रश्चापि शत्रुघ्नमुत्थाप्याभिप्रसाद्य च}
{श्रावयामास तत्त्वज्ञः सर्वभूतभवाभवौ} %||2-71-23||

\twolineshloka
{उत्थितौ तौ नरव्याघ्रौ प्रकाशेते यशस्विनौ}
{वर्षातपपरिक्लिन्नौ पृथगिन्द्रध्वजाविव} %||2-71-24||

\twolineshloka
{अश्रूणि परिमृद्नन्तौ रक्ताक्षौ दीनभाषिणौ}
{अमात्यास्त्वरयन्ति स्म तनयौ चापराः क्रियाः} %||2-72-25||


{॥इत्यार्षे श्रीमद्रामायणे वाल्मीकीये आदिकाव्ये अयोध्याकाण्डे एकसप्ततितमः सर्गः॥}

\sect{द्विसप्ततितमः सर्गः}

\twolineshloka
{अत्र यात्रां समीहन्तं शत्रुघ्नो लक्ष्मणानुजः}
{भरतं शोकसन्तप्तमिदं वचनमब्रवीत्} %||2-72-1||

\twolineshloka
{गतिर्यः सर्वभूतानां दुःखे किं पुनरात्मनः}
{स रामः सत्त्व सम्पन्नः स्त्रिया प्रव्राजितो वनम्} %||2-72-2||

\twolineshloka
{बलवान्वीर्य सम्पन्नो लक्ष्मणो नाम योऽप्यसौ}
{किं न मोचयते रामं कृत्वापि पितृनिग्रहम्} %||2-72-3||

\twolineshloka
{पूर्वमेव तु निग्राह्यः समवेक्ष्य नयानयौ}
{उत्पथं यः समारूढो नार्या राजा वशं गतः} %||2-72-4||

\twolineshloka
{इति सम्भाषमाणे तु शत्रुघ्ने लक्ष्मणानुजे}
{प्राग्द्वारेऽभूत्तदा कुब्जा सर्वाभरणभूषिता} %||2-72-5||

\twolineshloka
{लिप्ता चन्दनसारेण राजवस्त्राणि बिभ्रती}
{मेखला दामभिश्चित्रै रज्जुबद्धेव वानरी} %||2-72-6||

\twolineshloka
{तां समीक्ष्य तदा द्वाःस्थो भृशं पापस्य कारिणीम्}
{गृहीत्वाकरुणं कुब्जां शत्रुघ्नाय न्यवेदयत्} %||2-72-7||

\twolineshloka
{यस्याः कृते वने रामो न्यस्तदेहश्च वः पिता}
{सेयं पापा नृशंसा च तस्याः कुरु यथामति} %||2-72-8||

\twolineshloka
{शत्रुघ्नश्च तदाज्ञाय वचनं भृशदुःखितः}
{अन्तःपुरचरान्सर्वानित्युवाच धृतव्रतः} %||2-72-9||

\twolineshloka
{तीव्रमुत्पादितं दुःखं भ्रातॄणां मे तथा पितुः}
{यया सेयं नृशंसस्य कर्मणः फलमश्नुताम्} %||2-72-10||

\twolineshloka
{एवमुक्ता च तेनाशु सखी जनसमावृता}
{गृहीता बलवत्कुब्जा सा तद्गृहमनादयत्} %||2-72-11||

\twolineshloka
{ततः सुभृश सन्तप्तस्तस्याः सर्वः सखीजनः}
{क्रुद्धमाज्ञाय शत्रुघ्नं व्यपलायत सर्वशः} %||2-72-12||

\twolineshloka
{अमन्त्रयत कृत्स्नश्च तस्याः सर्वसखीजनः}
{यथायं समुपक्रान्तो निःशेषं नः करिष्यति} %||2-72-13||

\twolineshloka
{सानुक्रोशां वदान्यां च धर्मज्ञां च यशस्विनीम्}
{कौसल्यां शरणं यामः सा हि नोऽस्तु ध्रुवा गतिः} %||2-72-14||

\twolineshloka
{स च रोषेण ताम्राक्षः शत्रुघ्नः शत्रुतापनः}
{विचकर्ष तदा कुब्जां क्रोशन्तीं पृथिवीतले} %||2-72-15||

\twolineshloka
{तस्या ह्याकृष्यमाणाया मन्थरायास्ततस्ततः}
{चित्रं बहुविधं भाण्डं पृथिव्यां तद्व्यशीर्यत} %||2-72-16||

\twolineshloka
{तेन भाण्डेन सङ्कीर्णं श्रीमद्राजनिवेशनम्}
{अशोभत तदा भूयः शारदं गगनं यथा} %||2-72-17||

\twolineshloka
{स बली बलवत्क्रोधाद्गृहीत्वा पुरुषर्षभः}
{कैकेयीमभिनिर्भर्त्स्य बभाषे परुषं वचः} %||2-72-18||

\twolineshloka
{तैर्वाक्यैः परुषैर्दुःखैः कैकेयी भृशदुःखिता}
{शत्रुघ्न भयसन्त्रस्ता पुत्रं शरणमागता} %||2-72-19||

\twolineshloka
{तां प्रेक्ष्य भरतः क्रुद्धं शत्रुघ्नमिदमब्रवीत्}
{अवध्याः सर्वभूतानां प्रमदाः क्षम्यतामिति} %||2-72-20||

\twolineshloka
{हन्यामहमिमां पापां कैकेयीं दुष्टचारिणीम्}
{यदि मां धार्मिको रामो नासूयेन्मातृघातकम्} %||2-72-21||

\twolineshloka
{इमामपि हतां कुब्जां यदि जानाति राघवः}
{त्वां च मां चैव धर्मात्मा नाभिभाषिष्यते ध्रुवम्} %||2-72-22||

\twolineshloka
{भरतस्य वचः श्रुत्वा शत्रुघ्नो लक्ष्मणानुजः}
{न्यवर्तत ततो रोषात्तां मुमोच च मन्थराम्} %||2-72-23||

\twolineshloka
{सा पादमूले कैकेय्या मन्थरा निपपात ह}
{निःश्वसन्ती सुदुःखार्ता कृपणं विललाप च} %||2-72-24||

\fourlineindentedshloka
{शत्रुघ्नविक्षेपविमूढसंज्ञाम्}
{ समीक्ष्य कुब्जां भरतस्य माता}
{शनैः समाश्वासयदार्तरूपाम्}
{ क्रौञ्चीं विलग्नामिव वीक्षमाणाम्} %||2-73-25||


{॥इत्यार्षे श्रीमद्रामायणे वाल्मीकीये आदिकाव्ये अयोध्याकाण्डे द्विसप्ततितमः सर्गः॥}

\sect{त्रिसप्ततितमः सर्गः}

\twolineshloka
{ततः प्रभातसमये दिवसेऽथ चतुर्दशे}
{समेत्य राजकर्तारो भरतं वाक्यमब्रुवन्} %||2-73-1||

\twolineshloka
{गतो दशरथः स्वर्गं यो नो गुरुतरो गुरुः}
{रामं प्रव्राज्य वै ज्येष्ठं लक्ष्मणं च महाबलम्} %||2-73-2||

\twolineshloka
{त्वमद्य भव नो राजा राजपुत्र महायशः}
{सङ्गत्या नापराध्नोति राज्यमेतदनायकम्} %||2-73-3||

\twolineshloka
{आभिषेचनिकं सर्वमिदमादाय राघव}
{प्रतीक्षते त्वां स्वजनः श्रेणयश्च नृपात्मज} %||2-73-4||

\twolineshloka
{राज्यं गृहाण भरत पितृपैतामहं महत्}
{अभिषेचय चात्मानं पाहि चास्मान्नरर्षभ} %||2-73-5||

\twolineshloka
{आभिषेचनिकं भाण्डं कृत्वा सर्वं प्रदक्षिणम्}
{भरतस्तं जनं सर्वं प्रत्युवाच धृतव्रतः} %||2-73-6||

\twolineshloka
{ज्येष्ठस्य राजता नित्यमुचिता हि कुलस्य नः}
{नैवं भवन्तो मां वक्तुमर्हन्ति कुशला जनाः} %||2-73-7||

\twolineshloka
{रामः पूर्वो हि नो भ्राता भविष्यति महीपतिः}
{अहं त्वरण्ये वत्स्यामि वर्षाणि नव पञ्च च} %||2-73-8||

\twolineshloka
{युज्यतां महती सेना चतुरङ्गमहाबला}
{आनयिष्याम्यहं ज्येष्ठं भ्रातरं राघवं वनात्} %||2-73-9||

\twolineshloka
{आभिषेचनिकं चैव सर्वमेतदुपस्कृतम्}
{पुरस्कृत्य गमिष्यामि रामहेतोर्वनं प्रति} %||2-73-10||

\twolineshloka
{तत्रैव तं नरव्याघ्रमभिषिच्य पुरस्कृतम्}
{आनेष्यामि तु वै रामं हव्यवाहमिवाध्वरात्} %||2-73-11||

\twolineshloka
{न सकामा करिष्यामि स्वमिमां मातृगन्धिनीम्}
{वने वत्स्याम्यहं दुर्गे रामो राजा भविष्यति} %||2-73-12||

\twolineshloka
{क्रियतां शिल्पिभिः पन्थाः समानि विषमाणि च}
{रक्षिणश्चानुसंयान्तु पथि दुर्गविचारकाः} %||2-73-13||

\twolineshloka
{एवं सम्भाषमाणं तं रामहेतोर्नृपात्मजम्}
{प्रत्युवाच जनः सर्वः श्रीमद्वाक्यमनुत्तमम्} %||2-73-14||

\twolineshloka
{एवं ते भाषमाणस्य पद्मा श्रीरुपतिष्ठताम्}
{यस्त्वं ज्येष्ठे नृपसुते पृथिवीं दातुमिच्छसि} %||2-73-15||

\fourlineindentedshloka
{अनुत्तमं तद्वचनं नृपात्मज}
{ प्रभाषितं संश्रवणे निशम्य च}
{प्रहर्षजास्तं प्रति बाष्पबिन्दवो}
{ निपेतुरार्यानननेत्रसम्भवाः} %||2-73-16||

\fourlineindentedshloka
{ऊचुस्ते वचनमिदं निशम्य हृष्टाः}
{ सामात्याः सपरिषदो वियातशोकाः}
{पन्थानं नरवरभक्तिमाञ्जनश्च}
{ व्यादिष्टस्तव वचनाच्च शिल्पिवर्गः} %||2-74-17||


{॥इत्यार्षे श्रीमद्रामायणे वाल्मीकीये आदिकाव्ये अयोध्याकाण्डे त्रिसप्ततितमः सर्गः॥}

\sect{चतुःसप्ततितमः सर्गः}

\twolineshloka
{अथ भूमिप्रदेशज्ञाः सूत्रकर्मविशारदाः}
{स्वकर्माभिरताः शूराः खनका यन्त्रकास्तथा} %||2-74-1||

\twolineshloka
{कर्मान्तिकाः स्थपतयः पुरुषा यन्त्रकोविदाः}
{तथा वर्धकयश्चैव मार्गिणो वृक्षतक्षकाः} %||2-74-2||

\twolineshloka
{कूपकाराः सुधाकारा वंशकर्मकृतस्तथा}
{समर्था ये च द्रष्टारः पुरतस्ते प्रतस्थिरे} %||2-74-3||

\twolineshloka
{स तु हर्षात्तमुद्देशं जनौघो विपुलः प्रयान्}
{अशोभत महावेगः सागरस्येव पर्वणि} %||2-74-4||

\twolineshloka
{ते स्ववारं समास्थाय वर्त्मकर्माणि कोविदाः}
{करणैर्विविधोपेतैः पुरस्तात्सम्प्रतस्थिरे} %||2-74-5||

\twolineshloka
{लतावल्लीश्च गुल्मांश्च स्थाणूनश्मन एव च}
{जनास्ते चक्रिरे मार्गं छिन्दन्तो विविधान्द्रुमान्} %||2-74-6||

\twolineshloka
{अवृक्षेषु च देशेषु केचिद्वृक्षानरोपयन्}
{केचित्कुठारैष्टङ्कैश्च दात्रैश्छिन्दन्क्वचित्क्वचित्} %||2-74-7||

\twolineshloka
{अपरे वीरणस्तम्बान्बलिनो बलवत्तराः}
{विधमन्ति स्म दुर्गाणि स्थलानि च ततस्ततः} %||2-74-8||

\twolineshloka
{अपरेऽपूरयन्कूपान्पांसुभिः श्वभ्रमायतम्}
{निम्नभागांस्तथा केचित्समांश्चक्रुः समन्ततः} %||2-74-9||

\twolineshloka
{बबन्धुर्बन्धनीयांश्च क्षोद्यान्सञ्चुक्षुदुस्तदा}
{बिभिदुर्भेदनीयांश्च तांस्तान्देशान्नरास्तदा} %||2-74-10||

\threelineshloka
{अचिरेणैव कालेन परिवाहान्बहूदकान्}
{चक्रुर्बहुविधाकारान्सागरप्रतिमान्बहून्}
{उदपानान्बहुविधान्वेदिका परिमण्डितान्} %||2-74-11||

\twolineshloka
{ससुधाकुट्टिमतलः प्रपुष्पितमहीरुहः}
{मत्तोद्घुष्टद्विजगणः पताकाभिरलङ्कृतः} %||2-74-12||

\twolineshloka
{चन्दनोदकसंसिक्तो नानाकुसुमभूषितः}
{बह्वशोभत सेनायाः पन्थाः स्वर्गपथोपमः} %||2-74-13||

\twolineshloka
{आज्ञाप्याथ यथाज्ञप्ति युक्तास्तेऽधिकृता नराः}
{रमणीयेषु देशेषु बहुस्वादुफलेषु च} %||2-74-14||

\twolineshloka
{यो निवेशस्त्वभिप्रेतो भरतस्य महात्मनः}
{भूयस्तं शोभयामासुर्भूषाभिर्भूषणोपमम्} %||2-74-15||

\twolineshloka
{नक्षत्रेषु प्रशस्तेषु मुहूर्तेषु च तद्विदः}
{निवेशं स्थापयामासुर्भरतस्य महात्मनः} %||2-74-16||

\twolineshloka
{बहुपांसुचयाश्चापि परिखापरिवारिताः}
{तत्रेन्द्रकीलप्रतिमाः प्रतोलीवरशोभिताः} %||2-74-17||

\twolineshloka
{प्रासादमालासंयुक्ताः सौधप्राकारसंवृताः}
{पताका शोभिताः सर्वे सुनिर्मितमहापथाः} %||2-74-18||

\twolineshloka
{विसर्पत्भिरिवाकाशे विटङ्काग्रविमानकैः}
{समुच्छ्रितैर्निवेशास्ते बभुः शक्रपुरोपमाः} %||2-74-19||

\twolineshloka
{जाह्नवीं तु समासाद्य विविधद्रुम काननाम्}
{शीतलामलपानीयां महामीनसमाकुलाम्} %||2-74-20||

\fourlineindentedshloka
{सचन्द्रतारागणमण्डितं यथा}
{ नभःक्षपायाममलं विराजते}
{नरेन्द्रमार्गः स तथा व्यराजत}
{ क्रमेण रम्यः शुभशिल्पिनिर्मितः} %||2-75-21||


{॥इत्यार्षे श्रीमद्रामायणे वाल्मीकीये आदिकाव्ये अयोध्याकाण्डे चतुःसप्ततितमः सर्गः॥}

\sect{पञ्चसप्ततितमः सर्गः}

\twolineshloka
{ततो नान्दीमुखीं रात्रिं भरतं सूतमागधाः}
{तुष्टुवुर्वाग्विशेषज्ञाः स्तवैर्मङ्गलसंहितैः} %||2-75-1||

\twolineshloka
{सुवर्णकोणाभिहतः प्राणदद्यामदुन्दुभिः}
{दध्मुः शङ्खांश्च शतशो वाद्यांश्चोच्चावचस्वरान्} %||2-75-2||

\twolineshloka
{स तूर्य घोषः सुमहान्दिवमापूरयन्निव}
{भरतं शोकसन्तप्तं भूयः शोकैररन्ध्रयत्} %||2-75-3||

\twolineshloka
{ततो प्रबुद्धो भरतस्तं घोषं संनिवर्त्य च}
{नाहं राजेति चाप्युक्त्वा शत्रुघ्नमिदमब्रवीत्} %||2-75-4||

\twolineshloka
{पश्य शत्रुघ्न कैकेय्या लोकस्यापकृतं महत्}
{विसृज्य मयि दुःखानि राजा दशरथो गतः} %||2-75-5||

\twolineshloka
{तस्यैषा धर्मराजस्य धर्ममूला महात्मनः}
{परिभ्रमति राजश्रीर्नौरिवाकर्णिका जले} %||2-75-6||

\twolineshloka
{इत्येवं भरतं प्रेक्ष्य विलपन्तं विचेतनम्}
{कृपणं रुरुदुः सर्वाः सस्वरं योषितस्तदा} %||2-75-7||

\twolineshloka
{तथा तस्मिन्विलपति वसिष्ठो राजधर्मवित्}
{सभामिक्ष्वाकुनाथस्य प्रविवेश महायशाः} %||2-75-8||

\twolineshloka
{शात कुम्भमयीं रम्यां मणिरत्नसमाकुलाम्}
{सुधर्मामिव धर्मात्मा सगणः प्रत्यपद्यत} %||2-75-9||

\twolineshloka
{स काञ्चनमयं पीठं परार्ध्यास्तरणावृतम्}
{अध्यास्त सर्ववेदज्ञो दूताननुशशास च} %||2-75-10||

\twolineshloka
{ब्राह्मणान्क्षत्रियान्योधानमात्यान्गणबल्लभान्}
{क्षिप्रमानयताव्यग्राः कृत्यमात्ययिकं हि नः} %||2-75-11||

\twolineshloka
{ततो हलहलाशब्दो महान्समुदपद्यत}
{रथैरश्वैर्गजैश्चापि जनानामुपगच्छताम्} %||2-75-12||

\twolineshloka
{ततो भरतमायान्तं शतक्रतुमिवामराः}
{प्रत्यनन्दन्प्रकृतयो यथा दशरथं तथा} %||2-75-13||

\fourlineindentedshloka
{ह्रद इव तिमिनागसंवृतः}
{ स्तिमितजलो मणिशङ्खशर्करः}
{दशरथसुतशोभिता सभा}
{ सदशरथेव बभौ यथा पुरा} %||2-76-14||


{॥इत्यार्षे श्रीमद्रामायणे वाल्मीकीये आदिकाव्ये अयोध्याकाण्डे पञ्चसप्ततितमः सर्गः॥}

\sect{षट्सप्ततितमः सर्गः}

\twolineshloka
{तामार्यगणसम्पूर्णां भरतः प्रग्रहां सभाम्}
{ददर्श बुद्धिसम्पन्नः पूर्णचन्द्रां निशामिव} %||2-76-1||

\twolineshloka
{आसनानि यथान्यायमार्याणां विशतां तदा}
{अदृश्यत घनापाये पूर्णचन्द्रेव शर्वरी} %||2-76-2||

\twolineshloka
{राज्ञस्तु प्रकृतीः सर्वाः समग्राः प्रेक्ष्य धर्मवित्}
{इदं पुरोहितो वाक्यं भरतं मृदु चाब्रवीत्} %||2-76-3||

\twolineshloka
{तात राजा दशरथः स्वर्गतो धर्ममाचरन्}
{धन धान्यवतीं स्फीतां प्रदाय पृथिवीं तव} %||2-76-4||

\twolineshloka
{रामस्तथा सत्यधृतिः सतां धर्ममनुस्मरन्}
{नाजहात्पितुरादेशं शशी ज्योत्स्नामिवोदितः} %||2-76-5||

\twolineshloka
{पित्रा भ्रात्रा च ते दत्तं राज्यं निहतकण्टकम्}
{तद्भुङ्क्ष्व मुदितामात्यः क्षिप्रमेवाभिषेचय} %||2-76-6||

\twolineshloka
{उदीच्याश्च प्रतीच्याश्च दाक्षिणात्याश्च केवलाः}
{कोट्यापरान्ताः सामुद्रा रत्नान्यभिहरन्तु ते} %||2-76-7||

\twolineshloka
{तच्छ्रुत्वा भरतो वाक्यं शोकेनाभिपरिप्लुतः}
{जगाम मनसा रामं धर्मज्ञो धर्मकाङ्क्षया} %||2-76-8||

\twolineshloka
{स बाष्पकलया वाचा कलहंसस्वरो युवा}
{विललाप सभामध्ये जगर्हे च पुरोहितम्} %||2-76-9||

\twolineshloka
{चरितब्रह्मचर्यस्य विद्या स्नातस्य धीमतः}
{धर्मे प्रयतमानस्य को राज्यं मद्विधो हरेत्} %||2-76-10||

\twolineshloka
{कथं दशरथाज्जातो भवेद्राज्यापहारकः}
{राज्यं चाहं च रामस्य धर्मं वक्तुमिहार्हसि} %||2-76-11||

\twolineshloka
{ज्येष्ठः श्रेष्ठश्च धर्मात्मा दिलीपनहुषोपमः}
{लब्धुमर्हति काकुत्स्थो राज्यं दशरथो यथा} %||2-76-12||

\twolineshloka
{अनार्यजुष्टमस्वर्ग्यं कुर्यां पापमहं यदि}
{इक्ष्वाकूणामहं लोके भवेयं कुलपांसनः} %||2-76-13||

\twolineshloka
{यद्धि मात्रा कृतं पापं नाहं तदभिरोचये}
{इहस्थो वनदुर्गस्थं नमस्यामि कृताञ्जलिः} %||2-76-14||

\twolineshloka
{राममेवानुगच्छामि स राजा द्विपदां वरः}
{त्रयाणामपि लोकानां राघवो राज्यमर्हति} %||2-76-15||

\twolineshloka
{तद्वाक्यं धर्मसंयुक्तं श्रुत्वा सर्वे सभासदः}
{हर्षान्मुमुचुरश्रूणि रामे निहितचेतसः} %||2-76-16||

\twolineshloka
{यदि त्वार्यं न शक्ष्यामि विनिवर्तयितुं वनात्}
{वने तत्रैव वत्स्यामि यथार्यो लक्ष्मणस्तथा} %||2-76-17||

\twolineshloka
{सर्वोपायं तु वर्तिष्ये विनिवर्तयितुं बलात्}
{समक्षमार्य मिश्राणां साधूनां गुणवर्तिनाम्} %||2-76-18||

\twolineshloka
{एवमुक्त्वा तु धर्मात्मा भरतो भ्रातृवत्सलः}
{समीपस्थमुवाचेदं सुमन्त्रं मन्त्रकोविदम्} %||2-76-19||

\twolineshloka
{तूर्णमुत्थाय गच्छ त्वं सुमन्त्र मम शासनात्}
{यात्रामाज्ञापय क्षिप्रं बलं चैव समानय} %||2-76-20||

\twolineshloka
{एवमुक्तः सुमन्त्रस्तु भरतेन महात्मना}
{प्रहृष्टः सोऽदिशत्सर्वं यथा सन्दिष्टमिष्टवत्} %||2-76-21||

\twolineshloka
{ताः प्रहृष्टाः प्रकृतयो बलाध्यक्षा बलस्य च}
{श्रुत्वा यात्रां समाज्ञप्तां राघवस्य निवर्तने} %||2-76-22||

\twolineshloka
{ततो योधाङ्गनाः सर्वा भर्तॄन्सर्वान्गृहेगृहे}
{यात्रा गमनमाज्ञाय त्वरयन्ति स्म हर्षिताः} %||2-76-23||

\twolineshloka
{ते हयैर्गोरथैः शीघ्रैः स्यन्दनैश्च मनोजवैः}
{सह योधैर्बलाध्यक्षा बलं सर्वमचोदयन्} %||2-76-24||

\twolineshloka
{सज्जं तु तद्बलं दृष्ट्वा भरतो गुरुसंनिधौ}
{रथं मे त्वरयस्वेति सुमन्त्रं पार्श्वतोऽब्रवीत्} %||2-76-25||

\twolineshloka
{भरतस्य तु तस्याज्ञां प्रतिगृह्य प्रहर्षितः}
{रथं गृहीत्वा प्रययौ युक्तं परमवाजिभिः} %||2-76-26||

\fourlineindentedshloka
{स राघवः सत्यधृतिः प्रतापवा}
{न्ब्रुवन्सुयुक्तं दृढसत्यविक्रमः}
{गुरुं महारण्यगतं यशस्विनम्}
{ प्रसादयिष्यन्भरतोऽब्रवीत्तदा} %||2-76-27||

\fourlineindentedshloka
{तूण समुत्थाय सुमन्त्र गच्छ}
{ बलस्य योगाय बलप्रधानान्}
{आनेतुमिच्छामि हि तं वनस्थम्}
{ प्रसाद्य रामं जगतो हिताय} %||2-76-28||

\fourlineindentedshloka
{स सूतपुत्रो भरतेन सम्य}
{गाज्ञापितः सम्परिपूर्णकामः}
{शशास सर्वान्प्रकृतिप्रधाना}
{न्बलस्य मुख्यांश्च सुहृज्जनं च} %||2-76-29||

\fourlineindentedshloka
{ततः समुत्थाय कुले कुले ते}
{ राजन्यवैश्या वृषलाश्च विप्राः}
{अयूयुजन्नुष्ट्ररथान्खरांश्च}
{ नागान्हयांश्चैव कुलप्रसूतान्} %||2-77-30||


{॥इत्यार्षे श्रीमद्रामायणे वाल्मीकीये आदिकाव्ये अयोध्याकाण्डे षट्सप्ततितमः सर्गः॥}

\sect{सप्तसप्ततितमः सर्गः}

\twolineshloka
{ततः समुत्थितः काल्यमास्थाय स्यन्दनोत्तमम्}
{प्रययौ भरतः शीघ्रं रामदर्शनकाङ्क्षया} %||2-77-1||

\twolineshloka
{अग्रतः प्रययुस्तस्य सर्वे मन्त्रिपुरोधसः}
{अधिरुह्य हयैर्युक्तान्रथान्सूर्यरथोपमान्} %||2-77-2||

\twolineshloka
{नवनागसहस्राणि कल्पितानि यथाविधि}
{अन्वयुर्भरतं यान्तमिक्ष्वाकु कुलनन्दनम्} %||2-77-3||

\twolineshloka
{षष्ठी रथसहस्राणि धन्विनो विविधायुधाः}
{अन्वयुर्भरतं यान्तं राजपुत्रं यशस्विनम्} %||2-77-4||

\twolineshloka
{शतं सहस्राण्यश्वानां समारूढानि राघवम्}
{अन्वयुर्भरतं यान्तं राजपुत्रं यशस्विनम्} %||2-77-5||

\twolineshloka
{कैकेयी च सुमित्रा च कौसल्या च यशस्विनी}
{रामानयनसंहृष्टा ययुर्यानेन भास्वता} %||2-77-6||

\twolineshloka
{प्रयाताश्चार्यसङ्घाता रामं द्रष्टुं सलक्ष्मणम्}
{तस्यैव च कथाश्चित्राः कुर्वाणा हृष्टमानसाः} %||2-77-7||

\twolineshloka
{मेघश्यामं महाबाहुं स्थिरसत्त्वं दृढव्रतम्}
{कदा द्रक्ष्यामहे रामं जगतः शोकनाशनम्} %||2-77-8||

\twolineshloka
{दृष्ट एव हि नः शोकमपनेष्यति राघवः}
{तमः सर्वस्य लोकस्य समुद्यन्निव भास्करः} %||2-77-9||

\twolineshloka
{इत्येवं कथयन्तस्ते सम्प्रहृष्टाः कथाः शुभाः}
{परिष्वजानाश्चान्योन्यं ययुर्नागरिकास्तदा} %||2-77-10||

\twolineshloka
{ये च तत्रापरे सर्वे सम्मता ये च नैगमाः}
{रामं प्रति ययुर्हृष्टाः सर्वाः प्रकृतयस्तदा} %||2-77-11||

\twolineshloka
{मणि काराश्च ये केचित्कुम्भकाराश्च शोभनाः}
{सूत्रकर्मकृतश्चैव ये च शस्त्रोपजीविनः} %||2-77-12||

\twolineshloka
{मायूरकाः क्राकचिका रोचका वेधकास्तथा}
{दन्तकाराः सुधाकारास्तथा गन्धोपजीविनः} %||2-77-13||

\twolineshloka
{सुवर्णकाराः प्रख्यातास्तथा कम्बलधावकाः}
{स्नापकाच्छादका वैद्या धूपकाः शौण्डिकास्तथा} %||2-77-14||

\twolineshloka
{रजकास्तुन्नवायाश्च ग्रामघोषमहत्तराः}
{शैलूषाश्च सह स्त्रीभिर्यान्ति कैवर्तकास्तथा} %||2-77-15||

\twolineshloka
{समाहिता वेदविदो ब्राह्मणा वृत्तसम्मताः}
{गोरथैर्भरतं यान्तमनुजग्मुः सहस्रशः} %||2-77-16||

\twolineshloka
{सुवेषाः शुद्धवसनास्ताम्रमृष्टानुलेपनाः}
{सर्वे ते विविधैर्यानैः शनैर्भरतमन्वयुः} %||2-77-17||

\twolineshloka
{प्रहृष्टमुदिता सेना सान्वयात्कैकयीसुतम्}
{व्यवतिष्ठत सा सेना भरतस्यानुयायिनी} %||2-77-18||

\twolineshloka
{निरीक्ष्यानुगतां सेनां तां च गङ्गां शिवोदकाम्}
{भरतः सचिवान्सर्वानब्रवीद्वाक्यकोविदः} %||2-77-19||

\twolineshloka
{निवेशयत मे सैन्यमभिप्रायेण सर्वशः}
{विश्रान्तः प्रतरिष्यामः श्व इदानीं महानदीम्} %||2-77-20||

\twolineshloka
{दातुं च तावदिच्छामि स्वर्गतस्य महीपतेः}
{और्ध्वदेह निमित्तार्थमवतीर्योदकं नदीम्} %||2-77-21||

\twolineshloka
{तस्यैवं ब्रुवतोऽमात्यास्तथेत्युक्त्वा समाहिताः}
{न्यवेशयंस्तांश्छन्देन स्वेन स्वेन पृथक्पृथक्} %||2-77-22||

\fourlineindentedshloka
{निवेश्य गङ्गामनु तां महानदीम्}
{ चमूं विधानैः परिबर्ह शोभिनीम्}
{उवास रामस्य तदा महात्मनो}
{ विचिन्तयानो भरतो निवर्तनम्} %||2-78-23||


{॥इत्यार्षे श्रीमद्रामायणे वाल्मीकीये आदिकाव्ये अयोध्याकाण्डे सप्तसप्ततितमः सर्गः॥}

\sect{अष्टसप्ततितमः सर्गः}

\twolineshloka
{ततो निविष्टां ध्वजिनीं गङ्गामन्वाश्रितां नदीम्}
{निषादराजो दृष्ट्वैव ज्ञातीन्सन्त्वरितोऽब्रवीत्} %||2-78-1||

\twolineshloka
{महतीयमतः सेना सागराभा प्रदृश्यते}
{नास्यान्तमवगच्छामि मनसापि विचिन्तयन्} %||2-78-2||

\twolineshloka
{स एष हि महाकायः कोविदारध्वजो रथे}
{बन्धयिष्यति वा दाशानथ वास्मान्वधिष्यति} %||2-78-3||

\twolineshloka
{अथ दाशरथिं रामं पित्रा राज्याद्विवासितम्}
{भरतः कैकेयीपुत्रो हन्तुं समधिगच्छति} %||2-78-4||

\twolineshloka
{भर्ता चैव सखा चैव रामो दाशरथिर्मम}
{तस्यार्थकामाः संनद्धा गङ्गानूपेऽत्र तिष्ठत} %||2-78-5||

\twolineshloka
{तिष्ठन्तु सर्वदाशाश्च गङ्गामन्वाश्रिता नदीम्}
{बलयुक्ता नदीरक्षा मांसमूलफलाशनाः} %||2-78-6||

\twolineshloka
{नावां शतानां पञ्चानां कैवर्तानां शतं शतम्}
{संनद्धानां तथा यूनां तिष्ठन्त्वत्यभ्यचोदयत्} %||2-78-7||

\twolineshloka
{यदा तुष्टस्तु भरतो रामस्येह भविष्यति}
{सेयं स्वस्तिमयी सेना गङ्गामद्य तरिष्यति} %||2-78-8||

\twolineshloka
{इत्युक्त्वोपायनं गृह्य मत्स्यमांसमधूनि च}
{अभिचक्राम भरतं निषादाधिपतिर्गुहः} %||2-78-9||

\twolineshloka
{तमायान्तं तु सम्प्रेक्ष्य सूतपुत्रः प्रतापवान्}
{भरतायाचचक्षेऽथ विनयज्ञो विनीतवत्} %||2-78-10||

\twolineshloka
{एष ज्ञातिसहस्रेण स्थपतिः परिवारितः}
{कुशलो दण्डकारण्ये वृद्धो भ्रातुश्च ते सखा} %||2-78-11||

\twolineshloka
{तस्मात्पश्यतु काकुत्स्थ त्वां निषादाधिपो गुहः}
{असंशयं विजानीते यत्र तौ रामलक्ष्मणौ} %||2-78-12||

\twolineshloka
{एतत्तु वचनं श्रुत्वा सुमन्त्राद्भरतः शुभम्}
{उवाच वचनं शीघ्रं गुहः पश्यतु मामिति} %||2-78-13||

\twolineshloka
{लब्ध्वाभ्यनुज्ञां संहृष्टो ज्ञातिभिः परिवारितः}
{आगम्य भरतं प्रह्वो गुहो वचनमब्रवीत्} %||2-78-14||

\twolineshloka
{निष्कुटश्चैव देशोऽयं वञ्चिताश्चापि ते वयम्}
{निवेदयामस्ते सर्वे स्वके दाशकुले वस} %||2-78-15||

\twolineshloka
{अस्ति मूलं फलं चैव निषादैः समुपाहृतम्}
{आर्द्रं च मांसं शुष्कं च वन्यं चोच्चावचं महत्} %||2-78-16||

\twolineshloka
{आशंसे स्वाशिता सेना वत्स्यतीमां विभावरीम्}
{अर्चितो विविधैः कामैः श्वः ससैन्यो गमिष्यसि} %||2-79-17||


{॥इत्यार्षे श्रीमद्रामायणे वाल्मीकीये आदिकाव्ये अयोध्याकाण्डे अष्टसप्ततितमः सर्गः॥}

\sect{एकोनाशीतितमः सर्गः}

\twolineshloka
{एवमुक्तस्तु भरतो निषादाधिपतिं गुहम्}
{प्रत्युवाच महाप्राज्ञो वाक्यं हेत्वर्थसंहितम्} %||2-79-1||

\twolineshloka
{ऊर्जितः खलु ते कामः कृतो मम गुरोः सखे}
{यो मे त्वमीदृशीं सेनामेकोऽभ्यर्चितुमिच्छसि} %||2-79-2||

\twolineshloka
{इत्युक्त्वा तु महातेजा गुहं वचनमुत्तमम्}
{अब्रवीद्भरतः श्रीमान्निषादाधिपतिं पुनः} %||2-79-3||

\twolineshloka
{कतरेण गमिष्यामि भरद्वाजाश्रमं गुह}
{गहनोऽयं भृशं देशो गङ्गानूपो दुरत्ययः} %||2-79-4||

\twolineshloka
{तस्य तद्वचनं श्रुत्वा राजपुत्रस्य धीमतः}
{अब्रवीत्प्राञ्जलिर्वाक्यं गुहो गहनगोचरः} %||2-79-5||

\twolineshloka
{दाशास्त्वनुगमिष्यन्ति धन्विनः सुसमाहिताः}
{अहं चानुगमिष्यामि राजपुत्र महायशः} %||2-79-6||

\twolineshloka
{कच्चिन्न दुष्टो व्रजसि रामस्याक्लिष्टकर्मणः}
{इयं ते महती सेना शङ्कां जनयतीव मे} %||2-79-7||

\twolineshloka
{तमेवमभिभाषन्तमाकाश इव निर्मलः}
{भरतः श्लक्ष्णया वाचा गुहं वचनमब्रवीत्} %||2-79-8||

\twolineshloka
{मा भूत्स कालो यत्कष्टं न मां शङ्कितुमर्हसि}
{राघवः स हि मे भ्राता ज्येष्ठः पितृसमो मम} %||2-79-9||

\twolineshloka
{तं निवर्तयितुं यामि काकुत्स्थं वनवासिनम्}
{बुद्धिरन्या न ते कार्या गुह सत्यं ब्रवीमि ते} %||2-79-10||

\twolineshloka
{स तु संहृष्टवदनः श्रुत्वा भरतभाषितम्}
{पुनरेवाब्रवीद्वाक्यं भरतं प्रति हर्षितः} %||2-79-11||

\twolineshloka
{धन्यस्त्वं न त्वया तुल्यं पश्यामि जगतीतले}
{अयत्नादागतं राज्यं यस्त्वं त्यक्तुमिहेच्छसि} %||2-79-12||

\twolineshloka
{शाश्वती खलु ते कीर्तिर्लोकाननुचरिष्यति}
{यस्त्वं कृच्छ्रगतं रामं प्रत्यानयितुमिच्छसि} %||2-79-13||

\twolineshloka
{एवं सम्भाषमाणस्य गुहस्य भरतं तदा}
{बभौ नष्टप्रभः सूर्यो रजनी चाभ्यवर्तत} %||2-79-14||

\twolineshloka
{संनिवेश्य स तां सेनां गुहेन परितोषितः}
{शत्रुघ्नेन सह श्रीमाञ्शयनं पुनरागमत्} %||2-79-15||

\twolineshloka
{रामचिन्तामयः शोको भरतस्य महात्मनः}
{उपस्थितो ह्यनर्हस्य धर्मप्रेक्षस्य तादृशः} %||2-79-16||

\twolineshloka
{अन्तर्दाहेन दहनः सन्तापयति राघवम्}
{वनदाहाभिसन्तप्तं गूढोऽग्निरिव पादपम्} %||2-79-17||

\twolineshloka
{प्रस्रुतः सर्वगात्रेभ्यः स्वेदः शोकाग्निसम्भवः}
{यथा सूर्यांशुसन्तप्तो हिमवान्प्रस्रुतो हिमम्} %||2-79-18||

\twolineshloka
{ध्याननिर्दरशैलेन विनिःश्वसितधातुना}
{दैन्यपादपसङ्घेन शोकायासाधिशृङ्गिणा} %||2-79-19||

\twolineshloka
{प्रमोहानन्तसत्त्वेन सन्तापौषधिवेणुना}
{आक्रान्तो दुःखशैलेन महता कैकयीसुतः} %||2-79-20||

\fourlineindentedshloka
{गुहेन सार्धं भरतः समागतो}
{ महानुभावः सजनः समाहितः}
{सुदुर्मनास्तं भरतं तदा पुन}
{र्गुहः समाश्वासयदग्रजं प्रति} %||2-80-21||


{॥इत्यार्षे श्रीमद्रामायणे वाल्मीकीये आदिकाव्ये अयोध्याकाण्डे एकोनाशीतितमः सर्गः॥}

\sect{अशीतितमः सर्गः}

\twolineshloka
{आचचक्षेऽथ सद्भावं लक्ष्मणस्य महात्मनः}
{भरतायाप्रमेयाय गुहो गहनगोचरः} %||2-80-1||

\twolineshloka
{तं जाग्रतं गुणैर्युक्तं वरचापेषुधारिणम्}
{भ्रातृ गुप्त्यर्थमत्यन्तमहं लक्ष्मणमब्रवम्} %||2-80-2||

\twolineshloka
{इयं तात सुखा शय्या त्वदर्थमुपकल्पिता}
{प्रत्याश्वसिहि शेष्वास्यां सुखं राघवनन्दन} %||2-80-3||

\twolineshloka
{उचितोऽयं जनः सर्वे दुःखानां त्वं सुखोचितः}
{धर्मात्मंस्तस्य गुप्त्यर्थं जागरिष्यामहे वयम्} %||2-80-4||

\twolineshloka
{न हि रामात्प्रियतरो ममास्ति भुवि कश्चन}
{मोत्सुको भूर्ब्रवीम्येतदप्यसत्यं तवाग्रतः} %||2-80-5||

\twolineshloka
{अस्य प्रसादादाशंसे लोकेऽस्मिन्सुमहद्यशः}
{धर्मावाप्तिं च विपुलामर्थावाप्तिं च केवलाम्} %||2-80-6||

\twolineshloka
{सोऽहं प्रियसखं रामं शयानं सह सीतया}
{रक्षिष्यामि धनुष्पाणिः सर्वैः स्वैर्ज्ञातिभिः सह} %||2-80-7||

\twolineshloka
{न हि मेऽविदितं किञ्चिद्वनेऽस्मिंश्चरतः सदा}
{चतुरङ्गं ह्यपि बलं प्रसहेम वयं युधि} %||2-80-8||

\twolineshloka
{एवमस्माभिरुक्तेन लक्ष्मणेन महात्मना}
{अनुनीता वयं सर्वे धर्ममेवानुपश्यता} %||2-80-9||

\twolineshloka
{कथं दाशरथौ भूमौ शयाने सह सीतया}
{शक्या निद्रामया लब्धुं जीवितं वा सुखानि वा} %||2-80-10||

\twolineshloka
{यो न देवासुरैः सर्वैः शक्यः प्रसहितुं युधि}
{तं पश्य गुह संविष्टं तृणेषु सह सीतया} %||2-80-11||

\twolineshloka
{महता तपसा लब्धो विविधैश्च परिश्रमैः}
{एको दशरथस्यैष पुत्रः सदृशलक्षणः} %||2-80-12||

\twolineshloka
{अस्मिन्प्रव्राजिते राजा न चिरं वर्तयिष्यति}
{विधवा मेदिनी नूनं क्षिप्रमेव भविष्यति} %||2-80-13||

\twolineshloka
{विनद्य सुमहानादं श्रमेणोपरताः स्त्रियः}
{निर्घोषोपरतं नूनमद्य राजनिवेशनम्} %||2-80-14||

\twolineshloka
{कौसल्या चैव राजा च तथैव जननी मम}
{नाशंसे यदि ते सर्वे जीवेयुः शर्वरीमिमाम्} %||2-80-15||

\twolineshloka
{जीवेदपि हि मे माता शत्रुघ्नस्यान्ववेक्षया}
{दुःखिता या तु कौसल्या वीरसूर्विनशिष्यति} %||2-80-16||

\twolineshloka
{अतिक्रान्तमतिक्रान्तमनवाप्य मनोरथम्}
{राज्ये राममनिक्षिप्य पिता मे विनशिष्यति} %||2-80-17||

\twolineshloka
{सिद्धार्थाः पितरं वृत्तं तस्मिन्काले ह्युपस्थिते}
{प्रेतकार्येषु सर्वेषु संस्करिष्यन्ति भूमिपम्} %||2-80-18||

\twolineshloka
{रम्यचत्वरसंस्थानां सुविभक्तमहापथाम्}
{हर्म्यप्रासादसम्पन्नां सर्वरत्नविभूषिताम्} %||2-80-19||

\twolineshloka
{गजाश्वरथसम्बाधां तूर्यनादविनादिताम्}
{सर्वकल्याणसम्पूर्णां हृष्टपुष्टजनाकुलाम्} %||2-80-20||

\twolineshloka
{आरामोद्यानसम्पूर्णां समाजोत्सवशालिनीम्}
{सुखिता विचरिष्यन्ति राजधानीं पितुर्मम} %||2-80-21||

\twolineshloka
{अपि सत्यप्रतिज्ञेन सार्धं कुशलिना वयम्}
{निवृत्ते समये ह्यस्मिन्सुखिताः प्रविशेमहि} %||2-80-22||

\twolineshloka
{परिदेवयमानस्य तस्यैवं सुमहात्मनः}
{तिष्ठतो राजपुत्रस्य शर्वरी सात्यवर्तत} %||2-80-23||

\twolineshloka
{प्रभाते विमले सूर्ये कारयित्वा जटा उभौ}
{अस्मिन्भागीरथी तीरे सुखं सन्तारितौ मया} %||2-80-24||

\fourlineindentedshloka
{जटाधरौ तौ द्रुमचीरवाससौ}
{ महाबलौ कुञ्जरयूथपोपमौ}
{वरेषुचापासिधरौ परन्तपौ}
{ व्यवेक्षमाणौ सह सीतया गतौ} %||2-81-25||


{॥इत्यार्षे श्रीमद्रामायणे वाल्मीकीये आदिकाव्ये अयोध्याकाण्डे अशीतितमः सर्गः॥}

\sect{एकाशीतितमः सर्गः}

\twolineshloka
{गुहस्य वचनं श्रुत्वा भरतो भृशमप्रियम्}
{ध्यानं जगाम तत्रैव यत्र तच्छ्रुतमप्रियम्} %||2-81-1||

\twolineshloka
{सुकुमारो महासत्त्वः सिंहस्कन्धो महाभुजः}
{पुण्डरीक विशालाक्षस्तरुणः प्रियदर्शनः} %||2-81-2||

\twolineshloka
{प्रत्याश्वस्य मुहूर्तं तु कालं परमदुर्मनाः}
{पपात सहसा तोत्रैर्हृदि विद्ध इव द्विपः} %||2-81-3||

\twolineshloka
{तदवस्थं तु भरतं शत्रुघ्नोऽनन्तर स्थितः}
{परिष्वज्य रुरोदोच्चैर्विसंज्ञः शोककर्शितः} %||2-81-4||

\twolineshloka
{ततः सर्वाः समापेतुर्मातरो भरतस्य ताः}
{उपवास कृशा दीना भर्तृव्यसनकर्शिताः} %||2-81-5||

\twolineshloka
{ताश्च तं पतितं भूमौ रुदन्त्यः पर्यवारयन्}
{कौसल्या त्वनुसृत्यैनं दुर्मनाः परिषस्वजे} %||2-81-6||

\twolineshloka
{वत्सला स्वं यथा वत्समुपगूह्य तपस्विनी}
{परिपप्रच्छ भरतं रुदन्ती शोकलालसा} %||2-81-7||

\twolineshloka
{पुत्रव्याधिर्न ते कच्चिच्छरीरं परिबाधते}
{अद्य राजकुलस्यास्य त्वदधीनं हि जीवितम्} %||2-81-8||

\twolineshloka
{त्वां दृष्ट्वा पुत्र जीवामि रामे सभ्रातृके गते}
{वृत्ते दशरथे राज्ञि नाथ एकस्त्वमद्य नः} %||2-81-9||

\twolineshloka
{कच्चिन्न लक्ष्मणे पुत्र श्रुतं ते किञ्चिदप्रियम्}
{पुत्र वा ह्येकपुत्रायाः सहभार्ये वनं गते} %||2-81-10||

\twolineshloka
{स मुहूर्तं समाश्वस्य रुदन्नेव महायशाः}
{कौसल्यां परिसान्त्व्येदं गुहं वचनमब्रवीत्} %||2-81-11||

\twolineshloka
{भ्राता मे क्वावसद्रात्रिं क्व सीता क्व च लक्ष्मणः}
{अस्वपच्छयने कस्मिन्किं भुक्त्वा गुह शंस मे} %||2-81-12||

\twolineshloka
{सोऽब्रवीद्भरतं पृष्टो निषादाधिपतिर्गुहः}
{यद्विधं प्रतिपेदे च रामे प्रियहितेऽतिथौ} %||2-81-13||

\twolineshloka
{अन्नमुच्चावचं भक्ष्याः फलानि विविधानि च}
{रामायाभ्यवहारार्थं बहुचोपहृतं मया} %||2-81-14||

\twolineshloka
{तत्सर्वं प्रत्यनुज्ञासीद्रामः सत्यपराक्रमः}
{न हि तत्प्रत्यगृह्णात्स क्षत्रधर्ममनुस्मरन्} %||2-81-15||

\twolineshloka
{न ह्यस्माभिः प्रतिग्राह्यं सखे देयं तु सर्वदा}
{इति तेन वयं राजन्ननुनीता महात्मना} %||2-81-16||

\twolineshloka
{लक्ष्मणेन समानीतं पीत्वा वारि महायशाः}
{औपवास्यं तदाकार्षीद्राघवः सह सीतया} %||2-81-17||

\twolineshloka
{ततस्तु जलशेषेण लक्ष्मणोऽप्यकरोत्तदा}
{वाग्यतास्ते त्रयः सन्ध्यामुपासत समाहिताः} %||2-81-18||

\twolineshloka
{सौमित्रिस्तु ततः पश्चादकरोत्स्वास्तरं शुभम्}
{स्वयमानीय बर्हींषि क्षिप्रं राघव कारणात्} %||2-81-19||

\twolineshloka
{तस्मिन्समाविशद्रामः स्वास्तरे सह सीतया}
{प्रक्षाल्य च तयोः पादावपचक्राम लक्ष्मणः} %||2-81-20||

\twolineshloka
{एतत्तदिङ्गुदीमूलमिदमेव च तत्तृणम्}
{यस्मिन्रामश्च सीता च रात्रिं तां शयितावुभौ} %||2-81-21||

\fourlineindentedshloka
{नियम्य पृष्ठे तु तलाङ्गुलित्रवा}
{ञ्शरैः सुपूर्णाविषुधी परन्तपः}
{महद्धनुः सज्यमुपोह्य लक्ष्मणो}
{ निशामतिष्ठत्परितोऽस्य केवलम्} %||2-81-22||

\fourlineindentedshloka
{ततस्त्वहं चोत्तमबाणचापधृ}
{क्स्थितोऽभवं तत्र स यत्र लक्ष्मणः}
{अतन्द्रिभिर्ज्ञातिभिरात्तकार्मुकै}
{र्महेन्द्रकल्पं परिपालयंस्तदा} %||2-82-23||


{॥इत्यार्षे श्रीमद्रामायणे वाल्मीकीये आदिकाव्ये अयोध्याकाण्डे एकाशीतितमः सर्गः॥}

\sect{द्व्यशीतितमः सर्गः}

\twolineshloka
{तच्छ्रुत्वा निपुणं सर्वं भरतः सह मन्त्रिभिः}
{इङ्गुदीमूलमागम्य रामशय्यामवेक्ष्य ताम्} %||2-82-1||

\twolineshloka
{अब्रवीज्जननीः सर्वा इह तेन महात्मना}
{शर्वरी शयिता भूमाविदमस्य विमर्दितम्} %||2-82-2||

\twolineshloka
{महाभागकुलीनेन महाभागेन धीमता}
{जातो दशरथेनोर्व्यां न रामः स्वप्तुमर्हति} %||2-82-3||

\twolineshloka
{अजिनोत्तरसंस्तीर्णे वरास्तरणसञ्चये}
{शयित्वा पुरुषव्याघ्रः कथं शेते महीतले} %||2-82-4||

\twolineshloka
{प्रासादाग्र विमानेषु वलभीषु च सर्वदा}
{हैमराजतभौमेषु वरास्तरणशालिषु} %||2-82-5||

\twolineshloka
{पुष्पसञ्चयचित्रेषु चन्दनागरुगन्धिषु}
{पाण्डुराभ्रप्रकाशेषु शुकसङ्घरुतेषु च} %||2-82-6||

\twolineshloka
{गीतवादित्रनिर्घोषैर्वराभरणनिःस्वनैः}
{मृदङ्गवरशब्दैश्च सततं प्रतिबोधितः} %||2-82-7||

\twolineshloka
{बन्दिभिर्वन्दितः काले बहुभिः सूतमागधैः}
{गाथाभिरनुरूपाभिः स्तुतिभिश्च परन्तपः} %||2-82-8||

\twolineshloka
{अश्रद्धेयमिदं लोके न सत्यं प्रतिभाति मा}
{मुह्यते खलु मे भावः स्वप्नोऽयमिति मे मतिः} %||2-82-9||

\twolineshloka
{न नूनं दैवतं किञ्चित्कालेन बलवत्तरम्}
{यत्र दाशरथी रामो भूमावेवं शयीत सः} %||2-82-10||

\twolineshloka
{विदेहराजस्य सुता सीता च प्रियदर्शना}
{दयिता शयिता भूमौ स्नुषा दशरथस्य च} %||2-82-11||

\twolineshloka
{इयं शय्या मम भ्रातुरिदं हि परिवर्तितम्}
{स्थण्डिले कठिने सर्वं गात्रैर्विमृदितं तृणम्} %||2-82-12||

\twolineshloka
{मन्ये साभरणा सुप्ता सीतास्मिञ्शयने तदा}
{तत्र तत्र हि दृश्यन्ते सक्ताः कनकबिन्दवः} %||2-82-13||

\twolineshloka
{उत्तरीयमिहासक्तं सुव्यक्तं सीतया तदा}
{तथा ह्येते प्रकाशन्ते सक्ताः कौशेयतन्तवः} %||2-82-14||

\twolineshloka
{मन्ये भर्तुः सुखा शय्या येन बाला तपस्विनी}
{सुकुमारी सती दुःखं न विजानाति मैथिली} %||2-82-15||

\twolineshloka
{सार्वभौम कुले जातः सर्वलोकसुखावहः}
{सर्वलोकप्रियस्त्यक्त्वा राज्यं प्रियमनुत्तमम्} %||2-82-16||

\twolineshloka
{कथमिन्दीवरश्यामो रक्ताक्षः प्रियदर्शनः}
{सुखभागी च दुःखार्हः शयितो भुवि राघवः} %||2-82-17||

\twolineshloka
{सिद्धार्था खलु वैदेही पतिं यानुगता वनम्}
{वयं संशयिताः सर्वे हीनास्तेन महात्मना} %||2-82-18||

\twolineshloka
{अकर्णधारा पृथिवी शून्येव प्रतिभाति मा}
{गते दशरथे स्वर्गे रामे चारण्यमाश्रिते} %||2-82-19||

\twolineshloka
{न च प्रार्थयते कश्चिन्मनसापि वसुन्धराम्}
{वनेऽपि वसतस्तस्य बाहुवीर्याभिरक्षिताम्} %||2-82-20||

\twolineshloka
{शून्यसंवरणारक्षामयन्त्रितहयद्विपाम्}
{अपावृतपुरद्वारां राजधानीमरक्षिताम्} %||2-82-21||

\twolineshloka
{अप्रहृष्टबलां न्यूनां विषमस्थामनावृताम्}
{शत्रवो नाभिमन्यन्ते भक्ष्यान्विषकृतानिव} %||2-82-22||

\twolineshloka
{अद्य प्रभृति भूमौ तु शयिष्येऽहं तृणेषु वा}
{फलमूलाशनो नित्यं जटाचीराणि धारयन्} %||2-82-23||

\twolineshloka
{तस्यार्थमुत्तरं कालं निवत्स्यामि सुखं वने}
{तं प्रतिश्रवमामुच्य नास्य मिथ्या भविष्यति} %||2-82-24||

\twolineshloka
{वसन्तं भ्रातुरर्थाय शत्रुघ्नो मानुवत्स्यति}
{लक्ष्मणेन सह त्वार्यो अयोध्यां पालयिष्यति} %||2-82-25||

\twolineshloka
{अभिषेक्ष्यन्ति काकुत्स्थमयोध्यायां द्विजातयः}
{अपि मे देवताः कुर्युरिमं सत्यं मनोरथम्} %||2-82-26||

\fourlineindentedshloka
{प्रसाद्यमानः शिरसा मया स्वयम्}
{ बहुप्रकारं यदि न प्रपत्स्यते}
{ततोऽनुवत्स्यामि चिराय राघवम्}
{ वने वसन्नार्हति मामुपेक्षितुम्} %||2-83-27||


{॥इत्यार्षे श्रीमद्रामायणे वाल्मीकीये आदिकाव्ये अयोध्याकाण्डे द्व्यशीतितमः सर्गः॥}

\sect{त्र्यशीतितमः सर्गः}

\twolineshloka
{व्युष्य रात्रिं तु तत्रैव गङ्गाकूले स राघवः}
{भरतः काल्यमुत्थाय शत्रुघ्नमिदमब्रवीत्} %||2-83-1||

\twolineshloka
{शत्रुघोत्तिष्ठ किं शेषे निषादाधिपतिं गुहम्}
{शीघ्रमानय भद्रं ते तारयिष्यति वाहिनीम्} %||2-83-2||

\twolineshloka
{जागर्मि नाहं स्वपिमि तथैवार्यं विचिन्तयन्}
{इत्येवमब्रवीद्भ्रात्रा शत्रुघ्नोऽपि प्रचोदितः} %||2-83-3||

\twolineshloka
{इति संवदतोरेवमन्योन्यं नरसिंहयोः}
{आगम्य प्राञ्जलिः काले गुहो भरतमब्रवीत्} %||2-83-4||

\twolineshloka
{कच्चित्सुखं नदीतीरेऽवात्सीः काकुत्स्थ शर्वरीम्}
{कच्चिच्च सह सैन्यस्य तव सर्वमनामयम्} %||2-83-5||

\twolineshloka
{गुहस्य तत्तु वचनं श्रुत्वा स्नेहादुदीरितम्}
{रामस्यानुवशो वाक्यं भरतोऽपीदमब्रवीत्} %||2-83-6||

\twolineshloka
{सुखा नः शर्वरी राजन्पूजिताश्चापि ते वयम्}
{गङ्गां तु नौभिर्बह्वीभिर्दाशाः सन्तारयन्तु नः} %||2-83-7||

\twolineshloka
{ततो गुहः सन्त्वरितः श्रुत्वा भरतशासनम्}
{प्रतिप्रविश्य नगरं तं ज्ञातिजनमब्रवीत्} %||2-83-8||

\twolineshloka
{उत्तिष्ठत प्रबुध्यध्वं भद्रमस्तु हि वः सदा}
{नावः समनुकर्षध्वं तारयिष्याम वाहिनीम्} %||2-83-9||

\twolineshloka
{ते तथोक्ताः समुत्थाय त्वरिता राजशासनात्}
{पञ्च नावां शतान्येव समानिन्युः समन्ततः} %||2-83-10||

\twolineshloka
{अन्याः स्वस्तिकविज्ञेया महाघण्डा धरा वराः}
{शोभमानाः पताकिन्यो युक्तवाताः सुसंहताः} %||2-83-11||

\twolineshloka
{ततः स्वस्तिकविज्ञेयां पाण्डुकम्बलसंवृताम्}
{सनन्दिघोषां कल्याणीं गुहो नावमुपाहरत्} %||2-83-12||

\twolineshloka
{तामारुरोह भरतः शत्रुघ्नश्च महाबलः}
{कौसल्या च सुमित्रा च याश्चान्या राजयोषितः} %||2-83-13||

\twolineshloka
{पुरोहितश्च तत्पूर्वं गुरवे ब्राह्मणाश्च ये}
{अनन्तरं राजदारास्तथैव शकटापणाः} %||2-83-14||

\twolineshloka
{आवासमादीपयतां तीर्थं चाप्यवगाहताम्}
{भाण्डानि चाददानानां घोषस्त्रिदिवमस्पृशत्} %||2-83-15||

\twolineshloka
{पताकिन्यस्तु ता नावः स्वयं दाशैरधिष्ठिताः}
{वहन्त्यो जनमारूढं तदा सम्पेतुराशुगाः} %||2-83-16||

\twolineshloka
{नारीणामभिपूर्णास्तु काश्चित्काश्चित्तु वाजिनाम्}
{कश्चित्तत्र वहन्ति स्म यानयुग्यं महाधनम्} %||2-83-17||

\twolineshloka
{ताः स्म गत्वा परं तीरमवरोप्य च तं जनम्}
{निवृत्ताः काण्डचित्राणि क्रियन्ते दाशबन्धुभिः} %||2-83-18||

\twolineshloka
{सवैजयन्तास्तु गजा गजारोहैः प्रचोदिताः}
{तरन्तः स्म प्रकाशन्ते सध्वजा इव पर्वताः} %||2-83-19||

\twolineshloka
{नावश्चारुरुहुस्त्वन्ये प्लवैस्तेरुस्तथापरे}
{अन्ये कुम्भघटैस्तेरुरन्ये तेरुश्च बाहुभिः} %||2-83-20||

\twolineshloka
{सा पुण्या ध्वजिनी गङ्गां दाशैः सन्तारिता स्वयम्}
{मैत्रे मुहूर्ते प्रययौ प्रयागवनमुत्तमम्} %||2-83-21||

\fourlineindentedshloka
{आश्वासयित्वा च चमूं महात्मा}
{ निवेशयित्वा च यथोपजोषम्}
{द्रष्टुं भरद्वाजमृषिप्रवर्य}
{मृत्विग्वृतः सन्भरतः प्रतस्थे} %||2-84-22||


{॥इत्यार्षे श्रीमद्रामायणे वाल्मीकीये आदिकाव्ये अयोध्याकाण्डे त्र्यशीतितमः सर्गः॥}

\sect{चतुरशीतितमः सर्गः}

\twolineshloka
{भरद्वाजाश्रमं दृष्ट्वा क्रोशादेव नरर्षभः}
{बलं सर्वमवस्थाप्य जगाम सह मन्त्रिभिः} %||2-84-1||

\twolineshloka
{पद्भ्यामेव हि धर्मज्ञो न्यस्तशस्त्रपरिच्छदः}
{वसानो वाससी क्षौमे पुरोधाय पुरोहितम्} %||2-84-2||

\twolineshloka
{ततः सन्दर्शने तस्य भरद्वाजस्य राघवः}
{मन्त्रिणस्तानवस्थाप्य जगामानु पुरोहितम्} %||2-84-3||

\twolineshloka
{वसिष्ठमथ दृष्ट्वैव भरद्वाजो महातपाः}
{सञ्चचालासनात्तूर्णं शिष्यानर्घ्यमिति ब्रुवन्} %||2-84-4||

\twolineshloka
{समागम्य वसिष्ठेन भरतेनाभिवादितः}
{अबुध्यत महातेजाः सुतं दशरथस्य तम्} %||2-84-5||

\twolineshloka
{ताभ्यामर्घ्यं च पाद्यं च दत्त्वा पश्चात्फलानि च}
{आनुपूर्व्याच्च धर्मज्ञः पप्रच्छ कुशलं कुले} %||2-84-6||

\twolineshloka
{अयोध्यायां बले कोशे मित्रेष्वपि च मन्त्रिषु}
{जानन्दशरथं वृत्तं न राजानमुदाहरत्} %||2-84-7||

\twolineshloka
{वसिष्ठो भरतश्चैनं पप्रच्छतुरनामयम्}
{शरीरेऽग्निषु वृक्षेषु शिष्येषु मृगपक्षिषु} %||2-84-8||

\twolineshloka
{तथेति च प्रतिज्ञाय भरद्वाजो महातपाः}
{भरतं प्रत्युवाचेदं राघवस्नेहबन्धनात्} %||2-84-9||

\twolineshloka
{किमिहागमने कार्यं तव राज्यं प्रशासतः}
{एतदाचक्ष्व मे सर्वं न हि मे शुध्यते मनः} %||2-84-10||

\twolineshloka
{सुषुवे यम मित्रघ्नं कौसल्यानन्दवर्धनम्}
{भ्रात्रा सह सभार्यो यश्चिरं प्रव्राजितो वनम्} %||2-84-11||

\twolineshloka
{नियुक्तः स्त्रीनियुक्तेन पित्रा योऽसौ महायशाः}
{वनवासी भवेतीह समाः किल चतुर्दश} %||2-84-12||

\twolineshloka
{कच्चिन्न तस्यापापस्य पापं कर्तुमिहेच्छसि}
{अकण्टकं भोक्तुमना राज्यं तस्यानुजस्य च} %||2-84-13||

\twolineshloka
{एवमुक्तो भरद्वाजं भरतः प्रत्युवाच ह}
{पर्यश्रु नयनो दुःखाद्वाचा संसज्जमानया} %||2-84-14||

\twolineshloka
{हतोऽस्मि यदि मामेवं भगवानपि मन्यते}
{मत्तो न दोषमाशङ्केर्नैवं मामनुशाधि हि} %||2-84-15||

\twolineshloka
{न चैतदिष्टं माता मे यदवोचन्मदन्तरे}
{नाहमेतेन तुष्टश्च न तद्वचनमाददे} %||2-84-16||

\twolineshloka
{अहं तु तं नरव्याघ्रमुपयातः प्रसादकः}
{प्रतिनेतुमयोध्यां च पादौ तस्याभिवन्दितुम्} %||2-84-17||

\twolineshloka
{त्वं मामेवं गतं मत्वा प्रसादं कर्तुमर्हसि}
{शंस मे भगवन्रामः क्व सम्प्रति महीपतिः} %||2-84-18||

\threelineshloka
{उवाच तं भरद्वाजः प्रसादाद्भरतं वचः}
{त्वय्येतत्पुरुषव्याघ्रं युक्तं राघववंशजे}
{गुरुवृत्तिर्दमश्चैव साधूनां चानुयायिता} %||2-84-19||

\twolineshloka
{जाने चैतन्मनःस्थं ते दृढीकरणमस्त्विति}
{अपृच्छं त्वां तवात्यर्थं कीर्तिं समभिवर्धयन्} %||2-84-20||

\threelineshloka
{असौ वसति ते भ्राता चित्रकूटे महागिरौ}
{श्वस्तु गन्तासि तं देशं वसाद्य सह मन्त्रिभिः}
{एतं मे कुरु सुप्राज्ञ कामं कामार्थकोविद} %||2-84-21||

\fourlineindentedshloka
{ततस्तथेत्येवमुदारदर्शनः}
{ प्रतीतरूपो भरतोऽब्रवीद्वचः}
{चकार बुद्धिं च तदा महाश्रमे}
{ निशानिवासाय नराधिपात्मजः} %||2-85-22||


{॥इत्यार्षे श्रीमद्रामायणे वाल्मीकीये आदिकाव्ये अयोध्याकाण्डे चतुरशीतितमः सर्गः॥}

\sect{पञ्चाशीतितमः सर्गः}

\twolineshloka
{कृतबुद्धिं निवासाय तथैव स मुनिस्तदा}
{भरतं कैकयी पुत्रमातिथ्येन न्यमन्त्रयत्} %||2-85-1||

\twolineshloka
{अब्रवीद्भरतस्त्वेनं नन्विदं भवता कृतम्}
{पाद्यमर्घ्यं तथातिथ्यं वने यदूपपद्यते} %||2-85-2||

\twolineshloka
{अथोवाच भरद्वाजो भरतं प्रहसन्निव}
{जाने त्वां प्रीति संयुक्तं तुष्येस्त्वं येन केनचित्} %||2-85-3||

\twolineshloka
{सेनायास्तु तवैतस्याः कर्तुमिच्छामि भोजनम्}
{मम प्रितिर्यथा रूपा त्वमर्हो मनुजर्षभ} %||2-85-4||

\twolineshloka
{किमर्थं चापि निक्षिप्य दूरे बलमिहागतः}
{कस्मान्नेहोपयातोऽसि सबलः पुरुषर्षभ} %||2-85-5||

\twolineshloka
{भरतः प्रत्युवाचेदं प्राञ्जलिस्तं तपोधनम्}
{ससैन्यो नोपयातोऽस्मि भगवन्भगवद्भयात्} %||2-85-6||

\twolineshloka
{वाजि मुख्या मनुष्याश्च मत्ताश्च वर वारणाः}
{प्रच्छाद्य महतीं भूमिं भगवन्ननुयान्ति माम्} %||2-85-7||

\twolineshloka
{ते वृक्षानुदकं भूमिमाश्रमेषूटजांस्तथा}
{न हिंस्युरिति तेनाहमेक एवागतस्ततः} %||2-85-8||

\twolineshloka
{आनीयतामितः सेनेत्याज्ञप्तः परमर्षिणा}
{तथा तु चक्रे भरतः सेनायाः समुपागमम्} %||2-85-9||

\twolineshloka
{अग्निशालां प्रविश्याथ पीत्वापः परिमृज्य च}
{आतिथ्यस्य क्रियाहेतोर्विश्वकर्माणमाह्वयत्} %||2-85-10||

\twolineshloka
{आह्वये विश्वकर्माणमहं त्वष्टारमेव च}
{आतिथ्यं कर्तुमिच्छामि तत्र मे संविधीयताम्} %||2-85-11||

\twolineshloka
{प्राक्स्रोतसश्च या नद्यः प्रत्यक्स्रोतस एव च}
{पृथिव्यामन्तरिक्षे च समायान्त्वद्य सर्वशः} %||2-85-12||

\twolineshloka
{अन्याः स्रवन्तु मैरेयं सुरामन्याः सुनिष्ठिताम्}
{अपराश्चोदकं शीतमिक्षुकाण्डरसोपमम्} %||2-85-13||

\twolineshloka
{आह्वये देवगन्धर्वान्विश्वावसुहहाहुहून्}
{तथैवाप्सरसो देवीर्गन्धर्वीश्चापि सर्वशः} %||2-85-14||

\threelineshloka
{घृताचीमथ विश्वाचीं मिश्रकेशीमलम्बुसाम्}
{शक्रं याश्चोपतिष्ठन्ति ब्रह्माणं याश्च भामिनीः}
{सर्वास्तुम्बुरुणा सार्धमाह्वये सपरिच्छदाः} %||2-85-15||

\twolineshloka
{वनं कुरुषु यद्दिव्यं वासो भूषणपत्रवत्}
{दिव्यनारीफलं शश्वत्तत्कौबेरमिहैव तु} %||2-85-16||

\twolineshloka
{इह मे भगवान्सोमो विधत्तामन्नमुत्तमम्}
{भक्ष्यं भोज्यं च चोष्यं च लेह्यं च विविधं बहु} %||2-85-17||

\twolineshloka
{विचित्राणि च माल्यानि पादपप्रच्युतानि च}
{सुरादीनि च पेयानि मांसानि विविधानि च} %||2-85-18||

\twolineshloka
{एवं समाधिना युक्तस्तेजसाप्रतिमेन च}
{शिक्षास्वरसमायुक्तं तपसा चाब्रवीन्मुनिः} %||2-85-19||

\twolineshloka
{मनसा ध्यायतस्तस्य प्राङ्मुखस्य कृताञ्जलेः}
{आजग्मुस्तानि सर्वाणि दैवतानि पृथक्पृथक्} %||2-85-20||

\twolineshloka
{मलयं दुर्दुरं चैव ततः स्वेदनुदोऽनिलः}
{उपस्पृश्य ववौ युक्त्या सुप्रियात्मा सुखः शिवः} %||2-85-21||

\twolineshloka
{ततोऽभ्यवर्तन्त घना दिव्याः कुसुमवृष्टयः}
{देवदुन्दुभिघोषश्च दिक्षु सर्वासु शुश्रुवे} %||2-85-22||

\twolineshloka
{प्रववुश्चोत्तमा वाता ननृतुश्चाप्सरोगणाः}
{प्रजगुर्देवगन्धर्वा वीणा प्रमुमुचुः स्वरान्} %||2-85-23||

\twolineshloka
{स शब्दो द्यां च भूमिं च प्राणिनां श्रवणानि च}
{विवेशोच्चारितः श्लक्ष्णः समो लयगुणान्वितः} %||2-85-24||

\twolineshloka
{तस्मिन्नुपरते शब्दे दिव्ये श्रोत्रसुखे नृणाम्}
{ददर्श भारतं सैन्यं विधानं विश्वकर्मणः} %||2-85-25||

\twolineshloka
{बभूव हि समा भूमिः समन्तात्पञ्चयोजनम्}
{शाद्वलैर्बहुभिश्छन्ना नीलवैदूर्यसंनिभैः} %||2-85-26||

\twolineshloka
{तस्मिन्बिल्वाः कपित्थाश्च पनसा बीजपूरकाः}
{आमलक्यो बभूवुश्च चूताश्च फलभूषणाः} %||2-85-27||

\twolineshloka
{उत्तरेभ्यः कुरुभ्यश्च वनं दिव्योपभोगवत्}
{आजगाम नदी दिव्या तीरजैर्बहुभिर्वृता} %||2-85-28||

\twolineshloka
{चतुःशालानि शुभ्राणि शालाश्च गजवाजिनाम्}
{हर्म्यप्रासादसङ्घातास्तोरणानि शुभानि च} %||2-85-29||

\twolineshloka
{सितमेघनिभं चापि राजवेश्म सुतोरणम्}
{शुक्लमाल्यकृताकारं दिव्यगन्धसमुक्षितम्} %||2-85-30||

\twolineshloka
{चतुरस्रमसम्बाधं शयनासनयानवत्}
{दिव्यैः सर्वरसैर्युक्तं दिव्यभोजनवस्त्रवत्} %||2-85-31||

\twolineshloka
{उपकल्पित सर्वान्नं धौतनिर्मलभाजनम्}
{कॢप्तसर्वासनं श्रीमत्स्वास्तीर्णशयनोत्तमम्} %||2-85-32||

\twolineshloka
{प्रविवेश महाबाहुरनुज्ञातो महर्षिणा}
{वेश्म तद्रत्नसम्पूर्णं भरतः कैकयीसुतः} %||2-85-33||

\twolineshloka
{अनुजग्मुश्च तं सर्वे मन्त्रिणः सपुरोहिताः}
{बभूवुश्च मुदा युक्ता तं दृष्ट्वा वेश्म संविधिम्} %||2-85-34||

\twolineshloka
{तत्र राजासनं दिव्यं व्यजनं छत्रमेव च}
{भरतो मन्त्रिभिः सार्धमभ्यवर्तत राजवत्} %||2-85-35||

\twolineshloka
{आसनं पूजयामास रामायाभिप्रणम्य च}
{वालव्यजनमादाय न्यषीदत्सचिवासने} %||2-85-36||

\twolineshloka
{आनुपूर्व्यान्निषेदुश्च सर्वे मन्त्रपुरोहिताः}
{ततः सेनापतिः पश्चात्प्रशास्ता च निषेदतुः} %||2-85-37||

\twolineshloka
{ततस्तत्र मुहूर्तेन नद्यः पायसकर्दमाः}
{उपातिष्ठन्त भरतं भरद्वाजस्य शासनत्} %||2-85-38||

\twolineshloka
{तासामुभयतः कूलं पाण्डुमृत्तिकलेपनाः}
{रम्याश्चावसथा दिव्या ब्रह्मणस्तु प्रसादजाः} %||2-85-39||

\twolineshloka
{तेनैव च मुहूर्तेन दिव्याभरणभूषिताः}
{आगुर्विंशतिसाहस्रा ब्राह्मणा प्रहिताः स्त्रियः} %||2-85-40||

\twolineshloka
{सुवर्णमणिमुक्तेन प्रवालेन च शोभिताः}
{आगुर्विंशतिसाहस्राः कुबेरप्रहिताः स्त्रियः} %||2-85-41||

\twolineshloka
{याभिर्गृहीतः पुरुषः सोन्माद इव लक्ष्यते}
{आगुर्विंशतिसाहस्रा नन्दनादप्सरोगणाः} %||2-85-42||

\twolineshloka
{नारदस्तुम्बुरुर्गोपः पर्वतः सूर्यवर्चसः}
{एते गन्धर्वराजानो भरतस्याग्रतो जगुः} %||2-85-43||

\twolineshloka
{अलम्बुसा मिश्रकेशी पुण्डरीकाथ वामना}
{उपानृत्यंस्तु भरतं भरद्वाजस्य शासनात्} %||2-85-44||

\twolineshloka
{यानि माल्यानि देवेषु यानि चैत्ररथे वने}
{प्रयागे तान्यदृश्यन्त भरद्वाजस्य शासनात्} %||2-85-45||

\twolineshloka
{बिल्वा मार्दङ्गिका आसञ्शम्या ग्राहा बिभीतकाः}
{अश्वत्था नर्तकाश्चासन्भरद्वाजस्य तेजसा} %||2-85-46||

\twolineshloka
{ततः सरलतालाश्च तिलका नक्तमालकाः}
{प्रहृष्टास्तत्र सम्पेतुः कुब्जाभूताथ वामनाः} %||2-85-47||

\twolineshloka
{शिंशपामलकी जम्बूर्याश्चान्याः कानने लताः}
{प्रमदा विग्रहं कृत्वा भरद्वाजाश्रमेऽवसन्} %||2-85-48||

\twolineshloka
{सुरां सुरापाः पिबत पायसं च बुभुक्शिताः}
{मांसनि च सुमेध्यानि भक्ष्यन्तां यावदिच्छथ} %||2-85-49||

\twolineshloka
{उत्साद्य स्नापयन्ति स्म नदीतीरेषु वल्गुषु}
{अप्येकमेकं पुरुषं प्रमदाः सत्प चाष्ट च} %||2-85-50||

\twolineshloka
{संवहन्त्यः समापेतुर्नार्यो रुचिरलोचनाः}
{परिमृज्य तथा न्यायं पाययन्ति वराङ्गनाः} %||2-85-51||

\threelineshloka
{हयान्गजान्खरानुष्ट्रांस्तथैव सुरभेः सुतान्}
{इक्षूंश्च मधुजालांश्च भोजयन्ति स्म वाहनान्}
{इक्ष्वाकुवरयोधानां चोदयन्तो महाबलाः} %||2-85-52||

\twolineshloka
{नाश्वबन्धोऽश्वमाजानान्न गजं कुञ्जरग्रहः}
{मत्तप्रमत्तमुदिता चमूः सा तत्र सम्बभौ} %||2-85-53||

\twolineshloka
{तर्पिता सर्वकामैस्ते रक्तचन्दनरूषिताः}
{अप्सरोगणसंयुक्ताः सैन्या वाचमुदैरयन्} %||2-85-54||

\twolineshloka
{नैवायोध्यां गमिष्यामो न गमिष्याम दण्डकान्}
{कुशलं भरतस्यास्तु रामस्यास्तु तथा सुखम्} %||2-85-55||

\twolineshloka
{इति पादातयोधाश्च हस्त्यश्वारोहबन्धकाः}
{अनाथास्तं विधिं लब्ध्वा वाचमेतामुदैरयन्} %||2-85-56||

\twolineshloka
{सम्प्रहृष्टा विनेदुस्ते नरास्तत्र सहस्रशः}
{भरतस्यानुयातारः स्वर्गेऽयमिति चाब्रुवन्} %||2-85-57||

\twolineshloka
{ततो भुक्तवतां तेषां तदन्नममृतोपमम्}
{दिव्यानुद्वीक्ष्य भक्ष्यांस्तानभवद्भक्षणे मतिः} %||2-85-58||

\twolineshloka
{प्रेष्याश्चेट्यश्च वध्वश्च बलस्थाश्चापि सर्वशः}
{बभूवुस्ते भृशं तृप्ताः सर्वे चाहतवाससः} %||2-85-59||

\twolineshloka
{कुञ्जराश्च खरोष्ट्रश्च गोऽश्वाश्च मृगपक्षिणः}
{बभूवुः सुभृतास्तत्र नान्यो ह्यन्यमकल्पयत्} %||2-85-60||

\twolineshloka
{नाशुक्लवासास्तत्रासीत्क्षुधितो मलिनोऽपि वा}
{रजसा ध्वस्तकेशो वा नरः कश्चिददृश्यत} %||2-85-61||

\twolineshloka
{आजैश्चापि च वाराहैर्निष्ठानवरसञ्चयैः}
{फलनिर्यूहसंसिद्धैः सूपैर्गन्धरसान्वितैः} %||2-85-62||

\twolineshloka
{पुष्पध्वजवतीः पूर्णाः शुक्लस्यान्नस्य चाभितः}
{ददृशुर्विस्मितास्तत्र नरा लौहीः सहस्रशः} %||2-85-63||

\twolineshloka
{बभूवुर्वनपार्श्वेषु कूपाः पायसकर्दमाः}
{ताश्च कामदुघा गावो द्रुमाश्चासन्मधुश्च्युतः} %||2-85-64||

\twolineshloka
{वाप्यो मैरेय पूर्णाश्च मृष्टमांसचयैर्वृताः}
{प्रतप्त पिठरैश्चापि मार्गमायूरकौक्कुटैः} %||2-85-65||

\threelineshloka
{पात्रीणां च सहस्राणि शातकुम्भमयानि च}
{स्थाल्यः कुम्भ्यः करम्भ्यश्च दधिपूर्णाः सुसंस्कृताः}
{यौवनस्थस्य गौरस्य कपित्थस्य सुगन्धिनः} %||2-85-66||

\twolineshloka
{ह्रदाः पूर्णा रसालस्य दध्नः श्वेतस्य चापरे}
{बभूवुः पायसस्यान्ते शर्करायाश्च सञ्चयाः} %||2-85-67||

\twolineshloka
{कल्कांश्चूर्णकषायांश्च स्नानानि विविधानि च}
{ददृशुर्भाजनस्थानि तीर्थेषु सरितां नराः} %||2-85-68||

\twolineshloka
{शुक्लानंशुमतश्चापि दन्तधावनसञ्चयान्}
{शुक्लांश्चन्दनकल्कांश्च समुद्गेष्ववतिष्ठतः} %||2-85-69||

\twolineshloka
{दर्पणान्परिमृष्टांश्च वाससां चापि सञ्चयान्}
{पादुकोपानहां चैव युग्मान्यत्र सहस्रशः} %||2-85-70||

\twolineshloka
{आञ्जनीः कङ्कतान्कूर्चांश्छत्राणि च धनूंषि च}
{मर्मत्राणानि चित्राणि शयनान्यासनानि च} %||2-85-71||

\twolineshloka
{प्रतिपानह्रदान्पूर्णान्खरोष्ट्रगजवाजिनाम्}
{अवगाह्य सुतीर्थांश्च ह्रदान्सोत्पल पुष्करान्} %||2-85-72||

\twolineshloka
{नीलवैदूर्यवर्णांश्च मृदून्यवससञ्चयान्}
{निर्वापार्थं पशूनां ते ददृशुस्तत्र सर्वशः} %||2-85-73||

\twolineshloka
{व्यस्मयन्त मनुष्यास्ते स्वप्नकल्पं तदद्भुतम्}
{दृष्ट्वातिथ्यं कृतं तादृग्भरतस्य महर्षिणा} %||2-85-74||

\twolineshloka
{इत्येवं रममाणानां देवानामिव नन्दने}
{भरद्वाजाश्रमे रम्ये सा रात्रिर्व्यत्यवर्तत} %||2-85-75||

\twolineshloka
{प्रतिजग्मुश्च ता नद्यो गन्धर्वाश्च यथागतम्}
{भरद्वाजमनुज्ञाप्य ताश्च सर्वा वराङ्गनाः} %||2-85-76||

\fourlineindentedshloka
{तथैव मत्ता मदिरोत्कटा नरा}
{स्तथैव दिव्यागुरुचन्दनोक्षिताः}
{तथैव दिव्या विविधाः स्रगुत्तमाः}
{ पृथक्प्रकीर्णा मनुजैः प्रमर्दिताः} %||2-86-77||


{॥इत्यार्षे श्रीमद्रामायणे वाल्मीकीये आदिकाव्ये अयोध्याकाण्डे पञ्चाशीतितमः सर्गः॥}

\sect{षडशीतितमः सर्गः}

\twolineshloka
{ततस्तां रजनीमुष्य भरतः सपरिच्छदः}
{कृतातिथ्यो भरद्वाजं कामादभिजगाम ह} %||2-86-1||

\twolineshloka
{तमृषिः पुरुषव्याघ्रं प्रेक्ष्य प्राञ्जलिमागतम्}
{हुताग्निहोत्रो भरतं भरद्वाजोऽभ्यभाषत} %||2-86-2||

\twolineshloka
{कच्चिदत्र सुखा रात्रिस्तवास्मद्विषये गता}
{समग्रस्ते जनः कच्चिदातिथ्ये शंस मेऽनघ} %||2-86-3||

\twolineshloka
{तमुवाचाञ्जलिं कृत्वा भरतोऽभिप्रणम्य च}
{आश्रमादभिनिष्क्रन्तमृषिमुत्तम तेजसम्} %||2-86-4||

\twolineshloka
{सुखोषितोऽस्मि भगवन्समग्रबलवाहनः}
{तर्पितः सर्वकामैश्च सामात्यो बलवत्त्वया} %||2-86-5||

\twolineshloka
{अपेतक्लमसन्तापाः सुभक्ष्याः सुप्रतिश्रयाः}
{अपि प्रेष्यानुपादाय सर्वे स्म सुसुखोषिताः} %||2-86-6||

\twolineshloka
{आमन्त्रयेऽहं भगवन्कामं त्वामृषिसत्तम}
{समीपं प्रस्थितं भ्रातुर्मैरेणेक्षस्व चक्षुषा} %||2-86-7||

\twolineshloka
{आश्रमं तस्य धर्मज्ञ धार्मिकस्य महात्मनः}
{आचक्ष्व कतमो मार्गः कियानिति च शंस मे} %||2-86-8||

\twolineshloka
{इति पृष्टस्तु भरतं भ्रातृदर्शनलालसम्}
{प्रत्युवाच महातेजा भरद्वाजो महातपाः} %||2-86-9||

\twolineshloka
{भरतार्धतृतीयेषु योजनेष्वजने वने}
{चित्रकूटो गिरिस्तत्र रम्यनिर्दरकाननः} %||2-86-10||

\twolineshloka
{उत्तरं पार्श्वमासाद्य तस्य मन्दाकिनी नदी}
{पुष्पितद्रुमसञ्छन्ना रम्यपुष्पितकानना} %||2-86-11||

\twolineshloka
{अनन्तरं तत्सरितश्चित्रकूटश्च पर्वतः}
{ततो पर्णकुटी तात तत्र तौ वसतो ध्रुवम्} %||2-86-12||

\threelineshloka
{दक्षिणेनैव मार्गेण सव्यदक्षिणमेव च}
{गजवाजिरथाकीर्णां वाहिनीं वाहिनीपते}
{वाहयस्व महाभाग ततो द्रक्ष्यसि राघवम्} %||2-86-13||

\twolineshloka
{प्रयाणमिति च श्रुत्वा राजराजस्य योषितः}
{हित्वा यानानि यानार्हा ब्राह्मणं पर्यवारयन्} %||2-86-14||

\twolineshloka
{वेपमाना कृशा दीना सह देव्या सुमन्त्रिया}
{कौसल्या तत्र जग्राह कराभ्यां चरणौ मुनेः} %||2-86-15||

\twolineshloka
{असमृद्धेन कामेन सर्वलोकस्य गर्हिता}
{कैकेयी तस्य जग्राह चरणौ सव्यपत्रपा} %||2-86-16||

\twolineshloka
{तं प्रदक्षिणमागम्य भगवन्तं महामुनिम्}
{अदूराद्भरतस्यैव तस्थौ दीनमनास्तदा} %||2-86-17||

\twolineshloka
{ततः पप्रच्छ भरतं भरद्वाजो दृढव्रतः}
{विशेषं ज्ञातुमिच्छामि मातॄणां तव राघव} %||2-86-18||

\twolineshloka
{एवमुक्तस्तु भरतो भरद्वाजेन धार्मिकः}
{उवाच प्राञ्जलिर्भूत्वा वाक्यं वचनकोविदः} %||2-86-19||

\twolineshloka
{यामिमां भगवन्दीनां शोकानशनकर्शिताम्}
{पितुर्हि महिषीं देवीं देवतामिव पश्यसि} %||2-86-20||

\twolineshloka
{एषा तं पुरुषव्याघ्रं सिंहविक्रान्तगामिनम्}
{कौसल्या सुषुवे रामं धातारमदितिर्यथा} %||2-86-21||

\twolineshloka
{अस्या वामभुजं श्लिष्टा यैषा तिष्ठति दुर्मनाः}
{कर्णिकारस्य शाखेव शीर्णपुष्पा वनान्तरे} %||2-86-22||

\twolineshloka
{एतस्यास्तौ सुतौ देव्याः कुमारौ देववर्णिनौ}
{उभौ लक्ष्मणशत्रुघ्नौ वीरौ सत्यपराक्रमौ} %||2-86-23||

\twolineshloka
{यस्याः कृते नरयाघ्रौ जीवनाशमितो गतौ}
{राजा पुत्रविहीनश्च स्वर्गं दशरथो गतः} %||2-86-24||

\threelineshloka
{ऐश्वर्यकामां कैकेयीमनार्यामार्यरूपिणीम्}
{ममैतां मातरं विद्धि नृशंसां पापनिश्चयाम्}
{यतोमूलं हि पश्यामि व्यसनं महदात्मनः} %||2-86-25||

\twolineshloka
{इत्युक्त्वा नरशार्दूलो बाष्पगद्गदया गिरा}
{स निशश्वास ताम्राक्षो क्रुद्धो नाग इवासकृत्} %||2-86-26||

\twolineshloka
{भरद्वाजो महर्षिस्तं ब्रुवन्तं भरतं तदा}
{प्रत्युवाच महाबुद्धिरिदं वचनमर्थवत्} %||2-86-27||

\twolineshloka
{न दोषेणावगन्तव्या कैकेयी भरत त्वया}
{रामप्रव्राजनं ह्येतत्सुखोदर्कं भविष्यति} %||2-86-28||

\twolineshloka
{अभिवाद्य तु संसिद्धः कृत्वा चैनं प्रदक्षिणम्}
{आमन्त्र्य भरतः सैन्यं युज्यतामित्यचोदयत्} %||2-86-29||

\twolineshloka
{ततो वाजिरथान्युक्त्वा दिव्यान्हेमपरिष्क्रितान्}
{अध्यारोहत्प्रयाणार्थी बहून्बहुविधो जनः} %||2-86-30||

\twolineshloka
{गजकन्यागजाश्चैव हेमकक्ष्याः पताकिनः}
{जीमूता इव घर्मान्ते सघोषाः सम्प्रतस्थिरे} %||2-86-31||

\twolineshloka
{विविधान्यपि यानानि महानि च लघूनि च}
{प्रययुः सुमहार्हाणि पादैरेव पदातयः} %||2-86-32||

\twolineshloka
{अथ यानप्रवेकैस्तु कौसल्याप्रमुखाः स्त्रियः}
{रामदर्शनकाङ्क्षिण्यः प्रययुर्मुदितास्तदा} %||2-86-33||

\twolineshloka
{स चार्कतरुणाभासां नियुक्तां शिबिकां शुभाम्}
{आस्थाय प्रययौ श्रीमान्भरतः सपरिच्छदः} %||2-86-34||

\threelineshloka
{सा प्रयाता महासेना गजवाजिरथाकुला}
{दक्षिणां दिशमावृत्य महामेघ इवोत्थितः}
{वनानि तु व्यतिक्रम्य जुष्टानि मृगपक्षिभिः} %||2-86-35||

\fourlineindentedshloka
{सा सम्प्रहृष्टद्विपवाजियोधा}
{ वित्रासयन्ती मृगपक्षिसङ्घान्}
{महद्वनं तत्प्रविगाहमाना}
{ रराज सेना भरतस्य तत्र} %||2-87-36||


{॥इत्यार्षे श्रीमद्रामायणे वाल्मीकीये आदिकाव्ये अयोध्याकाण्डे षडशीतितमः सर्गः॥}

\sect{सप्ताशीतितमः सर्गः}

\twolineshloka
{तया महत्या यायिन्या ध्वजिन्या वनवासिनः}
{अर्दिता यूथपा मत्ताः सयूथाः सम्प्रदुद्रुवुः} %||2-87-1||

\twolineshloka
{ऋक्षाः पृषतसङ्घाश्च रुरवश्च समन्ततः}
{दृश्यन्ते वनराजीषु गिरिष्वपि नदीषु च} %||2-87-2||

\twolineshloka
{स सम्प्रतस्थे धर्मात्मा प्रीतो दशरथात्मजः}
{वृतो महत्या नादिन्या सेनया चतुरङ्गया} %||2-87-3||

\twolineshloka
{सागरौघनिभा सेना भरतस्य महात्मनः}
{महीं सञ्छादयामास प्रावृषि द्यामिवाम्बुदः} %||2-87-4||

\twolineshloka
{तुरङ्गौघैरवतता वारणैश्च महाजवैः}
{अनालक्ष्या चिरं कालं तस्मिन्काले बभूव भूः} %||2-87-5||

\twolineshloka
{स यात्वा दूरमध्वानं सुपरिश्रान्त वाहनः}
{उवाच भरतः श्रीमान्वसिष्ठं मन्त्रिणां वरम्} %||2-87-6||

\twolineshloka
{यादृशं लक्ष्यते रूपं यथा चैव श्रुतं मया}
{व्यक्तं प्राप्ताः स्म तं देशं भरद्वाजो यमब्रवीत्} %||2-87-7||

\twolineshloka
{अयं गिरिश्चित्रकूटस्तथा मन्दाकिनी नदी}
{एतत्प्रकाशते दूरान्नीलमेघनिभं वनम्} %||2-87-8||

\twolineshloka
{गिरेः सानूनि रम्याणि चित्रकूटस्य सम्प्रति}
{वारणैरवमृद्यन्ते मामकैः पर्वतोपमैः} %||2-87-9||

\twolineshloka
{मुञ्चन्ति कुसुमान्येते नगाः पर्वतसानुषु}
{नीला इवातपापाये तोयं तोयधरा घनाः} %||2-87-10||

\twolineshloka
{किन्नराचरितोद्देशं पश्य शत्रुघ्न पर्वतम्}
{हयैः समन्तादाकीर्णं मकरैरिव सागरम्} %||2-87-11||

\twolineshloka
{एते मृगगणा भान्ति शीघ्रवेगाः प्रचोदिताः}
{वायुप्रविद्धाः शरदि मेघराज्य इवाम्बरे} %||2-87-12||

\twolineshloka
{कुर्वन्ति कुसुमापीडाञ्शिरःसु सुरभीनमी}
{मेघप्रकाशैः फलकैर्दाक्षिणात्या यथा नराः} %||2-87-13||

\twolineshloka
{निष्कूजमिव भूत्वेदं वनं घोरप्रदर्शनम्}
{अयोध्येव जनाकीर्णा सम्प्रति प्रतिभाति मा} %||2-87-14||

\twolineshloka
{खुरैरुदीरितो रेणुर्दिवं प्रच्छाद्य तिष्ठति}
{तं वहत्यनिलः शीघ्रं कुर्वन्निव मम प्रियम्} %||2-87-15||

\twolineshloka
{स्यन्दनांस्तुरगोपेतान्सूतमुख्यैरधिष्ठितान्}
{एतान्सम्पततः शीघ्रं पश्य शत्रुघ्न कानने} %||2-87-16||

\twolineshloka
{एतान्वित्रासितान्पश्य बर्हिणः प्रियदर्शनान्}
{एतमाविशतः शैलमधिवासं पतत्रिणाम्} %||2-87-17||

\twolineshloka
{अतिमात्रमयं देशो मनोज्ञः प्रतिभाति मा}
{तापसानां निवासोऽयं व्यक्तं स्वर्गपथो यथा} %||2-87-18||

\twolineshloka
{मृगा मृगीभिः सहिता बहवः पृषता वने}
{मनोज्ञ रूपा लक्ष्यन्ते कुसुमैरिव चित्रितः} %||2-87-19||

\twolineshloka
{साधु सैन्याः प्रतिष्ठन्तां विचिन्वन्तु च काननम्}
{यथा तौ पुरुषव्याघ्रौ दृश्येते रामलक्ष्मणौ} %||2-87-20||

\twolineshloka
{भरतस्य वचः श्रुत्वा पुरुषाः शस्त्रपाणयः}
{विविशुस्तद्वनं शूरा धूमं च ददृशुस्ततः} %||2-87-21||

\twolineshloka
{ते समालोक्य धूमाग्रमूचुर्भरतमागताः}
{नामनुष्ये भवत्यग्निर्व्यक्तमत्रैव राघवौ} %||2-87-22||

\twolineshloka
{अथ नात्र नरव्याघ्रौ राजपुत्रौ परन्तपौ}
{अन्ये रामोपमाः सन्ति व्यक्तमत्र तपस्विनः} %||2-87-23||

\twolineshloka
{तच्छ्रुत्वा भरतस्तेषां वचनं साधु सम्मतम्}
{सैन्यानुवाच सर्वांस्तानमित्रबलमर्दनः} %||2-87-24||

\twolineshloka
{यत्ता भवन्तस्तिष्ठन्तु नेतो गन्तव्यमग्रतः}
{अहमेव गमिष्यामि सुमन्त्रो गुरुरेव च} %||2-87-25||

\twolineshloka
{एवमुक्तास्ततः सर्वे तत्र तस्थुः समन्ततः}
{भरतो यत्र धूमाग्रं तत्र दृष्टिं समादधत्} %||2-87-26||

\fourlineindentedshloka
{व्यवस्थिता या भरतेन सा चमू}
{र्निरीक्षमाणापि च धूममग्रतः}
{बभूव हृष्टा नचिरेण जानती}
{ प्रियस्य रामस्य समागमं तदा} %||2-88-27||


{॥इत्यार्षे श्रीमद्रामायणे वाल्मीकीये आदिकाव्ये अयोध्याकाण्डे सप्ताशीतितमः सर्गः॥}

\sect{अष्टाशीतितमः सर्गः}

\twolineshloka
{दीर्घकालोषितस्तस्मिन्गिरौ गिरिवनप्रियः}
{विदेह्याः प्रियमाकाङ्क्षन्स्वं च चित्तं विलोभयन्} %||2-88-1||

\twolineshloka
{अथ दाशरथिश्चित्रं चित्रकूटमदर्शयत्}
{भार्याममरसङ्काशः शचीमिव पुरन्दरः} %||2-88-2||

\twolineshloka
{न राज्याद्भ्रंशनं भद्रे न सुहृद्भिर्विनाभवः}
{मनो मे बाधते दृष्ट्वा रमणीयमिमं गिरिम्} %||2-88-3||

\twolineshloka
{पश्येममचलं भद्रे नानाद्विजगणायुतम्}
{शिखरैः खमिवोद्विद्धैर्धातुमद्भिर्विभूषितम्} %||2-88-4||

\twolineshloka
{केचिद्रजतसङ्काशाः केचित्क्षतजसंनिभाः}
{पीतमाञ्जिष्ठवर्णाश्च केचिन्मणिवरप्रभाः} %||2-88-5||

\twolineshloka
{पुष्यार्ककेतुकाभाश्च केचिज्ज्योती रसप्रभाः}
{विराजन्तेऽचलेन्द्रस्य देशा धातुविभूषिताः} %||2-88-6||

\twolineshloka
{नानामृगगणद्वीपितरक्ष्वृक्षगणैर्वृतः}
{अदुष्टैर्भात्ययं शैलो बहुपक्षिसमाकुलः} %||2-88-7||

\twolineshloka
{आम्रजम्ब्वसनैर्लोध्रैः प्रियालैः पनसैर्धवैः}
{अङ्कोलैर्भव्यतिनिशैर्बिल्वतिन्दुकवेणुभिः} %||2-88-8||

\twolineshloka
{काश्मर्यरिष्टवरणैर्मधूकैस्तिलकैस्तथा}
{बदर्यामलकैर्नीपैर्वेत्रधन्वनबीजकैः} %||2-88-9||

\twolineshloka
{पुष्पवद्भिः फलोपेतैश्छायावद्भिर्मनोरमैः}
{एवमादिभिराकीर्णः श्रियं पुष्यत्ययं गिरिः} %||2-88-10||

\twolineshloka
{शैलप्रस्थेषु रम्येषु पश्येमान्कामहर्षणान्}
{किन्नरान्द्वन्द्वशो भद्रे रममाणान्मनस्विनः} %||2-88-11||

\twolineshloka
{शाखावसक्तान्खड्गांश्च प्रवराण्यम्बराणि च}
{पश्य विद्याधरस्त्रीणां क्रीडेद्देशान्मनोरमान्} %||2-88-12||

\twolineshloka
{जलप्रपातैरुद्भेदैर्निष्यन्दैश्च क्वचित्क्वचित्}
{स्रवद्भिर्भात्ययं शैलः स्रवन्मद इव द्विपः} %||2-88-13||

\twolineshloka
{गुहासमीरणो गन्धान्नानापुष्पभवान्वहन्}
{घ्राणतर्पणमभ्येत्य कं नरं न प्रहर्षयेत्} %||2-88-14||

\twolineshloka
{यदीह शरदोऽनेकास्त्वया सार्धमनिन्दिते}
{लक्ष्मणेन च वत्स्यामि न मां शोकः प्रधक्ष्यति} %||2-88-15||

\twolineshloka
{बहुपुष्पफले रम्ये नानाद्विजगणायुते}
{विचित्रशिखरे ह्यस्मिन्रतवानस्मि भामिनि} %||2-88-16||

\twolineshloka
{अनेन वनवासेन मया प्राप्तं फलद्वयम्}
{पितुश्चानृणता धर्मे भरतस्य प्रियं तथा} %||2-88-17||

\twolineshloka
{वैदेहि रमसे कच्चिच्चित्रकूटे मया सह}
{पश्यन्ती विविधान्भावान्मनोवाक्कायसंयतान्} %||2-88-18||

\twolineshloka
{इदमेवामृतं प्राहू राज्ञां राजर्षयः परे}
{वनवासं भवार्थाय प्रेत्य मे प्रपितामहाः} %||2-88-19||

\twolineshloka
{शिलाः शैलस्य शोभन्ते विशालाः शतशोऽभितः}
{बहुला बहुलैर्वर्णैर्नीलपीतसितारुणैः} %||2-88-20||

\twolineshloka
{निशि भान्त्यचलेन्द्रस्य हुताशनशिखा इव}
{ओषध्यः स्वप्रभा लक्ष्म्या भ्राजमानाः सहस्रशः} %||2-88-21||

\twolineshloka
{केचित्क्षयनिभा देशाः केचिदुद्यानसंनिभाः}
{केचिदेकशिला भान्ति पर्वतस्यास्य भामिनि} %||2-88-22||

\twolineshloka
{भित्त्वेव वसुधां भाति चित्रकूटः समुत्थितः}
{चित्रकूटस्य कूटोऽसौ दृश्यते सर्वतः शिवः} %||2-88-23||

\twolineshloka
{कुष्ठपुंनागतगरभूर्जपत्रोत्तरच्छदान्}
{कामिनां स्वास्तरान्पश्य कुशेशयदलायुतान्} %||2-88-24||

\twolineshloka
{मृदिताश्चापविद्धाश्च दृश्यन्ते कमलस्रजः}
{कामिभिर्वनिते पश्य फलानि विविधानि च} %||2-88-25||

\twolineshloka
{वस्वौकसारां नलिनीमत्येतीवोत्तरान्कुरून्}
{पर्वतश्चित्रकूटोऽसौ बहुमूलफलोदकः} %||2-88-26||

\fourlineindentedshloka
{इमं तु कालं वनिते विजह्रिवाम्}
{स्त्वया च सीते सह लक्ष्मणेन च}
{रतिं प्रपत्स्ये कुलधर्मवर्धिनीम्}
{ सतां पथि स्वैर्नियमैः परैः स्थितः} %||2-89-27||


{॥इत्यार्षे श्रीमद्रामायणे वाल्मीकीये आदिकाव्ये अयोध्याकाण्डे अष्टाशीतितमः सर्गः॥}

\sect{एकोननवतितमः सर्गः}

\twolineshloka
{अथ शैलाद्विनिष्क्रम्य मैथिलीं कोसलेश्वरः}
{अदर्शयच्छुभजलां रम्यां मन्दाकिनीं नदीम्} %||2-89-1||

\twolineshloka
{अब्रवीच्च वरारोहां चारुचन्द्रनिभाननाम्}
{विदेहराजस्य सुतां रामो राजीवलोचनः} %||2-89-2||

\twolineshloka
{विचित्रपुलिनां रम्यां हंससारससेविताम्}
{कुसुमैरुपसम्पन्नां पश्य मन्दाकिनीं नदीम्} %||2-89-3||

\twolineshloka
{नानाविधैस्तीररुहैर्वृतां पुष्पफलद्रुमैः}
{राजन्तीं राजराजस्य नलिनीमिव सर्वतः} %||2-89-4||

\twolineshloka
{मृगयूथनिपीतानि कलुषाम्भांसि साम्प्रतम्}
{तीर्थानि रमणीयानि रतिं संजनयन्ति मे} %||2-89-5||

\twolineshloka
{जटाजिनधराः काले वल्कलोत्तरवाससः}
{ऋषयस्त्ववगाहन्ते नदीं मन्दाकिनीं प्रिये} %||2-89-6||

\twolineshloka
{आदित्यमुपतिष्ठन्ते नियमादूर्ध्वबाहवः}
{एतेऽपरे विशालाक्षि मुनयः संशितव्रताः} %||2-89-7||

\twolineshloka
{मारुतोद्धूत शिखरैः प्रनृत्त इव पर्वतः}
{पादपैः पत्रपुष्पाणि सृजद्भिरभितो नदीम्} %||2-89-8||

\twolineshloka
{कच्चिन्मणिनिकाशोदां कच्चित्पुलिनशालिनीम्}
{कच्चित्सिद्धजनाकीर्णां पश्य मन्दाकिनीं नदीम्} %||2-89-9||

\twolineshloka
{निर्धूतान्वायुना पश्य विततान्पुष्पसञ्चयान्}
{पोप्लूयमानानपरान्पश्य त्वं जलमध्यगान्} %||2-89-10||

\twolineshloka
{तांश्चातिवल्गु वचसो रथाङ्गाह्वयना द्विजाः}
{अधिरोहन्ति कल्याणि निष्कूजन्तः शुभा गिरः} %||2-89-11||

\twolineshloka
{दर्शनं चित्रकूटस्य मन्दाकिन्याश्च शोभने}
{अधिकं पुरवासाच्च मन्ये च तव दर्शनात्} %||2-89-12||

\twolineshloka
{विधूतकलुषैः सिद्धैस्तपोदमशमान्वितैः}
{नित्यविक्षोभित जलां विहाहस्व मया सह} %||2-89-13||

\twolineshloka
{सखीवच्च विगाहस्व सीते मन्दकिनीमिमाम्}
{कमलान्यवमज्जन्ती पुष्कराणि च भामिनि} %||2-89-14||

\twolineshloka
{त्वं पौरजनवद्व्यालानयोध्यामिव पर्वतम्}
{मन्यस्व वनिते नित्यं सरयूवदिमां नदीम्} %||2-89-15||

\twolineshloka
{लक्ष्मणश्चैव धर्मात्मा मन्निदेशे व्यवस्थितः}
{त्वं चानुकूला वैदेहि प्रीतिं जनयथो मम} %||2-89-16||

\twolineshloka
{उपस्पृशंस्त्रिषवणं मधुमूलफलाशनः}
{नायोध्यायै न राज्याय स्पृहयेऽद्य त्वया सह} %||2-89-17||

\fourlineindentedshloka
{इमां हि रम्यां गजयूथलोलिताम्}
{ निपीततोयां गजसिंहवानरैः}
{सुपुष्पितैः पुष्पधरैरलङ्कृताम्}
{ न सोऽस्ति यः स्यान्न गतक्रमः सुखी} %||2-89-18||

\fourlineindentedshloka
{इतीव रामो बहुसङ्गतं वचः}
{ प्रिया सहायः सरितं प्रति ब्रुवन्}
{चचार रम्यं नयनाञ्जनप्रभम्}
{ स चित्रकूटं रघुवंशवर्धनः} %||2-90-19||


{॥इत्यार्षे श्रीमद्रामायणे वाल्मीकीये आदिकाव्ये अयोध्याकाण्डे एकोननवतितमः सर्गः॥}

\sect{नवतितमः सर्गः}

\twolineshloka
{तथा तत्रासतस्तस्य भरतस्योपयायिनः}
{सैन्य रेणुश्च शब्दश्च प्रादुरास्तां नभः स्पृशौ} %||2-90-1||

\twolineshloka
{एतस्मिन्नन्तरे त्रस्ताः शब्देन महता ततः}
{अर्दिता यूथपा मत्ताः सयूथा दुद्रुवुर्दिशः} %||2-90-2||

\twolineshloka
{स तं सैन्यसमुद्भूतं शब्दं शुश्रव राघवः}
{तांश्च विप्रद्रुतान्सर्वान्यूथपानन्ववैक्षत} %||2-90-3||

\twolineshloka
{तांश्च विद्रवतो दृष्ट्वा तं च श्रुत्वा स निःस्वनम्}
{उवाच रामः सौमित्रिं लक्ष्मणं दीप्ततेजसम्} %||2-90-4||

\twolineshloka
{हन्त लक्ष्मण पश्येह सुमित्रा सुप्रजास्त्वया}
{भीमस्तनितगम्भीरस्तुमुलः श्रूयते स्वनः} %||2-90-5||

\threelineshloka
{राजा वा राजमात्रो वा मृगयामटते वने}
{अन्यद्वा श्वापदं किञ्चित्सौमित्रे ज्ञातुमर्हसि}
{सर्वमेतद्यथातत्त्वमचिराज्ज्ञातुमर्हसि} %||2-90-6||

\twolineshloka
{स लक्ष्मणः सन्त्वरितः सालमारुह्य पुष्पितम्}
{प्रेक्षमाणो दिशः सर्वाः पूर्वां दिशमवैक्षत} %||2-90-7||

\twolineshloka
{उदङ्मुखः प्रेक्षमाणो ददर्श महतीं चमूम्}
{रथाश्वगजसम्बाधां यत्तैर्युक्तां पदातिभिः} %||2-90-8||

\twolineshloka
{तामश्वगजसम्पूर्णां रथध्वजविभूषिताम्}
{शशंस सेनां रामाय वचनं चेदमब्रवीत्} %||2-90-9||

\twolineshloka
{अग्निं संशमयत्वार्यः सीता च भजतां गुहाम्}
{सज्यं कुरुष्व चापं च शरांश्च कवचं तथा} %||2-90-10||

\twolineshloka
{तं रामः पुरुषव्याघ्रो लक्ष्मणं प्रत्युवाच ह}
{अङ्गावेक्षस्व सौमित्रे कस्यैतां मन्यसे चमूम्} %||2-90-11||

\twolineshloka
{एवमुक्क्तस्तु रामेण लक्ष्माणो वाक्यमब्रवीत्}
{दिधक्षन्निव तां सेनां रुषितः पावको यथा} %||2-90-12||

\twolineshloka
{सम्पन्नं राज्यमिच्छंस्तु व्यक्तं प्राप्याभिषेचनम्}
{आवां हन्तुं समभ्येति कैकेय्या भरतः सुतः} %||2-90-13||

\twolineshloka
{एष वै सुमहाञ्श्रीमान्विटपी सम्प्रकाशते}
{विराजत्युद्गतस्कन्धः कोविदार ध्वजो रथे} %||2-90-14||

\twolineshloka
{भजन्त्येते यथाकाममश्वानारुह्य शीघ्रगान्}
{एते भ्राजन्ति संहृष्टा जगानारुह्य सादिनः} %||2-90-15||

\twolineshloka
{गृहीतधनुषौ चावां गिरिं वीर श्रयावहे}
{अपि नौ वशमागच्छेत्कोविदारध्वजो रणे} %||2-90-16||

\twolineshloka
{अपि द्रक्ष्यामि भरतं यत्कृते व्यसनं महत्}
{त्वया राघव सम्प्राप्तं सीतया च मया तथा} %||2-90-17||

\twolineshloka
{यन्निमित्तं भवान्राज्याच्च्युतो राघव शाश्वतीम्}
{सम्प्राप्तोऽयमरिर्वीर भरतो वध्य एव मे} %||2-90-18||

\threelineshloka
{भरतस्य वधे दोषं नाहं पश्यामि राघव}
{पूर्वापकरिणां त्यागे न ह्यधर्मो विधीयते}
{एतस्मिन्निहते कृत्स्नामनुशाधि वसुन्धराम्} %||2-90-19||

\twolineshloka
{अद्य पुत्रं हतं सङ्ख्ये कैकेयी राज्यकामुका}
{मया पश्येत्सुदुःखार्ता हस्तिभग्नमिव द्रुमम्} %||2-90-20||

\twolineshloka
{कैकेयीं च वधिष्यामि सानुबन्धां सबान्धवाम्}
{कलुषेणाद्य महता मेदिनी परिमुच्यताम्} %||2-90-21||

\twolineshloka
{अद्येमं संयतं क्रोधमसत्कारं च मानद}
{मोक्ष्यामि शत्रुसैन्येषु कक्षेष्विव हुताशनम्} %||2-90-22||

\twolineshloka
{अद्यैतच्चित्रकूटस्य काननं निशितैः शरैः}
{भिन्दञ्शत्रुशरीराणि करिष्ये शोणितोक्षितम्} %||2-90-23||

\twolineshloka
{शरैर्निर्भिन्नहृदयान्कुञ्जरांस्तुरगांस्तथा}
{श्वापदाः परिकर्षन्तु नराश्च निहतान्मया} %||2-90-24||

\twolineshloka
{शराणां धनुषश्चाहमनृणोऽस्मि महावने}
{ससैन्यं भरतं हत्वा भविष्यामि न संशयः} %||2-91-25||


{॥इत्यार्षे श्रीमद्रामायणे वाल्मीकीये आदिकाव्ये अयोध्याकाण्डे नवतितमः सर्गः॥}

\sect{एकनवतितमः सर्गः}

\twolineshloka
{सुसंरब्धं तु सौमित्रिं लक्ष्मणं क्रोधमूर्छितम्}
{रामस्तु परिसान्त्व्याथ वचनं चेदमब्रवीत्} %||2-91-1||

\twolineshloka
{किमत्र धनुषा कार्यमसिना वा सचर्मणा}
{महेष्वासे महाप्राज्ञे भरते स्वयमागते} %||2-91-2||

\twolineshloka
{प्राप्तकालं यदेषोऽस्मान्भरतो द्रष्टुमिच्छति}
{अस्मासु मनसाप्येष नाहितं किञ्चिदाचरेत्} %||2-91-3||

\twolineshloka
{विप्रियं कृतपूर्वं ते भरतेन कदा न किम्}
{ईदृशं वा भयं तेऽद्य भरतं योऽत्र शङ्कसे} %||2-91-4||

\twolineshloka
{न हि ते निष्ठुरं वाच्यो भरतो नाप्रियं वचः}
{अहं ह्यप्रियमुक्तः स्यां भरतस्याप्रिये कृते} %||2-91-5||

\twolineshloka
{कथं नु पुत्राः पितरं हन्युः कस्याञ्चिदापदि}
{भ्राता वा भ्रातरं हन्यात्सौमित्रे प्राणमात्मनः} %||2-91-6||

\twolineshloka
{यदि राज्यस्य हेतोस्त्वमिमां वाचं प्रभाषसे}
{वक्ष्यामि भरतं दृष्ट्वा राज्यमस्मै प्रदीयताम्} %||2-91-7||

\twolineshloka
{उच्यमानो हि भरतो मया लक्ष्मण तत्त्वतः}
{राज्यमस्मै प्रयच्छेति बाढमित्येव वक्ष्यति} %||2-91-8||

\twolineshloka
{तथोक्तो धर्मशीलेन भ्रात्रा तस्य हिते रतः}
{लक्ष्मणः प्रविवेशेव स्वानि गात्राणि लज्जया} %||2-91-9||

\twolineshloka
{व्रीडितं लक्ष्मणं दृष्ट्वा राघवः प्रत्युवाच ह}
{एष मन्ये महाबाहुरिहास्मान्द्रष्टुमागतः} %||2-91-10||

\twolineshloka
{वनवासमनुध्याय गृहाय प्रतिनेष्यति}
{इमां वाप्येश वैदेहीमत्यन्तसुखसेविनीम्} %||2-91-11||

\twolineshloka
{एतौ तौ सम्प्रकाशेते गोत्रवन्तौ मनोरमौ}
{वायुवेगसमौ वीर जवनौ तुरगोत्तमौ} %||2-91-12||

\twolineshloka
{स एष सुमहाकायः कम्पते वाहिनीमुखे}
{नागः शत्रुंजयो नाम वृद्धस्तातस्य धीमतः} %||2-91-13||

\twolineshloka
{अवतीर्य तु सालाग्रात्तस्मात्स समितिंजयः}
{लक्ष्मणः प्राञ्जलिर्भूत्वा तस्थौ रामस्य पार्श्वतः} %||2-91-14||

\twolineshloka
{भरतेनाथ सन्दिष्टा सम्मर्दो न भवेदिति}
{समन्तात्तस्य शैलस्य सेनावासमकल्पयत्} %||2-91-15||

\twolineshloka
{अध्यर्धमिक्ष्वाकुचमूर्योजनं पर्वतस्य सा}
{पार्श्वे न्यविशदावृत्य गजवाजिरथाकुला} %||2-91-16||

\fourlineindentedshloka
{सा चित्रकूटे भरतेन सेना}
{ धर्मं पुरस्कृत्य विधूय दर्पम्}
{प्रसादनार्थं रघुनन्दनस्य}
{ विरोचते नीतिमता प्रणीता} %||2-92-17||


{॥इत्यार्षे श्रीमद्रामायणे वाल्मीकीये आदिकाव्ये अयोध्याकाण्डे एकनवतितमः सर्गः॥}

\sect{द्विनवतितमः सर्गः}

\twolineshloka
{निवेश्य सेनां तु विभुः पद्भ्यां पादवतां वरः}
{अभिगन्तुं स काकुत्स्थमियेष गुरुवर्तकम्} %||2-92-1||

\twolineshloka
{निविष्ट मात्रे सैन्ये तु यथोद्देशं विनीतवत्}
{भरतो भ्रातरं वाक्यं शत्रुघ्नमिदमब्रवीत्} %||2-92-2||

\twolineshloka
{क्षिप्रं वनमिदं सौम्य नरसङ्घैः समन्ततः}
{लुब्धैश्च सहितैरेभिस्त्वमन्वेषितुमर्हसि} %||2-92-3||

\twolineshloka
{यावन्न रामं द्रक्ष्यामि लक्ष्मणं वा महाबलम्}
{वैदेहीं वा महाभागां न मे शान्तिर्भविष्यति} %||2-92-4||

\twolineshloka
{यावन्न चन्द्रसङ्काशं द्रक्ष्यामि शुभमाननम्}
{भ्रातुः पद्मपलाशाक्षं न मे शान्तिर्भविष्यति} %||2-92-5||

\twolineshloka
{यावन्न चरणौ भ्रातुः पार्थिव व्यञ्जनान्वितौ}
{शिरसा धारयिष्यामि न मे शान्तिर्भविष्यति} %||2-92-6||

\twolineshloka
{यावन्न राज्ये राज्यार्हः पितृपैतामहे स्थितः}
{अभिषेकजलक्लिन्नो न मे शान्तिर्भविष्यति} %||2-92-7||

\twolineshloka
{कृतकृत्या महाभागा वैदेही जनकात्मजा}
{भर्तारं सागरान्तायाः पृथिव्या यानुगच्छति} %||2-92-8||

\twolineshloka
{सुभगश्चित्रकूटोऽसौ गिरिराजोपमो गिरिः}
{यस्मिन्वसति काकुत्स्थः कुबेर इवनन्दने} %||2-92-9||

\twolineshloka
{कृतकार्यमिदं दुर्गं वनं व्यालनिषेवितम्}
{यदध्यास्ते महातेजा रामः शस्त्रभृतां वरः} %||2-92-10||

\twolineshloka
{एवमुक्त्वा महातेजा भरतः पुरुषर्षभः}
{पद्भ्यामेव महातेजाः प्रविवेश महद्वनम्} %||2-92-11||

\twolineshloka
{स तानि द्रुमजालानि जातानि गिरिसानुषु}
{पुष्पिताग्राणि मध्येन जगाम वदतां वरः} %||2-92-12||

\twolineshloka
{स गिरेश्चित्रकूटस्य सालमासाद्य पुष्पितम्}
{रामाश्रमगतस्याग्नेर्ददर्श ध्वजमुच्छ्रितम्} %||2-92-13||

\twolineshloka
{तं दृष्ट्वा भरतः श्रीमान्मुमोद सहबान्धवः}
{अत्र राम इति ज्ञात्वा गतः पारमिवाम्भसः} %||2-92-14||

\fourlineindentedshloka
{स चित्रकूटे तु गिरौ निशाम्य}
{ रामाश्रमं पुण्यजनोपपन्नम्}
{गुहेन सार्धं त्वरितो जगाम}
{ पुनर्निवेश्यैव चमूं महात्मा} %||2-93-15||


{॥इत्यार्षे श्रीमद्रामायणे वाल्मीकीये आदिकाव्ये अयोध्याकाण्डे द्विनवतितमः सर्गः॥}

\sect{त्रिनवतितमः सर्गः}

\twolineshloka
{निविष्टायां तु सेनायामुत्सुको भरतस्तदा}
{जगाम भ्रातरं द्रष्टुं शत्रुघ्नमनुदर्शयन्} %||2-93-1||

\twolineshloka
{ऋषिं वसिष्ठं सन्दिश्य मातॄर्मे शीघ्रमानय}
{इति तरितमग्रे स जागम गुरुवत्सलः} %||2-93-2||

\twolineshloka
{सुमन्त्रस्त्वपि शतुघ्नमदूरादन्वपद्यत}
{रामदार्शनजस्तर्षो भरतस्येव तस्य च} %||2-93-3||

\twolineshloka
{गच्छन्नेवाथ भरतस्तापसालयसंस्थिताम्}
{भ्रातुः पर्णकुटीं श्रीमानुटजं च ददर्श ह} %||2-93-4||

\twolineshloka
{शालायास्त्वग्रतस्तस्या ददर्श भरतस्तदा}
{काष्टानि चावभग्नानि पुष्पाण्यवचितानि च} %||2-93-5||

\twolineshloka
{ददर्श च वने तस्मिन्महतः सञ्चयान्कृतान्}
{मृगाणां महिषाणां च करीषैः शीतकारणात्} %||2-93-6||

\twolineshloka
{गच्छनेव महाबाहुर्द्युतिमान्भरतस्तदा}
{शत्रुघ्नं चाब्रवीद्धृष्टस्तानमात्यांश्च सर्वशः} %||2-93-7||

\twolineshloka
{मन्ये प्राप्ताः स्म तं देशं भरद्वाजो यमब्रवीत्}
{नातिदूरे हि मन्येऽहं नदीं मन्दाकिनीमितः} %||2-93-8||

\twolineshloka
{उच्चैर्बद्धानि चीराणि लक्ष्मणेन भवेदयम्}
{अभिज्ञानकृतः पन्था विकाले गन्तुमिच्छता} %||2-93-9||

\twolineshloka
{इदं चोदात्तदन्तानां कुञ्जराणां तरस्विनाम्}
{शैलपार्श्वे परिक्रान्तमन्योन्यमभिगर्जताम्} %||2-93-10||

\twolineshloka
{यमेवाधातुमिच्छन्ति तापसाः सततं वने}
{तस्यासौ दृश्यते धूमः सङ्कुलः कृष्टवर्त्मनः} %||2-93-11||

\twolineshloka
{अत्राहं पुरुषव्याघ्रं गुरुसत्कारकारिणम्}
{आर्यं द्रक्ष्यामि संहृष्टो महर्षिमिव राघवम्} %||2-93-12||

\twolineshloka
{अथ गत्वा मुहूर्तं तु चित्रकूटं स राघवः}
{मन्दाकिनीमनुप्राप्तस्तं जनं चेदमब्रवीत्} %||2-93-13||

\twolineshloka
{जगत्यां पुरुषव्याघ्र आस्ते वीरासने रतः}
{जनेन्द्रो निर्जनं प्राप्य धिन्मे जन्म सजीवितम्} %||2-93-14||

\twolineshloka
{मत्कृते व्यसनं प्राप्तो लोकनाथो महाद्युतिः}
{सरान्कामान्परित्यज्य वने वसति राघवः} %||2-93-15||

\twolineshloka
{इति लोकसमाक्रुष्टः पादेष्वद्य प्रसादयन्}
{रामस्य निपतिष्यामि सीतायाश्च पुनः पुनः} %||2-93-16||

\twolineshloka
{एवं स विलपंस्तस्मिन्वने दशरथात्मजः}
{ददर्श महतीं पुण्यां पर्णशालां मनोरमाम्} %||2-93-17||

\twolineshloka
{सालतालाश्वकर्णानां पर्णैर्बहुभिरावृताम्}
{विशालां मृदुभिस्तीर्णां कुशैर्वेदिमिवाध्वरे} %||2-93-18||

\twolineshloka
{शक्रायुध निकाशैश्च कार्मुकैर्भारसाधनैः}
{रुक्मपृष्ठैर्महासारैः शोभितां शत्रुबाधकैः} %||2-93-19||

\twolineshloka
{अर्करश्मिप्रतीकाशैर्घोरैस्तूणीगतैः शरैः}
{शोभितां दीप्तवदनैः सर्पैर्भोगवतीमिव} %||2-93-20||

\twolineshloka
{महारजतवासोभ्यामसिभ्यां च विराजिताम्}
{रुक्मबिन्दुविचित्राभ्यां चर्मभ्यां चापि शोभिताम्} %||2-93-21||

\twolineshloka
{गोधाङ्गुलित्रैरासाक्तैश्चित्रैः काञ्चनभूषितैः}
{अरिसङ्घैरनाधृष्यां मृगैः सिंहगुहामिव} %||2-93-22||

\twolineshloka
{प्रागुदक्स्रवणां वेदिं विशालां दीप्तपावकाम्}
{ददर्श भरतस्तत्र पुण्यां रामनिवेशने} %||2-93-23||

\twolineshloka
{निरीक्ष्य स मुहूर्तं तु ददर्श भरतो गुरुम्}
{उटजे राममासीनां जटामण्डलधारिणम्} %||2-93-24||

\twolineshloka
{तं तु कृष्णाजिनधरं चीरवल्कलवाससम्}
{ददर्श राममासीनमभितः पावकोपमम्} %||2-93-25||

\twolineshloka
{सिंहस्कन्धं महाबाहुं पुण्डरीकनिभेक्षणम्}
{पृथिव्याः सगरान्ताया भर्तारं धर्मचारिणम्} %||2-93-26||

\twolineshloka
{उपविष्टं महाबाहुं ब्रह्माणमिव शाश्वतम्}
{स्थण्डिले दर्भसस्म्तीर्णे सीतया लक्ष्मणेन च} %||2-93-27||

\twolineshloka
{तं दृष्ट्वा भरतः श्रीमान्दुःखमोहपरिप्लुतः}
{अभ्यधावत धर्मात्मा भरतः कैकयीसुतः} %||2-93-28||

\twolineshloka
{दृष्ट्वा च विललापार्तो बाष्पसन्दिग्धया गिरा}
{अशक्नुवन्धारयितुं धैर्याद्वचनमब्रवीत्} %||2-93-29||

\twolineshloka
{यः संसदि प्रकृतिभिर्भवेद्युक्त उपासितुम्}
{वन्यैर्मृगैरुपासीनः सोऽयमास्ते ममाग्रजः} %||2-93-30||

\twolineshloka
{वासोभिर्बहुसाहस्रैर्यो महात्मा पुरोचितः}
{मृगाजिने सोऽयमिह प्रवस्ते धर्ममाचरन्} %||2-93-31||

\twolineshloka
{अधारयद्यो विविधाश्चित्राः सुमनसस्तदा}
{सोऽयं जटाभारमिमं सहते राघवः कथम्} %||2-93-32||

\twolineshloka
{यस्य यज्ञैर्यथादिष्टैर्युक्तो धर्मस्य सञ्चयः}
{शरीर क्लेशसम्भूतं स धर्मं परिमार्गते} %||2-93-33||

\twolineshloka
{चन्दनेन महार्हेण यस्याङ्गमुपसेवितम्}
{मलेन तस्याङ्गमिदं कथमार्यस्य सेव्यते} %||2-93-34||

\twolineshloka
{मन्निमित्तमिदं दुःखं प्राप्तो रामः सुखोचितः}
{धिग्जीवितं नृशंसस्य मम लोकविगर्हितम्} %||2-93-35||

\twolineshloka
{इत्येवं विलपन्दीनः प्रस्विन्नमुखपङ्कजः}
{पादावप्राप्य रामस्य पपात भरतो रुदन्} %||2-93-36||

\twolineshloka
{दुःखाभितप्तो भरतो राजपुत्रो महाबलः}
{उक्त्वार्येति सकृद्दीनं पुनर्नोवाच किञ्चन} %||2-93-37||

\twolineshloka
{बाष्पापिहित कण्ठश्च प्रेक्ष्य रामं यशस्विनम्}
{आर्येत्येवाभिसङ्क्रुश्य व्याहर्तुं नाशकत्ततः} %||2-93-38||

\twolineshloka
{शत्रुघ्नश्चापि रामस्य ववन्दे चरणौ रुदन्}
{तावुभौ स समालिङ्ग्य रामोऽप्यश्रूण्यवर्तयत्} %||2-93-39||

\fourlineindentedshloka
{ततः सुमन्त्रेण गुहेन चैव}
{ समीयतू राजसुतावरण्ये}
{दिवाकरश्चैव निशाकरश्च}
{ यथाम्बरे शुक्रबृहस्पतिभ्याम्} %||2-93-40||

\fourlineindentedshloka
{तान्पार्थिवान्वारणयूथपाभा}
{न्समागतांस्तत्र महत्यरण्ये}
{वनौकसस्तेऽपि समीक्ष्य सर्वे}
{ऽप्यश्रूण्यमुञ्चन्प्रविहाय हर्षम्} %||2-94-41||


{॥इत्यार्षे श्रीमद्रामायणे वाल्मीकीये आदिकाव्ये अयोध्याकाण्डे त्रिनवतितमः सर्गः॥}

\sect{चतुर्नवतितमः सर्गः}

\twolineshloka
{आघ्राय रामस्तं मूर्ध्नि परिष्वज्य च राघवः}
{अङ्के भरतमारोप्य पर्यपृच्छत्समाहितः} %||2-94-1||

\twolineshloka
{क्व नु तेऽभूत्पिता तात यदरण्यं त्वमागतः}
{न हि त्वं जीवतस्तस्य वनमागन्तुमर्हसि} %||2-94-2||

\twolineshloka
{चिरस्य बत पश्यामि दूराद्भरतमागतम्}
{दुष्प्रतीकमरण्येऽस्मिन्किं तात वनमागतः} %||2-94-3||

\twolineshloka
{कच्चिद्दशरथो राजा कुशली सत्यसङ्गरः}
{राजसूयाश्वमेधानामाहर्ता धर्मनिश्चयः} %||2-94-4||

\twolineshloka
{स कच्चिद्ब्राह्मणो विद्वान्धर्मनित्यो महाद्युतिः}
{इक्ष्वाकूणामुपाध्यायो यथावत्तात पूज्यते} %||2-94-5||

\twolineshloka
{तात कच्चिच्च कौसल्या सुमित्रा च प्रजावती}
{सुखिनी कच्चिदार्या च देवी नन्दति कैकयी} %||2-94-6||

\twolineshloka
{कच्चिद्विनय सम्पन्नः कुलपुत्रो बहुश्रुतः}
{अनसूयुरनुद्रष्टा सत्कृतस्ते पुरोहितः} %||2-94-7||

\twolineshloka
{कच्चिदग्निषु ते युक्तो विधिज्ञो मतिमानृजुः}
{हुतं च होष्यमाणं च काले वेदयते सदा} %||2-94-8||

\twolineshloka
{इष्वस्त्रवरसम्पन्नमर्थशास्त्रविशारदम्}
{सुधन्वानमुपाध्यायं कच्चित्त्वं तात मन्यसे} %||2-94-9||

\twolineshloka
{कच्चिदात्म समाः शूराः श्रुतवन्तो जितेन्द्रियाः}
{कुलीनाश्चेङ्गितज्ञाश्च कृतास्ते तात मन्त्रिणः} %||2-94-10||

\twolineshloka
{मन्त्रो विजयमूलं हि राज्ञां भवति राघव}
{सुसंवृतो मन्त्रधरैरमात्यैः शास्त्रकोविदैः} %||2-94-11||

\twolineshloka
{कच्चिन्निद्रावशं नैषि कच्चित्काले विबुध्यसे}
{कच्चिंश्चापररात्रिषु चिन्तयस्यर्थनैपुणम्} %||2-94-12||

\twolineshloka
{कच्चिन्मन्त्रयसे नैकः कच्चिन्न बहुभिः सह}
{कच्चित्ते मन्त्रितो मन्त्रो राष्ट्रं न परिधावति} %||2-94-13||

\twolineshloka
{कच्चिदर्थं विनिश्चित्य लघुमूलं महोदयम्}
{क्षिप्रमारभसे कर्तुं न दीर्घयसि राघव} %||2-94-14||

\twolineshloka
{कच्चित्तु सुकृतान्येव कृतरूपाणि वा पुनः}
{विदुस्ते सर्वकार्याणि न कर्तव्यानि पार्थिवाः} %||2-94-15||

\twolineshloka
{कच्चिन्न तर्कैर्युक्त्वा वा ये चाप्यपरिकीर्तिताः}
{त्वया वा तव वामात्यैर्बुध्यते तात मन्त्रितम्} %||2-94-16||

\twolineshloka
{कच्चित्सहस्रान्मूर्खाणामेकमिच्छसि पण्डितम्}
{पण्डितो ह्यर्थकृच्छ्रेषु कुर्यान्निःश्रेयसं महत्} %||2-94-17||

\twolineshloka
{सहस्राण्यपि मूर्खाणां यद्युपास्ते महीपतिः}
{अथ वाप्ययुतान्येव नास्ति तेषु सहायता} %||2-94-18||

\twolineshloka
{एकोऽप्यमात्यो मेधावी शूरो दक्षो विचक्षणः}
{राजानं राजमात्रं वा प्रापयेन्महतीं श्रियम्} %||2-94-19||

\twolineshloka
{कच्चिन्मुख्या महत्स्वेव मध्यमेषु च मध्यमाः}
{जघन्याश्च जघन्येषु भृत्याः कर्मसु योजिताः} %||2-94-20||

\twolineshloka
{अमात्यानुपधातीतान्पितृपैतामहाञ्शुचीन्}
{श्रेष्ठाञ्श्रेष्ठेषु कच्चित्त्वं नियोजयसि कर्मसु} %||2-94-21||

\twolineshloka
{कच्चित्त्वां नावजानन्ति याजकाः पतितं यथा}
{उग्रप्रतिग्रहीतारं कामयानमिव स्त्रियः} %||2-94-22||

\twolineshloka
{उपायकुशलं वैद्यं भृत्यसन्दूषणे रतम्}
{शूरमैश्वर्यकामं च यो न हन्ति स वध्यते} %||2-94-23||

\twolineshloka
{कच्चिद्धृष्टश्च शूरश्च धृतिमान्मतिमाञ्शुचिः}
{कुलीनश्चानुरक्तश्च दक्षः सेनापतिः कृतः} %||2-94-24||

\twolineshloka
{बलवन्तश्च कच्चित्ते मुख्या युद्धविशारदाः}
{दृष्टापदाना विक्रान्तास्त्वया सत्कृत्य मानिताः} %||2-94-25||

\twolineshloka
{कचिद्बलस्य भक्तं च वेतनं च यथोचितम्}
{सम्प्राप्तकालं दातव्यं ददासि न विलम्बसे} %||2-94-26||

\twolineshloka
{कालातिक्रमणे ह्येव भक्त वेतनयोर्भृताः}
{भर्तुः कुप्यन्ति दुष्यन्ति सोऽनर्थः सुमहान्स्मृतः} %||2-94-27||

\twolineshloka
{कच्चित्सर्वेऽनुरक्तास्त्वां कुलपुत्राः प्रधानतः}
{कच्चित्प्राणांस्तवार्थेषु सन्त्यजन्ति समाहिताः} %||2-94-28||

\twolineshloka
{कच्चिज्जानपदो विद्वान्दक्षिणः प्रतिभानवान्}
{यथोक्तवादी दूतस्ते कृतो भरत पण्डितः} %||2-94-29||

\twolineshloka
{कच्चिदष्टादशान्येषु स्वपक्षे दश पञ्च च}
{त्रिभिस्त्रिभिरविज्ञातैर्वेत्सि तीर्थानि चारकैः} %||2-94-30||

\twolineshloka
{कच्चिद्व्यपास्तानहितान्प्रतियातांश्च सर्वदा}
{दुर्बलाननवज्ञाय वर्तसे रिपुसूदन} %||2-94-31||

\twolineshloka
{कच्चिन्न लोकायतिकान्ब्राह्मणांस्तात सेवसे}
{अनर्थ कुशला ह्येते बालाः पण्डितमानिनः} %||2-94-32||

\twolineshloka
{धर्मशास्त्रेषु मुख्येषु विद्यमानेषु दुर्बुधाः}
{बुद्धिमान्वीक्षिकीं प्राप्य निरर्थं प्रवदन्ति ते} %||2-94-33||

\twolineshloka
{वीरैरध्युषितां पूर्वमस्माकं तात पूर्वकैः}
{सत्यनामां दृढद्वारां हस्त्यश्वरथसङ्कुलाम्} %||2-94-34||

\twolineshloka
{ब्राह्मणैः क्षत्रियैर्वैश्यैः स्वकर्मनिरतैः सदा}
{जितेन्द्रियैर्महोत्साहैर्वृतामात्यैः सहस्रशः} %||2-94-35||

\twolineshloka
{प्रासादैर्विविधाकारैर्वृतां वैद्यजनाकुलाम्}
{कच्चित्समुदितां स्फीतामयोध्यां परिरक्षसि} %||2-94-36||

\twolineshloka
{कच्चिच्चैत्यशतैर्जुष्टः सुनिविष्टजनाकुलः}
{देवस्थानैः प्रपाभिश्च तडागैश्चोपशोभितः} %||2-94-37||

\twolineshloka
{प्रहृष्टनरनारीकः समाजोत्सवशोभितः}
{सुकृष्टसीमा पशुमान्हिंसाभिरभिवर्जितः} %||2-94-38||

\twolineshloka
{अदेवमातृको रम्यः श्वापदैः परिवर्जितः}
{कच्चिज्जनपदः स्फीतः सुखं वसति राघव} %||2-94-39||

\twolineshloka
{कच्चित्ते दयिताः सर्वे कृषिगोरक्षजीविनः}
{वार्तायां संश्रितस्तात लोको हि सुखमेधते} %||2-94-40||

\twolineshloka
{तेषां गुप्तिपरीहारैः कच्चित्ते भरणं कृतम्}
{रक्ष्या हि राज्ञा धर्मेण सर्वे विषयवासिनः} %||2-94-41||

\twolineshloka
{कच्चित्स्त्रियः सान्त्वयसि कच्चित्ताश्च सुरक्षिताः}
{कच्चिन्न श्रद्दधास्यासां कच्चिद्गुह्यं न भाषसे} %||2-94-42||

\threelineshloka
{कच्चिन्नाग वनं गुप्तं कुञ्जराणं च तृप्यसि}
{कच्चिद्दर्शयसे नित्यं मनुष्याणां विभूषितम्}
{उत्थायोत्थाय पूर्वाह्णे राजपुत्रो महापथे} %||2-94-43||

\twolineshloka
{कच्चित्सर्वाणि दुर्गाणि धनधान्यायुधोदकैः}
{यन्त्रैश्च परिपूर्णानि तथा शिल्पिधनुर्धरैः} %||2-94-44||

\twolineshloka
{आयस्ते विपुलः कच्चित्कच्चिदल्पतरो व्ययः}
{अपात्रेषु न ते कच्चित्कोशो गच्छति राघव} %||2-94-45||

\twolineshloka
{देवतार्थे च पित्रर्थे ब्राह्मणाभ्यागतेषु च}
{योधेषु मित्रवर्गेषु कच्चिद्गच्छति ते व्ययः} %||2-94-46||

\twolineshloka
{कच्चिदार्यो विशुद्धात्मा क्षारितश्चोरकर्मणा}
{अपृष्टः शास्त्रकुशलैर्न लोभाद्बध्यते शुचिः} %||2-94-47||

\twolineshloka
{गृहीतश्चैव पृष्टश्च काले दृष्टः सकारणः}
{कच्चिन्न मुच्यते चोरो धनलोभान्नरर्षभ} %||2-94-48||

\twolineshloka
{व्यसने कच्चिदाढ्यस्य दुगतस्य च राघव}
{अर्थं विरागाः पश्यन्ति तवामात्या बहुश्रुताः} %||2-94-49||

\twolineshloka
{यानि मिथ्याभिशस्तानां पतन्त्यस्राणि राघव}
{तानि पुत्रपशून्घ्नन्ति प्रीत्यर्थमनुशासतः} %||2-94-50||

\twolineshloka
{कच्चिद्वृधांश्च बालांश्च वैद्यमुख्यांश्च राघव}
{दानेन मनसा वाचा त्रिभिरेतैर्बुभूषसे} %||2-94-51||

\twolineshloka
{कच्चिद्गुरूंश्च वृद्धांश्च तापसान्देवतातिथीन्}
{चैत्यांश्च सर्वान्सिद्धार्थान्ब्राह्मणांश्च नमस्यसि} %||2-94-52||

\twolineshloka
{कच्चिदर्थेन वा धर्मं धर्मं धर्मेण वा पुनः}
{उभौ वा प्रीतिलोभेन कामेन न विबाधसे} %||2-94-53||

\twolineshloka
{कच्चिदर्थं च धर्मं च कामं च जयतां वर}
{विभज्य काले कालज्ञ सर्वान्भरत सेवसे} %||2-94-54||

\twolineshloka
{कच्चित्ते ब्राह्मणाः शर्म सर्वशास्त्रार्थकोविदः}
{आशंसन्ते महाप्राज्ञ पौरजानपदैः सह} %||2-94-55||

\twolineshloka
{नास्तिक्यमनृतं क्रोधं प्रमादं दीर्घसूत्रताम्}
{अदर्शनं ज्ञानवतामालस्यं पञ्चवृत्तिताम्} %||2-94-56||

\twolineshloka
{एकचिन्तनमर्थानामनर्थज्ञैश्च मन्त्रणम्}
{निश्चितानामनारम्भं मन्त्रस्यापरिलक्षणम्} %||2-94-57||

\twolineshloka
{मङ्गलस्याप्रयोगं च प्रत्युत्थानं च सर्वशः}
{कच्चित्त्वं वर्जयस्येतान्राजदोषांश्चतुर्दश} %||2-94-58||

\twolineshloka
{कच्चित्स्वादुकृतं भोज्यमेको नाश्नासि राघव}
{कच्चिदाशंसमानेभ्यो मित्रेभ्यः सम्प्रयच्छसि} %||2-95-59||


{॥इत्यार्षे श्रीमद्रामायणे वाल्मीकीये आदिकाव्ये अयोध्याकाण्डे चतुर्नवतितमः सर्गः॥}

\sect{पञ्चनवतितमः सर्गः}

\twolineshloka
{रामस्य वचनं श्रुत्वा भरतः प्रत्युवाच ह}
{किं मे धर्माद्विहीनस्य राजधर्मः करिष्यति} %||2-95-1||

\twolineshloka
{शाश्वतोऽयं सदा धर्मः स्थितोऽस्मासु नरर्षभ}
{ज्येष्ठ पुत्रे स्थिते राजन्न कनीयान्भवेन्नृपः} %||2-95-2||

\twolineshloka
{स समृद्धां मया सार्धमयोध्यां गच्छ राघव}
{अभिषेचय चात्मानं कुलस्यास्य भवाय नः} %||2-95-3||

\twolineshloka
{राजानं मानुषं प्राहुर्देवत्वे सम्मतो मम}
{यस्य धर्मार्थसहितं वृत्तमाहुरमानुषम्} %||2-95-4||

\twolineshloka
{केकयस्थे च मयि तु त्वयि चारण्यमाश्रिते}
{दिवमार्य गतो राजा यायजूकः सतां मतः} %||2-95-5||

\twolineshloka
{उत्तिष्ठ पुरुषव्याघ्र क्रियतामुदकं पितुः}
{अहं चायं च शत्रुघ्नः पूर्वमेव कृतोदकौ} %||2-95-6||

\twolineshloka
{प्रियेण किल दत्तं हि पितृलोकेषु राघव}
{अक्षय्यं भवतीत्याहुर्भवांश्चैव पितुः प्रियः} %||2-95-7||

\twolineshloka
{तां श्रुत्वा करुणां वाचं पितुर्मरणसंहिताम्}
{राघवो भरतेनोक्तां बभूव गतचेतनः} %||2-95-8||

\threelineshloka
{वाग्वज्रं भरतेनोक्तममनोज्ञं परन्तपः}
{प्रगृह्य बाहू रामो वै पुष्पिताग्रो यथा द्रुमः}
{वने परशुना कृत्तस्तथा भुवि पपात ह} %||2-95-9||

\twolineshloka
{तथा हि पतितं रामं जगत्यां जगतीपतिम्}
{कूलघातपरिश्रान्तं प्रसुप्तमिव कुञ्जरम्} %||2-95-10||

\twolineshloka
{भ्रातरस्ते महेष्वासं सर्वतः शोककर्शितम्}
{रुदन्तः सह वैदेह्या सिषिचुः सलिलेन वै} %||2-95-11||

\twolineshloka
{स तु संज्ञां पुनर्लब्ध्वा नेत्राभ्यामास्रमुत्सृजन्}
{उपाक्रामत काकुत्स्थः कृपणं बहुभाषितुम्} %||2-95-12||

\twolineshloka
{किं नु तस्य मया कार्यं दुर्जातेन महात्मना}
{यो मृतो मम शोकेन न मया चापि संस्कृतः} %||2-95-13||

\twolineshloka
{अहो भरत सिद्धार्थो येन राजा त्वयानघ}
{शत्रुघेण च सर्वेषु प्रेतकृत्येषु सत्कृतः} %||2-95-14||

\twolineshloka
{निष्प्रधानामनेकाग्रं नरेन्द्रेण विनाकृताम्}
{निवृत्तवनवासोऽपि नायोध्यां गन्तुमुत्सहे} %||2-95-15||

\twolineshloka
{समाप्तवनवासं मामयोध्यायां परन्तप}
{को नु शासिष्यति पुनस्ताते लोकान्तरं गते} %||2-95-16||

\twolineshloka
{पुरा प्रेक्ष्य सुवृत्तं मां पिता यान्याह सान्त्वयन्}
{वाक्यानि तानि श्रोष्यामि कुतः कर्णसुखान्यहम्} %||2-95-17||

\twolineshloka
{एवमुक्त्वा स भरतं भार्यामभ्येत्य राघवः}
{उवाच शोकसन्तप्तः पूर्णचन्द्रनिभाननाम्} %||2-95-18||

\twolineshloka
{सीते मृतस्ते श्वशुरः पित्रा हीनोऽसि लक्ष्मण}
{भरतो दुःखमाचष्टे स्वर्गतं पृथिवीपतिम्} %||2-95-19||

\twolineshloka
{सान्त्वयित्वा तु तां रामो रुदन्तीं जनकात्मजाम्}
{उवाच लक्ष्मणं तत्र दुःखितो दुःखितं वचः} %||2-95-20||

\twolineshloka
{आनयेङ्गुदिपिण्याकं चीरमाहर चोत्तरम्}
{जलक्रियार्थं तातस्य गमिष्यामि महात्मनः} %||2-95-21||

\twolineshloka
{सीता पुरस्ताद्व्रजतु त्वमेनामभितो व्रज}
{अहं पश्चाद्गमिष्यामि गतिर्ह्येषा सुदारुणा} %||2-95-22||

\twolineshloka
{ततो नित्यानुगस्तेषां विदितात्मा महामतिः}
{मृदुर्दान्तश्च शान्तश्च रामे च दृढ भक्तिमान्} %||2-95-23||

\twolineshloka
{सुमन्त्रस्तैर्नृपसुतैः सार्धमाश्वास्य राघवम्}
{अवातारयदालम्ब्य नदीं मन्दाकिनीं शिवाम्} %||2-95-24||

\twolineshloka
{ते सुतीर्थां ततः कृच्छ्रादुपागम्य यशस्विनः}
{नदीं मन्दाकिनीं रम्यां सदा पुष्पितकाननाम्} %||2-95-25||

\twolineshloka
{शीघ्रस्रोतसमासाद्य तीर्थं शिवमकर्दमम्}
{सिषिचुस्तूदकं राज्ञे तत एतद्भवत्विति} %||2-95-26||

\twolineshloka
{प्रगृह्य च महीपालो जलपूरितमञ्जलिम्}
{दिशं याम्यामभिमुखो रुदन्वचनमब्रवीत्} %||2-95-27||

\twolineshloka
{एतत्ते राजशार्दूल विमलं तोयमक्षयम्}
{पितृलोकगतस्याद्य मद्दत्तमुपतिष्ठतु} %||2-95-28||

\twolineshloka
{ततो मन्दाकिनी तीरात्प्रत्युत्तीर्य स राघवः}
{पितुश्चकार तेजस्वी निवापं भ्रातृभिः सह} %||2-95-29||

\twolineshloka
{ऐङ्गुदं बदरीमिश्रं पिण्याकं दर्भसंस्तरे}
{न्यस्य रामः सुदुःखार्तो रुदन्वचनमब्रवीत्} %||2-95-30||

\twolineshloka
{इदं भुङ्क्ष्व महाराजप्रीतो यदशना वयम्}
{यदन्नः पुरुषो भवति तदन्नास्तस्य देवताः} %||2-95-31||

\twolineshloka
{ततस्तेनैव मार्गेण प्रत्युत्तीर्य नदीतटात्}
{आरुरोह नरव्याघ्रो रम्यसानुं महीधरम्} %||2-95-32||

\twolineshloka
{ततः पर्णकुटीद्वारमासाद्य जगतीपतिः}
{परिजग्राह पाणिभ्यामुभौ भरतलक्ष्मणौ} %||2-95-33||

\twolineshloka
{तेषां तु रुदतां शब्दात्प्रतिश्रुत्काभवद्गिरौ}
{भ्रातॄणां सह वैदेह्या सिंहानां नर्दतामिव} %||2-95-34||

\threelineshloka
{विज्ञाय तुमुलं शब्दं त्रस्ता भरतसैनिकाः}
{अब्रुवंश्चापि रामेण भरतः सङ्गतो ध्रुवम्}
{तेषामेव महाञ्शब्दः शोचतां पितरं मृतम्} %||2-95-35||

\twolineshloka
{अथ वासान्परित्यज्य तं सर्वेऽभिमुखाः स्वनम्}
{अप्येक मनसो जग्मुर्यथास्थानं प्रधाविताः} %||2-95-36||

\twolineshloka
{हयैरन्ये गजैरन्ये रथैरन्ये स्वलङ्कृतैः}
{सुकुमारास्तथैवान्ये पद्भिरेव नरा ययुः} %||2-95-37||

\twolineshloka
{अचिरप्रोषितं रामं चिरविप्रोषितं यथा}
{द्रष्टुकामो जनः सर्वो जगाम सहसाश्रमम्} %||2-95-38||

\twolineshloka
{भ्रातॄणां त्वरितास्ते तु द्रष्टुकामाः समागमम्}
{ययुर्बहुविधैर्यानैः खुरनेमिसमाकुलैः} %||2-95-39||

\twolineshloka
{सा भूमिर्बहुभिर्यानैः खुरनेमिसमाहता}
{मुमोच तुमुलं शब्दं द्यौरिवाभ्रसमागमे} %||2-95-40||

\twolineshloka
{तेन वित्रासिता नागाः करेणुपरिवारिताः}
{आवासयन्तो गन्धेन जग्मुरन्यद्वनं ततः} %||2-95-41||

\twolineshloka
{वराहमृगसिंहाश्च महिषाः सर्क्षवानराः}
{व्याघ्र गोकर्णगवया वित्रेषुः पृषतैः सह} %||2-95-42||

\twolineshloka
{रथाङ्गसाह्वा नत्यूहा हंसाः कारण्डवाः प्लवाः}
{तथा पुंस्कोकिलाः क्रौञ्चा विसंज्ञा भेजिरे दिशः} %||2-95-43||

\twolineshloka
{तेन शब्देन वित्रस्तैराकाशं पक्षिभिर्वृतम्}
{मनुष्यैरावृता भूमिरुभयं प्रबभौ तदा} %||2-95-44||

\twolineshloka
{तान्नरान्बाष्पपूर्णाक्षान्समीक्ष्याथ सुदुःखितान्}
{पर्यष्वजत धर्मज्ञः पितृवन्मातृवच्च सः} %||2-95-45||

\fourlineindentedshloka
{स तत्र कांश्चित्परिषस्वजे नरा}
{न्नराश्च केचित्तु तमभ्यवादयन्}
{चकार सर्वान्सवयस्यबान्धवा}
{न्यथार्हमासाद्य तदा नृपात्मजः} %||2-95-46||

\fourlineindentedshloka
{ततः स तेषां रुदतां महात्मनाम्}
{ भुवं च खं चानुविनादयन्स्वनः}
{गुहा गिरीणां च दिशश्च सन्ततम्}
{ मृदङ्गघोषप्रतिमो विशुश्रुवे} %||2-96-47||


{॥इत्यार्षे श्रीमद्रामायणे वाल्मीकीये आदिकाव्ये अयोध्याकाण्डे पञ्चनवतितमः सर्गः॥}

\sect{षण्णवतितमः सर्गः}

\twolineshloka
{वसिष्ठः पुरतः कृत्वा दारान्दशरथस्य च}
{अभिचक्राम तं देशं रामदर्शनतर्षितः} %||2-96-1||

\twolineshloka
{राजपत्न्यश्च गच्छन्त्यो मन्दं मन्दाकिनीं प्रति}
{ददृशुस्तत्र तत्तीर्थं रामलक्ष्मणसेवितम्} %||2-96-2||

\twolineshloka
{कौसल्या बाष्पपूर्णेन मुखेन परिशुष्यता}
{सुमित्रामब्रवीद्दीना याश्चान्या राजयोषितः} %||2-96-3||

\twolineshloka
{इदं तेषामनाथानां क्लिष्टमक्लिष्ट कर्मणाम्}
{वने प्राक्केवलं तीर्थं ये ते निर्विषयी कृताः} %||2-96-4||

\twolineshloka
{इतः सुमित्रे पुत्रस्ते सदा जलमतन्द्रितः}
{स्वयं हरति सौमित्रिर्मम पुत्रस्य कारणात्} %||2-96-5||

\twolineshloka
{दक्षिणाग्रेषु दर्भेषु सा ददर्श महीतले}
{पितुरिङ्गुदिपिण्याकं न्यस्तमायतलोचना} %||2-96-6||

\twolineshloka
{तं भूमौ पितुरार्तेन न्यस्तं रामेण वीक्ष्य सा}
{उवाच देवी कौसल्या सर्वा दशरथस्त्रियः} %||2-96-7||

\twolineshloka
{इदमिक्ष्वाकुनाथस्य राघवस्य महात्मनः}
{राघवेण पितुर्दत्तं पश्यतैतद्यथाविधि} %||2-96-8||

\twolineshloka
{तस्य देवसमानस्य पार्थिवस्य महात्मनः}
{नैतदौपयिकं मन्ये भुक्तभोगस्य भोजनम्} %||2-96-9||

\twolineshloka
{चतुरन्तां महीं भुक्त्वा महेन्द्र सदृशो भुवि}
{कथमिङ्गुदिपिण्याकं स भुङ्क्ते वसुधाधिपः} %||2-96-10||

\twolineshloka
{अतो दुःखतरं लोके न किञ्चित्प्रतिभाति मा}
{यत्र रामः पितुर्दद्यादिङ्गुदीक्षोदमृद्धिमान्} %||2-96-11||

\twolineshloka
{रामेणेङ्गुदिपिण्याकं पितुर्दत्तं समीक्ष्य मे}
{कथं दुःखेन हृदयं न स्फोटति सहस्रधा} %||2-96-12||

\twolineshloka
{एवमार्तां सपत्न्यस्ता जग्मुराश्वास्य तां तदा}
{ददृशुश्चाश्रमे रामं स्वर्गाच्च्युतमिवामरम्} %||2-96-13||

\twolineshloka
{सर्वभोगैः परित्यक्तं राम सम्प्रेक्ष्य मातरः}
{आर्ता मुमुचुरश्रूणि सस्वरं शोककर्शिताः} %||2-96-14||

\twolineshloka
{तासां रामः समुत्थाय जग्राह चरणाञ्शुभान्}
{मातॄणां मनुजव्याघ्रः सर्वासां सत्यसङ्गरः} %||2-96-15||

\twolineshloka
{ताः पाणिभिः सुखस्पर्शैर्मृद्वङ्गुलितलैः शुभैः}
{प्रममार्जू रजः पृष्ठाद्रामस्यायतलोचनाः} %||2-96-16||

\twolineshloka
{सौमित्रिरपि ताः सर्वा मातॄह्सम्प्रेक्ष्य दुःखितः}
{अभ्यवादयतासक्तं शनै रामादनन्तरम्} %||2-96-17||

\twolineshloka
{यथा रामे तथा तस्मिन्सर्वा ववृतिरे स्त्रियः}
{वृत्तिं दशरथाज्जाते लक्ष्मणे शुभलक्षणे} %||2-96-18||

\twolineshloka
{सीतापि चरणांस्तासामुपसङ्गृह्य दुःखिता}
{श्वश्रूणामश्रुपूर्णाक्षी सा बभूवाग्रतः स्थिता} %||2-96-19||

\twolineshloka
{तां परिष्वज्य दुःखार्तां माता दुहितरं यथा}
{वनवासकृशां दीनां कौसल्या वाक्यमब्रवीत्} %||2-96-20||

\twolineshloka
{विदेहराजस्य सुता स्नुषा दशरथस्य च}
{रामपत्नी कथं दुःखं सम्प्राप्ता निर्जने वने} %||2-96-21||

\twolineshloka
{पद्ममातपसन्तप्तं परिक्लिष्टमिवोत्पलम्}
{काञ्चनं रजसा ध्वस्तं क्लिष्टं चन्द्रमिवाम्बुदैः} %||2-96-22||

\twolineshloka
{मुखं ते प्रेक्ष्य मां शोको दहत्यग्निरिवाश्रयम्}
{भृशं मनसि वैदेहि व्यसनारणिसम्भवः} %||2-96-23||

\twolineshloka
{ब्रुवन्त्यामेवमार्तायां जनन्यां भरताग्रजः}
{पादावासाद्य जग्राह वसिष्ठस्य स राघवः} %||2-96-24||

\fourlineindentedshloka
{पुरोहितस्याग्निसमस्य तस्य वै}
{ बृहस्पतेरिन्द्र इवामराधिपः}
{प्रगृह्य पादौ सुसमृद्धतेजसः}
{ सहैव तेनोपविवेश राघवः} %||2-96-25||

\fourlineindentedshloka
{ततो जघन्यं सहितैः स मन्त्रिभिः}
{ पुरप्रधानैश्च सहैव सैनिकैः}
{जनेन धर्मज्ञतमेन धर्मवा}
{नुपोपविष्टो भरतस्तदाग्रजम्} %||2-96-26||

\fourlineindentedshloka
{उपोपविष्टस्तु तदा स वीर्यवाम्}
{स्तपस्विवेषेण समीक्ष्य राघवम्}
{श्रिया ज्वलन्तं भरतः कृताञ्जलि}
{र्यथा महेन्द्रः प्रयतः प्रजापतिम्} %||2-96-27||

\fourlineindentedshloka
{किमेष वाक्यं भरतोऽद्य राघवम्}
{ प्रणम्य सत्कृत्य च साधु वक्ष्यति}
{इतीव तस्यार्यजनस्य तत्त्वतो}
{ बभूव कौतूहलमुत्तमं तदा} %||2-96-28||

\fourlineindentedshloka
{स राघवः सत्यधृतिश्च लक्ष्मणो}
{ महानुभावो भरतश्च धार्मिकः}
{वृताः सुहृद्भिश्च विरेजुरध्वरे}
{ यथा सदस्यैः सहितास्त्रयोऽग्नयः} %||2-97-29||


{॥इत्यार्षे श्रीमद्रामायणे वाल्मीकीये आदिकाव्ये अयोध्याकाण्डे षण्णवतितमः सर्गः॥}

\sect{सप्तनवतितमः सर्गः}

\twolineshloka
{तं तु रामः समाश्वास्य भ्रातरं गुरुवत्सलम्}
{लक्ष्मणेन सह भ्रात्रा प्रष्टुं समुपचक्रमे} %||2-97-1||

\twolineshloka
{किमेतदिच्छेयमहं श्रोतुं प्रव्याहृतं त्वया}
{यस्मात्त्वमागतो देशमिमं चीरजटाजिनी} %||2-97-2||

\twolineshloka
{यन्निमित्तमिमं देशं कृष्णाजिनजटाधरः}
{हित्वा राज्यं प्रविष्टस्त्वं तत्सर्वं वक्तुमर्हसि} %||2-97-3||

\twolineshloka
{इत्युक्तः केकयीपुत्रः काकुत्स्थेन महात्मना}
{प्रगृह्य बलवद्भूयः प्राञ्जलिर्वाक्यमब्रवीत्} %||2-97-4||

\twolineshloka
{आर्यं तातः परित्यज्य कृत्वा कर्म सुदुष्करम्}
{गतः स्वर्गं महाबाहुः पुत्रशोकाभिपीडितः} %||2-97-5||

\twolineshloka
{स्त्रिया नियुक्तः कैकेय्या मम मात्रा परन्तप}
{चकार सुमहत्पापमिदमात्मयशोहरम्} %||2-97-6||

\twolineshloka
{सा राज्यफलमप्राप्य विधवा शोककर्शिता}
{पतिष्यति महाघोरे निरये जननी मम} %||2-97-7||

\twolineshloka
{तस्य मे दासभूतस्य प्रसादं कर्तुमर्हसि}
{अभिषिञ्चस्व चाद्यैव राज्येन मघवानिव} %||2-97-8||

\twolineshloka
{इमाः प्रकृतयः सर्वा विधवा मातुरश्च याः}
{त्वत्सकाशमनुप्राप्ताः प्रसादं कर्तुमर्हसि} %||2-97-9||

\twolineshloka
{तदानुपूर्व्या युक्तं च युक्तं चात्मनि मानद}
{राज्यं प्राप्नुहि धर्मेण सकामान्सुहृदः कुरु} %||2-97-10||

\twolineshloka
{भवत्वविधवा भूमिः समग्रा पतिना त्वया}
{शशिना विमलेनेव शारदी रजनी यथा} %||2-97-11||

\twolineshloka
{एभिश्च सचिवैः सार्धं शिरसा याचितो मया}
{भ्रातुः शिष्यस्य दासस्य प्रसादं कर्तुमर्हसि} %||2-97-12||

\twolineshloka
{तदिदं शाश्वतं पित्र्यं सर्वं सचिवमण्डलम्}
{पूजितं पुरुषव्याघ्र नातिक्रमितुमुत्सहे} %||2-97-13||

\twolineshloka
{एवमुक्त्वा महाबाहुः सबाष्पः केकयीसुतः}
{रामस्य शिरसा पादौ जग्राह भरतः पुनः} %||2-97-14||

\twolineshloka
{तं मत्तमिव मातङ्गं निःश्वसन्तं पुनः पुनः}
{भ्रातरं भरतं रामः परिष्वज्येदमब्रवीत्} %||2-97-15||

\twolineshloka
{कुलीनः सत्त्वसम्पन्नस्तेजस्वी चरितव्रतः}
{राज्यहेतोः कथं पापमाचरेत्त्वद्विधो जनः} %||2-97-16||

\twolineshloka
{न दोषं त्वयि पश्यामि सूक्ष्ममप्यरि सूदन}
{न चापि जननीं बाल्यात्त्वं विगर्हितुमर्हसि} %||2-97-17||

\twolineshloka
{यावत्पितरि धर्मज्ञ गौरवं लोकसत्कृते}
{तावद्धर्मभृतां श्रेष्ठ जनन्यामपि गौरवम्} %||2-97-18||

\twolineshloka
{एताभ्यां धर्मशीलाभ्यां वनं गच्छेति राघव}
{माता पितृभ्यामुक्तोऽहं कथमन्यत्समाचरे} %||2-97-19||

\twolineshloka
{त्वया राज्यमयोध्यायां प्राप्तव्यं लोकसत्कृतम्}
{वस्तव्यं दण्डकारण्ये मया वल्कलवाससा} %||2-97-20||

\twolineshloka
{एवं कृत्वा महाराजो विभागं लोकसंनिधौ}
{व्यादिश्य च महातेजा दिवं दशरथो गतः} %||2-97-21||

\twolineshloka
{स च प्रमाणं धर्मात्मा राजा लोकगुरुस्तव}
{पित्रा दत्तं यथाभागमुपभोक्तुं त्वमर्हसि} %||2-97-22||

\twolineshloka
{चतुर्दश समाः सौम्य दण्डकारण्यमाश्रितः}
{उपभोक्ष्ये त्वहं दत्तं भागं पित्रा महात्मना} %||2-97-23||

\fourlineindentedshloka
{यदब्रवीन्मां नरलोकसत्कृतः}
{ पिता महात्मा विबुधाधिपोपमः}
{तदेव मन्ये परमात्मनो हितम्}
{ न सर्वलोकेश्वरभावमव्ययम्} %||2-98-24||


{॥इत्यार्षे श्रीमद्रामायणे वाल्मीकीये आदिकाव्ये अयोध्याकाण्डे सप्तनवतितमः सर्गः॥}

\sect{अष्टनवतितमः सर्गः}

\twolineshloka
{ततः पुरुषसिंहानां वृतानां तैः सुहृद्गणैः}
{शोचतामेव रजनी दुःखेन व्यत्यवर्तत} %||2-98-1||

\twolineshloka
{रजन्यां सुप्रभातायां भ्रातरस्ते सुहृद्वृताः}
{मन्दाकिन्यां हुतं जप्यं कृत्वा राममुपागमन्} %||2-98-2||

\twolineshloka
{तूष्णीं ते समुपासीना न कश्चित्किञ्चिदब्रवीत्}
{भरतस्तु सुहृन्मध्ये रामवचनमब्रवीत्} %||2-98-3||

\twolineshloka
{सान्त्विता मामिका माता दत्तं राज्यमिदं मम}
{तद्ददामि तवैवाहं भुङ्क्ष्व राज्यमकण्टकम्} %||2-98-4||

\twolineshloka
{महतेवाम्बुवेगेन भिन्नः सेतुर्जलागमे}
{दुरावारं त्वदन्येन राज्यखण्डमिदं महत्} %||2-98-5||

\twolineshloka
{गतिं खर इवाश्वस्य तार्क्ष्यस्येव पतत्रिणः}
{अनुगन्तुं न शक्तिर्मे गतिं तव महीपते} %||2-98-6||

\twolineshloka
{सुजीवं नित्यशस्तस्य यः परैरुपजीव्यते}
{राम तेन तु दुर्जीवं यः परानुपजीवति} %||2-98-7||

\twolineshloka
{यथा तु रोपितो वृक्षः पुरुषेण विवर्धितः}
{ह्रस्वकेन दुरारोहो रूढस्कन्धो महाद्रुमः} %||2-98-8||

\twolineshloka
{स यदा पुष्पितो भूत्वा फलानि न विदर्शयेत्}
{स तां नानुभवेत्प्रीतिं यस्य हेतोः प्रभावितः} %||2-98-9||

\twolineshloka
{एषोपमा महाबाहो त्वमर्थं वेत्तुमर्हसि}
{यदि त्वमस्मानृषभो भर्ता भृत्यान्न शाधि हि} %||2-98-10||

\twolineshloka
{श्रेणयस्त्वां महाराज पश्यन्त्वग्र्याश्च सर्वशः}
{प्रतपन्तमिवादित्यं राज्ये स्थितमरिन्दमम्} %||2-98-11||

\twolineshloka
{तवानुयाने काकुत्ष्ठ मत्ता नर्दन्तु कुञ्जराः}
{अन्तःपुर गता नार्यो नन्दन्तु सुसमाहिताः} %||2-98-12||

\twolineshloka
{तस्य साध्वित्यमन्यन्त नागरा विविधा जनाः}
{भरतस्य वचः श्रुत्वा रामं प्रत्यनुयाचतः} %||2-98-13||

\twolineshloka
{तमेवं दुःखितं प्रेक्ष्य विलपन्तं यशस्विनम्}
{रामः कृतात्मा भरतं समाश्वासयदात्मवान्} %||2-98-14||

\twolineshloka
{नात्मनः कामकारोऽस्ति पुरुषोऽयमनीश्वरः}
{इतश्चेतरतश्चैनं कृतान्तः परिकर्षति} %||2-98-15||

\twolineshloka
{सर्वे क्षयान्ता निचयाः पतनान्ताः समुच्छ्रयाः}
{संयोगा विप्रयोगान्ता मरणान्तं च जीवितम्} %||2-98-16||

\twolineshloka
{यथा फलानं पक्वानां नान्यत्र पतनाद्भयम्}
{एवं नरस्य जातस्य नान्यत्र मरणाद्भयम्} %||2-98-17||

\twolineshloka
{यथागारं दृढस्थूणं जीर्णं भूत्वावसीदति}
{तथावसीदन्ति नरा जरामृत्युवशं गताः} %||2-98-18||

\twolineshloka
{अहोरात्राणि गच्छन्ति सर्वेषां प्राणिनामिह}
{आयूंषि क्षपयन्त्याशु ग्रीष्मे जलमिवांशवः} %||2-98-19||

\twolineshloka
{आत्मानमनुशोच त्वं किमन्यमनुशोचसि}
{आयुस्ते हीयते यस्य स्थितस्य च गतस्य च} %||2-98-20||

\twolineshloka
{सहैव मृत्युर्व्रजति सह मृत्युर्निषीदति}
{गत्वा सुदीर्घमध्वानं सह मृत्युर्निवर्तते} %||2-98-21||

\twolineshloka
{गात्रेषु वलयः प्राप्ताः श्वेताश्चैव शिरोरुहाः}
{जरया पुरुषो जीर्णः किं हि कृत्वा प्रभावयेत्} %||2-98-22||

\twolineshloka
{नन्दन्त्युदित आदित्ये नन्दन्त्यस्तमिते रवौ}
{आत्मनो नावबुध्यन्ते मनुष्या जीवितक्षयम्} %||2-98-23||

\twolineshloka
{हृष्यन्त्यृतुमुखं दृष्ट्वा नवं नवमिहागतम्}
{ऋतूनां परिवर्तेन प्राणिनां प्राणसङ्क्षयः} %||2-98-24||

\twolineshloka
{यथा काष्ठं च काष्ठं च समेयातां महार्णवे}
{समेत्य च व्यपेयातां कालमासाद्य कञ्चन} %||2-98-25||

\twolineshloka
{एवं भार्याश्च पुत्राश्च ज्ञातयश्च वसूनि च}
{समेत्य व्यवधावन्ति ध्रुवो ह्येषां विनाभवः} %||2-98-26||

\twolineshloka
{नात्र कश्चिद्यथा भावं प्राणी समभिवर्तते}
{तेन तस्मिन्न सामर्थ्यं प्रेतस्यास्त्यनुशोचतः} %||2-98-27||

\twolineshloka
{यथा हि सार्थं गच्छन्तं ब्रूयात्कश्चित्पथि स्थितः}
{अहमप्यागमिष्यामि पृष्ठतो भवतामिति} %||2-98-28||

\twolineshloka
{एवं पूर्वैर्गतो मार्गः पितृपैतामहो ध्रुवः}
{तमापन्नः कथं शोचेद्यस्य नास्ति व्यतिक्रमः} %||2-98-29||

\twolineshloka
{वयसः पतमानस्य स्रोतसो वानिवर्तिनः}
{आत्मा सुखे नियोक्तव्यः सुखभाजः प्रजाः स्मृताः} %||2-98-30||

\twolineshloka
{धर्मात्मा स शुभैः कृत्स्नैः क्रतुभिश्चाप्तदक्षिणैः}
{धूतपापो गतः स्वर्गं पिता नः पृथिवीपतिः} %||2-98-31||

\twolineshloka
{भृत्यानां भरणात्सम्यक्प्रजानां परिपालनात्}
{अर्थादानाच्च धार्मेण पिता नस्त्रिदिवं गतः} %||2-98-32||

\twolineshloka
{इष्ट्वा बहुविधैर्यज्ञैर्भोगांश्चावाप्य पुष्कलान्}
{उत्तमं चायुरासाद्य स्वर्गतः पृथिवीपतिः} %||2-98-33||

\twolineshloka
{स जीर्णं मानुषं देहं परित्यज्य पिता हि नः}
{दैवीमृद्धिमनुप्राप्तो ब्रह्मलोकविहारिणीम्} %||2-98-34||

\twolineshloka
{तं तु नैवं विधः कश्चित्प्राज्ञः शोचितुमर्हति}
{त्वद्विधो यद्विधश्चापि श्रुतवान्बुद्धिमत्तरः} %||2-98-35||

\twolineshloka
{एते बहुविधाः शोका विलाप रुदिते तथा}
{वर्जनीया हि धीरेण सर्वावस्थासु धीमता} %||2-98-36||

\twolineshloka
{स स्वस्थो भव मा शोचो यात्वा चावस तां पुरीम्}
{तथा पित्रा नियुक्तोऽसि वशिना वदताम्व्वर} %||2-98-37||

\twolineshloka
{यत्राहमपि तेनैव नियुक्तः पुण्यकर्मणा}
{तत्रैवाहं करिष्यामि पितुरार्यस्य शासनम्} %||2-98-38||

\twolineshloka
{न मया शासनं तस्य त्यक्तुं न्याय्यमरिन्दम}
{तत्त्वयापि सदा मान्यं स वै बन्धुः स नः पिता} %||2-98-39||

\twolineshloka
{एवमुक्त्वा तु विरते रामे वचनमर्थवत्}
{उवाच भरतश्चित्रं धार्मिको धार्मिकं वचः} %||2-98-40||

\twolineshloka
{को हि स्यादीदृशो लोके यादृशस्त्वमरिन्दम}
{न त्वां प्रव्यथयेद्दुःखं प्रीतिर्वा न प्रहर्षयेत्} %||2-98-41||

\twolineshloka
{सम्मतश्चासि वृद्धानां तांश्च पृच्छसि संशयान्}
{यथा मृतस्तथा जीवन्यथासति तथा सति} %||2-98-42||

\twolineshloka
{यस्यैष बुद्धिलाभः स्यात्परितप्येत केन सः}
{स एवं व्यसनं प्राप्य न विषीदितुमर्हति} %||2-98-43||

\twolineshloka
{अमरोपमसत्त्वस्त्वं महात्मा सत्यसङ्गरः}
{सर्वज्ञः सर्वदर्शी च बुद्धिमांश्चासि राघव} %||2-98-44||

\twolineshloka
{न त्वामेवं गुणैर्युक्तं प्रभवाभवकोविदम्}
{अविषह्यतमं दुःखमासादयितुमर्हति} %||2-98-45||

\twolineshloka
{प्रोषिते मयि यत्पापं मात्रा मत्कारणात्कृतम्}
{क्षुद्रया तदनिष्टं मे प्रसीदतु भवान्मम} %||2-98-46||

\twolineshloka
{धर्मबन्धेन बद्धोऽस्मि तेनेमां नेह मातरम्}
{हन्मि तीव्रेण दण्डेन दण्डार्हां पापकारिणीम्} %||2-98-47||

\twolineshloka
{कथं दशरथाज्जातः शुद्धाभिजनकर्मणः}
{जानन्धर्ममधर्मिष्ठं कुर्यां कर्म जुगुप्सितम्} %||2-98-48||

\twolineshloka
{गुरुः क्रियावान्वृद्धश्च राजा प्रेतः पितेति च}
{तातं न परिगर्हेयं दैवतं चेति संसदि} %||2-98-49||

\twolineshloka
{को हि धर्मार्थयोर्हीनमीदृशं कर्म किल्बिषम्}
{स्त्रियाः प्रियचिकीर्षुः सन्कुर्याद्धर्मज्ञ धर्मवित्} %||2-98-50||

\twolineshloka
{अन्तकाले हि भूतानि मुह्यन्तीति पुराश्रुतिः}
{राज्ञैवं कुर्वता लोके प्रत्यक्षा सा श्रुतिः कृता} %||2-98-51||

\twolineshloka
{साध्वर्थमभिसन्धाय क्रोधान्मोहाच्च साहसात्}
{तातस्य यदतिक्रान्तं प्रत्याहरतु तद्भवान्} %||2-98-52||

\twolineshloka
{पितुर्हि समतिक्रान्तं पुत्रो यः साधु मन्यते}
{तदपत्यं मतं लोके विपरीतमतोऽन्यथा} %||2-98-53||

\twolineshloka
{तदपत्यं भवानस्तु मा भवान्दुष्कृतं पितुः}
{अभिपत्तत्कृतं कर्म लोके धीरविगर्हितम्} %||2-98-54||

\twolineshloka
{कैकेयीं मां च तातं च सुहृदो बान्धवांश्च नः}
{पौरजानपदान्सर्वांस्त्रातु सर्वमिदं भवान्} %||2-98-55||

\twolineshloka
{क्व चारण्यं क्व च क्षात्रं क्व जटाः क्व च पालनम्}
{ईदृशं व्याहतं कर्म न भवान्कर्तुमर्हति} %||2-98-56||

\twolineshloka
{अथ क्लेशजमेव त्वं धर्मं चरितुमिच्छसि}
{धर्मेण चतुरो वर्णान्पालयन्क्लेशमाप्नुहि} %||2-98-57||

\twolineshloka
{चतुर्णामाश्रमाणां हि गार्हस्थ्यं श्रेष्ठमाश्रमम्}
{आहुर्धर्मज्ञ धर्मज्ञास्तं कथं त्यक्तुमर्हसि} %||2-98-58||

\twolineshloka
{श्रुतेन बालः स्थानेन जन्मना भवतो ह्यहम्}
{स कथं पालयिष्यामि भूमिं भवति तिष्ठति} %||2-98-59||

\twolineshloka
{हीनबुद्धिगुणो बालो हीनः स्थानेन चाप्यहम्}
{भवता च विना भूतो न वर्तयितुमुत्सहे} %||2-98-60||

\twolineshloka
{इदं निखिलमव्यग्रं पित्र्यं राज्यमकण्टकम्}
{अनुशाधि स्वधर्मेण धर्मज्ञ सह बान्धवैः} %||2-98-61||

\twolineshloka
{इहैव त्वाभिषिञ्चन्तु धर्मज्ञ सह बान्धवैः}
{ऋत्विजः सवसिष्ठाश्च मन्त्रवन्मन्त्रकोविदाः} %||2-98-62||

\twolineshloka
{अभिषिक्तस्त्वमस्माभिरयोध्यां पालने व्रज}
{विजित्य तरसा लोकान्मरुद्भिरिव वासवः} %||2-98-63||

\twolineshloka
{ऋणानि त्रीण्यपाकुर्वन्दुर्हृदः साधु निर्दहन्}
{सुहृदस्तर्पयन्कामैस्त्वमेवात्रानुशाधि माम्} %||2-98-64||

\twolineshloka
{अद्यार्य मुदिताः सन्तु सुहृदस्तेऽभिषेचने}
{अद्य भीताः पालयन्तां दुर्हृदस्ते दिशो दश} %||2-98-65||

\twolineshloka
{आक्रोशं मम मातुश्च प्रमृज्य पुरुषर्षभ}
{अद्य तत्र भवन्तं च पितरं रक्ष किल्बिषात्} %||2-98-66||

\twolineshloka
{शिरसा त्वाभियाचेऽहं कुरुष्व करुणां मयि}
{बान्धवेषु च सर्वेषु भूतेष्विव महेश्वरः} %||2-98-67||

\twolineshloka
{अथ वा पृष्ठतः कृत्वा वनमेव भवानितः}
{गमिष्यति गमिष्यामि भवता सार्धमप्यहम्} %||2-98-68||

\fourlineindentedshloka
{तथापि रामो भरतेन ताम्यत}
{ प्रसाद्यमानः शिरसा महीपतिः}
{न चैव चक्रे गमनाय सत्त्ववा}
{न्मतिं पितुस्तद्वचने प्रतिष्ठितः} %||2-98-69||

\fourlineindentedshloka
{तदद्भुतं स्थैर्यमवेक्ष्य राघवे}
{ समं जनो हर्षमवाप दुःखितः}
{न यात्ययोध्यामिति दुःखितोऽभव}
{त्स्थिरप्रतिज्ञत्वमवेक्ष्य हर्षितः} %||2-98-70||

\fourlineindentedshloka
{तमृत्विजो नैगमयूथवल्लभा}
{स्तथा विसंज्ञाश्रुकलाश्च मातरः}
{तथा ब्रुवाणं भरतं प्रतुष्टुवुः}
{ प्रणम्य रामं च ययाचिरे सह} %||2-99-71||


{॥इत्यार्षे श्रीमद्रामायणे वाल्मीकीये आदिकाव्ये अयोध्याकाण्डे अष्टनवतितमः सर्गः॥}

\sect{एकोनशततमः सर्गः}

\twolineshloka
{पुनरेवं ब्रुवाणं तु भरतं लक्ष्मणाग्रजः}
{प्रत्युवच ततः श्रीमाञ्ज्ञातिमध्येऽतिसत्कृतः} %||2-99-1||

\twolineshloka
{उपपन्नमिदं वाक्यं यत्त्वमेवमभाषथाः}
{जातः पुत्रो दशरथात्कैकेय्यां राजसत्तमात्} %||2-99-2||

\twolineshloka
{पुरा भ्रातः पिता नः स मातरं ते समुद्वहन्}
{मातामहे समाश्रौषीद्राज्यशुल्कमनुत्तमम्} %||2-99-3||

\twolineshloka
{देवासुरे च सङ्ग्रामे जनन्यै तव पार्थिवः}
{सम्प्रहृष्टो ददौ राजा वरमाराधितः प्रभुः} %||2-99-4||

\twolineshloka
{ततः सा सम्प्रतिश्राव्य तव माता यशस्विनी}
{अयाचत नरश्रेष्ठं द्वौ वरौ वरवर्णिनी} %||2-99-5||

\twolineshloka
{तव राज्यं नरव्याघ्र मम प्रव्राजनं तथा}
{तच्च राजा तथा तस्यै नियुक्तः प्रददौ वरम्} %||2-99-6||

\twolineshloka
{तेन पित्राहमप्यत्र नियुक्तः पुरुषर्षभ}
{चतुर्दश वने वासं वर्षाणि वरदानिकम्} %||2-99-7||

\twolineshloka
{सोऽहं वनमिदं प्राप्तो निर्जनं लक्ष्मणान्वितः}
{शीतया चाप्रतिद्वन्द्वः सत्यवादे स्थितः पितुः} %||2-99-8||

\twolineshloka
{भवानपि तथेत्येव पितरं सत्यवादिनम्}
{कर्तुमर्हति राजेन्द्रं क्षिप्रमेवाभिषेचनात्} %||2-99-9||

\twolineshloka
{ऋणान्मोचय राजानं मत्कृते भरत प्रभुम्}
{पितरं त्राहि धर्मज्ञ मातरं चाभिनन्दय} %||2-99-10||

\twolineshloka
{श्रूयते हि पुरा तात श्रुतिर्गीता यशस्विनी}
{गयेन यजमानेन गयेष्वेव पितॄन्प्रति} %||2-99-11||

\twolineshloka
{पुं नाम्ना नरकाद्यस्मात्पितरं त्रायते सुतः}
{तस्मात्पुत्र इति प्रोक्तः पितॄन्यत्पाति वा सुतः} %||2-99-12||

\twolineshloka
{एष्टव्या बहवः पुत्रा गुणवन्तो बहुश्रुताः}
{तेषां वै समवेतानामपि कश्चिद्गयां व्रजेत्} %||2-99-13||

\twolineshloka
{एवं राजर्षयः सर्वे प्रतीता राजनन्दन}
{तस्मात्त्राहि नरश्रेष्ठ पितरं नरकात्प्रभो} %||2-99-14||

\twolineshloka
{अयोध्यां गच्छ भरत प्रकृतीरनुरञ्जय}
{शत्रुघ्न सहितो वीर सह सर्वैर्द्विजातिभिः} %||2-99-15||

\twolineshloka
{प्रवेक्ष्ये दण्डकारण्यमहमप्यविलम्बयन्}
{आभ्यां तु सहितो राजन्वैदेह्या लक्ष्मणेन च} %||2-99-16||

\fourlineindentedshloka
{त्वं राजा भव भरत स्वयं नराणाम्}
{ वन्यानामहमपि राजराण्मृगाणाम्}
{गच्छ त्वं पुरवरमद्य सम्प्रहृष्टः}
{ संहृष्टस्त्वहमपि दण्डकान्प्रवेक्ष्ये} %||2-99-17||

\fourlineindentedshloka
{छायां ते दिनकरभाः प्रबाधमानम्}
{ वर्षत्रं भरत करोतु मूर्ध्नि शीताम्}
{एतेषामहमपि काननद्रुमाणाम्}
{ छायां तामतिशयिनीं सुखं श्रयिष्ये} %||2-99-18||

\fourlineindentedshloka
{शत्रुघ्नः कुशलमतिस्तु ते सहायः}
{ सौमित्रिर्मम विदितः प्रधानमित्रम्}
{चत्वारस्तनयवरा वयं नरेन्द्रम्}
{ सत्यस्थं भरत चराम मा विषादम्} %||2-100-19||


{॥इत्यार्षे श्रीमद्रामायणे वाल्मीकीये आदिकाव्ये अयोध्याकाण्डे एकोनशततमः सर्गः॥}

\sect{शततमः सर्गः}

\twolineshloka
{आश्वासयन्तं भरतं जाबालिर्ब्राह्मणोत्तमः}
{उवाच रामं धर्मज्ञं धर्मापेतमिदं वचः} %||2-100-1||

\twolineshloka
{साधु राघव मा भूत्ते बुद्धिरेवं निरर्थका}
{प्राकृतस्य नरस्येव आर्य बुद्धेस्तपस्विनः} %||2-100-2||

\twolineshloka
{कः कस्य पुरुषो बन्धुः किमाप्यं कस्य केनचित्}
{यदेको जायते जन्तुरेक एव विनश्यति} %||2-100-3||

\twolineshloka
{तस्मान्माता पिता चेति राम सज्जेत यो नरः}
{उन्मत्त इव स ज्ञेयो नास्ति काचिद्धि कस्यचित्} %||2-100-4||

\twolineshloka
{यथा ग्रामान्तरं गच्छन्नरः कश्चित्क्वचिद्वसेत्}
{उत्सृज्य च तमावासं प्रतिष्ठेतापरेऽहनि} %||2-100-5||

\twolineshloka
{एवमेव मनुष्याणां पिता माता गृहं वसु}
{आवासमात्रं काकुत्स्थ सज्जन्ते नात्र सज्जनाः} %||2-100-6||

\twolineshloka
{पित्र्यं राज्यं समुत्सृज्य स नार्हति नरोत्तम}
{आस्थातुं कापथं दुःखं विषमं बहुकण्टकम्} %||2-100-7||

\twolineshloka
{समृद्धायामयोध्यायामात्मानमभिषेचय}
{एकवेणीधरा हि त्वां नगरी सम्प्रतीक्षते} %||2-100-8||

\twolineshloka
{राजभोगाननुभवन्महार्हान्पार्थिवात्मज}
{विहर त्वमयोध्यायां यथा शक्रस्त्रिविष्टपे} %||2-100-9||

\twolineshloka
{न ते कश्चिद्दशरतःस्त्वं च तस्य न कश्चन}
{अन्यो राजा त्वमन्यश्च तस्मात्कुरु यदुच्यते} %||2-100-10||

\twolineshloka
{गतः स नृपतिस्तत्र गन्तव्यं यत्र तेन वै}
{प्रवृत्तिरेषा मर्त्यानां त्वं तु मिथ्या विहन्यसे} %||2-100-11||

\twolineshloka
{अर्थधर्मपरा ये ये तांस्ताञ्शोचामि नेतरान्}
{ते हि दुःखमिह प्राप्य विनाशं प्रेत्य भेजिरे} %||2-100-12||

\twolineshloka
{अष्टका पितृदैवत्यमित्ययं प्रसृतो जनः}
{अन्नस्योपद्रवं पश्य मृतो हि किमशिष्यति} %||2-100-13||

\twolineshloka
{यदि भुक्तमिहान्येन देहमन्यस्य गच्छति}
{दद्यात्प्रवसतः श्राद्धं न तत्पथ्यशनं भवेत्} %||2-100-14||

\twolineshloka
{दानसंवनना ह्येते ग्रन्था मेधाविभिः कृताः}
{यजस्व देहि दीक्षस्व तपस्तप्यस्व सन्त्यज} %||2-100-15||

\twolineshloka
{स नास्ति परमित्येव कुरु बुद्धिं महामते}
{प्रत्यक्षं यत्तदातिष्ठ परोक्षं पृष्ठतः कुरु} %||2-100-16||

\twolineshloka
{सतां बुद्धिं पुरस्कृत्य सर्वलोकनिदर्शिनीम्}
{राज्यं त्वं प्रतिगृह्णीष्व भरतेन प्रसादितः} %||2-101-17||


{॥इत्यार्षे श्रीमद्रामायणे वाल्मीकीये आदिकाव्ये अयोध्याकाण्डे शततमः सर्गः॥}

\sect{एकाधिकशततमः सर्गः}

\twolineshloka
{जाबालेस्तु वचः श्रुत्वा रामः सत्यात्मनां वरः}
{उवाच परया युक्त्या स्वबुद्ध्या चाविपन्नया} %||2-101-1||

\twolineshloka
{भवान्मे प्रियकामार्थं वचनं यदिहोक्तवान्}
{अकार्यं कार्यसङ्काशमपथ्यं पथ्यसम्मितम्} %||2-101-2||

\twolineshloka
{निर्मर्यादस्तु पुरुषः पापाचारसमन्वितः}
{मानं न लभते सत्सु भिन्नचारित्रदर्शनः} %||2-101-3||

\twolineshloka
{कुलीनमकुलीनं वा वीरं पुरुषमानिनम्}
{चारित्रमेव व्याख्याति शुचिं वा यदि वाशुचिम्} %||2-101-4||

\twolineshloka
{अनारय्स्त्वार्य सङ्काशः शौचाद्धीनस्तथा शुचिः}
{लक्षण्यवदलक्षण्यो दुःशीलः शीलवानिव} %||2-101-5||

\twolineshloka
{अधर्मं धर्मवेषेण यदीमं लोकसङ्करम्}
{अभिपत्स्ये शुभं हित्वा क्रियाविधिविवर्जितम्} %||2-101-6||

\twolineshloka
{कश्चेतयानः पुरुषः कार्याकार्यविचक्षणः}
{बहु मंस्यति मां लोके दुर्वृत्तं लोकदूषणम्} %||2-101-7||

\twolineshloka
{कस्य यास्याम्यहं वृत्तं केन वा स्वर्गमाप्नुयाम्}
{अनया वर्तमानोऽहं वृत्त्या हीनप्रतिज्ञया} %||2-101-8||

\twolineshloka
{कामवृत्तस्त्वयं लोकः कृत्स्नः समुपवर्तते}
{यद्वृत्ताः सन्ति राजानस्तद्वृत्ताः सन्ति हि प्रजाः} %||2-101-9||

\twolineshloka
{सत्यमेवानृशंस्यं च राजवृत्तं सनातनम्}
{तस्मात्सत्यात्मकं राज्यं सत्ये लोकः प्रतिष्ठितः} %||2-101-10||

\twolineshloka
{ऋषयश्चैव देवाश्च सत्यमेव हि मेनिरे}
{सत्यवादी हि लोकेऽस्मिन्परमं गच्छति क्षयम्} %||2-101-11||

\twolineshloka
{उद्विजन्ते यथा सर्पान्नरादनृतवादिनः}
{धर्मः सत्यं परो लोके मूलं स्वर्गस्य चोच्यते} %||2-101-12||

\twolineshloka
{सत्यमेवेश्वरो लोके सत्यं पद्मा समाश्रिता}
{सत्यमूलानि सर्वाणि सत्यान्नास्ति परं पदम्} %||2-101-13||

\twolineshloka
{दत्तमिष्टं हुतं चैव तप्तानि च तपांसि च}
{वेदाः सत्यप्रतिष्ठानास्तस्मात्सत्यपरो भवेत्} %||2-101-14||

\twolineshloka
{एकः पालयते लोकमेकः पालयते कुलम्}
{मज्जत्येको हि निरय एकः स्वर्गे महीयते} %||2-101-15||

\twolineshloka
{सोऽहं पितुर्निदेशं तु किमर्थं नानुपालये}
{सत्यप्रतिश्रवः सत्यं सत्येन समयीकृतः} %||2-101-16||

\twolineshloka
{नैव लोभान्न मोहाद्वा न चाज्ञानात्तमोऽन्वितः}
{सेतुं सत्यस्य भेत्स्यामि गुरोः सत्यप्रतिश्रवः} %||2-101-17||

\twolineshloka
{असत्यसन्धस्य सतश्चलस्यास्थिरचेतसः}
{नैव देवा न पितरः प्रतीच्छन्तीति नः श्रुतम्} %||2-101-18||

\twolineshloka
{प्रत्यगात्ममिमं धर्मं सत्यं पश्याम्यहं स्वयम्}
{भारः सत्पुरुषाचीर्णस्तदर्थमभिनन्द्यते} %||2-101-19||

\twolineshloka
{क्षात्रं धर्ममहं त्यक्ष्ये ह्यधर्मं धर्मसंहितम्}
{क्षुद्रौर्नृशंसैर्लुब्धैश्च सेवितं पापकर्मभिः} %||2-101-20||

\twolineshloka
{कायेन कुरुते पापं मनसा सम्प्रधार्य च}
{अनृतं जिह्वया चाह त्रिविधं कर्म पातकम्} %||2-101-21||

\twolineshloka
{भूमिः कीर्तिर्यशो लक्ष्मीः पुरुषं प्रार्थयन्ति हि}
{स्वर्गस्थं चानुबध्नन्ति सत्यमेव भजेत तत्} %||2-101-22||

\twolineshloka
{श्रेष्ठं ह्यनार्यमेव स्याद्यद्भवानवधार्य माम्}
{आह युक्तिकरैर्वाक्यैरिदं भद्रं कुरुष्व ह} %||2-101-23||

\twolineshloka
{कथं ह्यहं प्रतिज्ञाय वनवासमिमं गुरोः}
{भरतस्य करिष्यामि वचो हित्वा गुरोर्वचः} %||2-101-24||

\twolineshloka
{स्थिरा मया प्रतिज्ञाता प्रतिज्ञा गुरुसंनिधौ}
{प्रहृष्टमानसा देवी कैकेयी चाभवत्तदा} %||2-101-25||

\twolineshloka
{वनवासं वसन्नेवं शुचिर्नियतभोजनः}
{मूलैः पुष्पैः फलैः पुण्यैः पितॄन्देवांश्च तर्पयन्} %||2-101-26||

\twolineshloka
{सन्तुष्टपञ्चवर्गोऽहं लोकयात्रां प्रवर्तये}
{अकुहः श्रद्दधानः सन्कार्याकार्यविचक्षणः} %||2-101-27||

\twolineshloka
{कर्मभूमिमिमां प्राप्य कर्तव्यं कर्म यच्छुभम्}
{अग्निर्वायुश्च सोमश्च कर्मणां फलभागिनः} %||2-101-28||

\twolineshloka
{शतं क्रतूनामाहृत्य देवराट्त्रिदिवं गतः}
{तपांस्युग्राणि चास्थाय दिवं याता महर्षयः} %||2-101-29||

\fourlineindentedshloka
{सत्यं च धर्मं च पराक्रमं च}
{ भूतानुकम्पां प्रियवादितां च}
{द्विजातिदेवातिथिपूजनं च}
{ पन्थानमाहुस्त्रिदिवस्य सन्तः} %||2-101-30||

\fourlineindentedshloka
{धर्मे रताः सत्पुरुषैः समेता}
{स्तेजस्विनो दानगुणप्रधानाः}
{अहिंसका वीतमलाश्च लोके}
{ भवन्ति पूज्या मुनयः प्रधानाः} %||2-102-31||


{॥इत्यार्षे श्रीमद्रामायणे वाल्मीकीये आदिकाव्ये अयोध्याकाण्डे एकाधिकशततमः सर्गः॥}

\sect{द्व्यधिकशततमः सर्गः}

\threelineshloka
{क्रुद्धमाज्ञाय राम तु वसिष्ठः प्रत्युवाच ह}
{जाबालिरपि जानीते लोकस्यास्य गतागतिम्}
{निवर्तयितु कामस्तु त्वामेतद्वाक्यमब्रवीत्} %||2-102-1||

\threelineshloka
{इमां लोकसमुत्पत्तिं लोकनाथ निबोध मे}
{सर्वं सलिलमेवासीत्पृथिवी यत्र निर्मिता}
{ततः समभवद्ब्रह्मा स्वयम्भूर्दैवतैः सह} %||2-102-2||

\twolineshloka
{स वराहस्ततो भूत्वा प्रोज्जहार वसुन्धराम्}
{असृजच्च जगत्सर्वं सह पुत्रैः कृतात्मभिः} %||2-102-3||

\twolineshloka
{आकाशप्रभवो ब्रह्मा शाश्वतो नित्य अव्ययः}
{तस्मान्मरीचिः संजज्ञे मरीचेः कश्यपः सुतः} %||2-102-4||

\twolineshloka
{विवस्वान्कश्यपाज्जज्ञे मनुर्वैवस्वतः स्मृतः}
{स तु प्रजापतिः पूर्वमिक्ष्वाकुस्तु मनोः सुतः} %||2-102-5||

\twolineshloka
{यस्येयं प्रथमं दत्ता समृद्धा मनुना मही}
{तमिक्ष्वाकुमयोध्यायां राजानं विद्धि पूर्वकम्} %||2-102-6||

\twolineshloka
{इक्ष्वाकोस्तु सुतः श्रीमान्कुक्षिरेवेति विश्रुतः}
{कुक्षेरथात्मजो वीरो विकुक्षिरुदपद्यत} %||2-102-7||

\twolineshloka
{विकुक्षेस्तु महातेजा बाणः पुत्रः प्रतापवान्}
{बाणस्य तु महाबाहुरनरण्यो महायशाः} %||2-102-8||

\twolineshloka
{नाना वृष्टिर्बभूवास्मिन्न दुर्भिक्षं सतां वरे}
{अनरण्ये महाराजे तस्करो वापि कश्चन} %||2-102-9||

\threelineshloka
{अनरण्यान्महाबाहुः पृथू राजा बभूव ह}
{तस्मात्पृथोर्महाराजस्त्रिशङ्कुरुदपद्यत}
{स सत्यवचनाद्वीरः सशरीरो दिवं गतः} %||2-102-10||

\twolineshloka
{त्रिशङ्कोरभवत्सूनुर्धुन्धुमारो महायशाः}
{धुन्धुमारान्महातेजा युवनाश्वो व्यजायत} %||2-102-11||

\twolineshloka
{युवनाश्व सुतः श्रीमान्मान्धाता समपद्यत}
{मान्धातुस्तु महातेजाः सुसन्धिरुदपद्यत} %||2-102-12||

\twolineshloka
{सुसन्धेरपि पुत्रौ द्वौ ध्रुवसन्धिः प्रसेनजित्}
{यशस्वी ध्रुवसन्धेस्तु भरतो रिपुसूदनः} %||2-102-13||

\threelineshloka
{भरतात्तु महाबाहोरसितो नाम जायत}
{यस्यैते प्रतिराजान उदपद्यन्त शत्रवः}
{हैहयास्तालजङ्घाश्च शूराश्च शशबिन्दवः} %||2-102-14||

\threelineshloka
{तांस्तु सर्वान्प्रतिव्यूह्य युद्धे राजा प्रवासितः}
{स च शैलवरे रम्ये बभूवाभिरतो मुनिः}
{द्वे चास्य भार्ये गर्भिण्यौ बभूवतुरिति श्रुतिः} %||2-102-15||

\twolineshloka
{भार्गवश्च्यवनो नाम हिमवन्तमुपाश्रितः}
{तमृषिं समुपागम्य कालिन्दी त्वभ्यवादयत्} %||2-102-16||

\twolineshloka
{स तामभ्यवदद्विप्रो वरेप्सुं पुत्रजन्मनि}
{ततः सा गृहमागम्य देवी पुत्रं व्यजायत} %||2-102-17||

\twolineshloka
{सपत्न्या तु गरस्तस्यै दत्तो गर्भजिघांसया}
{गरेण सह तेनैव जातः स सगरोऽभवत्} %||2-102-18||

\twolineshloka
{स राजा सगरो नाम यः समुद्रमखानयत्}
{इष्ट्वा पर्वणि वेगेन त्रासयन्तमिमाः प्रजाः} %||2-102-19||

\twolineshloka
{असमञ्जस्तु पुत्रोऽभूत्सगरस्येति नः श्रुतम्}
{जीवन्नेव स पित्रा तु निरस्तः पापकर्मकृत्} %||2-102-20||

\twolineshloka
{अंशुमानिति पुत्रोऽभूदसमञ्जस्य वीर्यवान्}
{दिलीपोंऽशुमतः पुत्रो दिलीपस्य भगीरथः} %||2-102-21||

\twolineshloka
{भगीरथात्ककुत्स्थस्तु काकुत्स्था येन तु स्मृताः}
{ककुत्स्थस्य तु पुत्रोऽभूद्रघुर्येन तु राघवः} %||2-102-22||

\twolineshloka
{रघोस्तु पुत्रस्तेजस्वी प्रवृद्धः पुरुषादकः}
{कल्माषपादः सौदास इत्येवं प्रथितो भुवि} %||2-102-23||

\twolineshloka
{कल्माषपादपुत्रोऽभूच्छङ्खणस्त्विति विश्रुतः}
{यस्तु तद्वीर्यमासाद्य सहसेनो व्यनीनशत्} %||2-102-24||

\twolineshloka
{शङ्खणस्य तु पुत्रोऽभूच्छूरः श्रीमान्सुदर्शनः}
{सुदर्शनस्याग्निवर्ण अग्निवर्षस्य शीघ्रगः} %||2-102-25||

\twolineshloka
{शीघ्रगस्य मरुः पुत्रो मरोः पुत्रः प्रशुश्रुकः}
{प्रशुश्रुकस्य पुत्रोऽभूदम्बरीषो महाद्युतिः} %||2-102-26||

\twolineshloka
{अम्बरीषस्य पुत्रोऽभून्नहुषः सत्यविक्रमः}
{नहुषस्य च नाभागः पुत्रः परमधार्मिकः} %||2-102-27||

\twolineshloka
{अजश्च सुव्रतश्चैव नाभागस्य सुतावुभौ}
{अजस्य चैव धर्मात्मा राजा दशरथः सुतः} %||2-102-28||

\twolineshloka
{तस्य ज्येष्ठोऽसि दायादो राम इत्यभिविश्रुतः}
{तद्गृहाण स्वकं राज्यमवेक्षस्व जगन्नृप} %||2-102-29||

\twolineshloka
{इक्ष्वाकूणां हि सर्वेषां राजा भवति पूर्वजः}
{पूर्वजेनावरः पुत्रो ज्येष्ठो राज्येऽभिषिच्यते} %||2-102-30||

\fourlineindentedshloka
{स राघवाणां कुलधर्ममात्मनः}
{ सनातनं नाद्य विहातुमर्हसि}
{प्रभूतरत्नामनुशाधि मेदिनीम्}
{ प्रभूतराष्ट्रां पितृवन्महायशाः} %||2-103-31||


{॥इत्यार्षे श्रीमद्रामायणे वाल्मीकीये आदिकाव्ये अयोध्याकाण्डे द्व्यधिकशततमः सर्गः॥}

\sect{त्र्यधिकशततमः सर्गः}

\twolineshloka
{वसिष्ठस्तु तदा राममुक्त्वा राजपुरोहितः}
{अब्रवीद्धर्मसंयुक्तं पुनरेवापरं वचः} %||2-103-1||

\twolineshloka
{पुरुषस्येह जातस्य भवन्ति गुरवस्त्रयः}
{आचार्यश्चैव काकुत्स्थ पिता माता च राघव} %||2-103-2||

\twolineshloka
{पिता ह्येनं जनयति पुरुषं पुरुषर्षभ}
{प्रज्ञां ददाति चाचार्यस्तस्मात्स गुरुरुच्यते} %||2-103-3||

\twolineshloka
{स तेऽहं पितुराचार्यस्तव चैव परन्तप}
{मम त्वं वचनं कुर्वन्नातिवर्तेः सतां गतिम्} %||2-103-4||

\twolineshloka
{इमा हि ते परिषदः श्रेणयश्च समागताः}
{एषु तात चरन्धर्मं नातिवर्तेः सतां गतिम्} %||2-103-5||

\twolineshloka
{वृद्धाया धर्मशीलाया मातुर्नार्हस्यवर्तितुम्}
{अस्यास्तु वचनं कुर्वन्नातिवर्तेः सतां गतिम्} %||2-103-6||

\twolineshloka
{भरतस्य वचः कुर्वन्याचमानस्य राघव}
{आत्मानं नातिवर्तेस्त्वं सत्यधर्मपराक्रम} %||2-103-7||

\twolineshloka
{एवं मधुरमुक्तस्तु गुरुणा राघवः स्वयम्}
{प्रत्युवाच समासीनं वसिष्ठं पुरुषर्षभः} %||2-103-8||

\twolineshloka
{यन्मातापितरौ वृत्तं तनये कुरुतः सदा}
{न सुप्रतिकरं तत्तु मात्रा पित्रा च यत्कृतम्} %||2-103-9||

\twolineshloka
{यथाशक्ति प्रदानेन स्नापनाच्छादनेन च}
{नित्यं च प्रियवादेन तथा संवर्धनेन च} %||2-103-10||

\twolineshloka
{स हि राजा जनयिता पिता दशरथो मम}
{आज्ञातं यन्मया तस्य न तन्मिथ्या भविष्यति} %||2-103-11||

\twolineshloka
{एवमुक्तस्तु रामेण भरतः प्रत्यनन्तरम्}
{उवाच परमोदारः सूतं परमदुर्मनाः} %||2-103-12||

\twolineshloka
{इह मे स्थण्डिले शीघ्रं कुशानास्तर सारथे}
{आर्यं प्रत्युपवेक्ष्यामि यावन्मे न प्रसीदति} %||2-103-13||

\twolineshloka
{अनाहारो निरालोको धनहीनो यथा द्विजः}
{शेष्ये पुरस्ताच्छालाया यावन्न प्रतियास्यति} %||2-103-14||

\twolineshloka
{स तु राममवेक्षन्तं सुमन्त्रं प्रेक्ष्य दुर्मनाः}
{कुशोत्तरमुपस्थाप्य भूमावेवास्तरत्स्वयम्} %||2-103-15||

\twolineshloka
{तमुवाच महातेजा रामो राजर्षिसत्तमाः}
{किं मां भरत कुर्वाणं तात प्रत्युपवेक्ष्यसि} %||2-103-16||

\twolineshloka
{ब्राह्मणो ह्येकपार्श्वेन नरान्रोद्धुमिहार्हति}
{न तु मूर्धावसिक्तानां विधिः प्रत्युपवेशने} %||2-103-17||

\twolineshloka
{उत्तिष्ठ नरशार्दूल हित्वैतद्दारुणं व्रतम्}
{पुरवर्यामितः क्षिप्रमयोध्यां याहि राघव} %||2-103-18||

\twolineshloka
{आसीनस्त्वेव भरतः पौरजानपदं जनम्}
{उवाच सर्वतः प्रेक्ष्य किमार्यं नानुशासथ} %||2-103-19||

\twolineshloka
{ते तमूचुर्महात्मानं पौरजानपदा जनाः}
{काकुत्स्थमभिजानीमः सम्यग्वदति राघवः} %||2-103-20||

\twolineshloka
{एषोऽपि हि महाभागः पितुर्वचसि तिष्ठति}
{अत एव न शक्ताः स्मो व्यावर्तयितुमञ्जसा} %||2-103-21||

\twolineshloka
{तेषामाज्ञाय वचनं रामो वचनमब्रवीत्}
{एवं निबोध वचनं सुहृदां धर्मचक्षुषाम्} %||2-103-22||

\twolineshloka
{एतच्चैवोभयं श्रुत्वा सम्यक्सम्पश्य राघव}
{उत्तिष्ठ त्वं महाबाहो मां च स्पृश तथोदकम्} %||2-103-23||

\twolineshloka
{अथोत्थाय जलं स्पृष्ट्वा भरतो वाक्यमब्रवीत्}
{शृण्वन्तु मे परिषदो मन्त्रिणः श्रेणयस्तथा} %||2-103-24||

\twolineshloka
{न याचे पितरं राज्यं नानुशासामि मातरम्}
{आर्यं परमधर्मज्ञमभिजानामि राघवम्} %||2-103-25||

\twolineshloka
{यदि त्ववश्यं वस्तव्यं कर्तव्यं च पितुर्वचः}
{अहमेव निवत्स्यामि चतुर्दश वने समाः} %||2-103-26||

\twolineshloka
{धर्मात्मा तस्य तथ्येन भ्रातुर्वाक्येन विस्मितः}
{उवाच रामः सम्प्रेक्ष्य पौरजानपदं जनम्} %||2-103-27||

\twolineshloka
{विक्रीतमाहितं क्रीतं यत्पित्रा जीवता मम}
{न तल्लोपयितुं शक्यं मया वा भरतेन वा} %||2-103-28||

\twolineshloka
{उपधिर्न मया कार्यो वनवासे जुगुप्सितः}
{युक्तमुक्तं च कैकेय्या पित्रा मे सुकृतं कृतम्} %||2-103-29||

\twolineshloka
{जानामि भरतं क्षान्तं गुरुसत्कारकारिणम्}
{सर्वमेवात्र कल्याणं सत्यसन्धे महात्मनि} %||2-103-30||

\twolineshloka
{अनेन धर्मशीलेन वनात्प्रत्यागतः पुनः}
{भ्रात्रा सह भविष्यामि पृथिव्याः पतिरुत्तमः} %||2-103-31||

\twolineshloka
{वृतो राजा हि कैकेय्या मया तद्वचनं कृतम्}
{अनृतान्मोचयानेन पितरं तं महीपतिम्} %||2-104-32||


{॥इत्यार्षे श्रीमद्रामायणे वाल्मीकीये आदिकाव्ये अयोध्याकाण्डे त्र्यधिकशततमः सर्गः॥}

\sect{चतुरधिकशततमः सर्गः}

\twolineshloka
{तमप्रतिमतेजोभ्यां भ्रातृभ्यां रोमहर्षणम्}
{विस्मिताः सङ्गमं प्रेक्ष्य समवेता महर्षयः} %||2-104-1||

\twolineshloka
{अन्तर्हितास्त्वृषिगणाः सिद्धाश्च परमर्षयः}
{तौ भ्रातरौ महात्मानौ काकुत्स्थौ प्रशशंसिरे} %||2-104-2||

\twolineshloka
{स धन्यो यस्य पुत्रौ द्वौ धर्मज्ञौ धर्मविक्रमौ}
{श्रुत्वा वयं हि सम्भाषामुभयोः स्पृहयामहे} %||2-104-3||

\twolineshloka
{ततस्त्वृषिगणाः क्षिप्रं दशग्रीववधैषिणः}
{भरतं राजशार्दूलमित्यूचुः सङ्गता वचः} %||2-104-4||

\twolineshloka
{कुले जात महाप्राज्ञ महावृत्त महायशः}
{ग्राह्यं रामस्य वाक्यं ते पितरं यद्यवेक्षसे} %||2-104-5||

\twolineshloka
{सदानृणमिमं रामं वयमिच्छामहे पितुः}
{अनृणत्वाच्च कैकेय्याः स्वर्गं दशरथो गतः} %||2-104-6||

\twolineshloka
{एतावदुक्त्वा वचनं गन्धर्वाः समहर्षयः}
{राजर्षयश्चैव तथा सर्वे स्वां स्वां गतिं गताः} %||2-104-7||

\twolineshloka
{ह्लादितस्तेन वाक्येन शुभेन शुभदर्शनः}
{रामः संहृष्टवदनस्तानृषीनभ्यपूजयत्} %||2-104-8||

\twolineshloka
{स्रस्तगात्रस्तु भरतः स वाचा सज्जमानया}
{कृताञ्जलिरिदं वाक्यं राघवं पुनरब्रवीत्} %||2-104-9||

\twolineshloka
{राजधर्ममनुप्रेक्ष्य कुलधर्मानुसन्ततिम्}
{कर्तुमर्हसि काकुत्स्थ मम मातुश्च याचनाम्} %||2-104-10||

\twolineshloka
{रक्षितुं सुमहद्राज्यमहमेकस्तु नोत्सहे}
{पौरजानपदांश्चापि रक्तान्रञ्जयितुं तथा} %||2-104-11||

\twolineshloka
{ज्ञातयश्च हि योधाश्च मित्राणि सुहृदश्च नः}
{त्वामेव प्रतिकाङ्क्षन्ते पर्जन्यमिव कर्षकाः} %||2-104-12||

\twolineshloka
{इदं राज्यं महाप्राज्ञ स्थापय प्रतिपद्य हि}
{शक्तिमानसि काकुत्स्थ लोकस्य परिपालने} %||2-104-13||

\twolineshloka
{इत्युक्त्वा न्यपतद्भ्रातुः पादयोर्भरतस्तदा}
{भृशं सम्प्रार्थयामास राममेवं प्रियं वदः} %||2-104-14||

\twolineshloka
{तमङ्के भ्रातरं कृत्वा रामो वचनमब्रवीत्}
{श्यामं नलिनपत्राक्षं मत्तहंसस्वरः स्वयम्} %||2-104-15||

\twolineshloka
{आगता त्वामियं बुद्धिः स्वजा वैनयिकी च या}
{भृशमुत्सहसे तात रक्षितुं पृथिवीमपि} %||2-104-16||

\twolineshloka
{अमात्यैश्च सुहृद्भिश्च बुद्धिमद्भिश्च मन्त्रिभिः}
{सर्वकार्याणि सम्मन्त्र्य सुमहान्त्यपि कारय} %||2-104-17||

\twolineshloka
{लक्ष्मीश्चन्द्रादपेयाद्वा हिमवान्वा हिमं त्यजेत्}
{अतीयात्सागरो वेलां न प्रतिज्ञामहं पितुः} %||2-104-18||

\twolineshloka
{कामाद्वा तात लोभाद्वा मात्रा तुभ्यमिदं कृतम्}
{न तन्मनसि कर्तव्यं वर्तितव्यं च मातृवत्} %||2-104-19||

\twolineshloka
{एवं ब्रुवाणं भरतः कौसल्यासुतमब्रवीत्}
{तेजसादित्यसङ्काशं प्रतिपच्चन्द्रदर्शनम्} %||2-104-20||

\twolineshloka
{अधिरोहार्य पादाभ्यां पादुके हेमभूषिते}
{एते हि सर्वलोकस्य योगक्षेमं विधास्यतः} %||2-104-21||

\twolineshloka
{सोऽधिरुह्य नरव्याघ्रः पादुके ह्यवरुह्य च}
{प्रायच्छत्सुमहातेजा भरताय महात्मने} %||2-104-22||

\fourlineindentedshloka
{स पादुके ते भरतः प्रतापवा}
{न्स्वलङ्कृते सम्परिगृह्य धर्मवित्}
{प्रदक्षिणं चैव चकार राघवम्}
{ चकार चैवोत्तमनागमूर्धनि} %||2-104-23||

\fourlineindentedshloka
{अथानुपूर्व्यात्प्रतिपूज्य तं जनम्}
{ गुरूंश्च मन्त्रिप्रकृतीस्तथानुजौ}
{व्यसर्जयद्राघववंशवर्धनः}
{ स्थितः स्वधर्मे हिमवानिवाचलः} %||2-104-24||

\fourlineindentedshloka
{तं मातरो बाष्पगृहीतकण्ठो}
{ दुःखेन नामन्त्रयितुं हि शेकुः}
{स त्वेव मातॄरभिवाद्य सर्वा}
{ रुदन्कुटीं स्वां प्रविवेश रामः} %||2-105-25||


{॥इत्यार्षे श्रीमद्रामायणे वाल्मीकीये आदिकाव्ये अयोध्याकाण्डे चतुरधिकशततमः सर्गः॥}

\sect{पञ्चाधिकशततमः सर्गः}

\twolineshloka
{ततः शिरसि कृत्वा तु पादुके भरतस्तदा}
{आरुरोह रथं हृष्टः शत्रुघ्नेन समन्वितः} %||2-105-1||

\twolineshloka
{वसिष्ठो वामदेवश्च जाबालिश्च दृढव्रतः}
{अग्रतः प्रययुः सर्वे मन्त्रिणो मन्त्रपूजिताः} %||2-105-2||

\twolineshloka
{मन्दाकिनीं नदीं रम्यां प्राङ्मुखास्ते ययुस्तदा}
{प्रदक्षिणं च कुर्वाणाश्चित्रकूटं महागिरिम्} %||2-105-3||

\twolineshloka
{पश्यन्धातुसहस्राणि रम्याणि विविधानि च}
{प्रययौ तस्य पार्श्वेन ससैन्यो भरतस्तदा} %||2-105-4||

\twolineshloka
{अदूराच्चित्रकूटस्य ददर्श भरतस्तदा}
{आश्रमं यत्र स मुनिर्भरद्वाजः कृतालयः} %||2-105-5||

\twolineshloka
{स तमाश्रममागम्य भरद्वाजस्य बुद्धिमान्}
{अवतीर्य रथात्पादौ ववन्दे कुलनन्दनः} %||2-105-6||

\twolineshloka
{ततो हृष्टो भरद्वाजो भरतं वाक्यमब्रवीत्}
{अपि कृत्यं कृतं तात रामेण च समागतम्} %||2-105-7||

\twolineshloka
{एवमुक्तस्तु भरतो भरद्वाजेन धीमता}
{प्रत्युवाच भरद्वाजं भरतो धर्मवत्सलः} %||2-105-8||

\twolineshloka
{स याच्यमानो गुरुणा मया च दृढविक्रमः}
{राघवः परमप्रीतो वसिष्ठं वाक्यमब्रवीत्} %||2-105-9||

\twolineshloka
{पितुः प्रतिज्ञां तामेव पालयिष्यामि तत्त्वतः}
{चतुर्दश हि वर्षाणि य प्रतिज्ञा पितुर्मम} %||2-105-10||

\twolineshloka
{एवमुक्तो महाप्राज्ञो वसिष्ठः प्रत्युवाच ह}
{वाक्यज्ञो वाक्यकुशलं राघवं वचनं महत्} %||2-105-11||

\twolineshloka
{एते प्रयच्छ संहृष्टः पादुके हेमभूषिते}
{अयोध्यायां महाप्राज्ञ योगक्षेमकरे तव} %||2-105-12||

\twolineshloka
{एवमुक्तो वसिष्ठेन राघवः प्राङ्मुखः स्थितः}
{पादुके हेमविकृते मम राज्याय ते ददौ} %||2-105-13||

\twolineshloka
{निवृत्तोऽहमनुज्ञातो रामेण सुमहात्मना}
{अयोध्यामेव गच्छामि गृहीत्वा पादुके शुभे} %||2-105-14||

\twolineshloka
{एतच्छ्रुत्वा शुभं वाक्यं भरतस्य महात्मनः}
{भरद्वाजः शुभतरं मुनिर्वाक्यमुदाहरत्} %||2-105-15||

\twolineshloka
{नैतच्चित्रं नरव्याघ्र शीलवृत्तवतां वर}
{यदार्यं त्वयि तिष्ठेत्तु निम्ने वृष्टिमिवोदकम्} %||2-105-16||

\twolineshloka
{अमृतः स महाबाहुः पिता दशरथस्तव}
{यस्य त्वमीदृशः पुत्रो धर्मात्मा धर्मवत्सलः} %||2-105-17||

\twolineshloka
{तमृषिं तु महात्मानमुक्तवाक्यं कृताञ्जलिः}
{आमन्त्रयितुमारेभे चरणावुपगृह्य च} %||2-105-18||

\twolineshloka
{ततः प्रदक्षिणं कृत्वा भरद्वाजं पुनः पुनः}
{भरतस्तु ययौ श्रीमानयोध्यां सह मन्त्रिभिः} %||2-105-19||

\twolineshloka
{यानैश्च शकटैश्चैव हयैश्नागैश्च सा चमूः}
{पुनर्निवृत्ता विस्तीर्णा भरतस्यानुयायिनी} %||2-105-20||

\twolineshloka
{ततस्ते यमुनां दिव्यां नदीं तीर्त्वोर्मिमालिनीम्}
{ददृशुस्तां पुनः सर्वे गङ्गां शिवजलां नदीम्} %||2-105-21||

\twolineshloka
{तां रम्यजलसम्पूर्णां सन्तीर्य सह बान्धवः}
{शृङ्गवेरपुरं रम्यं प्रविवेश ससैनिकः} %||2-105-22||

\twolineshloka
{शृङ्गवेरपुराद्भूय अयोध्यां सन्ददर्श ह}
{भरतो दुःखसन्तप्तः सारथिं चेदमब्रवीत्} %||2-105-23||

\twolineshloka
{सारथे पश्य विध्वस्ता अयोध्या न प्रकाशते}
{निराकारा निरानन्दा दीना प्रतिहतस्वना} %||2-106-24||


{॥इत्यार्षे श्रीमद्रामायणे वाल्मीकीये आदिकाव्ये अयोध्याकाण्डे पञ्चाधिकशततमः सर्गः॥}

\sect{षष्ठाधिकशततमः सर्गः}

\twolineshloka
{स्निग्धगम्भीरघोषेण स्यन्दनेनोपयान्प्रभुः}
{अयोध्यां भरतः क्षिप्रं प्रविवेश महायशाः} %||2-106-1||

\twolineshloka
{बिडालोलूकचरितामालीननरवारणाम्}
{तिमिराभ्याहतां कालीमप्रकाशां निशामिव} %||2-106-2||

\twolineshloka
{राहुशत्रोः प्रियां पत्नीं श्रिया प्रज्वलितप्रभाम्}
{ग्रहेणाभ्युत्थितेनैकां रोहिणीमिव पीडिताम्} %||2-106-3||

\twolineshloka
{अल्पोष्णक्षुब्धसलिलां घर्मोत्तप्तविहङ्गमाम्}
{लीनमीनझषग्राहां कृशां गिरिनदीमिव} %||2-106-4||

\twolineshloka
{विधूमामिव हेमाभामध्वराग्निसमुत्थिताम्}
{हविरभ्युक्षितां पश्चाच्छिखां विप्रलयं गताम्} %||2-106-5||

\twolineshloka
{विध्वस्तकवचां रुग्णगजवाजिरथध्वजाम्}
{हतप्रवीरामापन्नां चमूमिव महाहवे} %||2-106-6||

\twolineshloka
{सफेनां सस्वनां भूत्वा सागरस्य समुत्थिताम्}
{प्रशान्तमारुतोद्धूतां जलोर्मिमिव निःस्वनाम्} %||2-106-7||

\twolineshloka
{त्यक्तां यज्ञायुधैः सर्वैरभिरूपैश्च याजकैः}
{सुत्याकाले विनिर्वृत्ते वेदिं गतरवामिव} %||2-106-8||

\twolineshloka
{गोष्ठमध्ये स्थितामार्तामचरन्तीं नवं तृणम्}
{गोवृषेण परित्यक्तां गवां पत्नीमिवोत्सुकाम्} %||2-106-9||

\twolineshloka
{प्रभाकरालैः सुस्निग्धैः प्रज्वलद्भिरिवोत्तमैः}
{वियुक्तां मणिभिर्जात्यैर्नवां मुक्तावलीमिव} %||2-106-10||

\twolineshloka
{सहसा चलितां स्थानान्महीं पुण्यक्षयाद्गताम्}
{संहृतद्युतिविस्तारां तारामिव दिवश्च्युताम्} %||2-106-11||

\twolineshloka
{पुष्पनद्धां वसन्तान्ते मत्तभ्रमरशालिनीम्}
{द्रुतदावाग्निविप्लुष्टां क्लान्तां वनलतामिव} %||2-106-12||

\twolineshloka
{सम्मूढनिगमां सर्वां सङ्क्षिप्तविपणापणाम्}
{प्रच्छन्नशशिनक्षत्रां द्यामिवाम्बुधरैर्वृताम्} %||2-106-13||

\twolineshloka
{क्षीणपानोत्तमैर्भिन्नैः शरावैरभिसंवृताम्}
{हतशौण्डामिवाकाशे पानभूमिमसंस्कृताम्} %||2-106-14||

\twolineshloka
{वृक्णभूमितलां निम्नां वृक्णपात्रैः समावृताम्}
{उपयुक्तोदकां भग्नां प्रपां निपतितामिव} %||2-106-15||

\twolineshloka
{विपुलां विततां चैव युक्तपाशां तरस्विनाम्}
{भूमौ बाणैर्विनिष्कृत्तां पतितां ज्यामिवायुधात्} %||2-106-16||

\twolineshloka
{सहसा युद्धशौण्डेन हयारोहेण वाहिताम्}
{निक्षिप्तभाण्डामुत्सृष्टां किशोरीमिव दुर्बलाम्} %||2-106-17||

\twolineshloka
{प्रावृषि प्रविगाढायां प्रविष्टस्याभ्र मण्डलम्}
{प्रच्छन्नां नीलजीमूतैर्भास्करस्य प्रभामिव} %||2-106-18||

\twolineshloka
{भरतस्तु रथस्थः सञ्श्रीमान्दशरथात्मजः}
{वाहयन्तं रथश्रेष्ठं सारथिं वाक्यमब्रवीत्} %||2-106-19||

\twolineshloka
{किं नु खल्वद्य गम्भीरो मूर्छितो न निशम्यते}
{यथापुरमयोध्यायां गीतवादित्रनिःस्वनः} %||2-106-20||

\twolineshloka
{वारुणीमदगन्धाश्च माल्यगन्धश्च मूर्छितः}
{धूपितागरुगन्धश्च न प्रवाति समन्ततः} %||2-106-21||

\threelineshloka
{यानप्रवरघोषश्च स्निग्धश्च हयनिःस्वनः}
{प्रमत्तगजनादश्च महांश्च रथनिःस्वनः}
{नेदानीं श्रूयते पुर्यामस्यां रामे विवासिते} %||2-106-22||

\twolineshloka
{तरुणैश्चारु वेषैश्च नरैरुन्नतगामिभिः}
{सम्पतद्भिरयोध्यायां न विभान्ति महापथाः} %||2-106-23||

\twolineshloka
{एवं बहुविधं जल्पन्विवेश वसतिं पितुः}
{तेन हीनां नरेन्द्रेण सिंहहीनां गुहामिव} %||2-107-24||


{॥इत्यार्षे श्रीमद्रामायणे वाल्मीकीये आदिकाव्ये अयोध्याकाण्डे षष्ठाधिकशततमः सर्गः॥}

\sect{सप्तमाधिकशततमः सर्गः}

\twolineshloka
{ततो निक्षिप्य मातॄह्स अयोध्यायां दृढव्रतः}
{भरतः शोकसन्तप्तो गुरूनिदमथाब्रवीत्} %||2-107-1||

\twolineshloka
{नन्दिग्रामं गमिष्यामि सर्वानामन्त्रयेऽद्य वः}
{तत्र दुःखमिदं सर्वं सहिष्ये राघवं विना} %||2-107-2||

\twolineshloka
{गतश्च हि दिवं राजा वनस्थश्च गुरुर्मम}
{रामं प्रतीक्षे राज्याय स हि राजा महायशाः} %||2-107-3||

\twolineshloka
{एतच्छ्रुत्वा शुभं वाक्यं भरतस्य महात्मनः}
{अब्रुवन्मन्त्रिणः सर्वे वसिष्ठश्च पुरोहितः} %||2-107-4||

\twolineshloka
{सदृशं श्लाघनीयं च यदुक्तं भरत त्वया}
{वचनं भ्रातृवात्सल्यादनुरूपं तवैव तत्} %||2-107-5||

\twolineshloka
{नित्यं ते बन्धुलुब्धस्य तिष्ठतो भ्रातृसौहृदे}
{आर्यमार्गं प्रपन्नस्य नानुमन्येत कः पुमान्} %||2-107-6||

\twolineshloka
{मन्त्रिणां वचनं श्रुत्वा यथाभिलषितं प्रियम्}
{अब्रवीत्सारथिं वाक्यं रथो मे युज्यतामिति} %||2-107-7||

\twolineshloka
{प्रहृष्टवदनः सर्वा मातॄह्समभिवाद्य सः}
{आरुरोह रथं श्रीमाञ्शत्रुघ्नेन समन्वितः} %||2-107-8||

\twolineshloka
{आरुह्य तु रथं शीघ्रं शत्रुघ्नभरतावुभौ}
{ययतुः परमप्रीतौ वृतौ मन्त्रिपुरोहितैः} %||2-107-9||

\twolineshloka
{अग्रतो पुरवस्तत्र वसिष्ठ प्रमुखा द्विजाः}
{प्रययुः प्राङ्मुखाः सर्वे नन्दिग्रामो यतोऽभवत्} %||2-107-10||

\twolineshloka
{बलं च तदनाहूतं गजाश्वरथसङ्कुलम्}
{प्रययौ भरते याते सर्वे च पुरवासिनः} %||2-107-11||

\twolineshloka
{रथस्थः स तु धर्मात्मा भरतो भ्रातृवत्सलः}
{नन्दिग्रामं ययौ तूर्णं शिरस्याधाय पादुके} %||2-107-12||

\twolineshloka
{ततस्तु भरतः क्षिप्रं नन्दिग्रामं प्रविश्य सः}
{अवतीर्य रथात्तूर्णं गुरूनिदमुवाच ह} %||2-107-13||

\threelineshloka
{एतद्राज्यं मम भ्रात्रा दत्तं संन्यासवत्स्वयम्}
{योगक्षेमवहे चेमे पादुके हेमभूषिते}
{तमिमं पालयिष्यामि राघवागमनं प्रति} %||2-107-14||

\twolineshloka
{क्षिप्रं संयोजयित्वा तु राघवस्य पुनः स्वयम्}
{चरणौ तौ तु रामस्य द्रक्ष्यामि सहपादुकौ} %||2-107-15||

\twolineshloka
{ततो निक्षिप्तभारोऽहं राघवेण समागतः}
{निवेद्य गुरवे राज्यं भजिष्ये गुरुवृत्तिताम्} %||2-107-16||

\twolineshloka
{राघवाय च संन्यासं दत्त्वेमे वरपादुके}
{राज्यं चेदमयोध्यां च धूतपापो भवामि च} %||2-107-17||

\twolineshloka
{अभिषिक्ते तु काकुत्स्थे प्रहृष्टमुदिते जने}
{प्रीतिर्मम यशश्चैव भवेद्राज्याच्चतुर्गुणम्} %||2-107-18||

\twolineshloka
{एवं तु विलपन्दीनो भरतः स महायशाः}
{नन्दिग्रामेऽकरोद्राज्यं दुःखितो मन्त्रिभिः सह} %||2-107-19||

\twolineshloka
{स वल्कलजटाधारी मुनिवेषधरः प्रभुः}
{नन्दिग्रामेऽवसद्वीरः ससैन्यो भरतस्तदा} %||2-107-20||

\twolineshloka
{रामागमनमाकाङ्क्षन्भरतो भ्रातृवत्सलः}
{भ्रातुर्वचनकारी च प्रतिज्ञापारगस्तदा} %||2-107-21||

\twolineshloka
{पादुके त्वभिषिच्याथ नन्दिग्रामेऽवसत्तदा}
{भरतः शासनं सर्वं पादुकाभ्यां न्यवेदयत्} %||2-108-22||


{॥इत्यार्षे श्रीमद्रामायणे वाल्मीकीये आदिकाव्ये अयोध्याकाण्डे सप्तमाधिकशततमः सर्गः॥}

\sect{अष्टमाधिकशततमः सर्गः}

\twolineshloka
{प्रतिप्रयाते भरते वसन्रामस्तपोवने}
{लक्षयामास सोद्वेगमथौत्सुक्यं तपस्विनाम्} %||2-108-1||

\twolineshloka
{ये तत्र चित्रकूटस्य पुरस्तात्तापसाश्रमे}
{राममाश्रित्य निरतास्तानलक्षयदुत्सुकान्} %||2-108-2||

\twolineshloka
{नयनैर्भृकुटीभिश्च रामं निर्दिश्य शङ्किताः}
{अन्योन्यमुपजल्पन्तः शनैश्चक्रुर्मिथः कथाः} %||2-108-3||

\twolineshloka
{तेषामौत्सुक्यमालक्ष्य रामस्त्वात्मनि शङ्कितः}
{कृताञ्जलिरुवाचेदमृषिं कुलपतिं ततः} %||2-108-4||

\twolineshloka
{न कच्चिद्भगवन्किञ्चित्पूर्ववृत्तमिदं मयि}
{दृश्यते विकृतं येन विक्रियन्ते तपस्विनः} %||2-108-5||

\twolineshloka
{प्रमादाच्चरितं कच्चित्किञ्चिन्नावरजस्य मे}
{लक्ष्मणस्यर्षिभिर्दृष्टं नानुरूपमिवात्मनः} %||2-108-6||

\twolineshloka
{कच्चिच्छुश्रूषमाणा वः शुश्रूषणपरा मयि}
{प्रमदाभ्युचितां वृत्तिं सीता युक्तं न वर्तते} %||2-108-7||

\twolineshloka
{अथर्षिर्जरया वृद्धस्तपसा च जरां गतः}
{वेपमान इवोवाच रामं भूतदयापरम्} %||2-108-8||

\twolineshloka
{कुतः कल्याणसत्त्वायाः कल्याणाभिरतेस्तथा}
{चलनं तात वैदेह्यास्तपस्विषु विशेषतः} %||2-108-9||

\twolineshloka
{त्वन्निमित्तमिदं तावत्तापसान्प्रति वर्तते}
{रक्षोभ्यस्तेन संविग्नाः कथयन्ति मिथः कथाः} %||2-108-10||

\twolineshloka
{रावणावरजः कश्चित्खरो नामेह राक्षसः}
{उत्पाट्य तापसान्सर्वाञ्जनस्थाननिकेतनान्} %||2-108-11||

\twolineshloka
{धृष्टश्च जितकाशी च नृशंसः पुरुषादकः}
{अवलिप्तश्च पापश्च त्वां च तात न मृष्यते} %||2-108-12||

\twolineshloka
{त्वं यदा प्रभृति ह्यस्मिन्नाश्रमे तात वर्तसे}
{तदा प्रभृति रक्षांसि विप्रकुर्वन्ति तापसान्} %||2-108-13||

\twolineshloka
{दर्शयन्ति हि बीभत्सैः क्रूरैर्भीषणकैरपि}
{नाना रूपैर्विरूपैश्च रूपैरसुखदर्शनैः} %||2-108-14||

\twolineshloka
{अप्रशस्तैरशुचिभिः सम्प्रयोज्य च तापसान्}
{प्रतिघ्नन्त्यपरान्क्षिप्रमनार्याः पुरतः स्थितः} %||2-108-15||

\twolineshloka
{तेषु तेष्वाश्रमस्थानेष्वबुद्धमवलीय च}
{रमन्ते तापसांस्तत्र नाशयन्तोऽल्पचेतसः} %||2-108-16||

\twolineshloka
{अपक्षिपन्ति स्रुग्भाण्डानग्नीन्सिञ्चन्ति वारिणा}
{कलशांश्च प्रमृद्नन्ति हवने समुपस्थिते} %||2-108-17||

\twolineshloka
{तैर्दुरात्मभिराविष्टानाश्रमान्प्रजिहासवः}
{गमनायान्यदेशस्य चोदयन्त्यृषयोऽद्य माम्} %||2-108-18||

\twolineshloka
{तत्पुरा राम शारीरामुपहिंसां तपस्विषु}
{दर्शयति हि दुष्टास्ते त्यक्ष्याम इममाश्रमम्} %||2-108-19||

\twolineshloka
{बहुमूलफलं चित्रमविदूरादितो वनम्}
{पुराणाश्रममेवाहं श्रयिष्ये सगणः पुनः} %||2-108-20||

\twolineshloka
{खरस्त्वय्यपि चायुक्तं पुरा तात प्रवर्तते}
{सहास्माभिरितो गच्छ यदि बुद्धिः प्रवर्तते} %||2-108-21||

\twolineshloka
{सकलत्रस्य सन्देहो नित्यं यत्तस्य राघव}
{समर्थस्यापि हि सतो वासो दुःख इहाद्य ते} %||2-108-22||

\twolineshloka
{इत्युक्तवन्तं रामस्तं राजपुत्रस्तपस्विनम्}
{न शशाकोत्तरैर्वाक्यैरवरोद्धुं समुत्सुकम्} %||2-108-23||

\twolineshloka
{अभिनन्द्य समापृच्छ्य समाधाय च राघवम्}
{स जगामाश्रमं त्यक्त्वा कुलैः कुलपतिः सह} %||2-108-24||

\fourlineindentedshloka
{रामः संसाध्य त्वृषिगणमनुगमना}
{द्देशात्तस्माच्चित्कुलपतिमभिवाद्यर्षिम्}
{सम्यक्प्रीतैस्तैरनुमत उपदिष्टार्थः}
{ पुण्यं वासाय स्वनिलयमुपसम्पेदे} %||2-108-25||

\fourlineindentedshloka
{आश्रमं त्वृषिविरहितं प्रभुः}
{ क्षणमपि न जहौ स राघवः}
{राघवं हि सततमनुगता}
{स्तापसाश्चर्षिचरितधृतगुणाः} %||2-109-26||


{॥इत्यार्षे श्रीमद्रामायणे वाल्मीकीये आदिकाव्ये अयोध्याकाण्डे अष्टमाधिकशततमः सर्गः॥}

\sect{नवमाधिकशततमः सर्गः}

\twolineshloka
{राघवस्त्वपयातेषु तपस्विषु विचिन्तयन्}
{न तत्रारोचयद्वासं कारणैर्बहुभिस्तदा} %||2-109-1||

\twolineshloka
{इह मे भरतो दृष्टो मातरश्च सनागराः}
{सा च मे स्मृतिरन्वेति तान्नित्यमनुशोचतः} %||2-109-2||

\twolineshloka
{स्कन्धावारनिवेशेन तेन तस्य महात्मनः}
{हयहस्तिकरीषैश्च उपमर्दः कृतो भृशम्} %||2-109-3||

\twolineshloka
{तस्मादन्यत्र गच्छाम इति सञ्चिन्त्य राघवः}
{प्रातिष्ठत स वैदेह्या लक्ष्मणेन च सङ्गतः} %||2-109-4||

\twolineshloka
{सोऽत्रेराश्रममासाद्य तं ववन्दे महायशाः}
{तं चापि भगवानत्रिः पुत्रवत्प्रत्यपद्यत} %||2-109-5||

\twolineshloka
{स्वयमातिथ्यमादिश्य सर्वमस्य सुसत्कृतम्}
{सौमित्रिं च महाभागां सीतां च समसान्त्वयत्} %||2-109-6||

\twolineshloka
{पत्नीं च तमनुप्राप्तां वृद्धामामन्त्र्य सत्कृताम्}
{सान्त्वयामास धर्मज्ञः सर्वभूतहिते रतः} %||2-109-7||

\twolineshloka
{अनसूयां महाभागां तापसीं धर्मचारिणीम्}
{प्रतिगृह्णीष्व वैदेहीमब्रवीदृषिसत्तमः} %||2-109-8||

\twolineshloka
{रामाय चाचचक्षे तां तापसीं धर्मचारिणीम्}
{दश वर्षाण्यनावृष्ट्या दग्धे लोके निरन्तरम्} %||2-109-9||

\twolineshloka
{यया मूलफले सृष्टे जाह्नवी च प्रवर्तिता}
{उग्रेण तपसा युक्ता नियमैश्चाप्यलङ्कृता} %||2-109-10||

\twolineshloka
{दशवर्षसहस्राणि यया तप्तं महत्तपः}
{अनसूयाव्रतैस्तात प्रत्यूहाश्च निबर्हिताः} %||2-109-11||

\twolineshloka
{देवकार्यनिमित्तं च यया सन्त्वरमाणया}
{दशरात्रं कृत्वा रात्रिः सेयं मातेव तेऽनघ} %||2-109-12||

\twolineshloka
{तामिमां सर्वभूतानां नमस्कार्यां यशस्विनीम्}
{अभिगच्छतु वैदेही वृद्धामक्रोधनां सदा} %||2-109-13||

\twolineshloka
{एवं ब्रुवाणं तमृषिं तथेत्युक्त्वा स राघवः}
{सीतामुवाच धर्मज्ञामिदं वचनमुत्तमम्} %||2-109-14||

\twolineshloka
{राजपुत्रि श्रुतं त्वेतन्मुनेरस्य समीरितम्}
{श्रेयोऽर्थमात्मनः शीघ्रमभिगच्छ तपस्विनीम्} %||2-109-15||

\twolineshloka
{अनसूयेति या लोके कर्मभिः क्यातिमागता}
{तां शीघ्रमभिगच्छ त्वमभिगम्यां तपस्विनीम्} %||2-109-16||

\twolineshloka
{सीता त्वेतद्वचः श्रुत्वा राघवस्य हितैषिणी}
{तामत्रिपत्नीं धर्मज्ञामभिचक्राम मैथिली} %||2-109-17||

\twolineshloka
{शिथिलां वलितां वृद्धां जरापाण्डुरमूर्धजाम्}
{सततं वेपमानाङ्गीं प्रवाते कदली यथा} %||2-109-18||

\twolineshloka
{तां तु सीता महाभागामनसूयां पतिव्रताम्}
{अभ्यवादयदव्यग्रा स्वं नाम समुदाहरत्} %||2-109-19||

\twolineshloka
{अभिवाद्य च वैदेही तापसीं तामनिन्दिताम्}
{बद्धाञ्जलिपुटा हृष्टा पर्यपृच्छदनामयम्} %||2-109-20||

\twolineshloka
{ततः सीतां महाभागां दृष्ट्वा तां धर्मचारिणीम्}
{सान्त्वयन्त्यब्रवीद्धृष्टा दिष्ट्या धर्ममवेक्षसे} %||2-109-21||

\twolineshloka
{त्यक्त्वा ज्ञातिजनं सीते मानमृद्धिं च मानिनि}
{अवरुद्धं वने रामं दिष्ट्या त्वमनुगच्छसि} %||2-109-22||

\twolineshloka
{नगरस्थो वनस्थो वा पापो वा यदि वाशुभः}
{यासां स्त्रीणां प्रियो भर्ता तासां लोका महोदयाः} %||2-109-23||

\twolineshloka
{दुःशीलः कामवृत्तो वा धनैर्वा परिवर्जितः}
{स्त्रीणामार्य स्वभावानां परमं दैवतं पतिः} %||2-109-24||

\twolineshloka
{नातो विशिष्टं पश्यामि बान्धवं विमृशन्त्यहम्}
{सर्वत्र योग्यं वैदेहि तपः कृतमिवाव्ययम्} %||2-109-25||

\twolineshloka
{न त्वेवमवगच्छन्ति गुण दोषमसत्स्त्रियः}
{कामवक्तव्यहृदया भर्तृनाथाश्चरन्ति याः} %||2-109-26||

\twolineshloka
{प्राप्नुवन्त्ययशश्चैव धर्मभ्रंशं च मैथिलि}
{अकार्य वशमापन्नाः स्त्रियो याः खलु तद्विधाः} %||2-109-27||

\twolineshloka
{त्वद्विधास्तु गुणैर्युक्ता दृष्टलोकपरावराः}
{स्त्रियः स्वर्गे चरिष्यन्ति यथा पुण्यकृतस्तथा} %||2-110-28||


{॥इत्यार्षे श्रीमद्रामायणे वाल्मीकीये आदिकाव्ये अयोध्याकाण्डे नवमाधिकशततमः सर्गः॥}

\sect{दशमाधिकशततमः सर्गः}

\twolineshloka
{सा त्वेवमुक्ता वैदेही अनसूयानसूयया}
{प्रतिपूज्य वचो मन्दं प्रवक्तुमुपचक्रमे} %||2-110-1||

\twolineshloka
{नैतदाश्चर्यमार्याया यन्मां त्वमनुभाषसे}
{विदितं तु ममाप्येतद्यथा नार्याः पतिर्गुरुः} %||2-110-2||

\twolineshloka
{यद्यप्येष भवेद्भर्ता ममार्ये वृत्तवर्जितः}
{अद्वैधमुपवर्तव्यस्तथाप्येष मया भवेत्} %||2-110-3||

\twolineshloka
{किं पुनर्यो गुणश्लाघ्यः सानुक्रोशो जितेन्द्रियः}
{स्थिरानुरागो धर्मात्मा मातृवर्ती पितृ प्रियः} %||2-110-4||

\twolineshloka
{यां वृत्तिं वर्तते रामः कौसल्यायां महाबलः}
{तामेव नृपनारीणामन्यासामपि वर्तते} %||2-110-5||

\twolineshloka
{सकृद्दृष्टास्वपि स्त्रीषु नृपेण नृपवत्सलः}
{मातृवद्वर्तते वीरो मानमुत्सृज्य धर्मवित्} %||2-110-6||

\twolineshloka
{आगच्छन्त्याश्च विजनं वनमेवं भयावहम्}
{समाहितं हि मे श्वश्र्वा हृदये यत्स्थितं मम} %||2-110-7||

\twolineshloka
{प्राणिप्रदानकाले च यत्पुरा त्वग्निसंनिधौ}
{अनुशिष्टा जनन्यास्मि वाक्यं तदपि मे धृतम्} %||2-110-8||

\twolineshloka
{नवीकृतं तु तत्सर्वं वाक्यैस्ते धर्मचारिणि}
{पतिशुश्रूषणान्नार्यास्तपो नान्यद्विधीयते} %||2-110-9||

\twolineshloka
{सावित्री पतिशुश्रूषां कृत्वा स्वर्गे महीयते}
{तथा वृत्तिश्च याता त्वं पतिशुश्रूषया दिवम्} %||2-110-10||

\twolineshloka
{वरिष्ठा सर्वनारीणामेषा च दिवि देवता}
{रोहिणी च विना चन्द्रं मुहूर्तमपि दृश्यते} %||2-110-11||

\twolineshloka
{एवंविधाश्च प्रवराः स्त्रियो भर्तृदृढव्रताः}
{देवलोके महीयन्ते पुण्येन स्वेन कर्मणा} %||2-110-12||

\twolineshloka
{ततोऽनसूया संहृष्टा श्रुत्वोक्तं सीतया वचः}
{शिरस्याघ्राय चोवाच मैथिलीं हर्षयन्त्युत} %||2-110-13||

\twolineshloka
{नियमैर्विविधैराप्तं तपो हि महदस्ति मे}
{तत्संश्रित्य बलं सीते छन्दये त्वां शुचिव्रते} %||2-110-14||

\threelineshloka
{उपपन्नं च युक्तं च वचनं तव मैथिलि}
{प्रीता चास्म्युचितं किं ते करवाणि ब्रवीहि मे}
{कृतमित्यब्रवीत्सीता तपोबलसमन्विताम्} %||2-110-15||

\twolineshloka
{सा त्वेवमुक्ता धर्मज्ञा तया प्रीततराभवत्}
{सफलं च प्रहर्षं ते हन्त सीते करोम्यहम्} %||2-110-16||

\twolineshloka
{इदं दिव्यं वरं माल्यं वस्त्रमाभरणानि च}
{अङ्गरागं च वैदेहि महार्हमनुलेपनम्} %||2-110-17||

\twolineshloka
{मया दत्तमिदं सीते तव गात्राणि शोभयेत्}
{अनुरूपमसङ्क्लिष्टं नित्यमेव भविष्यति} %||2-110-18||

\twolineshloka
{अङ्गरागेण दिव्येन लिप्ताङ्गी जनकात्मजे}
{शोभयिष्यामि भर्तारं यथा श्रीर्विष्णुमव्ययम्} %||2-110-19||

\twolineshloka
{सा वस्त्रमङ्गरागं च भूषणानि स्रजस्तथा}
{मैथिली प्रतिजग्राह प्रीतिदानमनुत्तमम्} %||2-110-20||

\twolineshloka
{प्रतिगृह्य च तत्सीता प्रीतिदानं यशस्विनी}
{श्लिष्टाञ्जलिपुटा धीरा समुपास्त तपोधनाम्} %||2-110-21||

\twolineshloka
{तथा सीतामुपासीनामनसूया दृढव्रता}
{वचनं प्रष्टुमारेभे कथां काञ्चिदनुप्रियाम्} %||2-110-22||

\twolineshloka
{स्वयंवरे किल प्राप्ता त्वमनेन यशस्विना}
{राघवेणेति मे सीते कथा श्रुतिमुपागता} %||2-110-23||

\twolineshloka
{तां कथां श्रोतुमिच्छामि विस्तरेण च मैथिलि}
{यथानुभूतं कार्त्स्न्येन तन्मे त्वं वक्तुमर्हसि} %||2-110-24||

\twolineshloka
{एवमुक्ता तु सा सीता तां ततो धर्मचारिणीम्}
{श्रूयतामिति चोक्त्वा वै कथयामास तां कथाम्} %||2-110-25||

\twolineshloka
{मिथिलाधिपतिर्वीरो जनको नाम धर्मवित्}
{क्षत्रधर्मण्यभिरतो न्यायतः शास्ति मेदिनीम्} %||2-110-26||

\twolineshloka
{तस्य लाङ्गलहस्तस्य कर्षतः क्षेत्रमण्डलम्}
{अहं किलोत्थिता भित्त्वा जगतीं नृपतेः सुता} %||2-110-27||

\twolineshloka
{स मां दृष्ट्वा नरपतिर्मुष्टिविक्षेपतत्परः}
{पांशु गुण्ठित सर्वाङ्गीं विस्मितो जनकोऽभवत्} %||2-110-28||

\twolineshloka
{अनपत्येन च स्नेहादङ्कमारोप्य च स्वयम्}
{ममेयं तनयेत्युक्त्वा स्नेहो मयि निपातितः} %||2-110-29||

\twolineshloka
{अन्तरिक्षे च वागुक्ताप्रतिमा मानुषी किल}
{एवमेतन्नरपते धर्मेण तनया तव} %||2-110-30||

\twolineshloka
{ततः प्रहृष्टो धर्मात्मा पिता मे मिथिलाधिपः}
{अवाप्तो विपुलामृद्धिं मामवाप्य नराधिपः} %||2-110-31||

\twolineshloka
{दत्त्वा चास्मीष्टवद्देव्यै ज्येष्ठायै पुण्यकर्मणा}
{तया सम्भाविता चास्मि स्निग्धया मातृसौहृदात्} %||2-110-32||

\twolineshloka
{पतिसंयोगसुलभं वयो दृष्ट्वा तु मे पिता}
{चिन्तामभ्यगमद्दीनो वित्तनाशादिवाधनः} %||2-110-33||

\twolineshloka
{सदृशाच्चापकृष्टाच्च लोके कन्यापिता जनात्}
{प्रधर्षणामवाप्नोति शक्रेणापि समो भुवि} %||2-110-34||

\twolineshloka
{तां धर्षणामदूरस्थां सन्दृश्यात्मनि पार्थिवः}
{चिन्न्तार्णवगतः पारं नाससादाप्लवो यथ} %||2-110-35||

\twolineshloka
{अयोनिजां हि मां ज्ञात्वा नाध्यगच्छत्स चिन्तयन्}
{सदृशं चानुरूपं च महीपालः पतिं मम} %||2-110-36||

\twolineshloka
{तस्य बुद्धिरियं जाता चिन्तयानस्य सन्ततम्}
{स्वयं वरं तनूजायाः करिष्यामीति धीमतः} %||2-110-37||

\twolineshloka
{महायज्ञे तदा तस्य वरुणेन महात्मना}
{दत्तं धनुर्वरं प्रीत्या तूणी चाक्षय्य सायकौ} %||2-110-38||

\twolineshloka
{असञ्चाल्यं मनुष्यैश्च यत्नेनापि च गौरवात्}
{तन्न शक्ता नमयितुं स्वप्नेष्वपि नराधिपाः} %||2-110-39||

\twolineshloka
{तद्धनुः प्राप्य मे पित्रा व्याहृतं सत्यवादिना}
{समवाये नरेन्द्राणां पूर्वमामन्त्र्य पार्थिवान्} %||2-110-40||

\twolineshloka
{इदं च धनुरुद्यम्य सज्यं यः कुरुते नरः}
{तस्य मे दुहिता भार्या भविष्यति न संशयः} %||2-110-41||

\twolineshloka
{तच्च दृष्ट्वा धनुःश्रेष्ठं गौरवाद्गिरिसंनिभम्}
{अभिवाद्य नृपा जग्मुरशक्तास्तस्य तोलने} %||2-110-42||

\twolineshloka
{सुदीर्घस्य तु कालस्य राघवोऽयं महाद्युतिः}
{विश्वामित्रेण सहितो यज्ञं द्रष्टुं समागतः} %||2-110-43||

\twolineshloka
{लक्ष्मणेन सह भ्रात्रा रामः सत्यपराक्रमः}
{विश्वामित्रस्तु धर्मात्मा मम पित्रा सुपूजितः} %||2-110-44||

\threelineshloka
{प्रोवाच पितरं तत्र राघवो रामलक्ष्मणौ}
{सुतौ दशरथस्येमौ धनुर्दर्शनकाङ्क्षिणौ}
{इत्युक्तस्तेन विप्रेण तद्धनुः समुपानयत्} %||2-110-45||

\twolineshloka
{निमेषान्तरमात्रेण तदानम्य स वीर्यवान्}
{ज्यां समारोप्य झटिति पूरयामास वीर्यवान्} %||2-110-46||

\twolineshloka
{तेन पूरयता वेगान्मध्ये भग्नं द्विधा धनुः}
{तस्य शब्दोऽभवद्भीमः पतितस्याशनेरिव} %||2-110-47||

\twolineshloka
{ततोऽहं तत्र रामाय पित्रा सत्याभिसन्धिना}
{उद्यता दातुमुद्यम्य जलभाजनमुत्तमम्} %||2-110-48||

\twolineshloka
{दीयमानां न तु तदा प्रतिजग्राह राघवः}
{अविज्ञाय पितुश्छन्दमयोध्याधिपतेः प्रभोः} %||2-110-49||

\twolineshloka
{ततः श्वशुरमामन्त्र्य वृद्धं दशरथं नृपम्}
{मम पित्रा अहं दत्ता रामाय विदितात्मने} %||2-110-50||

\twolineshloka
{मम चैवानुजा साध्वी ऊर्मिला प्रियदर्शना}
{भार्यार्थे लक्ष्मणस्यापि दत्ता पित्रा मम स्वयम्} %||2-110-51||

\twolineshloka
{एवं दत्तास्मि रामाय तदा तस्मिन्स्वयं वरे}
{अनुरक्ता च धर्मेण पतिं वीर्यवतां वरम्} %||2-111-52||


{॥इत्यार्षे श्रीमद्रामायणे वाल्मीकीये आदिकाव्ये अयोध्याकाण्डे दशमाधिकशततमः सर्गः॥}

\sect{एकादशाधिकशततमः सर्गः}

\twolineshloka
{अनसूया तु धर्मज्ञा श्रुत्वा तां महतीं कथाम्}
{पर्यष्वजत बाहुभ्यां शिरस्याघ्राय मैथिलीम्} %||2-111-1||

\twolineshloka
{व्यक्ताक्षरपदं चित्रं भाषितं मधुरं त्वया}
{यथा स्वयंवरं वृत्तं तत्सर्वं हि श्रुतं मया} %||2-111-2||

\twolineshloka
{रमेऽहं कथया ते तु दृष्ढं मधुरभाषिणि}
{रविरस्तं गतः श्रीमानुपोह्य रजनीं शिवाम्} %||2-111-3||

\twolineshloka
{दिवसं प्रति कीर्णानामाहारार्थं पतत्रिणाम्}
{सन्ध्याकाले निलीनानां निद्रार्थं श्रूयते ध्वनिः} %||2-111-4||

\twolineshloka
{एते चाप्यभिषेकार्द्रा मुनयः फलशोधनाः}
{सहिता उपवर्तन्ते सलिलाप्लुतवल्कलाः} %||2-111-5||

\twolineshloka
{ऋषीणामग्निहोत्रेषु हुतेषु विधिपुर्वकम्}
{कपोताङ्गारुणो धूमो दृश्यते पवनोद्धतः} %||2-111-6||

\twolineshloka
{अल्पपर्णा हि तरवो घनीभूताः समन्ततः}
{विप्रकृष्टेऽपि ये देशे न प्रकाशन्ति वै दिशः} %||2-111-7||

\twolineshloka
{रजनी रससत्त्वानि प्रचरन्ति समन्ततः}
{तपोवनमृगा ह्येते वेदितीर्थेषु शेरते} %||2-111-8||

\twolineshloka
{सम्प्रवृत्ता निशा सीते नक्षत्रसमलङ्कृता}
{ज्योत्स्ना प्रावरणश्चन्द्रो दृश्यतेऽभ्युदितोऽम्बरे} %||2-111-9||

\twolineshloka
{गम्यतामनुजानामि रामस्यानुचरी भव}
{कथयन्त्या हि मधुरं त्वयाहं परितोषिता} %||2-111-10||

\twolineshloka
{अलङ्कुरु च तावत्त्वं प्रत्यक्षं मम मैथिलि}
{प्रीतिं जनय मे वत्स दिव्यालङ्कारशोभिनी} %||2-111-11||

\twolineshloka
{सा तदा समलङ्कृत्य सीता सुरसुतोपमा}
{प्रणम्य शिरसा तस्यै रामं त्वभिमुखी ययौ} %||2-111-12||

\twolineshloka
{तथा तु भूषितां सीतां ददर्श वदतां वरः}
{राघवः प्रीतिदानेन तपस्विन्या जहर्ष च} %||2-111-13||

\twolineshloka
{न्यवेदयत्ततः सर्वं सीता रामाय मैथिली}
{प्रीतिदानं तपस्विन्या वसनाभरणस्रजाम्} %||2-111-14||

\twolineshloka
{प्रहृष्टस्त्वभवद्रामो लक्ष्मणश्च महारथः}
{मैथिल्याः सत्क्रियां दृष्ट्वा मानुषेषु सुदुर्लभाम्} %||2-111-15||

\twolineshloka
{ततस्तां सर्वरीं प्रीतः पुण्यां शशिनिभाननः}
{अर्चितस्तापसैः सिद्धैरुवास रघुनन्दनः} %||2-111-16||

\twolineshloka
{तस्यां रात्र्यां व्यतीतायामभिषिच्य हुताग्निकान्}
{आपृच्छेतां नरव्याघ्रौ तापसान्वनगोचरान्} %||2-111-17||

\twolineshloka
{तावूचुस्ते वनचरास्तापसा धर्मचारिणः}
{वनस्य तस्य सञ्चारं राक्षसैः समभिप्लुतम्} %||2-111-18||

\twolineshloka
{एष पन्था महर्षीणां फलान्याहरतां वने}
{अनेन तु वनं दुर्गं गन्तुं राघव ते क्षमम्} %||2-111-19||

\fourlineindentedshloka
{इतीव तैः प्राञ्जलिभिस्तपस्विभि}
{र्द्विजैः कृतस्वस्त्ययनः परन्तपः}
{वनं सभार्यः प्रविवेश राघवः}
{ सलक्ष्मणः सूर्य इवाभ्रमण्डलम्} %||2-112-20||


{॥इत्यार्षे श्रीमद्रामायणे वाल्मीकीये आदिकाव्ये अयोध्याकाण्डे एकादशाधिकशततमः सर्गः॥}

